\hypertarget{nace_kartlar_10-12}{%
\chapter{Faz 3 --- NACE sınıf kartları, parça 1: 10 Gıda, 11 İçecekler, 12 Tütün ürünleri}\label{nace_kartlar_10-12}}

\textbf{Durum:} Faz 3 / 9, parça 1 / 16 tamamlandı (18 Ekim 2026); DEVAM bekleniyor. Önceki fazlar: \belgebag{01_teshis}{}{\kod{01\_\allowbreak{}teshis.\allowbreak{}md}}, \belgebag{02a_politika_uzayi}{}{\kod{02a\_\allowbreak{}politika\_\allowbreak{}uzayi.\allowbreak{}md}}, \belgebag{02b_secenekler_hukuk}{}{\kod{02b\_\allowbreak{}secenekler\_\allowbreak{}hukuk.\allowbreak{}md}}, \belgebag{02c_hazirlik}{}{\kod{02c\_\allowbreak{}hazirlik.\allowbreak{}md}}. Defter: \belgebag{defter}{}{\kod{defter.\allowbreak{}md}}. Ana matris (230 sınıfın tamamı): \href{../veri/nace_matris.csv}{\kod{veri/\allowbreak{}nace\_\allowbreak{}matris.\allowbreak{}csv}}.

\textbf{Bu belge ne yapıyor?}

\begin{itemize}
\tightlist
\item
  10--12 bölümlerinin 33 sınıfı için ana promptun 15 alanlı kartını dolduruyor.
\item
  Bölümlerin kısa genel görünümünü ve ana matris satırlarını veriyor.
\item
  Kartlara temel olan Faz 3 hesap altyapısını, bütün 230 C sınıfı için bir kez kuruyor. Bu altyapı ticaret, yapı, karmaşıklık, AB tarifesi, senaryolar ve skorları kapsıyor. Sonraki parçalar aynı tabloları kullanacak.
\item
  Kümülatif defteri güncelliyor.
\end{itemize}

\textbf{Takvim gösterimi.} Eylem tablolarında tarihler \kod{T}'ye göre yazılıyor. T2028'de \kod{T} = 1.1.2028, T2029'da \kod{T} = 1.1.2029; örneğin \kod{T−12} T2028'de Ocak 2027'ye, T2029'da Ocak 2028'e denk geliyor. Kanıt etiketleri (\etiketV{}, \etiketT{}, \etiketD{}, \etiketH{}) ve güven düzeyleri önceki fazlardaki gibi.

\textbf{Ana bulgu: 10--12 GB çıkışından en az etkilenen büyük sanayi grubu, ama üç ince nokta var.}

\begin{enumerate}
\def\labelenumi{\arabic{enumi}.}
\tightlist
\item
  \textbf{AB'ye ihracatın dörtte üçü tarım rejiminde (1/98).} Bu ihracat zaten menşe kuralına tabi ve GB'ye bağlı değil.

  \begin{itemize}
  \tightlist
  \item
    Rota A'da (1/98 sürer) S4'ün getirdiği ortalama AB tarifesi \%3,3. Bu oran yalnızca işlenmiş tarım ürünlerinin sanayi bileşeninden ve sigaradan geliyor.
  \item
    1/98 da kaybedilirse oran \%11,5'e çıkıyor.
  \item
    Riskin asıl kaynağı GB değil, hasmane bir AB'nin 1/98'i askıya alması. Bu durumda S4'ün P10 değerlerinde AB'ye ihracat −\$2.018 mn, istihdam −4.698 kişi oluyor. Kaybın yarıdan fazlası 10.39'da (kuru meyve, fındık, işlenmiş sebze).
  \end{itemize}
\item
  \textbf{İşlenmiş tarım ürünlerinde (İTÜ) S4'ün maliyeti sanayi bileşeni kadar.} Tarım bileşeni bugün de ödeniyor.

  \begin{itemize}
  \tightlist
  \item
    Sanayi bileşeni: bisküvi ve ekmek \%9; çikolata \%8,4; makarna \%7,7; glikoz \%14; gıda müstahzarları \%12,8; sigara \%57,6.
  \item
    En büyük mutlak kayıplar: bisküvi, gıda müstahzarları, glikoz ve sigara.
  \end{itemize}
\item
  \textbf{S2'de menşe kuralı tarifeden daha bağlayıcı olabilir.} Bisküvi ve ekmek ihracatçıları Dahilde İşleme Rejimi (DİR) ile ithal (çoğunlukla Rus) buğday kullanıyor. PEM'in Fasıl 11--19 kuralları ise unun menşeli olmasını istiyor.

  \begin{itemize}
  \tightlist
  \item
    Bisküvide S2'nin kaybı (−\$35 mn AB ihracatı) S4'ünkinin (−\$51 mn) üçte ikisinden fazla.
  \item
    Çözüm tarifede değil, DİR'deki ``eşdeğer eşya'' mekanizmasında ve ürüne özgü menşe kuralı (PSR) müzakeresinde (E-10-02).
  \end{itemize}
\end{enumerate}

\textbf{Olumsuz bulgular (dürüst rapor).}

\begin{itemize}
\tightlist
\item
  Sigarada (12.00) S2'de bile kayıp büyük. Kural ``tütünün \%70'i menşeli'' diyor ve Amerikan harmanı bunu karşılamıyor. AB'ye yönelik üretim büyük olasılıkla AB fabrikalarına kayar.
\item
  Fındık ve kuru meyvede Türkiye dünya ihracatının \%30--75'ini yapıyor. Bu, AB'ye karşı bir müzakere kaldıracı. Ama 1/98 kaybında bile fındık içinin AB MFN'si yalnızca \%3,2. Kaldıraç tarife tehdidinden değil, arz güvenliğinden geliyor.
\end{itemize}

\hypertarget{nace_kartlar_10-12-0-faz-3-yuxf6ntemi-buxfctuxfcn-paruxe7alar-iuxe7in}{%
\section{0. Faz 3 yöntemi (bütün parçalar için)}\label{nace_kartlar_10-12-0-faz-3-yuxf6ntemi-buxfctuxfcn-paruxe7alar-iuxe7in}}

\Needspace{7\baselineskip}

\textbf{Tablo 0.1 --- Faz 3 hesap zinciri}

\nopagebreak[4]

\begingroup\gbtablo\small

\begin{longtable}[]{@{}
  >{\raggedright\arraybackslash}p{(\columnwidth - 6\tabcolsep) * \real{0.1442}}
  >{\raggedright\arraybackslash}p{(\columnwidth - 6\tabcolsep) * \real{0.4430}}
  >{\raggedright\arraybackslash}p{(\columnwidth - 6\tabcolsep) * \real{0.2589}}
  >{\raggedright\arraybackslash}p{(\columnwidth - 6\tabcolsep) * \real{0.1540}}@{}}
\toprule\noalign{}
\begin{minipage}[b]{\linewidth}\raggedright
\textbf{Betik}
\end{minipage} & \begin{minipage}[b]{\linewidth}\raggedright
\textbf{Ne üretiyor}
\end{minipage} & \begin{minipage}[b]{\linewidth}\raggedright
\textbf{Ana kaynak}
\end{minipage} & \begin{minipage}[b]{\linewidth}\raggedright
\textbf{Etiket}
\end{minipage} \\
\midrule\noalign{}
\endhead
\bottomrule\noalign{}
\endlastfoot
\kod{kod/\allowbreak{}f3a\_\allowbreak{}ticaret.\allowbreak{}py} & Sınıf düzeyi ihracat, ithalat, AB ve Çin payları, RCA (dünya ve AB pazarı), AB pazar payı, birim değer, ilk 5 pazar/kaynak/rakip & BACI HS17 2020 ve 2024 (TR ikili; AB27 ithalatı; dünya ihracatı) & \etiketV{} veri, \etiketT{} eşleme \\
\kod{kod/\allowbreak{}f3b\_\allowbreak{}yapi.\allowbreak{}py} & Üretim, KD, istihdam, ölçek dağılımı, İBBS-2 bölge payları, AB'ye göre verimlilik & TÜİK YSHİ 2024, İş Kayıtları 2024, SGK 2025, Eurostat SBS 2023 & \etiketV{}; bölge ve gizli hücre \etiketT{} \\
\kod{kod/\allowbreak{}f3c\_\allowbreak{}karmasiklik.\allowbreak{}py} & PCI, ECI, ürün uzayı yoğunluğu, karmaşıklık kazancı, atlanabilir ürünler & BACI 2024 ihracatçı × HS6 (157 ülke, 4.858 ürün) & \etiketT{} \\
\kod{kod/\allowbreak{}f3d\_\allowbreak{}ab\_\allowbreak{}tarife.\allowbreak{}py} & KN8 düzeyinde AB MFN'si, spesifik vergilerin ad valorem eşdeğeri, rejime göre S4 artışı, tepe satırlar, EUDR/SKDM payları & AB KN 2026, Comext 2025 (değer ve miktar) & \etiketV{}; AVE \etiketT{} \\
\kod{kod/\allowbreak{}f3e\_\allowbreak{}senaryo.\allowbreak{}py} & S2/S3/S4 sınıf etkileri (P10--P50--P90), EDE, CKE & f3a--f3d, I-O 2023, ProTEE esneklikleri, Faz 2B GE & \etiketT{} \\
\kod{kod/\allowbreak{}f3f\_\allowbreak{}skor\_\allowbreak{}matris.\allowbreak{}py} & KS, SS, AME, EDE ve CKE puanları, arketip, duruş, ağırlık duyarlılığı, \kod{veri/\allowbreak{}nace\_\allowbreak{}matris.\allowbreak{}csv} & f3a--f3e, \kod{veri/\allowbreak{}f3\_\allowbreak{}ss\_\allowbreak{}boyutlar.\allowbreak{}csv} & \etiketT{}; SS \etiketH{} \\
\kod{kod/\allowbreak{}f3g\_\allowbreak{}kartlar.\allowbreak{}py} & Bu belge (kartların sayısal alanları; nitel alanlar \kod{kartlar/\allowbreak{}10-\allowbreak{}12*.\allowbreak{}yaml}) & --- & --- \\
\end{longtable}

\endgroup

\textbf{Ölçümün beş kuralı.}

\begin{enumerate}
\def\labelenumi{\arabic{enumi}.}
\tightlist
\item
  \textbf{Senaryolar genel dengeye bağlı.} Her Monte Carlo çekilişinde (N = 2.000) AB'ye toplam ihracat ve ithalat değişimi Faz 2B'nin genel denge dağılımından alınıyor. S4: −\%18,9 (GA90 −32/−5); S2: −\%11,1 (−27/+6). Kısmi denge modeli (\kod{ΔX =\allowbreak{} X[(\allowbreak{}1+\allowbreak{}τ\_\allowbreak{}etkin)\textasciicircum{}(\allowbreak{}−σ) − 1]}) yalnızca bu toplamı sınıflara dağıtıyor; S3 iki senaryonun kısmi denge yapısıyla ara değer alıyor. Kalibrasyon ölçeği S2'de 0,65, S4'te 0,54: ürün düzeyi esnekliklerle kısmi denge, genel dengenin yaklaşık iki katı kayıp üretiyor (\kod{V-\allowbreak{}30}).
\item
  \textbf{Esneklikler ürün düzeyinde.} CEPII ProTEE (Fontagné, Guimbard ve Orefice 2022) HS6 esneklikleri sınıf içinde AB ihracatıyla ağırlıklandırılıp {[}2, 15{]} aralığına kırpıldı; medyan 7,8 (\kod{V-\allowbreak{}32}).
\item
  \textbf{Tarife rejime göre ayrıştırılıyor.} S4'te GB satırına tam MFN, İTÜ satırına yalnızca sanayi bileşeni ekleniyor; tarım satırı (1/98) değişmiyor. Hasmane AB varsayımıyla (\kod{V-\allowbreak{}11}) S4'te 1/98'in askıya alınma olasılığı 0,25 (\kod{V-\allowbreak{}31}). 1/98 kaybı genel denge toplamına ekleniyor, yeniden dağıtılmıyor.
\item
  \textbf{Sürtünme ve menşe ayrı kalemler.} Sınır sürtünmesi \%0,5--2 ve yalnızca GB ve İTÜ satırlarına uygulanıyor; S2'de AEO karşılıklı tanıma ve tek pencereyle yarıya iniyor (\kod{V-\allowbreak{}33}). Menşe maliyeti \kod{c\_\allowbreak{}RoO} ve uyum olasılığı \kod{p\_\allowbreak{}uyum} sınıf varsayım tablosundan geliyor (\kod{veri/\allowbreak{}f3\_\allowbreak{}sinif\_\allowbreak{}varsayim.\allowbreak{}csv}, \etiketH{}). Tarım satırları 1/98'de bugün de menşe kuralına tabi olduğundan onlara ek menşe maliyeti eklenmiyor.
\item
  \textbf{TÜFE etkisi birinci tur ve I-O ile tutarlı.} Doğrudan kanal: COICOP 2018 sepetinden sınıfa eşlenen ağırlık × tüketimdeki ithal pay. Dolaylı kanal: I-O 2023 Leontief fiyat modeli (\kod{kod/\allowbreak{}f2a\_\allowbreak{}io\_\allowbreak{}fiyat.\allowbreak{}py}). 230 sınıfın EDE toplamı 22,0 bp; I-O ürün toplamı 23,5 bp. S4'te AB'ye uygulanacak c2 tarifesinin imalat toplamı birinci turda 69 bp (Faz 2A'daki 73 bp ile tutarlı).
\end{enumerate}

\textbf{Skorlar.}

\begin{itemize}
\tightlist
\item
  KS'nin yedi bileşeni, ağırlıkları ana prompttaki gibi, C sınıfları içinde yüzdelik sıraya çevrildi. Ar-Ge bileşeni OECD sınıflamasından geliyor (Türkiye'ye özgü sınıf düzeyi Ar-Ge verisi yok: VB-31).
\item
  Ağırlık duyarlılığı: varsayılan sıranın Spearman korelasyonu eşit ağırlıkta 0,93, ticaret ağırlıklı setle 0,92, yapı ağırlıklı setle 0,75; Dirichlet çekilişlerinin P10'unda 0,71. Sıralama ticaret ağırlıklarına karşı sağlam, yapısal ağırlıklara karşı orta derecede sağlam. Kartlarda her sınıfın sıra aralığı (P10--P90) veriliyor.
\item
  SS yedi boyutta, elle ve gerekçeli puanlandı \etiketH{}. SS = 0,7 × en yüksek boyut + 0,3 × ağırlıklı ortalama.
\item
  AME = AB ihracat payı × (beklenen \kod{τ\_\allowbreak{}AB} + \kod{c\_\allowbreak{}RoO} + sınır sürtünmesi) × değer zinciri çarpanı. Beklenen τ, 1/98 kaybını 0,25 olasılıkla içeriyor.
\item
  CKE: STA dışı ithalata 1 puanlık tarifenin ya da iç satışa \%1 üretim desteğinin 5 yıllık net cari denge getirisi. Getiri, sırasıyla TÜFE baz puanına ve mali maliyete bölünüyor; kapasite yatırımının ithalatı ve ithal girdi düşülüyor (\kod{V-\allowbreak{}36}).
\item
  Arketip ve duruş önce kurala göre atanıyor (\kod{V-\allowbreak{}34}), sonra kartta uzman muhakemesiyle düzeltiliyor. Düzeltmeler \kod{veri/\allowbreak{}nace\_\allowbreak{}matris.\allowbreak{}csv}'nin \kod{uzman\_\allowbreak{}duzeltmesi} sütununda.
\end{itemize}

\textbf{Düzeltme notu (Faz 0--2 sınıf tabloları).} Faz 3 hesapları iki veri hatası buldu.

\begin{itemize}
\tightlist
\item
  \textbf{D-1. KN8 → NACE ağırlıkları.} \kod{veri/\allowbreak{}cn2025\_\allowbreak{}nace\_\allowbreak{}rejim.\allowbreak{}csv} yazılırken birden çok sınıfa düşen 1.251 KN8 satırının ağırlığı tam sayıya yuvarlanmış ve sıfır görünüyordu. Bu satırlar TR→AB ihracatının \%12,8'ini oluşturuyor.

  \begin{itemize}
  \tightlist
  \item
    Faz 0--2'nin toplamları doğru: değerler eşit paylaştırılmış olarak zaten yazılıydı.
  \item
    Sınıf dağılımı ise bazı sınıflarda yanlış. Örnek: hibrit otomobil satırları (CPA 2.2 29.10.23) 29.10 ile 27.11 arasında eşit bölünmüştü. Faz 3 eşlemesiyle 29.32'nin AB ihracatı €2,1 mlr'dan €5,6 mlr'a, 29.10'unki €16,6 mlr'dan €18,7 mlr'a çıkıyor; 27.11 (€3,3 → €1,3 mlr) ve 28.25 (€2,9 → €1,3 mlr) düşüyor.
  \item
    Faz 3'te CPA 2.2 kodunun CPA 2.1'de aynen sürdüğü bağlantıya öncelik veriliyor (\kod{agirlik\_\allowbreak{}f3}, \kod{V-\allowbreak{}37}); eşleşmeyen ticaret \%6,0'dan \%0,3'e iniyor.
  \item
    Faz 0--2'deki sınıf düzeyi tablolar (\kod{nace\_\allowbreak{}kapsam.\allowbreak{}csv} ihracat sütunu, \kod{s4\_\allowbreak{}tarife\_\allowbreak{}duvari\_\allowbreak{}sinif.\allowbreak{}csv}, \kod{politika\_\allowbreak{}s4\_\allowbreak{}net\_\allowbreak{}istihdam\_\allowbreak{}sinif.\allowbreak{}csv}) bu sınıflar için Faz 3 kartlarıyla geçersiz kalıyor. Sınıf düzeyi Faz 2 sayılarının Faz 3 yöntemiyle yeniden hesaplanmış hâli matriste.
  \end{itemize}
\item
  \textbf{D-2. Kıbrıs akımları.} BACI'deki TR→Kıbrıs ihracatı (\$1,6 mlr, 2024; \%31'i akaryakıt) büyük ölçüde KKTC ticareti ve AB gümrük bölgesi dışında (Protokol 10).

  \begin{itemize}
  \tightlist
  \item
    Faz 3'te bu akımlar AB dışı sayılıyor ve AB pazarı verisinden düşülüyor. Etki: imalatın AB'ye ihracatının \%1,4'ü; bazı küçük gıda sınıflarında AB payı belirgin düşüyor.
  \item
    Faz 1--2'de AB tarafı maruziyet Comext'ten hesaplandığı için etkilenmiyor. Yalnız BACI'ye dayanan Türkiye ithalatı AB payları (EDE) çok az yüksek kalmış.
  \end{itemize}
\item
  \textbf{D-3. Faz 2'deki ``kaybı \%5'i aşan 21 sınıf''} σ = 5 ile, kalibrasyonsuz kısmi dengeden geliyordu. Faz 3'te genel denge kalibrasyonu ve iç ikameyle S4'te kaybı \%5'i aşan sınıf sayısı P50'de 9, P10'da 15 (\kod{veri/\allowbreak{}f3\_\allowbreak{}senaryo\_\allowbreak{}ozet.\allowbreak{}csv}).
\end{itemize}

\begin{landscape}

\textbf{Tablo 1.1 --- Bölümlerin genel görünümü (2024; senaryolar T+5)}

\nopagebreak[4]

\begingroup\gbtablo\footnotesize

\begin{longtable}[]{@{}
  >{\raggedright\arraybackslash}p{(\columnwidth - 24\tabcolsep) * \real{0.0961}}
  >{\raggedleft\arraybackslash}p{(\columnwidth - 24\tabcolsep) * \real{0.0453}}
  >{\raggedleft\arraybackslash}p{(\columnwidth - 24\tabcolsep) * \real{0.0564}}
  >{\raggedleft\arraybackslash}p{(\columnwidth - 24\tabcolsep) * \real{0.0674}}
  >{\raggedleft\arraybackslash}p{(\columnwidth - 24\tabcolsep) * \real{0.0916}}
  >{\raggedleft\arraybackslash}p{(\columnwidth - 24\tabcolsep) * \real{0.0567}}
  >{\raggedleft\arraybackslash}p{(\columnwidth - 24\tabcolsep) * \real{0.0888}}
  >{\raggedleft\arraybackslash}p{(\columnwidth - 24\tabcolsep) * \real{0.0531}}
  >{\raggedleft\arraybackslash}p{(\columnwidth - 24\tabcolsep) * \real{0.0732}}
  >{\raggedleft\arraybackslash}p{(\columnwidth - 24\tabcolsep) * \real{0.0732}}
  >{\raggedleft\arraybackslash}p{(\columnwidth - 24\tabcolsep) * \real{0.0964}}
  >{\raggedleft\arraybackslash}p{(\columnwidth - 24\tabcolsep) * \real{0.0964}}
  >{\raggedleft\arraybackslash}p{(\columnwidth - 24\tabcolsep) * \real{0.1054}}@{}}
\toprule\noalign{}
\begin{minipage}[b]{\linewidth}\raggedright
\textbf{Bölüm}
\end{minipage} & \begin{minipage}[b]{\linewidth}\raggedleft
\textbf{Sınıf}
\end{minipage} & \begin{minipage}[b]{\linewidth}\raggedleft
\textbf{Çalışan}
\end{minipage} & \begin{minipage}[b]{\linewidth}\raggedleft
\textbf{KD (\$ mn)}
\end{minipage} & \begin{minipage}[b]{\linewidth}\raggedleft
\textbf{İhracat (\$ mn)}
\end{minipage} & \begin{minipage}[b]{\linewidth}\raggedleft
\textbf{AB payı}
\end{minipage} & \begin{minipage}[b]{\linewidth}\raggedleft
\textbf{İthalat (\$ mn)}
\end{minipage} & \begin{minipage}[b]{\linewidth}\raggedleft
\textbf{Denge}
\end{minipage} & \begin{minipage}[b]{\linewidth}\raggedleft
\textbf{S2 cari P50}
\end{minipage} & \begin{minipage}[b]{\linewidth}\raggedleft
\textbf{S4 cari P50}
\end{minipage} & \begin{minipage}[b]{\linewidth}\raggedleft
\textbf{S4 istihdam P50}
\end{minipage} & \begin{minipage}[b]{\linewidth}\raggedleft
\textbf{S4 istihdam P10}
\end{minipage} & \begin{minipage}[b]{\linewidth}\raggedleft
\textbf{S4 TÜFE P50 (bp)}
\end{minipage} \\
\midrule\noalign{}
\endhead
\bottomrule\noalign{}
\endlastfoot
10 Gıda ürünleri & 25 & 610.282 & 21.311 & 22.655 & \%26 & 11.186 & +11.469 & −32 & +3 & −89 & −4.693 & 2,6 \\
11 İçecekler & 7 & 19.525 & 1.331 & 539 & \%30 & 1.242 & −704 & +3 & +16 & +54 & 0 & 0,9 \\
12 Tütün ürünleri & 1 & 5.769 & 656 & 876 & \%13 & 958 & −81 & −15 & +11 & +20 & −5 & 2,7 \\
\end{longtable}

\endgroup

\end{landscape}

\emph{Kaynak:} \etiketV{} TÜİK YSHİ, BACI; senaryolar \etiketT{} (f3e). Cari denge ve TÜFE etkileri sınıf toplamıdır; makro kanallar (kur, talep) Faz 6'da.

\hypertarget{nace_kartlar_10-12-1-genel-guxf6ruxfcnuxfcm}{%
\section{1. Genel görünüm}\label{nace_kartlar_10-12-1-genel-guxf6ruxfcnuxfcm}}

\textbf{10 Gıda ürünleri (25 sınıf).}

\begin{itemize}
\tightlist
\item
  İmalat istihdamının \%12'si (610 bin), AB'ye imalat ihracatının yalnızca \%5,6'sı. AB'ye ihracatın büyük kısmı tarım rejiminde (1/98): kuru meyve ve fındık, su ürünleri, meyve suyu, bitkisel yağ.
\item
  Türkiye birçok üründe dünya ihracatında ilk sırada (BACI 2024) \etiketV{}:

  \begin{itemize}
  \tightlist
  \item
    buğday unu (\%19,6), fındık içi (\%62), kuru kayısı (\%75), kuru incir (\%50), kuru üzüm (\%30);
  \item
    levrek (\%32), rafine ayçiçek yağı (\%19), bulgur (\%75), fruktoz (\%35);
  \item
    makarna ve çipurada ikinci.
  \end{itemize}
\item
  Bu ihracatın çoğu AB dışına gidiyor: un ve makarnada Irak, Somali, Gana; kanatlıda Irak.
\item
  Bölümün GB'ye bağlı tek parçası İTÜ: bisküvi, çikolata ve şekerleme, glikoz, gıda müstahzarları. AB'ye €1,1 mlr civarında ihracat ve S4'te \%7--16 sanayi bileşeni.
\end{itemize}

\textbf{11 İçecekler (7 sınıf).}

\begin{itemize}
\tightlist
\item
  AB'nin alkollü içki ve birada MFN'si sıfır ya da çok düşük, bu yüzden çıkışın AB tarafı etkisi küçük.
\item
  Alkolsüz içeceklerde (11.07) AB ihracatı \$111 mn ve S4 tarifesi \%6,6. İthalatta enerji içecekleri öne çıkıyor (Avusturya \%42).
\item
  11.05'in ithalatı (\$348 mn) aslında bira değil, ABD'den damıtma artığı yem (KN 2303 30): sınıflandırma CPA 11.05.20'ye bağlı.
\item
  ÖTV ile gümrük etkileşimi: alkollü içki ve kolada ÖTV maktu ya da yüksek oranlı. S4'te AB'ye MFN, ithal ürünün ÖTV öncesi fiyatına eklenir, ama TÜFE etkisi küçük (11 için 0,9 bp).
\end{itemize}

\textbf{12 Tütün ürünleri (1 sınıf).}

\begin{itemize}
\tightlist
\item
  Çok uluslu üretim üssü (PMI ve JTI İzmir; PMI Torbalı 2024'te üretiminin \%11'inden fazlasını ihraç etti) \etiketV{}.
\item
  AB'ye ihracat \$116 mn ve sigarada AB MFN'si \%57,6: S4'te AB pazarı fiilen kapanır.
\item
  S2'de menşe kuralı (\%70 menşeli yaprak) Amerikan harmanını dışarıda bırakıyor.
\item
  İthalatın \%48'i AB'den. Isıtılmış tütün ve sigara Almanya, Hollanda ve Belçika'dan geliyor. S4'te c2 tarifesi TÜFE'ye 2,7 bp ekliyor; bölümün S4 TÜFE etkisinin \%44'ü bu.
\end{itemize}

\hypertarget{nace_kartlar_10-12-2-ana-matris-satux131rlarux131}{%
\section{2. Ana matris satırları}\label{nace_kartlar_10-12-2-ana-matris-satux131rlarux131}}

Matrisin bütün sütunları ve 230 sınıfın tamamı \kod{veri/\allowbreak{}nace\_\allowbreak{}matris.\allowbreak{}csv}'de. Aşağıdaki iki tablo ana promptun matris sütunlarını 10--12 için gösteriyor.

\begin{landscape}

\textbf{Tablo 2.1a --- Ana matris: kimlik, ticaret, maruziyet}

\nopagebreak[4]

\begingroup\gbtablo\footnotesize

\begin{longtable}[]{@{}
  >{\raggedright\arraybackslash}p{(\columnwidth - 26\tabcolsep) * \real{0.0402}}
  >{\raggedright\arraybackslash}p{(\columnwidth - 26\tabcolsep) * \real{0.2105}}
  >{\raggedright\arraybackslash}p{(\columnwidth - 26\tabcolsep) * \real{0.0921}}
  >{\raggedleft\arraybackslash}p{(\columnwidth - 26\tabcolsep) * \real{0.0402}}
  >{\raggedleft\arraybackslash}p{(\columnwidth - 26\tabcolsep) * \real{0.0684}}
  >{\raggedleft\arraybackslash}p{(\columnwidth - 26\tabcolsep) * \real{0.0600}}
  >{\raggedleft\arraybackslash}p{(\columnwidth - 26\tabcolsep) * \real{0.0573}}
  >{\raggedleft\arraybackslash}p{(\columnwidth - 26\tabcolsep) * \real{0.0572}}
  >{\raggedleft\arraybackslash}p{(\columnwidth - 26\tabcolsep) * \real{0.0750}}
  >{\raggedleft\arraybackslash}p{(\columnwidth - 26\tabcolsep) * \real{0.0694}}
  >{\raggedright\arraybackslash}p{(\columnwidth - 26\tabcolsep) * \real{0.0814}}
  >{\raggedleft\arraybackslash}p{(\columnwidth - 26\tabcolsep) * \real{0.0684}}
  >{\raggedleft\arraybackslash}p{(\columnwidth - 26\tabcolsep) * \real{0.0383}}
  >{\raggedleft\arraybackslash}p{(\columnwidth - 26\tabcolsep) * \real{0.0415}}@{}}
\toprule\noalign{}
\begin{minipage}[b]{\linewidth}\raggedright
\textbf{Kod}
\end{minipage} & \begin{minipage}[b]{\linewidth}\raggedright
\textbf{Ad}
\end{minipage} & \begin{minipage}[b]{\linewidth}\raggedright
\textbf{GB}
\end{minipage} & \begin{minipage}[b]{\linewidth}\raggedleft
\textbf{KD}
\end{minipage} & \begin{minipage}[b]{\linewidth}\raggedleft
\textbf{İstihdam}
\end{minipage} & \begin{minipage}[b]{\linewidth}\raggedleft
\textbf{İhracat}
\end{minipage} & \begin{minipage}[b]{\linewidth}\raggedleft
\textbf{AB payı}
\end{minipage} & \begin{minipage}[b]{\linewidth}\raggedleft
\textbf{İthalat}
\end{minipage} & \begin{minipage}[b]{\linewidth}\raggedleft
\textbf{AB ith. payı}
\end{minipage} & \begin{minipage}[b]{\linewidth}\raggedleft
\textbf{τ\_AB (S4)}
\end{minipage} & \begin{minipage}[b]{\linewidth}\raggedright
\textbf{Menşe S2/S3}
\end{minipage} & \begin{minipage}[b]{\linewidth}\raggedleft
\textbf{İthal girdi}
\end{minipage} & \begin{minipage}[b]{\linewidth}\raggedleft
\textbf{RCA}
\end{minipage} & \begin{minipage}[b]{\linewidth}\raggedleft
\textbf{PCI}
\end{minipage} \\
\midrule\noalign{}
\endhead
\bottomrule\noalign{}
\endlastfoot
10.11 & Etin işlenmesi ve saklanması (kümes ha & KARMA & 804 & 10.835 & 146 & \%46 & 1.011 & \%64 & \%0,0 & D/D & \%20 & 0,09 & −0,89 \\
10.12 & Kümes hayvanları etlerinin işlenmesi v & TARIM-1/98 & 2.136 & 35.710 & 643 & \%0 & 86 & \%9 & \%0,0 & D/D & \%20 & 1,45 & −0,89 \\
10.13 & Et ve kümes hayvanları etlerinden üret & TARIM-1/98 & 379 & 11.726 & 214 & \%2 & 64 & \%56 & \%0,0 & D/D & \%20 & 0,49 & 0,17 \\
10.20 & Balık, kabuklu deniz hayvanları ve yum & TARIM-1/98 & 722 & 14.445 & 1.237 & \%40 & 514 & \%7 & \%0,0 & D/D & \%20 & 0,75 & −0,97 \\
10.31 & Patatesin işlenmesi ve saklanması & KARMA & 107 & 2.152 & 110 & \%3 & 22 & \%85 & \%0,0 & D/D & \%20 & 0,49 & 0,42 \\
10.32 & Sebze ve meyve suyu imalatı & TARIM-1/98 & 415 & 7.460 & 770 & \%27 & 144 & \%6 & \%0,0 & D/D & \%20 & 3,21 & −0,84 \\
10.39 & Meyve ve sebzelerin diğer şekillerde i & TARIM-1/98 & 2.174 & 66.354 & 6.472 & \%56 & 1.177 & \%8 & \%0,0 & D/D & \%20 & 5,05 & −1,03 \\
10.41 & Sıvı ve katı yağ imalatı & TARIM-1/98 & 910 & 17.859 & 3.183 & \%9 & 3.711 & \%2 & \%0,0 & D/O & \%20 & 1,53 & −1,10 \\
10.42 & Margarin ve benzeri yenilebilir katı y & TARIM-1/98 & 71 & 773 & 165 & \%7 & 79 & \%39 & \%0,0 & O/O & \%20 & 1,62 & −0,58 \\
10.51 & Süthane işletmelerinin faaliyetleri ve & TARIM-1/98 & 1.987 & 48.749 & 440 & \%6 & 112 & \%46 & \%0,2 & D/D & \%20 & 0,35 & 0,15 \\
10.52 & Dondurma imalatı & İTÜ & 489 & 7.589 & 72 & \%1 & 5 & \%54 & \%8,4 & D/O & \%20 & 0,99 & −0,13 \\
10.61 & Öğütülmüş hububat ve sebze ürünleri im & KARMA & 1.245 & 26.004 & 1.870 & \%7 & 739 & \%17 & \%4,9 & O/O & \%20 & 2,40 & −1,18 \\
10.62 & Nişasta ve nişastalı ürünlerin imalatı & KARMA & 357 & 3.117 & 528 & \%18 & 206 & \%68 & \%14,2 & D/O & \%20 & 1,84 & 0,05 \\
10.71 & Ekmek, taze pastane ürünleri ve taze k & İTÜ & 2.097 & 189.869 & 264 & \%19 & 29 & \%61 & \%9,0 & O/O & \%20 & 1,45 & −0,18 \\
10.72 & Peksimet ve bisküvi imalatı; dayanıklı & İTÜ & 1.548 & 42.160 & 1.174 & \%14 & 89 & \%75 & \%9,0 & O/O & \%20 & 3,10 & −0,49 \\
10.73 & Makarna, şehriye, kuskus ve benzeri un & İTÜ & 409 & 7.040 & 852 & \%2 & 17 & \%86 & \%7,7 & D/D & \%20 & 9,09 & −1,59 \\
10.81 & Şeker imalatı & TARIM-1/98 & 803 & 15.900 & 143 & \%10 & 458 & \%10 & \%0,0 & D/D & \%20 & 0,27 & −0,94 \\
10.82 & Kakao, çikolata ve şekerleme imalatı & İTÜ & 1.645 & 36.835 & 2.032 & \%14 & 963 & \%73 & \%8,4 & O/O & \%20 & 1,97 & −0,27 \\
10.83 & Kahve ve çayın işlenmesi & KARMA & 1.042 & 19.999 & 170 & \%23 & 364 & \%60 & \%3,2 & D/D & \%20 & 0,44 & −0,76 \\
10.84 & Baharat, sos, sirke ve diğer çeşni mad & KARMA & 296 & 7.151 & 242 & \%34 & 90 & \%40 & \%4,9 & D/D & \%20 & 0,78 & −0,92 \\
10.85 & Hazır yemeklerin imalatı & İTÜ & 221 & 10.010 & 100 & \%10 & 2 & \%65 & \%7,0 & O/O & \%20 & 1,00 & −1,08 \\
10.86 & Homojenize gıda müstahzarları ve diyet & KARMA & 61 & 1.321 & 214 & \%20 & 173 & \%91 & \%7,7 & O/O & \%20 & 0,80 & −0,17 \\
10.89 & Başka yerde sınıflandırılmamış diğer g & İTÜ & 467 & 10.858 & 1.114 & \%24 & 639 & \%62 & \%11,8 & D/O & \%20 & 0,99 & −0,13 \\
10.91 & Çiftlik hayvanları için hazır yem imal & TARIM-1/98 & 846 & 14.598 & 352 & \%1 & 279 & \%50 & \%0,0 & D/D & \%20 & 1,40 & −0,76 \\
10.92 & Ev hayvanları için hazır gıda imalatı & TARIM-1/98 & 81 & 1.768 & 147 & \%14 & 214 & \%87 & \%0,0 & D/D & \%20 & 0,49 & 0,03 \\
11.01 & Alkollü içeceklerin damıtılması, arıtı & İTÜ & 393 & 2.398 & 68 & \%48 & 425 & \%25 & \%0,0 & D/D & \%20 & 0,14 & −0,36 \\
11.02 & Üzümden şarap imalatı & TARIM-1/98 & 109 & 1.968 & 24 & \%61 & 53 & \%86 & \%0,0 & D/D & \%20 & 0,05 & −0,56 \\
11.03 & Elma şarabı ve diğer meyve şaraplarını & TARIM-1/98 & --- & --- & 0 & \%10 & 1 & \%47 & \%0,0 & D/D & \%20 & 0,01 & 0,07 \\
11.04 & Diğer damıtılmamış fermente içecekleri & İTÜ & --- & --- & 0 & \%15 & 2 & \%100 & \%0,0 & D/D & \%20 & 0,01 & −0,42 \\
11.05 & Bira imalatı & İTÜ & 65 & 394 & 61 & \%8 & 348 & \%10 & \%0,0 & D/D & \%20 & 0,23 & −0,64 \\
11.06 & Malt imalatı & TARIM-1/98 & --- & --- & 0 & \%2 & 83 & \%86 & \%0,0 & D/D & \%20 & 0,01 & 0,16 \\
11.07 & Alkolsüz içeceklerin imalatı; maden su & KARMA & 764 & 14.765 & 385 & \%29 & 330 & \%78 & \%6,6 & D/O & \%20 & 0,93 & −0,64 \\
12.00 & Tütün ürünleri imalatı & KARMA & 656 & 5.769 & 876 & \%13 & 958 & \%48 & \%51,3 & Y/Y & \%20 & 1,52 & −1,08 \\
\end{longtable}

\endgroup

\end{landscape}

\emph{Kaynak:} \kod{veri/\allowbreak{}nace\_\allowbreak{}matris.\allowbreak{}csv}. KD, ihracat, ithalat \$ mn (2024). τ\_AB: S4'te eklenen AB tarifesi (rota A).

\begin{landscape}

\textbf{Tablo 2.1b --- Ana matris: skorlar, strateji, senaryo}

\nopagebreak[4]

\begingroup\gbtablo\footnotesize

\begin{longtable}[]{@{}
  >{\raggedright\arraybackslash}p{(\columnwidth - 28\tabcolsep) * \real{0.0408}}
  >{\raggedleft\arraybackslash}p{(\columnwidth - 28\tabcolsep) * \real{0.0285}}
  >{\raggedleft\arraybackslash}p{(\columnwidth - 28\tabcolsep) * \real{0.0285}}
  >{\raggedleft\arraybackslash}p{(\columnwidth - 28\tabcolsep) * \real{0.0420}}
  >{\raggedleft\arraybackslash}p{(\columnwidth - 28\tabcolsep) * \real{0.0388}}
  >{\raggedleft\arraybackslash}p{(\columnwidth - 28\tabcolsep) * \real{0.0388}}
  >{\raggedright\arraybackslash}p{(\columnwidth - 28\tabcolsep) * \real{0.0580}}
  >{\raggedright\arraybackslash}p{(\columnwidth - 28\tabcolsep) * \real{0.0556}}
  >{\raggedright\arraybackslash}p{(\columnwidth - 28\tabcolsep) * \real{0.1608}}
  >{\raggedright\arraybackslash}p{(\columnwidth - 28\tabcolsep) * \real{0.2052}}
  >{\raggedright\arraybackslash}p{(\columnwidth - 28\tabcolsep) * \real{0.0586}}
  >{\raggedleft\arraybackslash}p{(\columnwidth - 28\tabcolsep) * \real{0.0582}}
  >{\raggedleft\arraybackslash}p{(\columnwidth - 28\tabcolsep) * \real{0.0608}}
  >{\raggedright\arraybackslash}p{(\columnwidth - 28\tabcolsep) * \real{0.0751}}
  >{\raggedright\arraybackslash}p{(\columnwidth - 28\tabcolsep) * \real{0.0503}}@{}}
\toprule\noalign{}
\begin{minipage}[b]{\linewidth}\raggedright
\textbf{Kod}
\end{minipage} & \begin{minipage}[b]{\linewidth}\raggedleft
\textbf{KS}
\end{minipage} & \begin{minipage}[b]{\linewidth}\raggedleft
\textbf{SS}
\end{minipage} & \begin{minipage}[b]{\linewidth}\raggedleft
\textbf{AME}
\end{minipage} & \begin{minipage}[b]{\linewidth}\raggedleft
\textbf{EDE}
\end{minipage} & \begin{minipage}[b]{\linewidth}\raggedleft
\textbf{CKE}
\end{minipage} & \begin{minipage}[b]{\linewidth}\raggedright
\textbf{Arketip}
\end{minipage} & \begin{minipage}[b]{\linewidth}\raggedright
\textbf{Duruş}
\end{minipage} & \begin{minipage}[b]{\linewidth}\raggedright
\textbf{Tarife yolu}
\end{minipage} & \begin{minipage}[b]{\linewidth}\raggedright
\textbf{Ana araçlar}
\end{minipage} & \begin{minipage}[b]{\linewidth}\raggedright
\textbf{AB müz.}
\end{minipage} & \begin{minipage}[b]{\linewidth}\raggedleft
\textbf{S2 cari P50}
\end{minipage} & \begin{minipage}[b]{\linewidth}\raggedleft
\textbf{S2 TÜFE P50}
\end{minipage} & \begin{minipage}[b]{\linewidth}\raggedright
\textbf{İstihdam riski}
\end{minipage} & \begin{minipage}[b]{\linewidth}\raggedright
\textbf{Güven}
\end{minipage} \\
\midrule\noalign{}
\endhead
\bottomrule\noalign{}
\endlastfoot
10.11 & 41 & 53 & 17 & 0,46 & 34 & --- & Koru & Değişiklik yok (tarım ürünü; 1/98) & ESK ithalat kaynağı çeşitlendirme; SPS; EUDR (sığır izlenebilirliği) & nötr & +0,2 & 0,01 & Düşük & Yüksek \\
10.12 & 46 & 60 & 11 & 0,05 & 25 & B/D & Geliştir & Değişiklik yok (tarım); yem girdilerinde mevcut rejim & Pazar çeşitlendirme; AB SPS listelenme talebi; yem arz güvenliği & hücum & +0,1 & 0,00 & Düşük & Yüksek \\
10.13 & 54 & 30 & 18 & 0,05 & 16 & --- & Koru & Değişiklik yok (tarım) & İzleme & nötr & 0,0 & 0,00 & Düşük & Yüksek \\
10.20 & 47 & 38 & 62 & 0,05 & 42 & A & Koru & 1/98 balıkçılık tavizleri korunur; S2'de kota yerine \%0 & 1/98 sürekliliği (E-10-01); ticaret savunması hazırlığı (E-10-04); yem ve yavru arzı & hücum & 0,0 & 0,00 & Orta & Yüksek \\
10.31 & 50 & 22 & 13 & 0,01 & 33 & --- & Koru & Değişiklik yok & İzleme & nötr & 0,0 & 0,00 & Düşük & Yüksek \\
10.32 & 52 & 15 & 62 & 0,02 & 32 & B/A & Geliştir & Değişiklik yok (tarım; 1/98) & 1/98 sürekliliği (E-10-01); AB dışı pazar; giriş fiyatı istisnası (E-10-06) & hücum & 0,0 & 0,00 & Orta & Yüksek \\
10.39 & 50 & 37 & 54 & 0,12 & 51 & A & Koru & 1/98 korunur; S2'de tavizler ve giriş fiyatı istisnası; TR tarifesi değişmez & 1/98 sürekliliği (E-10-01); arz güvenliği kaldıracı (fındık); SPS kontrol sıklığının azaltılması & hücum & −0,4 & 0,00 & Orta & Yüksek \\
10.41 & 46 & 63 & 18 & 0,39 & 63 & D/B & Geliştir & Tarım tarifeleri değişmez; DİR sürer & Yağlı tohum üretimi; stratejik ham yağ stoku; kaynak çeşitlendirme & nötr & −0,1 & 0,00 & Düşük & Orta \\
10.42 & 42 & 30 & 17 & 0,01 & 26 & B & Geliştir & Değişiklik yok & Pazar çeşitlendirme & nötr & 0,0 & 0,00 & Düşük & Orta \\
10.51 & 46 & 57 & 19 & 0,14 & 12 & --- & Koru & Değişiklik yok (tarım) & İzleme; S2'de süt ürünleri SPS listelenmesi & nötr & 0,0 & 0,00 & Düşük & Orta \\
10.52 & 52 & 4 & 15 & 0,00 & 15 & --- & Koru & Değişiklik yok & İzleme & nötr & 0,0 & 0,00 & Düşük & Orta \\
10.61 & 48 & 71 & 27 & 0,11 & 37 & B/D & Geliştir & Tarım tarifeleri değişmez; İTÜ satırlarında (bulgur, karışımlar) AB'de S4'te \%4,9 & Buğday arz güvenliği (TMO stoku, kaynak çeşitlendirme); DİR eşdeğer eşya (E-10-02); Afrika ve Körfez STA'ları & nötr & −6,5 & 0,00 & Düşük & Yüksek \\
10.62 & 53 & 30 & 84 & 0,02 & 50 & B & Geliştir & Değişiklik yok; S2'de AB'ye \%0 (İTÜ sanayi bileşeni) & S2'de İTÜ sanayi bileşeninin sıfırı; nişasta bazlı şeker kotası; mısır arzı & savunma & −4,2 & 0,01 & Düşük & Yüksek \\
10.71 & 56 & 67 & 72 & 0,05 & 13 & ---/D & Koru & Değişiklik yok & Buğday ve un arz güvenliği; menşe uyumlu DİR (E-10-02) & savunma & −5,2 & 0,04 & Düşük & Orta \\
10.72 & 55 & 22 & 63 & 0,03 & 25 & B & Geliştir & Değişiklik yok; S2'de AB'ye \%0 & Menşe uyumlu DİR (E-10-02); Fasıl 19 PSR toleransı; S4 tamponu (E-10-05) & savunma & −16,3 & 0,03 & Düşük & Yüksek \\
10.73 & 42 & 52 & 16 & 0,01 & 20 & B/D & Geliştir & Değişiklik yok & Afrika, Latin Amerika ve Japonya STA'ları; makarnalık buğday arzı & nötr & −0,4 & 0,01 & Düşük & Orta \\
10.81 & 34 & 53 & 10 & 0,09 & 20 & --- & Koru & Değişiklik yok (tarım; kota rejimi) & İzleme; S2'de şeker hassas listesi & savunma & 0,0 & 0,00 & Düşük & Yüksek \\
10.82 & 56 & 7 & 57 & 0,11 & 57 & B/G & Geliştir & Kakao ara ürünlerinde (1803--1805) MFN \%0, tüm menşeler; nihai üründe MFN korunur & Girdide sıfır (kakao); EUDR uyumu (E-10-03); menşe uyumlu DİR (şeker) & savunma & +2,7 & 0,11 & Düşük & Yüksek \\
10.83 & 44 & 22 & 39 & 0,18 & 25 & --- & Koru & Değişiklik yok & EUDR (kahve; E-10-03); izleme & nötr & +2,8 & 0,08 & Düşük & Yüksek \\
10.84 & 42 & 22 & 73 & 0,02 & 30 & A & Geliştir & S2'de AB'ye \%0 (İTÜ); baharatta 1/98 & S2'de İTÜ sıfırı; SPS kontrol sıklığı (2019/1793); izleme & savunma & −2,1 & 0,01 & Düşük & Yüksek \\
10.85 & 38 & 24 & 41 & 0,01 & 8 & --- & Koru & Değişiklik yok & İzleme; afet ve askerî tayın tedariki & nötr & −0,3 & 0,01 & Düşük & Yüksek \\
10.86 & 50 & 55 & 69 & 0,02 & 72 & D/C & Geliştir & S4'te bebek mamasında (1901 10) MFN \%0, tüm menşeler; diğerlerinde MFN & S4'te 1901 10 MFN \%0; yatırım teşviki (bebek maması tesisi); stok izleme & nötr & +1,0 & 0,03 & Düşük & Yüksek \\
10.89 & 53 & 38 & 86 & 0,05 & 36 & A/C & Geliştir & S2'de AB'ye \%0; S4'te 2106 90'da AB'ye MFN (c2), TÜFE izlenir & S2'de İTÜ sıfırı; maya ihracatında pazar çeşitlendirme; S4 tamponu (E-10-05) & savunma & −3,6 & 0,08 & Düşük & Yüksek \\
10.91 & 49 & 63 & 12 & 0,02 & 30 & D/B & Geliştir & Tarım tarifeleri değişmez & Soya küspesi ve mısır kaynak çeşitlendirme; EUDR (soya); stok & nötr & 0,0 & 0,00 & Düşük & Yüksek \\
10.92 & 48 & 4 & 15 & 0,04 & 41 & C & Geliştir & Değişiklik yok (tarım) & Yatırım teşviki (yerlileştirme); izleme & nötr & 0,0 & 0,00 & Düşük & Yüksek \\
11.01 & 41 & 0 & 45 & 0,08 & 23 & --- & Koru & Değişiklik yok (AB MFN ≈ \%0) & Coğrafi işaret (Türk Rakısı) korunması; izleme & hücum & +0,1 & 0,02 & Düşük & Düşük \\
11.02 & 50 & 0 & 13 & 0,01 & 33 & --- & Koru & Değişiklik yok & İzleme & nötr & 0,0 & 0,00 & Düşük & Yüksek \\
11.03 & 42 & 0 & 11 & 0,00 & 3 & --- & Koru & Değişiklik yok & İzleme & nötr & 0,0 & 0,00 & Düşük & Düşük \\
11.04 & 43 & 0 & 20 & 0,00 & --- & --- & Koru & Değişiklik yok & İzleme & nötr & 0,0 & 0,00 & Düşük & Düşük \\
11.05 & 42 & 0 & 16 & 0,08 & 33 & --- & Koru & Değişiklik yok (AB MFN \%0) & İzleme & nötr & 0,0 & 0,01 & Düşük & Orta \\
11.06 & 44 & 15 & 12 & 0,01 & 3 & --- & Koru & Değişiklik yok (tarım) & İzleme & nötr & 0,0 & 0,00 & Düşük & Düşük \\
11.07 & 41 & 30 & 74 & 0,13 & 29 & --- & Koru & S2'de AB'ye \%0; S4'te AB'ye c2 (\%8,6) & İzleme; S4'te enerji içeceklerinde TÜFE izleme & savunma & +2,9 & 0,14 & Düşük & Düşük \\
12.00 & 40 & 0 & 98 & 0,31 & 26 & B & Koru & Tütün ürünlerinde MFN korunur; S4'te ısıtılmış tütünde MFN \%0 & AB dışı pazar; S2'de PSR talebi (\%50 menşeli yaprak); S4'te ısıtılmış tütünde MFN \%0 & savunma & −14,7 & 0,29 & Düşük & Yüksek \\
\end{longtable}

\endgroup

\end{landscape}

\emph{Kaynak:} Skorlar 0--100 (C sınıfları içinde yüzdelik); EDE bp/puan; S2 cari \$ mn; TÜFE bp. Arketip ve duruş kural tabanlı öneri, uzman düzeltmesiyle (kartlar).

\hypertarget{nace_kartlar_10-12-3-sux131nux131f-kartlarux131}{%
\section{3. Sınıf kartları}\label{nace_kartlar_10-12-3-sux131nux131f-kartlarux131}}

Her kartta sayısal alanlar veriden, nitel alanlar \kod{kartlar/\allowbreak{}10-\allowbreak{}12\_\allowbreak{}siniflar.\allowbreak{}yaml}'dan geliyor. Bölümler arası ortak eylemler (E-10-01 \ldots{} E-10-06) §4'te; kartlar bunlara atıf yapıyor. Senaryo tablolarındaki değerler T+5'teki yıllık etki (P50, parantez içinde P10 / P90).

\hypertarget{nace_kartlar_10-12-31-1011-etin-iux15flenmesi-ve-saklanmasux131-kuxfcmes-hayvanlarux131-eti-hariuxe7--gb-statuxfcsuxfc-karma--birincilikincil-arketip---duruux15f-koru}{%
\subsection{3.1 10.11 Etin işlenmesi ve saklanması (kümes hayvanları eti hariç) --- GB statüsü: KARMA \textbar{} Birincil/ikincil arketip: --- \textbar{} Duruş: Koru}\label{nace_kartlar_10-12-31-1011-etin-iux15flenmesi-ve-saklanmasux131-kuxfcmes-hayvanlarux131-eti-hariuxe7--gb-statuxfcsuxfc-karma--birincilikincil-arketip---duruux15f-koru}}

\emph{İç pazar ve tarım politikası sınıfı; GB çıkışı doğrudan etkilemiyor. AB'den et ithalatı (\%64) 1/98'e bağlı.}

\textbf{1. Kimlik.} Rev. 2.1: 10.11 (1:1); CPA 10.11. Başlıca HS4 (TR ihracatı 2024): 0504 \%48, 0202 \%13, 0511 \%12, 4102 \%9. BEC payları (ticaret ağırlıklı): ara \%28 / sermaye \%0 / tüketim \%72.

\textbf{2. Ölçek}

\begingroup\gbtablo\small

\begin{longtable}[]{@{}
  >{\raggedright\arraybackslash}p{(\columnwidth - 4\tabcolsep) * \real{0.2849}}
  >{\raggedright\arraybackslash}p{(\columnwidth - 4\tabcolsep) * \real{0.3960}}
  >{\raggedright\arraybackslash}p{(\columnwidth - 4\tabcolsep) * \real{0.3191}}@{}}
\toprule\noalign{}
\begin{minipage}[b]{\linewidth}\raggedright
\textbf{Gösterge}
\end{minipage} & \begin{minipage}[b]{\linewidth}\raggedright
\textbf{Değer}
\end{minipage} & \begin{minipage}[b]{\linewidth}\raggedright
\textbf{Etiket ve kaynak}
\end{minipage} \\
\midrule\noalign{}
\endhead
\bottomrule\noalign{}
\endlastfoot
Üretim değeri; katma değer & \$4.069 mn; \$804 mn (GSYH'nin \%0,06) & TÜİK YSHİ 2024 \etiketV{} \\
Çalışan; girişim (0 / 1--9 / 10--49 / 50--249 / 250+) & 10.835; 608 (57 / 410 / 144 / 37 / 5) & TÜİK YSHİ ve İş Kayıtları 2024 \etiketV{}; gizli büyüklük hücreleri \etiketT{} \\
250+ girişimlerin istihdam payı; AB'ye göre kişi başı KD & \%23; 1,41 & \etiketT{} (AB27 SBS 2023, nominal kur) \\
İlk 3 İBBS-2 bölgesi & Kocaeli, Sakarya, Düzce, Bolu, Yalova \%34; Balıkesir, Çanakkale \%6; Konya, Karaman \%6 & \etiketT{} iş kayıtları + SGK 2025, RAS \\
Yabancı sermaye; çapa yatırımcılar & Namet, Pınar Et, Bonfilet ve ESK (Et ve Süt Kurumu; kamu ithalatı). Yabancı çapa yok & \etiketH{} (firma bilgisi; pay verisi yok: VB-30) \\
\end{longtable}

\endgroup

\textbf{3. Dış ticaret}

\begingroup\gbtablo\small

\begin{longtable}[]{@{}
  >{\raggedright\arraybackslash}p{(\columnwidth - 2\tabcolsep) * \real{0.4711}}
  >{\raggedright\arraybackslash}p{(\columnwidth - 2\tabcolsep) * \real{0.5289}}@{}}
\toprule\noalign{}
\begin{minipage}[b]{\linewidth}\raggedright
\textbf{Gösterge}
\end{minipage} & \begin{minipage}[b]{\linewidth}\raggedright
\textbf{Değer (BACI 2024, \$; \etiketV{}, eşleme \etiketT{})}
\end{minipage} \\
\midrule\noalign{}
\endhead
\bottomrule\noalign{}
\endlastfoot
İhracat (2024); 2020--24 değişimi & \$146 mn; \%26 \\
AB payı; ilk 5 AB pazarı & \%46; ITA \%21; DEU \%19; FRA \%16; NLD \%10; DNK \%6 \\
İlk 5 AB dışı pazar & VNM \%36; HKG \%10; USA \%10; LBR \%4; CIV \%4 \\
İthalat; AB ve Çin payı; ilk kaynaklar & \$1.011 mn; AB \%64, Çin \%0; POL \%30; BRA \%24; FRA \%15; BGR \%8; ITA \%7 \\
Dış ticaret dengesi & \$−865 mn \\
İthal girdi içeriği (I-O, toplam) / KD÷üretim & \%20 / \%20 \\
RCA dünya (2020 → 2024) / AB pazarında & 0,09 → 0,09 / 0,28 \\
Birim değer ÷ AB pazarı ortalaması & 2,03 \\
AB pazar payı (AB dışı tedarik içinde, 2020 → 2024); TR'nin sırası & \%0,8 → \%1,1; 14. \\
\end{longtable}

\endgroup

\textbf{4. AB tarafı maruziyet}

\begin{itemize}
\tightlist
\item
  \kod{τ\_\allowbreak{}AB} (S4 artışı, rota A, TR→AB 2025 ağırlıklı): \%0,0; 1/98 kaybıyla \%0,3; tam MFN ortalaması \%0,3. Rejim: GB \%20, İTÜ \%0, tarım \%80 \etiketV{}/\etiketT{}.
\item
  Tepe satırlar (S4 artışı ≥ \%10; en büyük 3): yok.
\item
  AB pazarında AB dışı rakipler (tercih durumu): GBR \%27; NZL \%13; ARG \%11; BRA \%10; CHN \%9. Tercihli rakip payı \%51.
\item
  Menşe (S2 / S3): risk D/D; \kod{c\_\allowbreak{}RoO} \%1,5 / \%2,0; \kod{p\_\allowbreak{}uyum} 0,90 / 0,85 \etiketH{}. Kural: Fasıl 2: Fasıl 1--2 malzemeleri tamamen elde edilmiş; ithal canlı hayvandan et menşe kazanmaz (1/98 rejiminde bugün de geçerli).
\item
  EUDR kapsamındaki AB ihracatı payı: \%0; SKDM: \%0.
\item
  AB'ye ihracat (\$68 mn) bağırsak, tuzlanmış deri artıkları ve yağ gibi yan ürünler; et ihracatı SPS izni olmadığı için sınırlı \etiketH{}.
\item
  AB'den ithalat (\$651 mn; Polonya, Fransa) ESK'nın fiyat istikrarı alımları. 1/98'deki Türk tavizleri kalkarsa bu ithalat MFN'ye (et \%40--225) döner; bu da TÜFE riski (E-10-01).
\end{itemize}

\textbf{5. Türkiye tarafı maruziyet}

\begin{itemize}
\tightlist
\item
  Uygulanan tarife (2026, ithalat ağırlıklı): MFN \%96,7, İGV dahil \%96,7 \etiketV{}. S4'te AB menşeline (c2): \%0,0.
\item
  Kritik girdiler (HS4 0102, 0201, 0202), Türkiye'nin toplam ithalatı 2024: \$1.745 mn (AB \%58; BRA \%28, POL \%18, BGR \%13) \etiketV{}.
\item
  Canlı hayvan ve et ithalatı TARIM rejiminde. GB çıkışı Türkiye'nin et tarifesini değiştirmiyor.
\end{itemize}

\textbf{6. Kabiliyet}

\begin{itemize}
\tightlist
\item
  KS 40,7 (sıra 202/230; ağırlık duyarlılığında P10--P90 130--220). Bileşenler: RCA 12, PCI 13 (PCI −0,89), yurt içi KD 52, Ar-Ge 25, tedarikçi 84, verimlilik 98, firma 59 \etiketT{}.
\item
  Ürün uzayında atlanabilir ürünler (sınıf içi; yoğunluk üst çeyrek, PCI \textgreater{} sınıf): 020120 (PCI +0.14, RCA 0.00); 020110 (PCI -0.07, RCA 0.00); 410210 (PCI -0.88, RCA 0.35) \etiketT{}.
\item
  \textbf{Kırmızı çizgi yetkinlikleri:} Kesimhane ve soğuk zincir hijyen sertifikasyonu (AB'ye bağırsak ve yan ürün ihracatının şartı).
\item
  \textbf{Geliştirme hedefi:} T+5: AB'ye ısıl işlem görmüş et ürünleri için SPS listelenmesi (E-10-06); 2035: hayvan hastalıklarından ari bölge (Trakya) üzerinden AB ihracatı.
\end{itemize}

\textbf{7. Stratejik önem.} SS 52,6 (en yüksek boyut: gıda güvenliği 70; ağırlıklı ortalama 12,0) \etiketH{}. Kırmızı et arzı; canlı hayvan ithalatına bağımlı

\textbf{8. Makro duyarlılık}

\begingroup\gbtablo\small

\begin{longtable}[]{@{}lr@{}}
\toprule\noalign{}
\textbf{Gösterge} & \textbf{Değer (\etiketT{})} \\
\midrule\noalign{}
\endhead
\bottomrule\noalign{}
\endlastfoot
TÜFE ağırlığı (doğrudan, COICOP 2018 eşlemesi); tüketimde ithal pay & \%2,10; \%18 \\
ÜFE vekili (imalat üretimindeki pay) & \%0,61 \\
EDE: 1 puan tarife, tüm menşe / yalnız AB (bp; ρ = 1) & 0,464 / 0,299 \\
Geçişkenlik ρ (senaryolarda) & 0,6--1,0 (P50 0,8) \\
Piyasa yoğunlaşması (250+ istihdam payı) & \%23 \\
CKE: tarife (\$ mn / bp) ve mali (\$ mn / mlr TL); 0--100 puan & 415,0 ve 15,42; 34 \\
\end{longtable}

\endgroup

\textbf{9. Senaryo etkileri, T+5, P50 (P10 / P90)}

\begingroup\gbtablo\small

\begin{longtable}[]{@{}
  >{\raggedright\arraybackslash}p{(\columnwidth - 6\tabcolsep) * \real{0.3056}}
  >{\raggedleft\arraybackslash}p{(\columnwidth - 6\tabcolsep) * \real{0.3188}}
  >{\raggedleft\arraybackslash}p{(\columnwidth - 6\tabcolsep) * \real{0.1878}}
  >{\raggedleft\arraybackslash}p{(\columnwidth - 6\tabcolsep) * \real{0.1878}}@{}}
\toprule\noalign{}
\begin{minipage}[b]{\linewidth}\raggedright
\textbf{Değişken}
\end{minipage} & \begin{minipage}[b]{\linewidth}\raggedleft
\textbf{S2}
\end{minipage} & \begin{minipage}[b]{\linewidth}\raggedleft
\textbf{S3}
\end{minipage} & \begin{minipage}[b]{\linewidth}\raggedleft
\textbf{S4}
\end{minipage} \\
\midrule\noalign{}
\endhead
\bottomrule\noalign{}
\endlastfoot
AB'ye ihracat (\$ mn) & −1 (−2 / +1) & −1 (−2 / 0) & −1 (−1 / 0) \\
Toplam ihracat (\$ mn) & −1 (−2 / +1) & −1 (−2 / 0) & 0 (−1 / 0) \\
Nihai mal ithalatı (\$ mn) & −1 (−3 / +43) & −1 (−4 / +25) & −1 (−2 / 0) \\
Cari denge (\$ mn) & 0 (−34 / +1) & 0 (−20 / +2) & 0 (0 / +1) \\
Üretim (\$ mn) & 0 (−42 / +2) & +1 (−26 / +2) & 0 (0 / +1) \\
İstihdam (kişi) & 0 (−78 / +3) & +1 (−47 / +4) & 0 (0 / +1) \\
TÜFE (bp, birinci tur) & +0,01 (+0,01 / +0,02) & +0,03 (+0,02 / +0,04) & +0,02 (+0,01 / +0,02) \\
Yatırım ve DYY (bütün senaryolar) & Etkisi yok; ESK alımları kamu kararı & & \\
\end{longtable}

\endgroup

\emph{Kaynak:} \etiketT{} f3e: kısmi denge, toplamlar Faz 2B GE dağılımına kalibre; S4'te 1/98 kaybı olasılığı 0,25.

\textbf{10. Politika paketi}

\begingroup\gbtablo\footnotesize

\begin{longtable}[]{@{}
  >{\raggedright\arraybackslash}p{(\columnwidth - 16\tabcolsep) * \real{0.1117}}
  >{\raggedright\arraybackslash}p{(\columnwidth - 16\tabcolsep) * \real{0.1325}}
  >{\raggedright\arraybackslash}p{(\columnwidth - 16\tabcolsep) * \real{0.1262}}
  >{\raggedright\arraybackslash}p{(\columnwidth - 16\tabcolsep) * \real{0.0895}}
  >{\raggedright\arraybackslash}p{(\columnwidth - 16\tabcolsep) * \real{0.0831}}
  >{\raggedright\arraybackslash}p{(\columnwidth - 16\tabcolsep) * \real{0.0668}}
  >{\raggedright\arraybackslash}p{(\columnwidth - 16\tabcolsep) * \real{0.1656}}
  >{\raggedright\arraybackslash}p{(\columnwidth - 16\tabcolsep) * \real{0.1258}}
  >{\raggedright\arraybackslash}p{(\columnwidth - 16\tabcolsep) * \real{0.0989}}@{}}
\toprule\noalign{}
\begin{minipage}[b]{\linewidth}\raggedright
\textbf{Eylem ID}
\end{minipage} & \begin{minipage}[b]{\linewidth}\raggedright
\textbf{Araç}
\end{minipage} & \begin{minipage}[b]{\linewidth}\raggedright
\textbf{Parametre}
\end{minipage} & \begin{minipage}[b]{\linewidth}\raggedright
\textbf{Sorumlu}
\end{minipage} & \begin{minipage}[b]{\linewidth}\raggedright
\textbf{Hukuki araç}
\end{minipage} & \begin{minipage}[b]{\linewidth}\raggedright
\textbf{Bütçe}
\end{minipage} & \begin{minipage}[b]{\linewidth}\raggedright
\textbf{Başlangıç--bitiş; sunset}
\end{minipage} & \begin{minipage}[b]{\linewidth}\raggedright
\textbf{Performans koşulu}
\end{minipage} & \begin{minipage}[b]{\linewidth}\raggedright
\textbf{KPI}
\end{minipage} \\
\midrule\noalign{}
\endhead
\bottomrule\noalign{}
\endlastfoot
E-10.11-01 & ESK ithalat kaynağı çeşitlendirme & AB payı \%64 → ≤ \%40; Brezilya, Uruguay ve Avustralya ile ihale çerçevesi & TOB, ESK & ESK yönetim kurulu kararı & İdari & T−12 → T+12; sunset: T+24 & Fiyat ≤ AB teklifi + \%5 & Kırmızı et TÜFE alt endeksi yıllık artışı ≤ genel TÜFE + 5 puan \\
\end{longtable}

\endgroup

\textbf{12. Erken uyarı ve tetikler}

\begingroup\gbtablo\small

\begin{longtable}[]{@{}
  >{\raggedright\arraybackslash}p{(\columnwidth - 4\tabcolsep) * \real{0.3238}}
  >{\raggedright\arraybackslash}p{(\columnwidth - 4\tabcolsep) * \real{0.3792}}
  >{\raggedright\arraybackslash}p{(\columnwidth - 4\tabcolsep) * \real{0.2970}}@{}}
\toprule\noalign{}
\begin{minipage}[b]{\linewidth}\raggedright
\textbf{Gösterge}
\end{minipage} & \begin{minipage}[b]{\linewidth}\raggedright
\textbf{Eşik}
\end{minipage} & \begin{minipage}[b]{\linewidth}\raggedright
\textbf{Tetiklenecek önlem}
\end{minipage} \\
\midrule\noalign{}
\endhead
\bottomrule\noalign{}
\endlastfoot
Kırmızı et TÜFE alt endeksi (aylık) & 3 ay üst üste genel TÜFE'nin 5 puan üstünde & ESK'nın ek ithalat ve stok satışı \\
\end{longtable}

\endgroup

\textbf{13. Kazanan ve kaybedenler; bölgesel ve sosyal etki; telafi.} Etkilenen kesim yok. 1/98 askıya alınırsa et ithalatçıları ve tüketiciler kaybeder; telafi ESK'nın kaynak çeşitlendirmesi.

\textbf{14. Riskler ve azaltım}

\begin{itemize}
\tightlist
\item
  AB'nin 1/98'i askıya alması ithalat maliyetini artırır → E-10-01 ve E-10.11-01.
\end{itemize}

\textbf{15. Güven düzeyi; veri boşlukları; \etiketD{}.} Yüksek (veri); ihracatın bileşimi \etiketH{}. \etiketD{}: AB'ye et ürünü ihracatındaki SPS listelenme durumu.

\hypertarget{nace_kartlar_10-12-32-1012-kuxfcmes-hayvanlarux131-etlerinin-iux15flenmesi-ve-saklanmasux131--gb-statuxfcsuxfc-tarim-198--birincilikincil-arketip-bd--duruux15f-geliux15ftir}{%
\subsection{3.2 10.12 Kümes hayvanları etlerinin işlenmesi ve saklanması --- GB statüsü: TARIM-1/98 \textbar{} Birincil/ikincil arketip: B/D \textbar{} Duruş: Geliştir}\label{nace_kartlar_10-12-32-1012-kuxfcmes-hayvanlarux131-etlerinin-iux15flenmesi-ve-saklanmasux131--gb-statuxfcsuxfc-tarim-198--birincilikincil-arketip-bd--duruux15f-geliux15ftir}}

\emph{Güçlü ihracatçı (\$643 mn), ama ihracatın \%70'i Irak'a gidiyor; AB'ye ihracat yok denecek kadar az. GB çıkışının etkisi sıfıra yakın, asıl risk pazar yoğunlaşması.}

\textbf{1. Kimlik.} Rev. 2.1: 10.12 (1:1); CPA 10.12. Başlıca HS4 (TR ihracatı 2024): 0207 \%100, 0505 \%0, 6701 \%0, 0209 \%0. BEC payları (ticaret ağırlıklı): ara \%2 / sermaye \%0 / tüketim \%98.

\textbf{2. Ölçek}

\begingroup\gbtablo\small

\begin{longtable}[]{@{}
  >{\raggedright\arraybackslash}p{(\columnwidth - 4\tabcolsep) * \real{0.2708}}
  >{\raggedright\arraybackslash}p{(\columnwidth - 4\tabcolsep) * \real{0.4283}}
  >{\raggedright\arraybackslash}p{(\columnwidth - 4\tabcolsep) * \real{0.3009}}@{}}
\toprule\noalign{}
\begin{minipage}[b]{\linewidth}\raggedright
\textbf{Gösterge}
\end{minipage} & \begin{minipage}[b]{\linewidth}\raggedright
\textbf{Değer}
\end{minipage} & \begin{minipage}[b]{\linewidth}\raggedright
\textbf{Etiket ve kaynak}
\end{minipage} \\
\midrule\noalign{}
\endhead
\bottomrule\noalign{}
\endlastfoot
Üretim değeri; katma değer & \$7.057 mn; \$2.136 mn (GSYH'nin \%0,16) & TÜİK YSHİ 2024 \etiketV{} \\
Çalışan; girişim (0 / 1--9 / 10--49 / 50--249 / 250+) & 35.710; 130 (12 / 77 / 24 / 0 / 68) & TÜİK YSHİ ve İş Kayıtları 2024 \etiketV{}; gizli büyüklük hücreleri \etiketT{} \\
250+ girişimlerin istihdam payı; AB'ye göre kişi başı KD & \%98; 0,96 & \etiketT{} (AB27 SBS 2023, nominal kur) \\
İlk 3 İBBS-2 bölgesi & Manisa, Afyon, Kütahya, Uşak \%16; Kocaeli, Sakarya, Düzce, Bolu, Yalova \%15; Balıkesir, Çanakkale \%13 & \etiketT{} iş kayıtları + SGK 2025, RAS \\
Yabancı sermaye; çapa yatırımcılar & Banvit (BRF ve Katar Yatırım Otoritesi), Beypi, Şenpiliç, Erpiliç, Keskinoğlu, Gedik Piliç, Abalıoğlu \etiketH{} & \etiketH{} (firma bilgisi; pay verisi yok: VB-30) \\
\end{longtable}

\endgroup

\textbf{3. Dış ticaret}

\begingroup\gbtablo\small

\begin{longtable}[]{@{}
  >{\raggedright\arraybackslash}p{(\columnwidth - 2\tabcolsep) * \real{0.4839}}
  >{\raggedright\arraybackslash}p{(\columnwidth - 2\tabcolsep) * \real{0.5161}}@{}}
\toprule\noalign{}
\begin{minipage}[b]{\linewidth}\raggedright
\textbf{Gösterge}
\end{minipage} & \begin{minipage}[b]{\linewidth}\raggedright
\textbf{Değer (BACI 2024, \$; \etiketV{}, eşleme \etiketT{})}
\end{minipage} \\
\midrule\noalign{}
\endhead
\bottomrule\noalign{}
\endlastfoot
İhracat (2024); 2020--24 değişimi & \$643 mn; \%18 \\
AB payı; ilk 5 AB pazarı & \%0; DEU \%45; AUT \%28; MLT \%9; HUN \%4; ITA \%4 \\
İlk 5 AB dışı pazar & IRQ \%69; UZB \%4; MRT \%3; COG \%2; BLR \%2 \\
İthalat; AB ve Çin payı; ilk kaynaklar & \$86 mn; AB \%9, Çin \%1; UKR \%38; BRA \%33; USA \%12; ITA \%4; IRQ \%4 \\
Dış ticaret dengesi & \$+557 mn \\
İthal girdi içeriği (I-O, toplam) / KD÷üretim & \%20 / \%30 \\
RCA dünya (2020 → 2024) / AB pazarında & 1,94 → 1,45 / 0,07 \\
Birim değer ÷ AB pazarı ortalaması & 1,95 \\
AB pazar payı (AB dışı tedarik içinde, 2020 → 2024); TR'nin sırası & \%0,1 → \%0,3; 14. \\
\end{longtable}

\endgroup

\textbf{4. AB tarafı maruziyet}

\begin{itemize}
\tightlist
\item
  \kod{τ\_\allowbreak{}AB} (S4 artışı, rota A, TR→AB 2025 ağırlıklı): \%0,0; 1/98 kaybıyla \%0,0; tam MFN ortalaması \%0,0. Rejim: GB \%100, İTÜ \%0, tarım \%0 \etiketV{}/\etiketT{}.
\item
  Tepe satırlar (S4 artışı ≥ \%10; en büyük 3): yok.
\item
  AB pazarında AB dışı rakipler (tercih durumu): UKR \%39; BRA \%22; GBR \%15; CHN \%9; THA \%4. Tercihli rakip payı \%57.
\item
  Menşe (S2 / S3): risk D/D; \kod{c\_\allowbreak{}RoO} \%1,5 / \%2,0; \kod{p\_\allowbreak{}uyum} 0,90 / 0,85 \etiketH{}. Kural: Fasıl 2 (0207): tamamen elde edilmiş.
\item
  EUDR kapsamındaki AB ihracatı payı: \%0; SKDM: \%0.
\item
  AB pazarına taze kanatlı eti ihracatı SPS listelenmesine bağlı; bugün yok denecek kadar az (\$3 mn) \etiketD{}.
\end{itemize}

\textbf{5. Türkiye tarafı maruziyet}

\begin{itemize}
\tightlist
\item
  Uygulanan tarife (2026, ithalat ağırlıklı): MFN \%56,1, İGV dahil \%56,1 \etiketV{}. S4'te AB menşeline (c2): \%0,4.
\item
  Kritik girdiler (HS4 1005, 1201, 2304), Türkiye'nin toplam ithalatı 2024: \$3.393 mn (AB \%4; BRA \%35, UKR \%25, ARG \%13) \etiketV{}.
\item
  Yem maliyetinin ana kalemleri mısır ve soya küspesi. Bunlar ithal (Ukrayna, Brezilya, ABD) ve tarım rejiminde, GB'den bağımsız. DİR ile ihracata yönelik kullanılıyor.
\end{itemize}

\textbf{6. Kabiliyet}

\begin{itemize}
\tightlist
\item
  KS 46,0 (sıra 158/230; ağırlık duyarlılığında P10--P90 83--199). Bileşenler: RCA 42, PCI 13 (PCI −0,89), yurt içi KD 73, Ar-Ge 25, tedarikçi 84, verimlilik 95, firma 24 \etiketT{}.
\item
  Ürün uzayında atlanabilir ürünler (sınıf içi; yoğunluk üst çeyrek, PCI \textgreater{} sınıf): 050510 (PCI +0.13, RCA 0.08); 020711 (PCI +0.15, RCA 0.21) \etiketT{}.
\item
  \textbf{Kırmızı çizgi yetkinlikleri:} Entegre üretim (damızlık--kuluçka--kesim) ve helal sertifikasyonu; Irak dışı pazarlara giriş için kalıntı izleme.
\item
  \textbf{Geliştirme hedefi:} T+5: Irak payı \%70 → ≤ \%50 (Körfez, Orta Asya, Afrika); 2035: AB'ye pişmiş kanatlı ürünü listelenmesi.
\end{itemize}

\textbf{7. Stratejik önem.} SS 60,0 (en yüksek boyut: gıda güvenliği 80; ağırlıklı ortalama 13,5) \etiketH{}. En ucuz hayvansal protein; yem (mısır/soya) ithalatına bağımlı

\textbf{8. Makro duyarlılık}

\begingroup\gbtablo\small

\begin{longtable}[]{@{}lr@{}}
\toprule\noalign{}
\textbf{Gösterge} & \textbf{Değer (\etiketT{})} \\
\midrule\noalign{}
\endhead
\bottomrule\noalign{}
\endlastfoot
TÜFE ağırlığı (doğrudan, COICOP 2018 eşlemesi); tüketimde ithal pay & \%1,65; \%2 \\
ÜFE vekili (imalat üretimindeki pay) & \%1,06 \\
EDE: 1 puan tarife, tüm menşe / yalnız AB (bp; ρ = 1) & 0,048 / 0,004 \\
Geçişkenlik ρ (senaryolarda) & 0,6--1,0 (P50 0,8) \\
Piyasa yoğunlaşması (250+ istihdam payı) & \%98 \\
CKE: tarife (\$ mn / bp) ve mali (\$ mn / mlr TL); 0--100 puan & 873,9 ve 2,45; 25 \\
\end{longtable}

\endgroup

\textbf{9. Senaryo etkileri, T+5, P50 (P10 / P90)}

\begingroup\gbtablo\small

\begin{longtable}[]{@{}
  >{\raggedright\arraybackslash}p{(\columnwidth - 6\tabcolsep) * \real{0.3125}}
  >{\raggedleft\arraybackslash}p{(\columnwidth - 6\tabcolsep) * \real{0.3155}}
  >{\raggedleft\arraybackslash}p{(\columnwidth - 6\tabcolsep) * \real{0.1904}}
  >{\raggedleft\arraybackslash}p{(\columnwidth - 6\tabcolsep) * \real{0.1816}}@{}}
\toprule\noalign{}
\begin{minipage}[b]{\linewidth}\raggedright
\textbf{Değişken}
\end{minipage} & \begin{minipage}[b]{\linewidth}\raggedleft
\textbf{S2}
\end{minipage} & \begin{minipage}[b]{\linewidth}\raggedleft
\textbf{S3}
\end{minipage} & \begin{minipage}[b]{\linewidth}\raggedleft
\textbf{S4}
\end{minipage} \\
\midrule\noalign{}
\endhead
\bottomrule\noalign{}
\endlastfoot
AB'ye ihracat (\$ mn) & 0 (−1 / 0) & −1 (−1 / 0) & 0 (0 / 0) \\
Toplam ihracat (\$ mn) & 0 (−1 / 0) & 0 (−1 / 0) & 0 (0 / 0) \\
Nihai mal ithalatı (\$ mn) & 0 (−2 / +1) & −1 (−2 / 0) & 0 (−1 / 0) \\
Cari denge (\$ mn) & 0 (0 / +1) & 0 (0 / +1) & 0 (0 / 0) \\
Üretim (\$ mn) & 0 (−1 / +1) & 0 (0 / +1) & 0 (0 / +1) \\
İstihdam (kişi) & 0 (−2 / +3) & +1 (−1 / +4) & +1 (0 / +2) \\
TÜFE (bp, birinci tur) & 0,00 (0,00 / 0,00) & +0,01 (+0,01 / +0,01) & +0,01 (0,00 / +0,01) \\
Yatırım ve DYY (bütün senaryolar) & GB'den bağımsız; Irak riskine bağlı & & \\
\end{longtable}

\endgroup

\emph{Kaynak:} \etiketT{} f3e: kısmi denge, toplamlar Faz 2B GE dağılımına kalibre; S4'te 1/98 kaybı olasılığı 0,25.

\textbf{10. Politika paketi}

\begingroup\gbtablo\scriptsize

\begin{longtable}[]{@{}
  >{\raggedright\arraybackslash}p{(\columnwidth - 16\tabcolsep) * \real{0.0983}}
  >{\raggedright\arraybackslash}p{(\columnwidth - 16\tabcolsep) * \real{0.1197}}
  >{\raggedright\arraybackslash}p{(\columnwidth - 16\tabcolsep) * \real{0.1403}}
  >{\raggedright\arraybackslash}p{(\columnwidth - 16\tabcolsep) * \real{0.1018}}
  >{\raggedright\arraybackslash}p{(\columnwidth - 16\tabcolsep) * \real{0.1284}}
  >{\raggedright\arraybackslash}p{(\columnwidth - 16\tabcolsep) * \real{0.0738}}
  >{\raggedright\arraybackslash}p{(\columnwidth - 16\tabcolsep) * \real{0.1497}}
  >{\raggedright\arraybackslash}p{(\columnwidth - 16\tabcolsep) * \real{0.1167}}
  >{\raggedright\arraybackslash}p{(\columnwidth - 16\tabcolsep) * \real{0.0713}}@{}}
\toprule\noalign{}
\begin{minipage}[b]{\linewidth}\raggedright
\textbf{Eylem ID}
\end{minipage} & \begin{minipage}[b]{\linewidth}\raggedright
\textbf{Araç}
\end{minipage} & \begin{minipage}[b]{\linewidth}\raggedright
\textbf{Parametre}
\end{minipage} & \begin{minipage}[b]{\linewidth}\raggedright
\textbf{Sorumlu}
\end{minipage} & \begin{minipage}[b]{\linewidth}\raggedright
\textbf{Hukuki araç}
\end{minipage} & \begin{minipage}[b]{\linewidth}\raggedright
\textbf{Bütçe}
\end{minipage} & \begin{minipage}[b]{\linewidth}\raggedright
\textbf{Başlangıç--bitiş; sunset}
\end{minipage} & \begin{minipage}[b]{\linewidth}\raggedright
\textbf{Performans koşulu}
\end{minipage} & \begin{minipage}[b]{\linewidth}\raggedright
\textbf{KPI}
\end{minipage} \\
\midrule\noalign{}
\endhead
\bottomrule\noalign{}
\endlastfoot
E-10.12-01 & Pazar çeşitlendirme & Körfez İşbirliği Konseyi, Orta Asya ve Batı Afrika'da helal kanatlı tescili ve kota & Ticaret Bakanlığı (İhracat GM), TOB & STA müzakereleri; SPS ikili protokolleri & 0,05 mlr TL/yıl \etiketT{} & T−12 → T+36; sunset: T+36 & Yeni pazarlarda tescil & Irak payı ≤ \%50 \\
E-10.12-02 & AB SPS listelenme talebi & Isıl işlem görmüş kanatlı ürünleri için AB üçüncü ülke listesi (S2 SPS bölümü) & TOB, Ticaret Bakanlığı & S2 anlaşma metni (SPS) & İdari & T−18 → T+24; sunset: --- & AB denetimi (DG SANTE) & Liste kararı \\
\end{longtable}

\endgroup

\textbf{12. Erken uyarı ve tetikler}

\begingroup\gbtablo\small

\begin{longtable}[]{@{}
  >{\raggedright\arraybackslash}p{(\columnwidth - 4\tabcolsep) * \real{0.2576}}
  >{\raggedright\arraybackslash}p{(\columnwidth - 4\tabcolsep) * \real{0.1162}}
  >{\raggedright\arraybackslash}p{(\columnwidth - 4\tabcolsep) * \real{0.6262}}@{}}
\toprule\noalign{}
\begin{minipage}[b]{\linewidth}\raggedright
\textbf{Gösterge}
\end{minipage} & \begin{minipage}[b]{\linewidth}\raggedright
\textbf{Eşik}
\end{minipage} & \begin{minipage}[b]{\linewidth}\raggedright
\textbf{Tetiklenecek önlem}
\end{minipage} \\
\midrule\noalign{}
\endhead
\bottomrule\noalign{}
\endlastfoot
Irak'a kanatlı ihracatı (aylık) & Yıllık −\%20 & Eximbank kredi sigortası limitinin artırılması; pazar çeşitlendirme hızlandırılır \\
\end{longtable}

\endgroup

\textbf{13. Kazanan ve kaybedenler; bölgesel ve sosyal etki; telafi.} GB çıkışından etkilenen yok. Bölgeler: Manisa--Afyon (\%16), Kocaeli--Sakarya--Bolu (\%15), Balıkesir--Çanakkale (\%13); risk Irak'a bağlı.

\textbf{14. Riskler ve azaltım}

\begin{itemize}
\tightlist
\item
  Irak'ın ithalat yasağı ya da kotası (2024 öncesinde de görüldü) → pazar çeşitlendirme.
\item
  Kuş gribi → bölgeselleştirme protokolleri.
\end{itemize}

\textbf{15. Güven düzeyi; veri boşlukları; \etiketD{}.} Yüksek. \etiketD{}: AB'ye kanatlı ürünü listelenme durumu (DG SANTE üçüncü ülke listeleri).

\hypertarget{nace_kartlar_10-12-33-1013-et-ve-kuxfcmes-hayvanlarux131-etlerinden-uxfcretilen-uxfcruxfcnlerin-imalatux131--gb-statuxfcsuxfc-tarim-198--birincilikincil-arketip---duruux15f-koru}{%
\subsection{3.3 10.13 Et ve kümes hayvanları etlerinden üretilen ürünlerin imalatı --- GB statüsü: TARIM-1/98 \textbar{} Birincil/ikincil arketip: --- \textbar{} Duruş: Koru}\label{nace_kartlar_10-12-33-1013-et-ve-kuxfcmes-hayvanlarux131-etlerinden-uxfcretilen-uxfcruxfcnlerin-imalatux131--gb-statuxfcsuxfc-tarim-198--birincilikincil-arketip---duruux15f-koru}}

\emph{Küçük ve iç pazara dönük. AB ihracatı Kıbrıs düzeltmesinden sonra \$5 mn. GB çıkışının etkisi yok; S4'te 1/98 kaybında AB MFN'si \%53, ama hacim küçük.}

\textbf{1. Kimlik.} Rev. 2.1: 10.13 (1:1); CPA 10.13. Başlıca HS4 (TR ihracatı 2024): 1602 \%64, 1601 \%32, 2301 \%2, 0210 \%2. BEC payları (ticaret ağırlıklı): ara \%23 / sermaye \%0 / tüketim \%77.

\textbf{2. Ölçek}

\begingroup\gbtablo\small

\begin{longtable}[]{@{}
  >{\raggedright\arraybackslash}p{(\columnwidth - 4\tabcolsep) * \real{0.3158}}
  >{\raggedright\arraybackslash}p{(\columnwidth - 4\tabcolsep) * \real{0.3249}}
  >{\raggedright\arraybackslash}p{(\columnwidth - 4\tabcolsep) * \real{0.3593}}@{}}
\toprule\noalign{}
\begin{minipage}[b]{\linewidth}\raggedright
\textbf{Gösterge}
\end{minipage} & \begin{minipage}[b]{\linewidth}\raggedright
\textbf{Değer}
\end{minipage} & \begin{minipage}[b]{\linewidth}\raggedright
\textbf{Etiket ve kaynak}
\end{minipage} \\
\midrule\noalign{}
\endhead
\bottomrule\noalign{}
\endlastfoot
Üretim değeri; katma değer & \$2.240 mn; \$379 mn (GSYH'nin \%0,03) & TÜİK YSHİ 2024 \etiketV{} \\
Çalışan; girişim (0 / 1--9 / 10--49 / 50--249 / 250+) & 11.726; 945 (62 / 778 / 135 / 31 / 7) & TÜİK YSHİ ve İş Kayıtları 2024 \etiketV{}; gizli büyüklük hücreleri \etiketT{} \\
250+ girişimlerin istihdam payı; AB'ye göre kişi başı KD & \%26; 0,64 & \etiketT{} (AB27 SBS 2023, nominal kur) \\
İlk 3 İBBS-2 bölgesi & İzmir \%16; İstanbul \%16; Manisa, Afyon, Kütahya, Uşak \%8 & \etiketT{} iş kayıtları + SGK 2025, RAS \\
Yabancı sermaye; çapa yatırımcılar & Pınar, Namet, Polonez, Apikoğlu \etiketH{} & \etiketH{} (firma bilgisi; pay verisi yok: VB-30) \\
\end{longtable}

\endgroup

\textbf{3. Dış ticaret}

\begingroup\gbtablo\small

\begin{longtable}[]{@{}
  >{\raggedright\arraybackslash}p{(\columnwidth - 2\tabcolsep) * \real{0.4802}}
  >{\raggedright\arraybackslash}p{(\columnwidth - 2\tabcolsep) * \real{0.5198}}@{}}
\toprule\noalign{}
\begin{minipage}[b]{\linewidth}\raggedright
\textbf{Gösterge}
\end{minipage} & \begin{minipage}[b]{\linewidth}\raggedright
\textbf{Değer (BACI 2024, \$; \etiketV{}, eşleme \etiketT{})}
\end{minipage} \\
\midrule\noalign{}
\endhead
\bottomrule\noalign{}
\endlastfoot
İhracat (2024); 2020--24 değişimi & \$214 mn; \%152 \\
AB payı; ilk 5 AB pazarı & \%2; AUT \%55; FRA \%20; BGR \%11; MLT \%5; NLD \%2 \\
İlk 5 AB dışı pazar & IRQ \%21; VEN \%15; SYR \%7; CYP \%5; KWT \%4 \\
İthalat; AB ve Çin payı; ilk kaynaklar & \$64 mn; AB \%56, Çin \%0; POL \%23; GBR \%16; DNK \%11; IRN \%9; TUN \%8 \\
Dış ticaret dengesi & \$+150 mn \\
İthal girdi içeriği (I-O, toplam) / KD÷üretim & \%20 / \%17 \\
RCA dünya (2020 → 2024) / AB pazarında & 0,28 → 0,49 / 0,06 \\
Birim değer ÷ AB pazarı ortalaması & 0,83 \\
AB pazar payı (AB dışı tedarik içinde, 2020 → 2024); TR'nin sırası & \%0,1 → \%0,2; 14. \\
\end{longtable}

\endgroup

\textbf{4. AB tarafı maruziyet}

\begin{itemize}
\tightlist
\item
  \kod{τ\_\allowbreak{}AB} (S4 artışı, rota A, TR→AB 2025 ağırlıklı): \%0,0; 1/98 kaybıyla \%53,0; tam MFN ortalaması \%53,0. Rejim: GB \%0, İTÜ \%0, tarım \%100 \etiketV{}/\etiketT{}.
\item
  Tepe satırlar (S4 artışı ≥ \%10; en büyük 3): yok.
\item
  AB pazarında AB dışı rakipler (tercih durumu): BRA \%33; THA \%31; GBR \%14; CHN \%13; CHE \%3. Tercihli rakip payı \%20.
\item
  Menşe (S2 / S3): risk D/D; \kod{c\_\allowbreak{}RoO} \%1,5 / \%2,0; \kod{p\_\allowbreak{}uyum} 0,90 / 0,85 \etiketH{}. Kural: 1601--1602: Fasıl 1, 2 ve 16 malzemeleri tamamen elde edilmiş.
\item
  EUDR kapsamındaki AB ihracatı payı: \%0; SKDM: \%0.
\end{itemize}

\textbf{5. Türkiye tarafı maruziyet}

\begin{itemize}
\tightlist
\item
  Uygulanan tarife (2026, ithalat ağırlıklı): MFN \%11,6, İGV dahil \%11,6 \etiketV{}. S4'te AB menşeline (c2): \%0,0.
\item
  Kritik girdiler (HS4 0201, 0202, 0207), Türkiye'nin toplam ithalatı 2024: \$923 mn (AB \%65; POL \%33, BRA \%28, FRA \%16) \etiketV{}.
\end{itemize}

\textbf{6. Kabiliyet}

\begin{itemize}
\tightlist
\item
  KS 54,1 (sıra 69/230; ağırlık duyarlılığında P10--P90 41--118). Bileşenler: RCA 45, PCI 60 (PCI 0,17), yurt içi KD 50, Ar-Ge 25, tedarikçi 84, verimlilik 81, firma 55 \etiketT{}.
\item
  Ürün uzayında atlanabilir ürünler (sınıf içi; yoğunluk üst çeyrek, PCI \textgreater{} sınıf): 160231 (PCI +0.43, RCA 0.38) \etiketT{}.
\end{itemize}

\textbf{7. Stratejik önem.} SS 29,8 (en yüksek boyut: gıda güvenliği 40; ağırlıklı ortalama 6,0) \etiketH{}. İşlenmiş et; temel tüketim

\textbf{8. Makro duyarlılık}

\begingroup\gbtablo\small

\begin{longtable}[]{@{}lr@{}}
\toprule\noalign{}
\textbf{Gösterge} & \textbf{Değer (\etiketT{})} \\
\midrule\noalign{}
\endhead
\bottomrule\noalign{}
\endlastfoot
TÜFE ağırlığı (doğrudan, COICOP 2018 eşlemesi); tüketimde ithal pay & \%0,85; \%5 \\
ÜFE vekili (imalat üretimindeki pay) & \%0,34 \\
EDE: 1 puan tarife, tüm menşe / yalnız AB (bp; ρ = 1) & 0,046 / 0,026 \\
Geçişkenlik ρ (senaryolarda) & 0,6--1,0 (P50 0,8) \\
Piyasa yoğunlaşması (250+ istihdam payı) & \%26 \\
CKE: tarife (\$ mn / bp) ve mali (\$ mn / mlr TL); 0--100 puan & 257,0 ve 2,09; 16 \\
\end{longtable}

\endgroup

\textbf{9. Senaryo etkileri, T+5, P50 (P10 / P90)}

\begingroup\gbtablo\small

\begin{longtable}[]{@{}lrrr@{}}
\toprule\noalign{}
\textbf{Değişken} & \textbf{S2} & \textbf{S3} & \textbf{S4} \\
\midrule\noalign{}
\endhead
\bottomrule\noalign{}
\endlastfoot
AB'ye ihracat (\$ mn) & 0 (0 / 0) & 0 (0 / 0) & 0 (−2 / 0) \\
Toplam ihracat (\$ mn) & 0 (0 / 0) & 0 (0 / 0) & 0 (−2 / 0) \\
Nihai mal ithalatı (\$ mn) & 0 (0 / +3) & 0 (0 / +2) & 0 (0 / 0) \\
Cari denge (\$ mn) & 0 (−2 / 0) & 0 (−1 / 0) & 0 (−2 / 0) \\
Üretim (\$ mn) & 0 (−2 / 0) & 0 (−2 / 0) & 0 (−2 / 0) \\
İstihdam (kişi) & 0 (−9 / 0) & 0 (−5 / 0) & 0 (−7 / 0) \\
TÜFE (bp, birinci tur) & 0,00 (0,00 / 0,00) & 0,00 (0,00 / 0,00) & 0,00 (0,00 / 0,00) \\
Yatırım ve DYY (bütün senaryolar) & --- & & \\
\end{longtable}

\endgroup

\emph{Kaynak:} \etiketT{} f3e: kısmi denge, toplamlar Faz 2B GE dağılımına kalibre; S4'te 1/98 kaybı olasılığı 0,25.

\textbf{13. Kazanan ve kaybedenler; bölgesel ve sosyal etki; telafi.} Etkilenen kesim yok.

\textbf{15. Güven düzeyi; veri boşlukları; \etiketD{}.} Yüksek.

\hypertarget{nace_kartlar_10-12-34-1020-balux131k-kabuklu-deniz-hayvanlarux131-ve-yumuux15fakuxe7alarux131n-iux15flenmesi-ve-saklanmasux131--gb-statuxfcsuxfc-tarim-198--birincilikincil-arketip-a--duruux15f-koru}{%
\subsection{3.4 10.20 Balık, kabuklu deniz hayvanları ve yumuşakçaların işlenmesi ve saklanması --- GB statüsü: TARIM-1/98 \textbar{} Birincil/ikincil arketip: A \textbar{} Duruş: Koru}\label{nace_kartlar_10-12-34-1020-balux131k-kabuklu-deniz-hayvanlarux131-ve-yumuux15fakuxe7alarux131n-iux15flenmesi-ve-saklanmasux131--gb-statuxfcsuxfc-tarim-198--birincilikincil-arketip-a--duruux15f-koru}}

\emph{AB'ye entegre bir tarım ihracatçısı: levrekte dünya birincisi (\%32), çipurada ikinci (\%31); AB'ye \$491 mn ihracat. GB'ye bağlı değil, 1/98'in balıkçılık tavizlerine bağlı. AB ticaret savunmasının tarihi var.}

\textbf{1. Kimlik.} Rev. 2.1: 10.20 (1:1); CPA 10.20. Başlıca HS4 (TR ihracatı 2024): 0304 \%42, 0303 \%39, 1604 \%8, 0305 \%4. BEC payları (ticaret ağırlıklı): ara \%16 / sermaye \%0 / tüketim \%79.

\textbf{2. Ölçek}

\begingroup\gbtablo\small

\begin{longtable}[]{@{}
  >{\raggedright\arraybackslash}p{(\columnwidth - 4\tabcolsep) * \real{0.2812}}
  >{\raggedright\arraybackslash}p{(\columnwidth - 4\tabcolsep) * \real{0.4044}}
  >{\raggedright\arraybackslash}p{(\columnwidth - 4\tabcolsep) * \real{0.3144}}@{}}
\toprule\noalign{}
\begin{minipage}[b]{\linewidth}\raggedright
\textbf{Gösterge}
\end{minipage} & \begin{minipage}[b]{\linewidth}\raggedright
\textbf{Değer}
\end{minipage} & \begin{minipage}[b]{\linewidth}\raggedright
\textbf{Etiket ve kaynak}
\end{minipage} \\
\midrule\noalign{}
\endhead
\bottomrule\noalign{}
\endlastfoot
Üretim değeri; katma değer & \$2.685 mn; \$722 mn (GSYH'nin \%0,05) & TÜİK YSHİ 2024 \etiketV{} \\
Çalışan; girişim (0 / 1--9 / 10--49 / 50--249 / 250+) & 14.445; 380 (42 / 296 / 51 / 23 / 9) & TÜİK YSHİ ve İş Kayıtları 2024 \etiketV{}; gizli büyüklük hücreleri \etiketT{} \\
250+ girişimlerin istihdam payı; AB'ye göre kişi başı KD & \%69; 0,85 & \etiketT{} (AB27 SBS 2023, nominal kur) \\
İlk 3 İBBS-2 bölgesi & İzmir \%19; Aydın, Denizli, Muğla \%18; Trabzon, Ordu, Giresun, Rize, Artvin, Gümüşhane \%16 & \etiketT{} iş kayıtları + SGK 2025, RAS \\
Yabancı sermaye; çapa yatırımcılar & Kılıç Deniz, Gümüşdoğa, Agromey, Özpekler, Fjord Marin (Norveç sermayesi) \etiketH{} & \etiketH{} (firma bilgisi; pay verisi yok: VB-30) \\
\end{longtable}

\endgroup

\textbf{3. Dış ticaret}

\begingroup\gbtablo\small

\begin{longtable}[]{@{}
  >{\raggedright\arraybackslash}p{(\columnwidth - 2\tabcolsep) * \real{0.4846}}
  >{\raggedright\arraybackslash}p{(\columnwidth - 2\tabcolsep) * \real{0.5154}}@{}}
\toprule\noalign{}
\begin{minipage}[b]{\linewidth}\raggedright
\textbf{Gösterge}
\end{minipage} & \begin{minipage}[b]{\linewidth}\raggedright
\textbf{Değer (BACI 2024, \$; \etiketV{}, eşleme \etiketT{})}
\end{minipage} \\
\midrule\noalign{}
\endhead
\bottomrule\noalign{}
\endlastfoot
İhracat (2024); 2020--24 değişimi & \$1.237 mn; \%102 \\
AB payı; ilk 5 AB pazarı & \%40; ITA \%29; NLD \%22; DEU \%15; GRC \%10; FRA \%5 \\
İlk 5 AB dışı pazar & RUS \%41; GBR \%14; JPN \%8; USA \%7; KOR \%5 \\
İthalat; AB ve Çin payı; ilk kaynaklar & \$514 mn; AB \%7, Çin \%4; MAR \%31; NOR \%8; MYS \%7; RUS \%6; GIN \%5 \\
Dış ticaret dengesi & \$+723 mn \\
İthal girdi içeriği (I-O, toplam) / KD÷üretim & \%20 / \%27 \\
RCA dünya (2020 → 2024) / AB pazarında & 0,48 → 0,75 / 0,53 \\
Birim değer ÷ AB pazarı ortalaması & 1,18 \\
AB pazar payı (AB dışı tedarik içinde, 2020 → 2024); TR'nin sırası & \%1,5 → \%2,1; 14. \\
\end{longtable}

\endgroup

\textbf{4. AB tarafı maruziyet}

\begin{itemize}
\tightlist
\item
  \kod{τ\_\allowbreak{}AB} (S4 artışı, rota A, TR→AB 2025 ağırlıklı): \%0,0; 1/98 kaybıyla \%16,3; tam MFN ortalaması \%16,3. Rejim: GB \%0, İTÜ \%0, tarım \%100 \etiketV{}/\etiketT{}.
\item
  Tepe satırlar (S4 artışı ≥ \%10; en büyük 3): yok.
\item
  AB pazarında AB dışı rakipler (tercih durumu): NOR \%15; ECU \%9; MAR \%8; CHN \%7; IND \%4. Tercihli rakip payı \%64.
\item
  Menşe (S2 / S3): risk D/D; \kod{c\_\allowbreak{}RoO} \%1,5 / \%2,0; \kod{p\_\allowbreak{}uyum} 0,90 / 0,85 \etiketH{}. Kural: Fasıl 3 ve 1604--1605: tamamen elde edilmiş (yetiştiricilikte yumurta/yavru TR'de).
\item
  EUDR kapsamındaki AB ihracatı payı: \%0; SKDM: \%0.
\item
  AB ticaret savunması: AB Türk gökkuşağı alabalığına 2015'te telafi edici vergi koydu (2015/309), vergiyi 2021'de beş yıl uzattı (2021/823), vergi 26.5.2026'da sona erdi. Levrek ve çipurada sübvansiyon soruşturması 2016'da önlemsiz kapandı \etiketV{}. Yunan ve İspanyol üreticiler hassas; hasmane AB'de yeni soruşturma olasılığı yüksek \etiketH{}.
\item
  Taze balıkta AB MFN'si \%15. 1/98 kaybında (P10) AB'ye ihracat −\$137 mn ve −417 kişi.
\item
  Değer zinciri: AB perakende zincirlerine taze ürün (çabuk bozulan, tam zamanında). Sınır gecikmesi maliyeti yüksek.
\end{itemize}

\textbf{5. Türkiye tarafı maruziyet}

\begin{itemize}
\tightlist
\item
  Uygulanan tarife (2026, ithalat ağırlıklı): MFN \%14,7, İGV dahil \%19,5 \etiketV{}. S4'te AB menşeline (c2): \%0,0.
\item
  Kritik girdiler (HS4 0303, 2301), Türkiye'nin toplam ithalatı 2024: \$432 mn (AB \%13; MAR \%31, NOR \%7, RUS \%7) \etiketV{}.
\item
  Balık yemi girdisi balık unu (Moritanya, Peru) ve soya. Yavru üretimi yurt içinde.
\end{itemize}

\textbf{6. Kabiliyet}

\begin{itemize}
\tightlist
\item
  KS 46,7 (sıra 153/230; ağırlık duyarlılığında P10--P90 79--192). Bileşenler: RCA 51, PCI 10 (PCI −0,97), yurt içi KD 65, Ar-Ge 25, tedarikçi 84, verimlilik 94, firma 31 \etiketT{}.
\item
  Ürün uzayında atlanabilir ürünler (sınıf içi; yoğunluk üst çeyrek, PCI \textgreater{} sınıf): 030539 (PCI -0.91, RCA 0.30); 030619 (PCI -0.96, RCA 0.08) \etiketT{}.
\item
  \textbf{Kırmızı çizgi yetkinlikleri:} Kuluçkahane ve genetik (yavru) kapasitesi; AB sağlık sertifikası ve izlenebilirlik (ASC/GlobalG.A.P.).
\item
  \textbf{Geliştirme hedefi:} T+5: AB payını koruyup ABD ve Körfez'e işlenmiş fileto; 2035: açık deniz yetiştiriciliği ve yeni türler (kuzey Ege).
\end{itemize}

\textbf{7. Stratejik önem.} SS 37,7 (en yüksek boyut: gıda güvenliği 50; ağırlıklı ortalama 9,0) \etiketH{}. Su ürünleri; yetiştiricilik ihracatı

\textbf{8. Makro duyarlılık}

\begingroup\gbtablo\small

\begin{longtable}[]{@{}lr@{}}
\toprule\noalign{}
\textbf{Gösterge} & \textbf{Değer (\etiketT{})} \\
\midrule\noalign{}
\endhead
\bottomrule\noalign{}
\endlastfoot
TÜFE ağırlığı (doğrudan, COICOP 2018 eşlemesi); tüketimde ithal pay & \%0,14; \%6 \\
ÜFE vekili (imalat üretimindeki pay) & \%0,40 \\
EDE: 1 puan tarife, tüm menşe / yalnız AB (bp; ρ = 1) & 0,048 / 0,003 \\
Geçişkenlik ρ (senaryolarda) & 0,6--1,0 (P50 0,8) \\
Piyasa yoğunlaşması (250+ istihdam payı) & \%69 \\
CKE: tarife (\$ mn / bp) ve mali (\$ mn / mlr TL); 0--100 puan & 1.287,1 ve 14,92; 42 \\
\end{longtable}

\endgroup

\textbf{9. Senaryo etkileri, T+5, P50 (P10 / P90)}

\begingroup\gbtablo\small

\begin{longtable}[]{@{}
  >{\raggedright\arraybackslash}p{(\columnwidth - 6\tabcolsep) * \real{0.2707}}
  >{\raggedleft\arraybackslash}p{(\columnwidth - 6\tabcolsep) * \real{0.4221}}
  >{\raggedleft\arraybackslash}p{(\columnwidth - 6\tabcolsep) * \real{0.1536}}
  >{\raggedleft\arraybackslash}p{(\columnwidth - 6\tabcolsep) * \real{0.1536}}@{}}
\toprule\noalign{}
\begin{minipage}[b]{\linewidth}\raggedright
\textbf{Değişken}
\end{minipage} & \begin{minipage}[b]{\linewidth}\raggedleft
\textbf{S2}
\end{minipage} & \begin{minipage}[b]{\linewidth}\raggedleft
\textbf{S3}
\end{minipage} & \begin{minipage}[b]{\linewidth}\raggedleft
\textbf{S4}
\end{minipage} \\
\midrule\noalign{}
\endhead
\bottomrule\noalign{}
\endlastfoot
AB'ye ihracat (\$ mn) & 0 (0 / +10) & 0 (0 / 0) & 0 (−137 / 0) \\
Toplam ihracat (\$ mn) & 0 (0 / +7) & 0 (0 / 0) & 0 (−112 / 0) \\
Nihai mal ithalatı (\$ mn) & 0 (0 / +2) & 0 (0 / +1) & 0 (0 / 0) \\
Cari denge (\$ mn) & 0 (0 / +4) & 0 (0 / 0) & 0 (−89 / 0) \\
Üretim (\$ mn) & 0 (−1 / +5) & 0 (−1 / 0) & 0 (−112 / 0) \\
İstihdam (kişi) & 0 (−2 / +19) & 0 (−2 / 0) & 0 (−417 / 0) \\
TÜFE (bp, birinci tur) & 0,00 (0,00 / 0,00) & 0,00 (0,00 / 0,00) & 0,00 (0,00 / 0,00) \\
Yatırım ve DYY (bütün senaryolar) & Norveç sermayeli yatırımlar 1/98 belirsizliğine duyarlı \etiketH{} & & \\
\end{longtable}

\endgroup

\emph{Kaynak:} \etiketT{} f3e: kısmi denge, toplamlar Faz 2B GE dağılımına kalibre; S4'te 1/98 kaybı olasılığı 0,25.

\textbf{10. Politika paketi}

\begingroup\gbtablo\scriptsize

\begin{longtable}[]{@{}
  >{\raggedright\arraybackslash}p{(\columnwidth - 16\tabcolsep) * \real{0.0983}}
  >{\raggedright\arraybackslash}p{(\columnwidth - 16\tabcolsep) * \real{0.1347}}
  >{\raggedright\arraybackslash}p{(\columnwidth - 16\tabcolsep) * \real{0.1483}}
  >{\raggedright\arraybackslash}p{(\columnwidth - 16\tabcolsep) * \real{0.0969}}
  >{\raggedright\arraybackslash}p{(\columnwidth - 16\tabcolsep) * \real{0.0846}}
  >{\raggedright\arraybackslash}p{(\columnwidth - 16\tabcolsep) * \real{0.0588}}
  >{\raggedright\arraybackslash}p{(\columnwidth - 16\tabcolsep) * \real{0.1457}}
  >{\raggedright\arraybackslash}p{(\columnwidth - 16\tabcolsep) * \real{0.1107}}
  >{\raggedright\arraybackslash}p{(\columnwidth - 16\tabcolsep) * \real{0.1220}}@{}}
\toprule\noalign{}
\begin{minipage}[b]{\linewidth}\raggedright
\textbf{Eylem ID}
\end{minipage} & \begin{minipage}[b]{\linewidth}\raggedright
\textbf{Araç}
\end{minipage} & \begin{minipage}[b]{\linewidth}\raggedright
\textbf{Parametre}
\end{minipage} & \begin{minipage}[b]{\linewidth}\raggedright
\textbf{Sorumlu}
\end{minipage} & \begin{minipage}[b]{\linewidth}\raggedright
\textbf{Hukuki araç}
\end{minipage} & \begin{minipage}[b]{\linewidth}\raggedright
\textbf{Bütçe}
\end{minipage} & \begin{minipage}[b]{\linewidth}\raggedright
\textbf{Başlangıç--bitiş; sunset}
\end{minipage} & \begin{minipage}[b]{\linewidth}\raggedright
\textbf{Performans koşulu}
\end{minipage} & \begin{minipage}[b]{\linewidth}\raggedright
\textbf{KPI}
\end{minipage} \\
\midrule\noalign{}
\endhead
\bottomrule\noalign{}
\endlastfoot
E-10.20-01 & Destek programlarının ASCM uyum dosyası & Yetiştiricilik desteklerinin (yem, yavru) ihracata bağlı olmadığının belgelenmesi; destek tasarımında ihracat koşulu yok & TOB, Ticaret Bakanlığı (İthalat GM) & Destek tebliğleri & İdari & T−12 → T; sunset: --- & Dosya hazır & AB soruşturması açılırsa ilk yanıt 30 gün \\
E-10.20-02 & S2'de balıkçılık tavizi & Levrek, çipura ve alabalıkta kota/tavan yerine \%0; menşe: yetiştiricilikte yavru TR'de & Ticaret Bakanlığı, TOB & S2 anlaşma metni (tarım eki) & İdari & T−18 → T; sunset: --- & --- & Tavizli satırların AB ihracatı ≥ T−12 düzeyi \\
\end{longtable}

\endgroup

\textbf{11. Müzakere pozisyonu}

\begin{itemize}
\tightlist
\item
  AB ile hücum: 1/98 balıkçılık tavizlerinin tam liberalizasyona dönüştürülmesi; karşılığında AB'ye işlenmiş ürün kotası (E-10-06).
\item
  Üçüncü ülke STA'ları: Körfez (taze levrek/çipura), Japonya (alabalık).
\end{itemize}

\textbf{12. Erken uyarı ve tetikler}

\begingroup\gbtablo\small

\begin{longtable}[]{@{}
  >{\raggedright\arraybackslash}p{(\columnwidth - 4\tabcolsep) * \real{0.3788}}
  >{\raggedright\arraybackslash}p{(\columnwidth - 4\tabcolsep) * \real{0.1168}}
  >{\raggedright\arraybackslash}p{(\columnwidth - 4\tabcolsep) * \real{0.5044}}@{}}
\toprule\noalign{}
\begin{minipage}[b]{\linewidth}\raggedright
\textbf{Gösterge}
\end{minipage} & \begin{minipage}[b]{\linewidth}\raggedright
\textbf{Eşik}
\end{minipage} & \begin{minipage}[b]{\linewidth}\raggedright
\textbf{Tetiklenecek önlem}
\end{minipage} \\
\midrule\noalign{}
\endhead
\bottomrule\noalign{}
\endlastfoot
AB'nin soruşturma başlatma ilanı (OJ C) & İlan & E-10-04: yanıt ekibi; AB ithalatçılarıyla ortak savunma \\
AB'ye taze balık ihracatı (aylık, ton) & Yıllık −\%15 & 1/98 ihlali iddiasıyla Ortaklık Konseyi ve DTÖ başvurusu \\
\end{longtable}

\endgroup

\textbf{13. Kazanan ve kaybedenler; bölgesel ve sosyal etki; telafi.} Kaybeden (yalnız 1/98 kaybında): Muğla (Milas, Bodrum) ve İzmir kıyı işletmeleri; 14 bin çalışan. Telafi: E-10-05 benzeri ihracat sigortası ve AB dışı pazar.

\textbf{14. Riskler ve azaltım}

\begin{itemize}
\tightlist
\item
  Yeni AB telafi edici vergi soruşturması → E-10-04.
\item
  1/98'in askıya alınması → E-10-01.
\end{itemize}

\textbf{15. Güven düzeyi; veri boşlukları; \etiketD{}.} Orta (1/98 tavizlerinin satır içeriği: VB-06). \etiketD{}: levrek ve çipurada bugünkü AB tercih oranı ve kota.

\hypertarget{nace_kartlar_10-12-35-1031-patatesin-iux15flenmesi-ve-saklanmasux131--gb-statuxfcsuxfc-karma--birincilikincil-arketip---duruux15f-koru}{%
\subsection{3.5 10.31 Patatesin işlenmesi ve saklanması --- GB statüsü: KARMA \textbar{} Birincil/ikincil arketip: --- \textbar{} Duruş: Koru}\label{nace_kartlar_10-12-35-1031-patatesin-iux15flenmesi-ve-saklanmasux131--gb-statuxfcsuxfc-karma--birincilikincil-arketip---duruux15f-koru}}

\emph{Küçük sınıf (patates ürünleri). AB ihracatı Kıbrıs düzeltmesinden sonra \$4 mn; GB etkisi yok.}

\textbf{1. Kimlik.} Rev. 2.1: 10.31 (1:1); CPA 10.31. Başlıca HS4 (TR ihracatı 2024): 2004 \%85, 2005 \%13, 0710 \%2, 1105 \%1. BEC payları (ticaret ağırlıklı): ara \%12 / sermaye \%0 / tüketim \%88.

\textbf{2. Ölçek}

\begingroup\gbtablo\small

\begin{longtable}[]{@{}
  >{\raggedright\arraybackslash}p{(\columnwidth - 4\tabcolsep) * \real{0.2861}}
  >{\raggedright\arraybackslash}p{(\columnwidth - 4\tabcolsep) * \real{0.3932}}
  >{\raggedright\arraybackslash}p{(\columnwidth - 4\tabcolsep) * \real{0.3207}}@{}}
\toprule\noalign{}
\begin{minipage}[b]{\linewidth}\raggedright
\textbf{Gösterge}
\end{minipage} & \begin{minipage}[b]{\linewidth}\raggedright
\textbf{Değer}
\end{minipage} & \begin{minipage}[b]{\linewidth}\raggedright
\textbf{Etiket ve kaynak}
\end{minipage} \\
\midrule\noalign{}
\endhead
\bottomrule\noalign{}
\endlastfoot
Üretim değeri; katma değer & \$385 mn; \$107 mn (GSYH'nin \%0,01) & TÜİK YSHİ 2024 \etiketV{} \\
Çalışan; girişim (0 / 1--9 / 10--49 / 50--249 / 250+) & 2.152; 145 (10 / 125 / 13 / 6 / 2) & TÜİK YSHİ ve İş Kayıtları 2024 \etiketV{}; gizli büyüklük hücreleri \etiketT{} \\
250+ girişimlerin istihdam payı; AB'ye göre kişi başı KD & \%38; 0,37 & \etiketT{} (AB27 SBS 2023, nominal kur) \\
İlk 3 İBBS-2 bölgesi & İstanbul \%46; Konya, Karaman \%16; Kırıkkale, Aksaray, Niğde, Nevşehir, Kırşehir \%12 & \etiketT{} iş kayıtları + SGK 2025, RAS \\
Yabancı sermaye; çapa yatırımcılar & Pepsico (Lay's; Tarsus ve Kocaeli \etiketD{}), Ülker \etiketH{} & \etiketH{} (firma bilgisi; pay verisi yok: VB-30) \\
\end{longtable}

\endgroup

\textbf{3. Dış ticaret}

\begingroup\gbtablo\small

\begin{longtable}[]{@{}
  >{\raggedright\arraybackslash}p{(\columnwidth - 2\tabcolsep) * \real{0.4839}}
  >{\raggedright\arraybackslash}p{(\columnwidth - 2\tabcolsep) * \real{0.5161}}@{}}
\toprule\noalign{}
\begin{minipage}[b]{\linewidth}\raggedright
\textbf{Gösterge}
\end{minipage} & \begin{minipage}[b]{\linewidth}\raggedright
\textbf{Değer (BACI 2024, \$; \etiketV{}, eşleme \etiketT{})}
\end{minipage} \\
\midrule\noalign{}
\endhead
\bottomrule\noalign{}
\endlastfoot
İhracat (2024); 2020--24 değişimi & \$110 mn; \%140 \\
AB payı; ilk 5 AB pazarı & \%3; DEU \%37; GRC \%23; BGR \%15; FRA \%9; NLD \%5 \\
İlk 5 AB dışı pazar & RUS \%54; IRQ \%15; CHN \%10; CYP \%6; GEO \%4 \\
İthalat; AB ve Çin payı; ilk kaynaklar & \$22 mn; AB \%85, Çin \%0; NLD \%38; DEU \%32; BEL \%7; RUS \%7; LBN \%5 \\
Dış ticaret dengesi & \$+88 mn \\
İthal girdi içeriği (I-O, toplam) / KD÷üretim & \%20 / \%28 \\
RCA dünya (2020 → 2024) / AB pazarında & 0,40 → 0,49 / 0,40 \\
Birim değer ÷ AB pazarı ortalaması & 0,76 \\
AB pazar payı (AB dışı tedarik içinde, 2020 → 2024); TR'nin sırası & \%1,3 → \%1,6; 11. \\
\end{longtable}

\endgroup

\textbf{4. AB tarafı maruziyet}

\begin{itemize}
\tightlist
\item
  \kod{τ\_\allowbreak{}AB} (S4 artışı, rota A, TR→AB 2025 ağırlıklı): \%0,0; 1/98 kaybıyla \%15,1; tam MFN ortalaması \%15,1. Rejim: GB \%0, İTÜ \%0, tarım \%100 \etiketV{}/\etiketT{}.
\item
  Tepe satırlar (S4 artışı ≥ \%10; en büyük 3): yok.
\item
  AB pazarında AB dışı rakipler (tercih durumu): GBR \%49; SRB \%13; EGY \%5; NOR \%4; CHE \%4. Tercihli rakip payı \%67.
\item
  Menşe (S2 / S3): risk D/D; \kod{c\_\allowbreak{}RoO} \%1,5 / \%2,0; \kod{p\_\allowbreak{}uyum} 0,90 / 0,85 \etiketH{}. Kural: 2004--2005: Fasıl 7 malzemesi tamamen elde edilmiş.
\item
  EUDR kapsamındaki AB ihracatı payı: \%0; SKDM: \%0.
\end{itemize}

\textbf{5. Türkiye tarafı maruziyet}

\begin{itemize}
\tightlist
\item
  Uygulanan tarife (2026, ithalat ağırlıklı): MFN \%25,7, İGV dahil \%25,7 \etiketV{}. S4'te AB menşeline (c2): \%0,0.
\end{itemize}

\textbf{6. Kabiliyet}

\begin{itemize}
\tightlist
\item
  KS 50,3 (sıra 111/230; ağırlık duyarlılığında P10--P90 61--169). Bileşenler: RCA 38, PCI 73 (PCI 0,42), yurt içi KD 68, Ar-Ge 25, tedarikçi 84, verimlilik 29, firma 29 \etiketT{}.
\item
  Ürün uzayında atlanabilir ürünler (sınıf içi; yoğunluk üst çeyrek, PCI \textgreater{} sınıf): yok \etiketT{}.
\end{itemize}

\textbf{7. Stratejik önem.} SS 22,4 (en yüksek boyut: gıda güvenliği 30; ağırlıklı ortalama 4,5) \etiketH{}.

\textbf{8. Makro duyarlılık}

\begingroup\gbtablo\small

\begin{longtable}[]{@{}lr@{}}
\toprule\noalign{}
\textbf{Gösterge} & \textbf{Değer (\etiketT{})} \\
\midrule\noalign{}
\endhead
\bottomrule\noalign{}
\endlastfoot
TÜFE ağırlığı (doğrudan, COICOP 2018 eşlemesi); tüketimde ithal pay & \%0,22; \%4 \\
ÜFE vekili (imalat üretimindeki pay) & \%0,06 \\
EDE: 1 puan tarife, tüm menşe / yalnız AB (bp; ρ = 1) & 0,009 / 0,008 \\
Geçişkenlik ρ (senaryolarda) & 0,6--1,0 (P50 0,8) \\
Piyasa yoğunlaşması (250+ istihdam payı) & \%38 \\
CKE: tarife (\$ mn / bp) ve mali (\$ mn / mlr TL); 0--100 puan & 835,3 ve 9,83; 33 \\
\end{longtable}

\endgroup

\textbf{9. Senaryo etkileri, T+5, P50 (P10 / P90)}

\begingroup\gbtablo\small

\begin{longtable}[]{@{}lrrr@{}}
\toprule\noalign{}
\textbf{Değişken} & \textbf{S2} & \textbf{S3} & \textbf{S4} \\
\midrule\noalign{}
\endhead
\bottomrule\noalign{}
\endlastfoot
AB'ye ihracat (\$ mn) & 0 (0 / 0) & 0 (0 / 0) & 0 (−2 / 0) \\
Toplam ihracat (\$ mn) & 0 (0 / 0) & 0 (0 / 0) & 0 (−1 / 0) \\
Nihai mal ithalatı (\$ mn) & 0 (0 / +1) & 0 (0 / +1) & 0 (0 / 0) \\
Cari denge (\$ mn) & 0 (−1 / 0) & 0 (−1 / 0) & 0 (−1 / 0) \\
Üretim (\$ mn) & 0 (−1 / 0) & 0 (−1 / 0) & 0 (−1 / 0) \\
İstihdam (kişi) & 0 (−5 / +1) & 0 (−3 / +1) & 0 (−5 / 0) \\
TÜFE (bp, birinci tur) & 0,00 (0,00 / 0,00) & 0,00 (0,00 / 0,00) & 0,00 (0,00 / 0,00) \\
Yatırım ve DYY (bütün senaryolar) & --- & & \\
\end{longtable}

\endgroup

\emph{Kaynak:} \etiketT{} f3e: kısmi denge, toplamlar Faz 2B GE dağılımına kalibre; S4'te 1/98 kaybı olasılığı 0,25.

\textbf{13. Kazanan ve kaybedenler; bölgesel ve sosyal etki; telafi.} Etkilenen kesim yok.

\textbf{15. Güven düzeyi; veri boşlukları; \etiketD{}.} Yüksek.

\hypertarget{nace_kartlar_10-12-36-1032-sebze-ve-meyve-suyu-imalatux131--gb-statuxfcsuxfc-tarim-198--birincilikincil-arketip-ba--duruux15f-geliux15ftir}{%
\subsection{3.6 10.32 Sebze ve meyve suyu imalatı --- GB statüsü: TARIM-1/98 \textbar{} Birincil/ikincil arketip: B/A \textbar{} Duruş: Geliştir}\label{nace_kartlar_10-12-36-1032-sebze-ve-meyve-suyu-imalatux131--gb-statuxfcsuxfc-tarim-198--birincilikincil-arketip-ba--duruux15f-geliux15ftir}}

\emph{Elma suyu konsantresinde dünya ikincisi (\%16). ABD ana pazar (\%67); AB'ye ihracatın büyük kısmı Hollanda'daki karıştırıcılara gidiyor. Tarım rejiminde; risk 1/98 kaybı ve ABD'nin ticaret savunması.}

\textbf{1. Kimlik.} Rev. 2.1: 10.32 (1:1); CPA 10.32. Başlıca HS4 (TR ihracatı 2024): 2009 \%100. BEC payları (ticaret ağırlıklı): ara \%0 / sermaye \%0 / tüketim \%100.

\textbf{2. Ölçek}

\begingroup\gbtablo\small

\begin{longtable}[]{@{}
  >{\raggedright\arraybackslash}p{(\columnwidth - 4\tabcolsep) * \real{0.3105}}
  >{\raggedright\arraybackslash}p{(\columnwidth - 4\tabcolsep) * \real{0.3372}}
  >{\raggedright\arraybackslash}p{(\columnwidth - 4\tabcolsep) * \real{0.3523}}@{}}
\toprule\noalign{}
\begin{minipage}[b]{\linewidth}\raggedright
\textbf{Gösterge}
\end{minipage} & \begin{minipage}[b]{\linewidth}\raggedright
\textbf{Değer}
\end{minipage} & \begin{minipage}[b]{\linewidth}\raggedright
\textbf{Etiket ve kaynak}
\end{minipage} \\
\midrule\noalign{}
\endhead
\bottomrule\noalign{}
\endlastfoot
Üretim değeri; katma değer & \$1.980 mn; \$415 mn (GSYH'nin \%0,03) & TÜİK YSHİ 2024 \etiketV{} \\
Çalışan; girişim (0 / 1--9 / 10--49 / 50--249 / 250+) & 7.460; 322 (19 / 264 / 42 / 13 / 9) & TÜİK YSHİ ve İş Kayıtları 2024 \etiketV{}; gizli büyüklük hücreleri \etiketT{} \\
250+ girişimlerin istihdam payı; AB'ye göre kişi başı KD & \%58; 0,66 & \etiketT{} (AB27 SBS 2023, nominal kur) \\
İlk 3 İBBS-2 bölgesi & Adana, Mersin \%16; İstanbul \%14; Bursa, Eskişehir, Bilecik \%13 & \etiketT{} iş kayıtları + SGK 2025, RAS \\
Yabancı sermaye; çapa yatırımcılar & Döhler, Konfrut, Elmasu, Aroma, Dimes, Tat \etiketH{} & \etiketH{} (firma bilgisi; pay verisi yok: VB-30) \\
\end{longtable}

\endgroup

\textbf{3. Dış ticaret}

\begingroup\gbtablo\small

\begin{longtable}[]{@{}
  >{\raggedright\arraybackslash}p{(\columnwidth - 2\tabcolsep) * \real{0.4766}}
  >{\raggedright\arraybackslash}p{(\columnwidth - 2\tabcolsep) * \real{0.5234}}@{}}
\toprule\noalign{}
\begin{minipage}[b]{\linewidth}\raggedright
\textbf{Gösterge}
\end{minipage} & \begin{minipage}[b]{\linewidth}\raggedright
\textbf{Değer (BACI 2024, \$; \etiketV{}, eşleme \etiketT{})}
\end{minipage} \\
\midrule\noalign{}
\endhead
\bottomrule\noalign{}
\endlastfoot
İhracat (2024); 2020--24 değişimi & \$770 mn; \%96 \\
AB payı; ilk 5 AB pazarı & \%27; NLD \%40; DEU \%12; ITA \%12; POL \%7; ESP \%7 \\
İlk 5 AB dışı pazar & USA \%67; GBR \%5; SAU \%4; JPN \%4; RUS \%3 \\
İthalat; AB ve Çin payı; ilk kaynaklar & \$144 mn; AB \%6, Çin \%47; CHN \%47; IRN \%24; CYP \%11; IRL \%2; PAK \%2 \\
Dış ticaret dengesi & \$+625 mn \\
İthal girdi içeriği (I-O, toplam) / KD÷üretim & \%20 / \%21 \\
RCA dünya (2020 → 2024) / AB pazarında & 2,55 → 3,21 / 1,43 \\
Birim değer ÷ AB pazarı ortalaması & 1,19 \\
AB pazar payı (AB dışı tedarik içinde, 2020 → 2024); TR'nin sırası & \%4,8 → \%5,7; 2. \\
\end{longtable}

\endgroup

\textbf{4. AB tarafı maruziyet}

\begin{itemize}
\tightlist
\item
  \kod{τ\_\allowbreak{}AB} (S4 artışı, rota A, TR→AB 2025 ağırlıklı): \%0,0; 1/98 kaybıyla \%24,0; tam MFN ortalaması \%24,0. Rejim: GB \%0, İTÜ \%0, tarım \%100 \etiketV{}/\etiketT{}.
\item
  Tepe satırlar (S4 artışı ≥ \%10; en büyük 3): yok.
\item
  AB pazarında AB dışı rakipler (tercih durumu): BRA \%48; UKR \%5; ZAF \%4; EGY \%4; CRI \%3. Tercihli rakip payı \%35.
\item
  Menşe (S2 / S3): risk D/D; \kod{c\_\allowbreak{}RoO} \%1,5 / \%2,0; \kod{p\_\allowbreak{}uyum} 0,90 / 0,85 \etiketH{}. Kural: 2009: aynı başlık dışı; meyve tamamen elde edilmiş.
\item
  EUDR kapsamındaki AB ihracatı payı: \%0; SKDM: \%0.
\item
  1/98 kaybında (P10) meyve sularında AB MFN'si \%24'e çıkar (giriş fiyatı sistemi dahil değil): AB'ye ihracat −\$74 mn.
\item
  Değer zinciri: AB içecek sanayisine ara girdi (konsantre).
\end{itemize}

\textbf{5. Türkiye tarafı maruziyet}

\begin{itemize}
\tightlist
\item
  Uygulanan tarife (2026, ithalat ağırlıklı): MFN \%58,5, İGV dahil \%58,5 \etiketV{}. S4'te AB menşeline (c2): \%0,0.
\item
  Kritik girdiler (HS4 2009), Türkiye'nin toplam ithalatı 2024: \$144 mn (AB \%18; CHN \%47, IRN \%24, CYP \%11) \etiketV{}.
\end{itemize}

\textbf{6. Kabiliyet}

\begin{itemize}
\tightlist
\item
  KS 51,6 (sıra 95/230; ağırlık duyarlılığında P10--P90 35--160). Bileşenler: RCA 84, PCI 16 (PCI −0,84), yurt içi KD 54, Ar-Ge 25, tedarikçi 84, verimlilik 83, firma 31 \etiketT{}.
\item
  Ürün uzayında atlanabilir ürünler (sınıf içi; yoğunluk üst çeyrek, PCI \textgreater{} sınıf): yok \etiketT{}.
\item
  \textbf{Kırmızı çizgi yetkinlikleri:} Konsantre üretiminde patulin ve kalıntı kontrolü (ABD ve AB eşikleri).
\item
  \textbf{Geliştirme hedefi:} T+5: konsantre yerine aseptik nihai ürün (NFC) payını artırmak.
\end{itemize}

\textbf{7. Stratejik önem.} SS 14,9 (en yüksek boyut: gıda güvenliği 20; ağırlıklı ortalama 3,0) \etiketH{}.

\textbf{8. Makro duyarlılık}

\begingroup\gbtablo\small

\begin{longtable}[]{@{}lr@{}}
\toprule\noalign{}
\textbf{Gösterge} & \textbf{Değer (\etiketT{})} \\
\midrule\noalign{}
\endhead
\bottomrule\noalign{}
\endlastfoot
TÜFE ağırlığı (doğrudan, COICOP 2018 eşlemesi); tüketimde ithal pay & \%0,18; \%3 \\
ÜFE vekili (imalat üretimindeki pay) & \%0,30 \\
EDE: 1 puan tarife, tüm menşe / yalnız AB (bp; ρ = 1) & 0,017 / 0,001 \\
Geçişkenlik ρ (senaryolarda) & 0,6--1,0 (P50 0,8) \\
Piyasa yoğunlaşması (250+ istihdam payı) & \%58 \\
CKE: tarife (\$ mn / bp) ve mali (\$ mn / mlr TL); 0--100 puan & 1.256,0 ve 6,32; 32 \\
\end{longtable}

\endgroup

\textbf{9. Senaryo etkileri, T+5, P50 (P10 / P90)}

\begingroup\gbtablo\small

\begin{longtable}[]{@{}
  >{\raggedright\arraybackslash}p{(\columnwidth - 6\tabcolsep) * \real{0.2908}}
  >{\raggedleft\arraybackslash}p{(\columnwidth - 6\tabcolsep) * \real{0.3919}}
  >{\raggedleft\arraybackslash}p{(\columnwidth - 6\tabcolsep) * \real{0.1586}}
  >{\raggedleft\arraybackslash}p{(\columnwidth - 6\tabcolsep) * \real{0.1586}}@{}}
\toprule\noalign{}
\begin{minipage}[b]{\linewidth}\raggedright
\textbf{Değişken}
\end{minipage} & \begin{minipage}[b]{\linewidth}\raggedleft
\textbf{S2}
\end{minipage} & \begin{minipage}[b]{\linewidth}\raggedleft
\textbf{S3}
\end{minipage} & \begin{minipage}[b]{\linewidth}\raggedleft
\textbf{S4}
\end{minipage} \\
\midrule\noalign{}
\endhead
\bottomrule\noalign{}
\endlastfoot
AB'ye ihracat (\$ mn) & 0 (0 / +4) & 0 (0 / 0) & 0 (−74 / 0) \\
Toplam ihracat (\$ mn) & 0 (0 / +3) & 0 (0 / 0) & 0 (−60 / 0) \\
Nihai mal ithalatı (\$ mn) & 0 (0 / +1) & 0 (0 / 0) & 0 (0 / 0) \\
Cari denge (\$ mn) & 0 (0 / +2) & 0 (0 / 0) & 0 (−48 / 0) \\
Üretim (\$ mn) & 0 (0 / +2) & 0 (0 / 0) & 0 (−60 / 0) \\
İstihdam (kişi) & 0 (0 / +6) & 0 (0 / 0) & 0 (−156 / 0) \\
TÜFE (bp, birinci tur) & 0,00 (0,00 / 0,00) & 0,00 (0,00 / 0,00) & 0,00 (0,00 / 0,00) \\
Yatırım ve DYY (bütün senaryolar) & Döhler (Almanya) yatırımları 1/98'e duyarlı \etiketH{} & & \\
\end{longtable}

\endgroup

\emph{Kaynak:} \etiketT{} f3e: kısmi denge, toplamlar Faz 2B GE dağılımına kalibre; S4'te 1/98 kaybı olasılığı 0,25.

\textbf{12. Erken uyarı ve tetikler}

\begingroup\gbtablo\small

\begin{longtable}[]{@{}
  >{\raggedright\arraybackslash}p{(\columnwidth - 4\tabcolsep) * \real{0.3405}}
  >{\raggedright\arraybackslash}p{(\columnwidth - 4\tabcolsep) * \real{0.1160}}
  >{\raggedright\arraybackslash}p{(\columnwidth - 4\tabcolsep) * \real{0.5435}}@{}}
\toprule\noalign{}
\begin{minipage}[b]{\linewidth}\raggedright
\textbf{Gösterge}
\end{minipage} & \begin{minipage}[b]{\linewidth}\raggedright
\textbf{Eşik}
\end{minipage} & \begin{minipage}[b]{\linewidth}\raggedright
\textbf{Tetiklenecek önlem}
\end{minipage} \\
\midrule\noalign{}
\endhead
\bottomrule\noalign{}
\endlastfoot
AB'ye ve ABD'ye konsantre ihracatı (aylık) & Yıllık −\%20 & E-10-05 benzeri kredi sigortası; ABD'de damping soruşturması hazırlığı \\
\end{longtable}

\endgroup

\textbf{13. Kazanan ve kaybedenler; bölgesel ve sosyal etki; telafi.} Kaybeden (yalnız 1/98 kaybında): Isparta, Karaman, Niğde elma bölgeleri.

\textbf{14. Riskler ve azaltım}

\begin{itemize}
\tightlist
\item
  1/98 askısı → E-10-01.
\item
  ABD'nin elma suyu ticaret savunması (geçmişte Çin'e uyguladı) \etiketH{}.
\end{itemize}

\textbf{15. Güven düzeyi; veri boşlukları; \etiketD{}.} Orta.

\hypertarget{nace_kartlar_10-12-37-1039-meyve-ve-sebzelerin-diux11fer-ux15fekillerde-iux15flenmesi-ve-korunmasux131--gb-statuxfcsuxfc-tarim-198--birincilikincil-arketip-a--duruux15f-koru}{%
\subsection{3.7 10.39 Meyve ve sebzelerin diğer şekillerde işlenmesi ve korunması --- GB statüsü: TARIM-1/98 \textbar{} Birincil/ikincil arketip: A \textbar{} Duruş: Koru}\label{nace_kartlar_10-12-37-1039-meyve-ve-sebzelerin-diux11fer-ux15fekillerde-iux15flenmesi-ve-korunmasux131--gb-statuxfcsuxfc-tarim-198--birincilikincil-arketip-a--duruux15f-koru}}

\emph{10--12'nin en büyük AB ihracatçısı: AB'ye \$3,6 mlr, AB payı \%56. Dünya birincisi olduğu ürünler: fındık içi (\%62), kuru kayısı (\%75), kuru incir (\%50), kuru üzüm (\%30). Tamamı 1/98'de. S4'ün P50 etkisi yok; P10'da (1/98 askısı) −\$1.141 mn ve −3.215 kişi.}

\textbf{1. Kimlik.} Rev. 2.1: 10.39 (1:1); CPA 10.39. Başlıca HS4 (TR ihracatı 2024): 0802 \%29, 2008 \%19, 0806 \%9, 0813 \%8. BEC payları (ticaret ağırlıklı): ara \%1 / sermaye \%0 / tüketim \%99. Kabuklu fındık ve kuru meyvenin bir kısmı A bölümünde (01.25); burada işlenmiş, kavrulmuş, ezme ve kuru meyve hazırlıkları var.

\textbf{2. Ölçek}

\begingroup\gbtablo\small

\begin{longtable}[]{@{}
  >{\raggedright\arraybackslash}p{(\columnwidth - 4\tabcolsep) * \real{0.2699}}
  >{\raggedright\arraybackslash}p{(\columnwidth - 4\tabcolsep) * \real{0.4304}}
  >{\raggedright\arraybackslash}p{(\columnwidth - 4\tabcolsep) * \real{0.2997}}@{}}
\toprule\noalign{}
\begin{minipage}[b]{\linewidth}\raggedright
\textbf{Gösterge}
\end{minipage} & \begin{minipage}[b]{\linewidth}\raggedright
\textbf{Değer}
\end{minipage} & \begin{minipage}[b]{\linewidth}\raggedright
\textbf{Etiket ve kaynak}
\end{minipage} \\
\midrule\noalign{}
\endhead
\bottomrule\noalign{}
\endlastfoot
Üretim değeri; katma değer & \$13.264 mn; \$2.174 mn (GSYH'nin \%0,16) & TÜİK YSHİ 2024 \etiketV{} \\
Çalışan; girişim (0 / 1--9 / 10--49 / 50--249 / 250+) & 66.354; 3.961 (426 / 3.333 / 426 / 221 / 49) & TÜİK YSHİ ve İş Kayıtları 2024 \etiketV{}; gizli büyüklük hücreleri \etiketT{} \\
250+ girişimlerin istihdam payı; AB'ye göre kişi başı KD & \%35; 0,51 & \etiketT{} (AB27 SBS 2023, nominal kur) \\
İlk 3 İBBS-2 bölgesi & Manisa, Afyon, Kütahya, Uşak \%16; İzmir \%14; Bursa, Eskişehir, Bilecik \%10 & \etiketT{} iş kayıtları + SGK 2025, RAS \\
Yabancı sermaye; çapa yatırımcılar & Ferrero (Oltan dahil), ofi (Olam), Balsu, Durak, Gürsoy; kuru meyvede Malatya ve Ege ihracatçıları \etiketH{} & \etiketH{} (firma bilgisi; pay verisi yok: VB-30) \\
\end{longtable}

\endgroup

\textbf{3. Dış ticaret}

\begingroup\gbtablo\small

\begin{longtable}[]{@{}
  >{\raggedright\arraybackslash}p{(\columnwidth - 2\tabcolsep) * \real{0.4731}}
  >{\raggedright\arraybackslash}p{(\columnwidth - 2\tabcolsep) * \real{0.5269}}@{}}
\toprule\noalign{}
\begin{minipage}[b]{\linewidth}\raggedright
\textbf{Gösterge}
\end{minipage} & \begin{minipage}[b]{\linewidth}\raggedright
\textbf{Değer (BACI 2024, \$; \etiketV{}, eşleme \etiketT{})}
\end{minipage} \\
\midrule\noalign{}
\endhead
\bottomrule\noalign{}
\endlastfoot
İhracat (2024); 2020--24 değişimi & \$6.472 mn; \%36 \\
AB payı; ilk 5 AB pazarı & \%56; DEU \%37; ITA \%17; FRA \%10; NLD \%8; POL \%7 \\
İlk 5 AB dışı pazar & GBR \%14; USA \%14; IRQ \%9; RUS \%5; CHE \%5 \\
İthalat; AB ve Çin payı; ilk kaynaklar & \$1.177 mn; AB \%8, Çin \%4; USA \%24; VNM \%12; IRN \%12; ECU \%8; CHN \%4 \\
Dış ticaret dengesi & \$+5.296 mn \\
İthal girdi içeriği (I-O, toplam) / KD÷üretim & \%20 / \%16 \\
RCA dünya (2020 → 2024) / AB pazarında & 4,86 → 5,05 / 4,68 \\
Birim değer ÷ AB pazarı ortalaması & 1,23 \\
AB pazar payı (AB dışı tedarik içinde, 2020 → 2024); TR'nin sırası & \%15,8 → \%18,6; 1. \\
\end{longtable}

\endgroup

\textbf{4. AB tarafı maruziyet}

\begin{itemize}
\tightlist
\item
  \kod{τ\_\allowbreak{}AB} (S4 artışı, rota A, TR→AB 2025 ağırlıklı): \%0,0; 1/98 kaybıyla \%9,5; tam MFN ortalaması \%9,5. Rejim: GB \%0, İTÜ \%0, tarım \%100 \etiketV{}/\etiketT{}.
\item
  Tepe satırlar (S4 artışı ≥ \%10; en büyük 3): 1 satır, AB ihracatının \%0'i: KN 2008 11 10 \%12,8 (€1,0 mn).
\item
  AB pazarında AB dışı rakipler (tercih durumu): USA \%12; CHN \%9; ECU \%5; VNM \%5; ARG \%5. Tercihli rakip payı \%44.
\item
  Menşe (S2 / S3): risk D/D; \kod{c\_\allowbreak{}RoO} \%1,5 / \%2,0; \kod{p\_\allowbreak{}uyum} 0,90 / 0,85 \etiketH{}. Kural: Fasıl 8 ve 20: meyve/fındık tamamen elde edilmiş; şeker ağırlık sınırı (2008).
\item
  EUDR kapsamındaki AB ihracatı payı: \%0; SKDM: \%0.
\item
  1/98 kaybında tarife satıra göre değişiyor: fındık içinde \%3,2, kuru kayısıda \%5,6, kuru incirde \%7,7, işlenmiş meyve ve sebzede (2001--2008) \%12--20. Ağırlıklı ortalama \%9,5.
\item
  Sınır dışı engel (bugün de): AB'nin artırılmış resmî kontroller listesinde (Uygulama Tüzüğü 2019/1793) Türk kuru incir, fındık ve antep fıstığı aflatoksin nedeniyle yer alıyor \etiketV{} (güncel satırlar \etiketD{}).
\item
  Değer zinciri: fındıkta AB şekerleme sanayisine (Ferrero, Lindt, Nestlé) ara girdi. Alıcı çok yoğunlaşmış, ama bağımlılık karşılıklı: Türkiye'nin dünya ihracatındaki payı \%62.
\end{itemize}

\textbf{5. Türkiye tarafı maruziyet}

\begin{itemize}
\tightlist
\item
  Uygulanan tarife (2026, ithalat ağırlıklı): MFN \%25,2, İGV dahil \%25,2 \etiketV{}. S4'te AB menşeline (c2): \%0,3.
\item
  Kritik girdiler (HS4 0801, 0802, 0813), Türkiye'nin toplam ithalatı 2024: \$1.161 mn (AB \%3; USA \%44, VNM \%12, IRN \%12) \etiketV{}.
\item
  İthalat (\$1,2 mlr) badem, ceviz, kaju ve antep fıstığı: ABD, Vietnam, İran. Tarım rejiminde; GB'den bağımsız.
\end{itemize}

\textbf{6. Kabiliyet}

\begin{itemize}
\tightlist
\item
  KS 49,8 (sıra 119/230; ağırlık duyarlılığında P10--P90 48--170). Bileşenler: RCA 79, PCI 9 (PCI −1,03), yurt içi KD 49, Ar-Ge 25, tedarikçi 84, verimlilik 61, firma 67 \etiketT{}.
\item
  Ürün uzayında atlanabilir ürünler (sınıf içi; yoğunluk üst çeyrek, PCI \textgreater{} sınıf): yok \etiketT{}.
\item
  \textbf{Kırmızı çizgi yetkinlikleri:} Aflatoksin kontrolü ve izlenebilirlik (bahçe → depo → fabrika); fındık kırma ve kavurma kapasitesi; Malatya kayısı kurutma standardı.
\item
  \textbf{Geliştirme hedefi:} T+5: ham fındık içi yerine fındık ezmesi ve kavrulmuş üründe payı +10 puan (atlanabilir ürünler: 2008 19); 2035: AB'nin yanında Asya şekerleme sanayisi (Çin, Hindistan) pazarı.
\end{itemize}

\textbf{7. Stratejik önem.} SS 37,2 (en yüksek boyut: gıda güvenliği 50; ağırlıklı ortalama 7,5) \etiketH{}. Gıda güvenliğinden çok tarımsal gelir ve ihracat: 400 binden fazla fındık üreticisi hane \etiketD{}.

\textbf{8. Makro duyarlılık}

\begingroup\gbtablo\small

\begin{longtable}[]{@{}lr@{}}
\toprule\noalign{}
\textbf{Gösterge} & \textbf{Değer (\etiketT{})} \\
\midrule\noalign{}
\endhead
\bottomrule\noalign{}
\endlastfoot
TÜFE ağırlığı (doğrudan, COICOP 2018 eşlemesi); tüketimde ithal pay & \%0,88; \%3 \\
ÜFE vekili (imalat üretimindeki pay) & \%1,99 \\
EDE: 1 puan tarife, tüm menşe / yalnız AB (bp; ρ = 1) & 0,118 / 0,010 \\
Geçişkenlik ρ (senaryolarda) & 0,6--1,0 (P50 0,8) \\
Piyasa yoğunlaşması (250+ istihdam payı) & \%35 \\
CKE: tarife (\$ mn / bp) ve mali (\$ mn / mlr TL); 0--100 puan & 2.626,2 ve 16,51; 51 \\
\end{longtable}

\endgroup

\textbf{9. Senaryo etkileri, T+5, P50 (P10 / P90)}

\begingroup\gbtablo\small

\begin{longtable}[]{@{}
  >{\raggedright\arraybackslash}p{(\columnwidth - 6\tabcolsep) * \real{0.2406}}
  >{\raggedleft\arraybackslash}p{(\columnwidth - 6\tabcolsep) * \real{0.4670}}
  >{\raggedleft\arraybackslash}p{(\columnwidth - 6\tabcolsep) * \real{0.1462}}
  >{\raggedleft\arraybackslash}p{(\columnwidth - 6\tabcolsep) * \real{0.1462}}@{}}
\toprule\noalign{}
\begin{minipage}[b]{\linewidth}\raggedright
\textbf{Değişken}
\end{minipage} & \begin{minipage}[b]{\linewidth}\raggedleft
\textbf{S2}
\end{minipage} & \begin{minipage}[b]{\linewidth}\raggedleft
\textbf{S3}
\end{minipage} & \begin{minipage}[b]{\linewidth}\raggedleft
\textbf{S4}
\end{minipage} \\
\midrule\noalign{}
\endhead
\bottomrule\noalign{}
\endlastfoot
AB'ye ihracat (\$ mn) & −1 (−2 / +72) & −1 (−2 / 0) & −3 (−1.141 / −1) \\
Toplam ihracat (\$ mn) & −1 (−2 / +49) & −1 (−2 / 0) & −2 (−940 / −1) \\
Nihai mal ithalatı (\$ mn) & 0 (−1 / +6) & 0 (−1 / +4) & −1 (−2 / 0) \\
Cari denge (\$ mn) & 0 (−1 / +34) & 0 (−1 / 0) & −1 (−748 / −1) \\
Üretim (\$ mn) & 0 (−2 / +43) & 0 (−2 / 0) & −1 (−939 / −1) \\
İstihdam (kişi) & −2 (−6 / +160) & −2 (−6 / −1) & −4 (−3.215 / −2) \\
TÜFE (bp, birinci tur) & 0,00 (0,00 / 0,00) & 0,00 (0,00 / 0,00) & 0,00 (0,00 / 0,00) \\
Yatırım ve DYY (bütün senaryolar) & Ferrero'nun Türkiye'deki işleme yatırımı 1/98'e değil fındık arzına bağlı \etiketH{} & & \\
\end{longtable}

\endgroup

\emph{Kaynak:} \etiketT{} f3e: kısmi denge, toplamlar Faz 2B GE dağılımına kalibre; S4'te 1/98 kaybı olasılığı 0,25.

\textbf{10. Politika paketi}

\begingroup\gbtablo\scriptsize

\begin{longtable}[]{@{}
  >{\raggedright\arraybackslash}p{(\columnwidth - 16\tabcolsep) * \real{0.0983}}
  >{\raggedright\arraybackslash}p{(\columnwidth - 16\tabcolsep) * \real{0.0896}}
  >{\raggedright\arraybackslash}p{(\columnwidth - 16\tabcolsep) * \real{0.1580}}
  >{\raggedright\arraybackslash}p{(\columnwidth - 16\tabcolsep) * \real{0.1377}}
  >{\raggedright\arraybackslash}p{(\columnwidth - 16\tabcolsep) * \real{0.0964}}
  >{\raggedright\arraybackslash}p{(\columnwidth - 16\tabcolsep) * \real{0.0714}}
  >{\raggedright\arraybackslash}p{(\columnwidth - 16\tabcolsep) * \real{0.1478}}
  >{\raggedright\arraybackslash}p{(\columnwidth - 16\tabcolsep) * \real{0.1139}}
  >{\raggedright\arraybackslash}p{(\columnwidth - 16\tabcolsep) * \real{0.0870}}@{}}
\toprule\noalign{}
\begin{minipage}[b]{\linewidth}\raggedright
\textbf{Eylem ID}
\end{minipage} & \begin{minipage}[b]{\linewidth}\raggedright
\textbf{Araç}
\end{minipage} & \begin{minipage}[b]{\linewidth}\raggedright
\textbf{Parametre}
\end{minipage} & \begin{minipage}[b]{\linewidth}\raggedright
\textbf{Sorumlu}
\end{minipage} & \begin{minipage}[b]{\linewidth}\raggedright
\textbf{Hukuki araç}
\end{minipage} & \begin{minipage}[b]{\linewidth}\raggedright
\textbf{Bütçe}
\end{minipage} & \begin{minipage}[b]{\linewidth}\raggedright
\textbf{Başlangıç--bitiş; sunset}
\end{minipage} & \begin{minipage}[b]{\linewidth}\raggedright
\textbf{Performans koşulu}
\end{minipage} & \begin{minipage}[b]{\linewidth}\raggedright
\textbf{KPI}
\end{minipage} \\
\midrule\noalign{}
\endhead
\bottomrule\noalign{}
\endlastfoot
E-10.39-01 & Fındık arz güvenliği diyaloğu & AB şekerleme sanayisi ve Komisyonla çok yıllı arz ve fiyat şeffaflığı çerçevesi; 1/98'in sürmesiyle bağlantılı & Ticaret Bakanlığı, TOB, Karadeniz Fındık İhracatçıları Birliği & İdari mutabakat & İdari & T−18 → T; sunset: --- & AB'nin 1/98'i sürdürmesi & 1/98 kesintisiz \\
E-10.39-02 & SPS kontrol sıklığını azaltma & 2019/1793 listesindeki Türk ürünlerinde kontrol oranının düşürülmesi; ön ihracat kontrolü (PEC) ve laboratuvar akreditasyonu & TOB, TÜRKAK & AB ile SPS teknik protokolü & 0,03 mlr TL/yıl \etiketT{} & T−12 → T+24; sunset: T+24 & RASFF bildirimleri −\%30 & Kontrol oranı ≤ \%10 \\
\end{longtable}

\endgroup

\textbf{11. Müzakere pozisyonu}

\begin{itemize}
\tightlist
\item
  AB ile hücum: S2'de 1/98'in tam aktarımı; işlenmiş meyve ve sebzede giriş fiyatı sistemi istisnası ve kotaların genişletilmesi (E-10-06).
\item
  Üçüncü ülke: Çin ve Hindistan'da fındık ve kuru meyve tarifesi (Çin'de fındık MFN \%10 \etiketD{}); Körfez STA'sı.
\end{itemize}

\textbf{12. Erken uyarı ve tetikler}

\begingroup\gbtablo\small

\begin{longtable}[]{@{}
  >{\raggedright\arraybackslash}p{(\columnwidth - 4\tabcolsep) * \real{0.2590}}
  >{\raggedright\arraybackslash}p{(\columnwidth - 4\tabcolsep) * \real{0.2970}}
  >{\raggedright\arraybackslash}p{(\columnwidth - 4\tabcolsep) * \real{0.4440}}@{}}
\toprule\noalign{}
\begin{minipage}[b]{\linewidth}\raggedright
\textbf{Gösterge}
\end{minipage} & \begin{minipage}[b]{\linewidth}\raggedright
\textbf{Eşik}
\end{minipage} & \begin{minipage}[b]{\linewidth}\raggedright
\textbf{Tetiklenecek önlem}
\end{minipage} \\
\midrule\noalign{}
\endhead
\bottomrule\noalign{}
\endlastfoot
AB'nin 1/98 tavizlerini askıya alma girişimi & Komisyon önerisi ya da Uygulama Tüzüğü 654/2014 bildirimi & E-10-01: Ortaklık Konseyi toplantı talebi ve DTÖ istişaresi; fındık arz diyaloğu devreye girer \\
AB'ye fındık ve kuru meyve ihracatı (aylık, ton) & Yıllık −\%10 & E-10-05 genişletilir; Asya pazarı hızlandırılır \\
\end{longtable}

\endgroup

\textbf{13. Kazanan ve kaybedenler; bölgesel ve sosyal etki; telafi.} Kaybeden (yalnız 1/98 askısında): Karadeniz (Ordu, Giresun, Trabzon; fındık) ve Malatya, Manisa, Aydın (kuru meyve). En büyük üç bölgenin istihdam payı \%40. Telafi: tarımsal gelir desteği değil ihracat sigortası ve pazar çeşitlendirme; fındıkta TMO alımı son çare.

\textbf{14. Riskler ve azaltım}

\begin{itemize}
\tightlist
\item
  1/98 askısı (olasılık 0,25 varsayıldı) → E-10-01 ve fındık arz diyaloğu.
\item
  Aflatoksin kaynaklı AB kontrol artışı (çıkıştan bağımsız) → SPS eylemi.
\item
  Ferrero'nun alım payı (yaklaşık üçte bir \etiketD{}): fiyat gücü yoğunlaşması.
\end{itemize}

\textbf{15. Güven düzeyi; veri boşlukları; \etiketD{}.} Orta (1/98 satır içeriği VB-06; 1/98 askı olasılığı \etiketH{}). \etiketD{}: Ferrero'nun Türkiye fındık alım payı; 2019/1793'ün güncel Türkiye satırları.

\hypertarget{nace_kartlar_10-12-38-1041-sux131vux131-ve-katux131-yaux11f-imalatux131--gb-statuxfcsuxfc-tarim-198--birincilikincil-arketip-db--duruux15f-geliux15ftir}{%
\subsection{3.8 10.41 Sıvı ve katı yağ imalatı --- GB statüsü: TARIM-1/98 \textbar{} Birincil/ikincil arketip: D/B \textbar{} Duruş: Geliştir}\label{nace_kartlar_10-12-38-1041-sux131vux131-ve-katux131-yaux11f-imalatux131--gb-statuxfcsuxfc-tarim-198--birincilikincil-arketip-db--duruux15f-geliux15ftir}}

\emph{Rafine ayçiçek yağında dünya birincisi (\%19), ama ithal ham yağa ve tohuma bağımlı: Rusya ve Ukrayna \%42, palm yağında Malezya \%21. AB'ye ihracat \%9. Gıda güvenliği sınıfı; GB çıkışının etkisi sıfıra yakın.}

\textbf{1. Kimlik.} Rev. 2.1: 10.41 (1:1); CPA 10.41. Başlıca HS4 (TR ihracatı 2024): 1512 \%41, 1509 \%21, 2304 \%17, 1507 \%11. BEC payları (ticaret ağırlıklı): ara \%71 / sermaye \%0 / tüketim \%29.

\textbf{2. Ölçek}

\begingroup\gbtablo\small

\begin{longtable}[]{@{}
  >{\raggedright\arraybackslash}p{(\columnwidth - 4\tabcolsep) * \real{0.3004}}
  >{\raggedright\arraybackslash}p{(\columnwidth - 4\tabcolsep) * \real{0.3603}}
  >{\raggedright\arraybackslash}p{(\columnwidth - 4\tabcolsep) * \real{0.3393}}@{}}
\toprule\noalign{}
\begin{minipage}[b]{\linewidth}\raggedright
\textbf{Gösterge}
\end{minipage} & \begin{minipage}[b]{\linewidth}\raggedright
\textbf{Değer}
\end{minipage} & \begin{minipage}[b]{\linewidth}\raggedright
\textbf{Etiket ve kaynak}
\end{minipage} \\
\midrule\noalign{}
\endhead
\bottomrule\noalign{}
\endlastfoot
Üretim değeri; katma değer & \$7.265 mn; \$910 mn (GSYH'nin \%0,07) & TÜİK YSHİ 2024 \etiketV{} \\
Çalışan; girişim (0 / 1--9 / 10--49 / 50--249 / 250+) & 17.859; 1.938 (107 / 1.664 / 264 / 43 / 5) & TÜİK YSHİ ve İş Kayıtları 2024 \etiketV{}; gizli büyüklük hücreleri \etiketT{} \\
250+ girişimlerin istihdam payı; AB'ye göre kişi başı KD & \%14; 0,44 & \etiketT{} (AB27 SBS 2023, nominal kur) \\
İlk 3 İBBS-2 bölgesi & Manisa, Afyon, Kütahya, Uşak \%13; İzmir \%11; Balıkesir, Çanakkale \%10 & \etiketT{} iş kayıtları + SGK 2025, RAS \\
Yabancı sermaye; çapa yatırımcılar & Küçükbay, Yonca, Zade, Oruçoğlu, Savola (Suudi), Cargill ve Bunge \etiketH{} & \etiketH{} (firma bilgisi; pay verisi yok: VB-30) \\
\end{longtable}

\endgroup

\textbf{3. Dış ticaret}

\begingroup\gbtablo\small

\begin{longtable}[]{@{}
  >{\raggedright\arraybackslash}p{(\columnwidth - 2\tabcolsep) * \real{0.4696}}
  >{\raggedright\arraybackslash}p{(\columnwidth - 2\tabcolsep) * \real{0.5304}}@{}}
\toprule\noalign{}
\begin{minipage}[b]{\linewidth}\raggedright
\textbf{Gösterge}
\end{minipage} & \begin{minipage}[b]{\linewidth}\raggedright
\textbf{Değer (BACI 2024, \$; \etiketV{}, eşleme \etiketT{})}
\end{minipage} \\
\midrule\noalign{}
\endhead
\bottomrule\noalign{}
\endlastfoot
İhracat (2024); 2020--24 değişimi & \$3.183 mn; \%110 \\
AB payı; ilk 5 AB pazarı & \%9; ESP \%44; ITA \%23; EST \%10; DEU \%6; ROU \%4 \\
İlk 5 AB dışı pazar & IRQ \%15; USA \%11; IRN \%11; DJI \%8; DZA \%7 \\
İthalat; AB ve Çin payı; ilk kaynaklar & \$3.711 mn; AB \%2, Çin \%0; RUS \%30; MYS \%21; UKR \%12; ARG \%11; SYR \%6 \\
Dış ticaret dengesi & \$−527 mn \\
İthal girdi içeriği (I-O, toplam) / KD÷üretim & \%20 / \%13 \\
RCA dünya (2020 → 2024) / AB pazarında & 1,10 → 1,53 / 0,29 \\
Birim değer ÷ AB pazarı ortalaması & 1,07 \\
AB pazar payı (AB dışı tedarik içinde, 2020 → 2024); TR'nin sırası & \%0,3 → \%1,2; 15. \\
\end{longtable}

\endgroup

\textbf{4. AB tarafı maruziyet}

\begin{itemize}
\tightlist
\item
  \kod{τ\_\allowbreak{}AB} (S4 artışı, rota A, TR→AB 2025 ağırlıklı): \%0,0; 1/98 kaybıyla \%12,4; tam MFN ortalaması \%12,4. Rejim: GB \%0, İTÜ \%0, tarım \%100 \etiketV{}/\etiketT{}.
\item
  Tepe satırlar (S4 artışı ≥ \%10; en büyük 3): yok.
\item
  AB pazarında AB dışı rakipler (tercih durumu): BRA \%18; UKR \%17; IDN \%13; ARG \%12; MYS \%8. Tercihli rakip payı \%37.
\item
  Menşe (S2 / S3): risk D/O; \kod{c\_\allowbreak{}RoO} \%1,5 / \%3,0; \kod{p\_\allowbreak{}uyum} 0,90 / 0,80 \etiketH{}. Kural: 1507--1515: aynı alt başlık dışı malzemeden; ham yağ ithalatından rafine yağ menşe kazanır.
\item
  EUDR kapsamındaki AB ihracatı payı: \%43; SKDM: \%0.
\item
  AB'ye ihracat (\$283 mn) ağırlıkla zeytinyağı (İspanya \%40, İtalya \%21). 1/98 kaybında spesifik vergi AVE'si \%12'ye çıkar (zeytinyağında €124,5/100 kg).
\item
  EUDR: AB ihracatının \%43'ü palm ve soya türevi satırlarda; ham yağın kökenine ilişkin coğrafi konum verisi gerekiyor (E-10-03).
\end{itemize}

\textbf{5. Türkiye tarafı maruziyet}

\begin{itemize}
\tightlist
\item
  Uygulanan tarife (2026, ithalat ağırlıklı): MFN \%18,9, İGV dahil \%18,9 \etiketV{}. S4'te AB menşeline (c2): \%0,0.
\item
  Kritik girdiler (HS4 1206, 1512, 1511, 1507, 1201), Türkiye'nin toplam ithalatı 2024: \$3.755 mn (AB \%2; BRA \%27, RUS \%20, MYS \%17) \etiketV{}.
\item
  2022 şoku: Ukrayna savaşında ham ayçiçek yağı fiyatı iki katına çıktı. Yerli tohum üretimi tüketimin yaklaşık yarısını karşılıyor \etiketD{}.
\end{itemize}

\textbf{6. Kabiliyet}

\begin{itemize}
\tightlist
\item
  KS 45,7 (sıra 162/230; ağırlık duyarlılığında P10--P90 91--193). Bileşenler: RCA 67, PCI 7 (PCI −1,10), yurt içi KD 47, Ar-Ge 25, tedarikçi 84, verimlilik 45, firma 71 \etiketT{}.
\item
  Ürün uzayında atlanabilir ürünler (sınıf içi; yoğunluk üst çeyrek, PCI \textgreater{} sınıf): yok \etiketT{}.
\item
  \textbf{Kırmızı çizgi yetkinlikleri:} Rafineri ve ekstraksiyon kapasitesi; DİR'deki ihracat ağı (Irak, Körfez).
\item
  \textbf{Geliştirme hedefi:} T+5: yağlı tohum (ayçiçek, kanola) ekim alanı +\%15; 2035: ham yağ ithalatında tek kaynak payı ≤ \%30.
\end{itemize}

\textbf{7. Stratejik önem.} SS 63,3 (en yüksek boyut: gıda güvenliği 85; ağırlıklı ortalama 12,8) \etiketH{}. Bitkisel yağ temel tüketim ürünü. Ham yağ ve tohumda ithalat bağımlılığı yüksek, 2022'de fiyat şoku yaşandı.

\textbf{8. Makro duyarlılık}

\begingroup\gbtablo\small

\begin{longtable}[]{@{}lr@{}}
\toprule\noalign{}
\textbf{Gösterge} & \textbf{Değer (\etiketT{})} \\
\midrule\noalign{}
\endhead
\bottomrule\noalign{}
\endlastfoot
TÜFE ağırlığı (doğrudan, COICOP 2018 eşlemesi); tüketimde ithal pay & \%1,00; \%11 \\
ÜFE vekili (imalat üretimindeki pay) & \%1,09 \\
EDE: 1 puan tarife, tüm menşe / yalnız AB (bp; ρ = 1) & 0,395 / 0,009 \\
Geçişkenlik ρ (senaryolarda) & 0,6--1,0 (P50 0,8) \\
Piyasa yoğunlaşması (250+ istihdam payı) & \%14 \\
CKE: tarife (\$ mn / bp) ve mali (\$ mn / mlr TL); 0--100 puan & 1.502,4 ve 52,71; 63 \\
\end{longtable}

\endgroup

\textbf{9. Senaryo etkileri, T+5, P50 (P10 / P90)}

\begingroup\gbtablo\small

\begin{longtable}[]{@{}lrrr@{}}
\toprule\noalign{}
\textbf{Değişken} & \textbf{S2} & \textbf{S3} & \textbf{S4} \\
\midrule\noalign{}
\endhead
\bottomrule\noalign{}
\endlastfoot
AB'ye ihracat (\$ mn) & 0 (0 / +6) & 0 (0 / 0) & 0 (−102 / 0) \\
Toplam ihracat (\$ mn) & 0 (0 / +4) & 0 (0 / 0) & 0 (−82 / 0) \\
Nihai mal ithalatı (\$ mn) & 0 (0 / +3) & 0 (0 / +2) & 0 (0 / 0) \\
Cari denge (\$ mn) & 0 (−1 / +1) & 0 (−1 / 0) & 0 (−65 / 0) \\
Üretim (\$ mn) & 0 (−1 / +1) & 0 (−1 / 0) & 0 (−82 / 0) \\
İstihdam (kişi) & 0 (−2 / +1) & 0 (−2 / 0) & 0 (−137 / 0) \\
TÜFE (bp, birinci tur) & 0,00 (0,00 / 0,00) & 0,00 (0,00 / 0,00) & 0,00 (0,00 / 0,00) \\
Yatırım ve DYY (bütün senaryolar) & Etkisi yok & & \\
\end{longtable}

\endgroup

\emph{Kaynak:} \etiketT{} f3e: kısmi denge, toplamlar Faz 2B GE dağılımına kalibre; S4'te 1/98 kaybı olasılığı 0,25.

\textbf{10. Politika paketi}

\begingroup\gbtablo\scriptsize

\begin{longtable}[]{@{}
  >{\raggedright\arraybackslash}p{(\columnwidth - 16\tabcolsep) * \real{0.0983}}
  >{\raggedright\arraybackslash}p{(\columnwidth - 16\tabcolsep) * \real{0.0928}}
  >{\raggedright\arraybackslash}p{(\columnwidth - 16\tabcolsep) * \real{0.1176}}
  >{\raggedright\arraybackslash}p{(\columnwidth - 16\tabcolsep) * \real{0.0797}}
  >{\raggedright\arraybackslash}p{(\columnwidth - 16\tabcolsep) * \real{0.1396}}
  >{\raggedright\arraybackslash}p{(\columnwidth - 16\tabcolsep) * \real{0.0958}}
  >{\raggedright\arraybackslash}p{(\columnwidth - 16\tabcolsep) * \real{0.1591}}
  >{\raggedright\arraybackslash}p{(\columnwidth - 16\tabcolsep) * \real{0.1128}}
  >{\raggedright\arraybackslash}p{(\columnwidth - 16\tabcolsep) * \real{0.1043}}@{}}
\toprule\noalign{}
\begin{minipage}[b]{\linewidth}\raggedright
\textbf{Eylem ID}
\end{minipage} & \begin{minipage}[b]{\linewidth}\raggedright
\textbf{Araç}
\end{minipage} & \begin{minipage}[b]{\linewidth}\raggedright
\textbf{Parametre}
\end{minipage} & \begin{minipage}[b]{\linewidth}\raggedright
\textbf{Sorumlu}
\end{minipage} & \begin{minipage}[b]{\linewidth}\raggedright
\textbf{Hukuki araç}
\end{minipage} & \begin{minipage}[b]{\linewidth}\raggedright
\textbf{Bütçe}
\end{minipage} & \begin{minipage}[b]{\linewidth}\raggedright
\textbf{Başlangıç--bitiş; sunset}
\end{minipage} & \begin{minipage}[b]{\linewidth}\raggedright
\textbf{Performans koşulu}
\end{minipage} & \begin{minipage}[b]{\linewidth}\raggedright
\textbf{KPI}
\end{minipage} \\
\midrule\noalign{}
\endhead
\bottomrule\noalign{}
\endlastfoot
E-10.41-01 & Stratejik ham yağ stoku & 60 günlük tüketim; TMO ve özel lisanslı depo & TOB, TMO & Cumhurbaşkanı kararı (stok programı) & 0,15 mlr TL/yıl \etiketT{} (taşıma maliyeti) & T−12 → T+36; sunset: T+36 (gözden geçirme) & Stok ≥ 60 gün & Fiyat şokunda TÜFE yağ alt endeksi oynaklığı −\%30 \\
\end{longtable}

\endgroup

\textbf{12. Erken uyarı ve tetikler}

\begingroup\gbtablo\small

\begin{longtable}[]{@{}
  >{\raggedright\arraybackslash}p{(\columnwidth - 4\tabcolsep) * \real{0.4137}}
  >{\raggedright\arraybackslash}p{(\columnwidth - 4\tabcolsep) * \real{0.1142}}
  >{\raggedright\arraybackslash}p{(\columnwidth - 4\tabcolsep) * \real{0.4721}}@{}}
\toprule\noalign{}
\begin{minipage}[b]{\linewidth}\raggedright
\textbf{Gösterge}
\end{minipage} & \begin{minipage}[b]{\linewidth}\raggedright
\textbf{Eşik}
\end{minipage} & \begin{minipage}[b]{\linewidth}\raggedright
\textbf{Tetiklenecek önlem}
\end{minipage} \\
\midrule\noalign{}
\endhead
\bottomrule\noalign{}
\endlastfoot
Ham ayçiçek yağı ithalat fiyatı (aylık) & 3 ayda +\%25 & Stoktan satış; kaynak çeşitlendirme ihaleleri \\
\end{longtable}

\endgroup

\textbf{13. Kazanan ve kaybedenler; bölgesel ve sosyal etki; telafi.} GB çıkışından etkilenen yok.

\textbf{14. Riskler ve azaltım}

\begin{itemize}
\tightlist
\item
  Karadeniz arz kesintisi (GB'den bağımsız) → stok.
\end{itemize}

\textbf{15. Güven düzeyi; veri boşlukları; \etiketD{}.} Yüksek (veri), Orta (stok bütçesi).

\hypertarget{nace_kartlar_10-12-39-1042-margarin-ve-benzeri-yenilebilir-katux131-yaux11flarux131n-imalatux131--gb-statuxfcsuxfc-tarim-198--birincilikincil-arketip-b--duruux15f-geliux15ftir}{%
\subsection{3.9 10.42 Margarin ve benzeri yenilebilir katı yağların imalatı --- GB statüsü: TARIM-1/98 \textbar{} Birincil/ikincil arketip: B \textbar{} Duruş: Geliştir}\label{nace_kartlar_10-12-39-1042-margarin-ve-benzeri-yenilebilir-katux131-yaux11flarux131n-imalatux131--gb-statuxfcsuxfc-tarim-198--birincilikincil-arketip-b--duruux15f-geliux15ftir}}

\emph{Küçük sınıf (margarin): ihracat \$165 mn, Irak ve Suriye ağırlıklı. GB etkisi yok; palm yağına bağımlı (Endonezya, Malezya).}

\textbf{1. Kimlik.} Rev. 2.1: 10.42 (1:1); CPA 10.42. Başlıca HS4 (TR ihracatı 2024): 1517 \%100. BEC payları (ticaret ağırlıklı): ara \%0 / sermaye \%0 / tüketim \%100.

\textbf{2. Ölçek}

\begingroup\gbtablo\small

\begin{longtable}[]{@{}
  >{\raggedright\arraybackslash}p{(\columnwidth - 4\tabcolsep) * \real{0.3362}}
  >{\raggedright\arraybackslash}p{(\columnwidth - 4\tabcolsep) * \real{0.2781}}
  >{\raggedright\arraybackslash}p{(\columnwidth - 4\tabcolsep) * \real{0.3857}}@{}}
\toprule\noalign{}
\begin{minipage}[b]{\linewidth}\raggedright
\textbf{Gösterge}
\end{minipage} & \begin{minipage}[b]{\linewidth}\raggedright
\textbf{Değer}
\end{minipage} & \begin{minipage}[b]{\linewidth}\raggedright
\textbf{Etiket ve kaynak}
\end{minipage} \\
\midrule\noalign{}
\endhead
\bottomrule\noalign{}
\endlastfoot
Üretim değeri; katma değer & \$380 mn; \$71 mn (GSYH'nin \%0,01) & TÜİK YSHİ 2024 \etiketV{} \\
Çalışan; girişim (0 / 1--9 / 10--49 / 50--249 / 250+) & 773; 13 (1 / 7 / 2 / 3 / 1) & TÜİK YSHİ ve İş Kayıtları 2024 \etiketV{}; gizli büyüklük hücreleri \etiketT{} \\
250+ girişimlerin istihdam payı; AB'ye göre kişi başı KD & \%49; 0,43 & \etiketT{} (AB27 SBS 2023, nominal kur) \\
İlk 3 İBBS-2 bölgesi & İstanbul \%55; İzmir \%19; Adana, Mersin \%16 & \etiketT{} iş kayıtları + SGK 2025, RAS \\
Yabancı sermaye; çapa yatırımcılar & Unilever (Sana), Yonca, Bunge \etiketH{} & \etiketH{} (firma bilgisi; pay verisi yok: VB-30) \\
\end{longtable}

\endgroup

\textbf{3. Dış ticaret}

\begingroup\gbtablo\small

\begin{longtable}[]{@{}
  >{\raggedright\arraybackslash}p{(\columnwidth - 2\tabcolsep) * \real{0.4733}}
  >{\raggedright\arraybackslash}p{(\columnwidth - 2\tabcolsep) * \real{0.5267}}@{}}
\toprule\noalign{}
\begin{minipage}[b]{\linewidth}\raggedright
\textbf{Gösterge}
\end{minipage} & \begin{minipage}[b]{\linewidth}\raggedright
\textbf{Değer (BACI 2024, \$; \etiketV{}, eşleme \etiketT{})}
\end{minipage} \\
\midrule\noalign{}
\endhead
\bottomrule\noalign{}
\endlastfoot
İhracat (2024); 2020--24 değişimi & \$165 mn; \%69 \\
AB payı; ilk 5 AB pazarı & \%7; ROU \%62; DEU \%15; BGR \%8; NLD \%3; GRC \%3 \\
İlk 5 AB dışı pazar & IRQ \%46; SYR \%20; GEO \%6; IRN \%4; LBY \%2 \\
İthalat; AB ve Çin payı; ilk kaynaklar & \$79 mn; AB \%39, Çin \%0; IDN \%35; MYS \%16; NLD \%15; SWE \%14; IND \%6 \\
Dış ticaret dengesi & \$+87 mn \\
İthal girdi içeriği (I-O, toplam) / KD÷üretim & \%20 / \%19 \\
RCA dünya (2020 → 2024) / AB pazarında & 1,61 → 1,62 / 0,78 \\
Birim değer ÷ AB pazarı ortalaması & 0,72 \\
AB pazar payı (AB dışı tedarik içinde, 2020 → 2024); TR'nin sırası & \%2,7 → \%3,1; 10. \\
\end{longtable}

\endgroup

\textbf{4. AB tarafı maruziyet}

\begin{itemize}
\tightlist
\item
  \kod{τ\_\allowbreak{}AB} (S4 artışı, rota A, TR→AB 2025 ağırlıklı): \%0,0; 1/98 kaybıyla \%14,9; tam MFN ortalaması \%14,9. Rejim: GB \%0, İTÜ \%0, tarım \%100 \etiketV{}/\etiketT{}.
\item
  Tepe satırlar (S4 artışı ≥ \%10; en büyük 3): yok.
\item
  AB pazarında AB dışı rakipler (tercih durumu): USA \%21; GBR \%20; IDN \%14; CHE \%10; NOR \%8. Tercihli rakip payı \%48.
\item
  Menşe (S2 / S3): risk O/O; \kod{c\_\allowbreak{}RoO} \%2,5 / \%3,5; \kod{p\_\allowbreak{}uyum} 0,85 / 0,75 \etiketH{}. Kural: 1517: Fasıl 15 ve 4 malzemeleri ağırlık sınırı.
\item
  EUDR kapsamındaki AB ihracatı payı: \%0; SKDM: \%0.
\end{itemize}

\textbf{5. Türkiye tarafı maruziyet}

\begin{itemize}
\tightlist
\item
  Uygulanan tarife (2026, ithalat ağırlıklı): MFN \%22,5, İGV dahil \%22,5 \etiketV{}. S4'te AB menşeline (c2): \%0,0.
\item
  Kritik girdiler (HS4 1511, 1513), Türkiye'nin toplam ithalatı 2024: \$825 mn (AB \%1; MYS \%87, IDN \%10, PHL \%2) \etiketV{}.
\end{itemize}

\textbf{6. Kabiliyet}

\begin{itemize}
\tightlist
\item
  KS 41,8 (sıra 196/230; ağırlık duyarlılığında P10--P90 145--209). Bileşenler: RCA 57, PCI 22 (PCI −0,58), yurt içi KD 51, Ar-Ge 25, tedarikçi 84, verimlilik 42, firma 18 \etiketT{}.
\item
  Ürün uzayında atlanabilir ürünler (sınıf içi; yoğunluk üst çeyrek, PCI \textgreater{} sınıf): yok \etiketT{}.
\end{itemize}

\textbf{7. Stratejik önem.} SS 29,8 (en yüksek boyut: gıda güvenliği 40; ağırlıklı ortalama 6,0) \etiketH{}.

\textbf{8. Makro duyarlılık}

\begingroup\gbtablo\small

\begin{longtable}[]{@{}lr@{}}
\toprule\noalign{}
\textbf{Gösterge} & \textbf{Değer (\etiketT{})} \\
\midrule\noalign{}
\endhead
\bottomrule\noalign{}
\endlastfoot
TÜFE ağırlığı (doğrudan, COICOP 2018 eşlemesi); tüketimde ithal pay & \%0,07; \%7 \\
ÜFE vekili (imalat üretimindeki pay) & \%0,06 \\
EDE: 1 puan tarife, tüm menşe / yalnız AB (bp; ρ = 1) & 0,011 / 0,004 \\
Geçişkenlik ρ (senaryolarda) & 0,6--1,0 (P50 0,8) \\
Piyasa yoğunlaşması (250+ istihdam payı) & \%49 \\
CKE: tarife (\$ mn / bp) ve mali (\$ mn / mlr TL); 0--100 puan & 419,2 ve 6,76; 26 \\
\end{longtable}

\endgroup

\textbf{9. Senaryo etkileri, T+5, P50 (P10 / P90)}

\begingroup\gbtablo\small

\begin{longtable}[]{@{}lrrr@{}}
\toprule\noalign{}
\textbf{Değişken} & \textbf{S2} & \textbf{S3} & \textbf{S4} \\
\midrule\noalign{}
\endhead
\bottomrule\noalign{}
\endlastfoot
AB'ye ihracat (\$ mn) & 0 (0 / 0) & 0 (0 / 0) & 0 (−2 / 0) \\
Toplam ihracat (\$ mn) & 0 (0 / 0) & 0 (0 / 0) & 0 (−1 / 0) \\
Nihai mal ithalatı (\$ mn) & 0 (0 / +2) & 0 (0 / +1) & 0 (0 / 0) \\
Cari denge (\$ mn) & 0 (−1 / 0) & 0 (−1 / 0) & 0 (−1 / 0) \\
Üretim (\$ mn) & 0 (−2 / 0) & 0 (−1 / 0) & 0 (−1 / 0) \\
İstihdam (kişi) & 0 (−2 / 0) & 0 (−1 / 0) & 0 (−2 / 0) \\
TÜFE (bp, birinci tur) & 0,00 (0,00 / 0,00) & 0,00 (0,00 / 0,00) & 0,00 (0,00 / 0,00) \\
Yatırım ve DYY (bütün senaryolar) & --- & & \\
\end{longtable}

\endgroup

\emph{Kaynak:} \etiketT{} f3e: kısmi denge, toplamlar Faz 2B GE dağılımına kalibre; S4'te 1/98 kaybı olasılığı 0,25.

\textbf{13. Kazanan ve kaybedenler; bölgesel ve sosyal etki; telafi.} Etkilenen kesim yok.

\textbf{15. Güven düzeyi; veri boşlukları; \etiketD{}.} Yüksek.

\hypertarget{nace_kartlar_10-12-310-1051-suxfcthane-iux15fletmelerinin-faaliyetleri-ve-peynir-imalatux131--gb-statuxfcsuxfc-tarim-198--birincilikincil-arketip---duruux15f-koru}{%
\subsection{3.10 10.51 Süthane işletmelerinin faaliyetleri ve peynir imalatı --- GB statüsü: TARIM-1/98 \textbar{} Birincil/ikincil arketip: --- \textbar{} Duruş: Koru}\label{nace_kartlar_10-12-310-1051-suxfcthane-iux15fletmelerinin-faaliyetleri-ve-peynir-imalatux131--gb-statuxfcsuxfc-tarim-198--birincilikincil-arketip---duruux15f-koru}}

\emph{İç pazara dönük büyük işveren (49 bin kişi). AB ihracatı \$26 mn, Almanya'daki diaspora pazarı. Tarım rejiminde; GB etkisi yok.}

\textbf{1. Kimlik.} Rev. 2.1: 10.51 (1:1); CPA 10.51. Başlıca HS4 (TR ihracatı 2024): 0406 \%52, 0402 \%16, 0404 \%15, 0403 \%7. BEC payları (ticaret ağırlıklı): ara \%28 / sermaye \%0 / tüketim \%72.

\textbf{2. Ölçek}

\begingroup\gbtablo\small

\begin{longtable}[]{@{}
  >{\raggedright\arraybackslash}p{(\columnwidth - 4\tabcolsep) * \real{0.2886}}
  >{\raggedright\arraybackslash}p{(\columnwidth - 4\tabcolsep) * \real{0.3876}}
  >{\raggedright\arraybackslash}p{(\columnwidth - 4\tabcolsep) * \real{0.3239}}@{}}
\toprule\noalign{}
\begin{minipage}[b]{\linewidth}\raggedright
\textbf{Gösterge}
\end{minipage} & \begin{minipage}[b]{\linewidth}\raggedright
\textbf{Değer}
\end{minipage} & \begin{minipage}[b]{\linewidth}\raggedright
\textbf{Etiket ve kaynak}
\end{minipage} \\
\midrule\noalign{}
\endhead
\bottomrule\noalign{}
\endlastfoot
Üretim değeri; katma değer & \$10.843 mn; \$1.987 mn (GSYH'nin \%0,15) & TÜİK YSHİ 2024 \etiketV{} \\
Çalışan; girişim (0 / 1--9 / 10--49 / 50--249 / 250+) & 48.749; 2.293 (180 / 1.802 / 374 / 89 / 37) & TÜİK YSHİ ve İş Kayıtları 2024 \etiketV{}; gizli büyüklük hücreleri \etiketT{} \\
250+ girişimlerin istihdam payı; AB'ye göre kişi başı KD & \%52; 0,48 & \etiketT{} (AB27 SBS 2023, nominal kur) \\
İlk 3 İBBS-2 bölgesi & Kocaeli, Sakarya, Düzce, Bolu, Yalova \%15; İzmir \%14; Bursa, Eskişehir, Bilecik \%13 & \etiketT{} iş kayıtları + SGK 2025, RAS \\
Yabancı sermaye; çapa yatırımcılar & Sütaş, Pınar, Danone, Ak Gıda (İçim), SEK, Torku \etiketH{} & \etiketH{} (firma bilgisi; pay verisi yok: VB-30) \\
\end{longtable}

\endgroup

\textbf{3. Dış ticaret}

\begingroup\gbtablo\small

\begin{longtable}[]{@{}
  >{\raggedright\arraybackslash}p{(\columnwidth - 2\tabcolsep) * \real{0.4802}}
  >{\raggedright\arraybackslash}p{(\columnwidth - 2\tabcolsep) * \real{0.5198}}@{}}
\toprule\noalign{}
\begin{minipage}[b]{\linewidth}\raggedright
\textbf{Gösterge}
\end{minipage} & \begin{minipage}[b]{\linewidth}\raggedright
\textbf{Değer (BACI 2024, \$; \etiketV{}, eşleme \etiketT{})}
\end{minipage} \\
\midrule\noalign{}
\endhead
\bottomrule\noalign{}
\endlastfoot
İhracat (2024); 2020--24 değişimi & \$440 mn; \%21 \\
AB payı; ilk 5 AB pazarı & \%6; DEU \%46; FRA \%18; BGR \%9; NLD \%6; CZE \%4 \\
İlk 5 AB dışı pazar & IRQ \%20; EGY \%8; DZA \%7; SAU \%6; ARE \%5 \\
İthalat; AB ve Çin payı; ilk kaynaklar & \$112 mn; AB \%46, Çin \%0; CYP \%24; NLD \%10; DEU \%7; IRL \%7; ITA \%6 \\
Dış ticaret dengesi & \$+328 mn \\
İthal girdi içeriği (I-O, toplam) / KD÷üretim & \%20 / \%18 \\
RCA dünya (2020 → 2024) / AB pazarında & 0,40 → 0,35 / 0,25 \\
Birim değer ÷ AB pazarı ortalaması & 1,22 \\
AB pazar payı (AB dışı tedarik içinde, 2020 → 2024); TR'nin sırası & \%0,3 → \%1,0; 7. \\
\end{longtable}

\endgroup

\textbf{4. AB tarafı maruziyet}

\begin{itemize}
\tightlist
\item
  \kod{τ\_\allowbreak{}AB} (S4 artışı, rota A, TR→AB 2025 ağırlıklı): \%0,2; 1/98 kaybıyla \%22,5; tam MFN ortalaması \%22,6. Rejim: GB \%0, İTÜ \%2, tarım \%98 \etiketV{}/\etiketT{}.
\item
  Tepe satırlar (S4 artışı ≥ \%10; en büyük 3): yok.
\item
  AB pazarında AB dışı rakipler (tercih durumu): GBR \%55; CHE \%24; NZL \%7; UKR \%4; USA \%2. Tercihli rakip payı \%93.
\item
  Menşe (S2 / S3): risk D/D; \kod{c\_\allowbreak{}RoO} \%1,5 / \%2,0; \kod{p\_\allowbreak{}uyum} 0,90 / 0,85 \etiketH{}. Kural: Fasıl 4: Fasıl 4 malzemeleri tamamen elde edilmiş.
\item
  EUDR kapsamındaki AB ihracatı payı: \%0; SKDM: \%0.
\end{itemize}

\textbf{5. Türkiye tarafı maruziyet}

\begin{itemize}
\tightlist
\item
  Uygulanan tarife (2026, ithalat ağırlıklı): MFN \%137,9, İGV dahil \%137,9 \etiketV{}. S4'te AB menşeline (c2): \%0,0.
\item
  Kritik girdiler (HS4 0402, 0404, 0405, 0406), Türkiye'nin toplam ithalatı 2024: \$90 mn (AB \%64; CYP \%24, IRL \%9, NLD \%8) \etiketV{}.
\item
  AB'den ithalat süt tozu, peynir altı suyu tozu ve tereyağı. Tarım rejiminde; 1/98'deki Türk tavizleri kalkarsa maliyet artar, ama ithalat küçük (\$52 mn).
\end{itemize}

\textbf{6. Kabiliyet}

\begin{itemize}
\tightlist
\item
  KS 45,7 (sıra 161/230; ağırlık duyarlılığında P10--P90 105--196). Bileşenler: RCA 17, PCI 58 (PCI 0,15), yurt içi KD 51, Ar-Ge 25, tedarikçi 84, verimlilik 54, firma 56 \etiketT{}.
\item
  Ürün uzayında atlanabilir ürünler (sınıf içi; yoğunluk üst çeyrek, PCI \textgreater{} sınıf): 040520 (PCI +0.16, RCA 0.07) \etiketT{}.
\end{itemize}

\textbf{7. Stratejik önem.} SS 56,8 (en yüksek boyut: gıda güvenliği 75; ağırlıklı ortalama 14,2) \etiketH{}. Süt ve süt ürünleri; çiğ süt üreticisi geliri

\textbf{8. Makro duyarlılık}

\begingroup\gbtablo\small

\begin{longtable}[]{@{}lr@{}}
\toprule\noalign{}
\textbf{Gösterge} & \textbf{Değer (\etiketT{})} \\
\midrule\noalign{}
\endhead
\bottomrule\noalign{}
\endlastfoot
TÜFE ağırlığı (doğrudan, COICOP 2018 eşlemesi); tüketimde ithal pay & \%3,01; \%4 \\
ÜFE vekili (imalat üretimindeki pay) & \%1,63 \\
EDE: 1 puan tarife, tüm menşe / yalnız AB (bp; ρ = 1) & 0,137 / 0,063 \\
Geçişkenlik ρ (senaryolarda) & 0,6--1,0 (P50 0,8) \\
Piyasa yoğunlaşması (250+ istihdam payı) & \%52 \\
CKE: tarife (\$ mn / bp) ve mali (\$ mn / mlr TL); 0--100 puan & 135,1 ve 0,64; 12 \\
\end{longtable}

\endgroup

\textbf{9. Senaryo etkileri, T+5, P50 (P10 / P90)}

\begingroup\gbtablo\small

\begin{longtable}[]{@{}lrrr@{}}
\toprule\noalign{}
\textbf{Değişken} & \textbf{S2} & \textbf{S3} & \textbf{S4} \\
\midrule\noalign{}
\endhead
\bottomrule\noalign{}
\endlastfoot
AB'ye ihracat (\$ mn) & 0 (0 / +1) & 0 (0 / 0) & 0 (−9 / 0) \\
Toplam ihracat (\$ mn) & 0 (0 / 0) & 0 (0 / 0) & 0 (−7 / 0) \\
Nihai mal ithalatı (\$ mn) & 0 (0 / +4) & 0 (0 / +2) & 0 (0 / 0) \\
Cari denge (\$ mn) & 0 (−3 / 0) & 0 (−2 / 0) & 0 (−6 / 0) \\
Üretim (\$ mn) & 0 (−3 / 0) & 0 (−2 / 0) & 0 (−7 / 0) \\
İstihdam (kişi) & 0 (−10 / 0) & 0 (−7 / 0) & 0 (−23 / 0) \\
TÜFE (bp, birinci tur) & 0,00 (0,00 / 0,00) & 0,00 (0,00 / 0,00) & 0,00 (0,00 / 0,00) \\
Yatırım ve DYY (bütün senaryolar) & --- & & \\
\end{longtable}

\endgroup

\emph{Kaynak:} \etiketT{} f3e: kısmi denge, toplamlar Faz 2B GE dağılımına kalibre; S4'te 1/98 kaybı olasılığı 0,25.

\textbf{13. Kazanan ve kaybedenler; bölgesel ve sosyal etki; telafi.} Etkilenen kesim yok.

\textbf{15. Güven düzeyi; veri boşlukları; \etiketD{}.} Yüksek.

\hypertarget{nace_kartlar_10-12-311-1052-dondurma-imalatux131--gb-statuxfcsuxfc-iux307tuxfc--birincilikincil-arketip---duruux15f-koru}{%
\subsection{3.11 10.52 Dondurma imalatı --- GB statüsü: İTÜ \textbar{} Birincil/ikincil arketip: --- \textbar{} Duruş: Koru}\label{nace_kartlar_10-12-311-1052-dondurma-imalatux131--gb-statuxfcsuxfc-iux307tuxfc--birincilikincil-arketip---duruux15f-koru}}

\emph{İç pazara dönük (dondurma). Kıbrıs düzeltmesinden sonra AB ihracatı \$1 mn; GB etkisi yok.}

\textbf{1. Kimlik.} Rev. 2.1: 10.52 (1:1); CPA 10.52. Başlıca HS4 (TR ihracatı 2024): 2105 \%100. BEC payları (ticaret ağırlıklı): ara \%0 / sermaye \%0 / tüketim \%100.

\textbf{2. Ölçek}

\begingroup\gbtablo\small

\begin{longtable}[]{@{}
  >{\raggedright\arraybackslash}p{(\columnwidth - 4\tabcolsep) * \real{0.2933}}
  >{\raggedright\arraybackslash}p{(\columnwidth - 4\tabcolsep) * \real{0.3768}}
  >{\raggedright\arraybackslash}p{(\columnwidth - 4\tabcolsep) * \real{0.3300}}@{}}
\toprule\noalign{}
\begin{minipage}[b]{\linewidth}\raggedright
\textbf{Gösterge}
\end{minipage} & \begin{minipage}[b]{\linewidth}\raggedright
\textbf{Değer}
\end{minipage} & \begin{minipage}[b]{\linewidth}\raggedright
\textbf{Etiket ve kaynak}
\end{minipage} \\
\midrule\noalign{}
\endhead
\bottomrule\noalign{}
\endlastfoot
Üretim değeri; katma değer & \$2.017 mn; \$489 mn (GSYH'nin \%0,04) & TÜİK YSHİ 2024 \etiketV{} \\
Çalışan; girişim (0 / 1--9 / 10--49 / 50--249 / 250+) & 7.589; 878 (22 / 841 / 42 / 10 / 4) & TÜİK YSHİ ve İş Kayıtları 2024 \etiketV{}; gizli büyüklük hücreleri \etiketT{} \\
250+ girişimlerin istihdam payı; AB'ye göre kişi başı KD & \%40; 1,35 & \etiketT{} (AB27 SBS 2023, nominal kur) \\
İlk 3 İBBS-2 bölgesi & İstanbul \%26; Hatay, Kahramanmaraş, Osmaniye \%23; Konya, Karaman \%5 & \etiketT{} iş kayıtları + SGK 2025, RAS \\
Yabancı sermaye; çapa yatırımcılar & Unilever (Algida; Konya), Ülker (Golf), Panda \etiketH{} & \etiketH{} (firma bilgisi; pay verisi yok: VB-30) \\
\end{longtable}

\endgroup

\textbf{3. Dış ticaret}

\begingroup\gbtablo\small

\begin{longtable}[]{@{}
  >{\raggedright\arraybackslash}p{(\columnwidth - 2\tabcolsep) * \real{0.4824}}
  >{\raggedright\arraybackslash}p{(\columnwidth - 2\tabcolsep) * \real{0.5176}}@{}}
\toprule\noalign{}
\begin{minipage}[b]{\linewidth}\raggedright
\textbf{Gösterge}
\end{minipage} & \begin{minipage}[b]{\linewidth}\raggedright
\textbf{Değer (BACI 2024, \$; \etiketV{}, eşleme \etiketT{})}
\end{minipage} \\
\midrule\noalign{}
\endhead
\bottomrule\noalign{}
\endlastfoot
İhracat (2024); 2020--24 değişimi & \$72 mn; \%56 \\
AB payı; ilk 5 AB pazarı & \%1; DEU \%32; BEL \%28; NLD \%23; FRA \%5; MLT \%3 \\
İlk 5 AB dışı pazar & USA \%16; IRQ \%16; CYP \%10; ARE \%7; SRB \%7 \\
İthalat; AB ve Çin payı; ilk kaynaklar & \$5 mn; AB \%54, Çin \%2; CHE \%30; FRA \%30; MDA \%10; POL \%7; NLD \%5 \\
Dış ticaret dengesi & \$+67 mn \\
İthal girdi içeriği (I-O, toplam) / KD÷üretim & \%20 / \%24 \\
RCA dünya (2020 → 2024) / AB pazarında & 0,96 → 0,99 / 0,10 \\
Birim değer ÷ AB pazarı ortalaması & 1,10 \\
AB pazar payı (AB dışı tedarik içinde, 2020 → 2024); TR'nin sırası & \%0,1 → \%0,4; 12. \\
\end{longtable}

\endgroup

\textbf{4. AB tarafı maruziyet}

\begin{itemize}
\tightlist
\item
  \kod{τ\_\allowbreak{}AB} (S4 artışı, rota A, TR→AB 2025 ağırlıklı): \%8,5; 1/98 kaybıyla \%8,5; tam MFN ortalaması \%15,7. Rejim: GB \%0, İTÜ \%100, tarım \%0 \etiketV{}/\etiketT{}.
\item
  Tepe satırlar (S4 artışı ≥ \%10; en büyük 3): yok.
\item
  AB pazarında AB dışı rakipler (tercih durumu): GBR \%45; SRB \%30; UKR \%10; CHE \%4; MKD \%3. Tercihli rakip payı \%63.
\item
  Menşe (S2 / S3): risk D/O; \kod{c\_\allowbreak{}RoO} \%1,5 / \%2,5; \kod{p\_\allowbreak{}uyum} 0,90 / 0,85 \etiketH{}. Kural: 2105: başlık değişikliği; Fasıl 4 ve şeker ağırlık sınırları.
\item
  EUDR kapsamındaki AB ihracatı payı: \%0; SKDM: \%0.
\end{itemize}

\textbf{5. Türkiye tarafı maruziyet}

\begin{itemize}
\tightlist
\item
  Uygulanan tarife (2026, ithalat ağırlıklı): MFN \%8,2, İGV dahil \%8,2 \etiketV{}. S4'te AB menşeline (c2): \%8,2.
\end{itemize}

\textbf{6. Kabiliyet}

\begin{itemize}
\tightlist
\item
  KS 51,8 (sıra 91/230; ağırlık duyarlılığında P10--P90 48--149). Bileşenler: RCA 43, PCI 43 (PCI −0,13), yurt içi KD 60, Ar-Ge 25, tedarikçi 84, verimlilik 97, firma 38 \etiketT{}.
\item
  Ürün uzayında atlanabilir ürünler (sınıf içi; yoğunluk üst çeyrek, PCI \textgreater{} sınıf): yok \etiketT{}.
\end{itemize}

\textbf{7. Stratejik önem.} SS 3,7 (en yüksek boyut: gıda güvenliği 5; ağırlıklı ortalama 0,8) \etiketH{}.

\textbf{8. Makro duyarlılık}

\begingroup\gbtablo\small

\begin{longtable}[]{@{}lr@{}}
\toprule\noalign{}
\textbf{Gösterge} & \textbf{Değer (\etiketT{})} \\
\midrule\noalign{}
\endhead
\bottomrule\noalign{}
\endlastfoot
TÜFE ağırlığı (doğrudan, COICOP 2018 eşlemesi); tüketimde ithal pay & \%0,21; \%2 \\
ÜFE vekili (imalat üretimindeki pay) & \%0,30 \\
EDE: 1 puan tarife, tüm menşe / yalnız AB (bp; ρ = 1) & 0,004 / 0,002 \\
Geçişkenlik ρ (senaryolarda) & 0,6--1,0 (P50 0,8) \\
Piyasa yoğunlaşması (250+ istihdam payı) & \%40 \\
CKE: tarife (\$ mn / bp) ve mali (\$ mn / mlr TL); 0--100 puan & 376,6 ve 0,31; 15 \\
\end{longtable}

\endgroup

\textbf{9. Senaryo etkileri, T+5, P50 (P10 / P90)}

\begingroup\gbtablo\small

\begin{longtable}[]{@{}
  >{\raggedright\arraybackslash}p{(\columnwidth - 6\tabcolsep) * \real{0.3834}}
  >{\raggedleft\arraybackslash}p{(\columnwidth - 6\tabcolsep) * \real{0.1816}}
  >{\raggedleft\arraybackslash}p{(\columnwidth - 6\tabcolsep) * \real{0.2175}}
  >{\raggedleft\arraybackslash}p{(\columnwidth - 6\tabcolsep) * \real{0.2175}}@{}}
\toprule\noalign{}
\begin{minipage}[b]{\linewidth}\raggedright
\textbf{Değişken}
\end{minipage} & \begin{minipage}[b]{\linewidth}\raggedleft
\textbf{S2}
\end{minipage} & \begin{minipage}[b]{\linewidth}\raggedleft
\textbf{S3}
\end{minipage} & \begin{minipage}[b]{\linewidth}\raggedleft
\textbf{S4}
\end{minipage} \\
\midrule\noalign{}
\endhead
\bottomrule\noalign{}
\endlastfoot
AB'ye ihracat (\$ mn) & 0 (0 / 0) & 0 (0 / 0) & 0 (−1 / 0) \\
Toplam ihracat (\$ mn) & 0 (0 / 0) & 0 (0 / 0) & 0 (0 / 0) \\
Nihai mal ithalatı (\$ mn) & 0 (−1 / 0) & 0 (−1 / 0) & 0 (−1 / 0) \\
Cari denge (\$ mn) & 0 (0 / 0) & 0 (0 / 0) & 0 (0 / 0) \\
Üretim (\$ mn) & 0 (0 / 0) & 0 (0 / 0) & 0 (0 / 0) \\
İstihdam (kişi) & 0 (0 / +1) & 0 (0 / +1) & +1 (0 / +1) \\
TÜFE (bp, birinci tur) & 0,00 (0,00 / 0,00) & +0,01 (+0,01 / +0,01) & +0,02 (+0,02 / +0,02) \\
Yatırım ve DYY (bütün senaryolar) & --- & & \\
\end{longtable}

\endgroup

\emph{Kaynak:} \etiketT{} f3e: kısmi denge, toplamlar Faz 2B GE dağılımına kalibre; S4'te 1/98 kaybı olasılığı 0,25.

\textbf{13. Kazanan ve kaybedenler; bölgesel ve sosyal etki; telafi.} Etkilenen kesim yok.

\textbf{15. Güven düzeyi; veri boşlukları; \etiketD{}.} Yüksek.

\hypertarget{nace_kartlar_10-12-312-1061-uxf6ux11fuxfctuxfclmuxfcux15f-hububat-ve-sebze-uxfcruxfcnleri-imalatux131--gb-statuxfcsuxfc-karma--birincilikincil-arketip-bd--duruux15f-geliux15ftir}{%
\subsection{3.12 10.61 Öğütülmüş hububat ve sebze ürünleri imalatı --- GB statüsü: KARMA \textbar{} Birincil/ikincil arketip: B/D \textbar{} Duruş: Geliştir}\label{nace_kartlar_10-12-312-1061-uxf6ux11fuxfctuxfclmuxfcux15f-hububat-ve-sebze-uxfcruxfcnleri-imalatux131--gb-statuxfcsuxfc-karma--birincilikincil-arketip-bd--duruux15f-geliux15ftir}}

\emph{Dünya un ihracatında birinci (\%19,6), bulgurda \%75. İhracatın \%92'si AB dışına (Irak, Suriye, Somali) ve DİR'de Rus buğdayıyla. GB etkisi küçük; asıl kırılganlık buğday arzı.}

\textbf{1. Kimlik.} Rev. 2.1: 10.61 (n:1); CPA 10.61. Başlıca HS4 (TR ihracatı 2024): 1101 \%62, 1904 \%12, 1006 \%11, 1103 \%5. BEC payları (ticaret ağırlıklı): ara \%68 / sermaye \%0 / tüketim \%32.

\textbf{2. Ölçek}

\begingroup\gbtablo\small

\begin{longtable}[]{@{}
  >{\raggedright\arraybackslash}p{(\columnwidth - 4\tabcolsep) * \real{0.2832}}
  >{\raggedright\arraybackslash}p{(\columnwidth - 4\tabcolsep) * \real{0.3999}}
  >{\raggedright\arraybackslash}p{(\columnwidth - 4\tabcolsep) * \real{0.3169}}@{}}
\toprule\noalign{}
\begin{minipage}[b]{\linewidth}\raggedright
\textbf{Gösterge}
\end{minipage} & \begin{minipage}[b]{\linewidth}\raggedright
\textbf{Değer}
\end{minipage} & \begin{minipage}[b]{\linewidth}\raggedright
\textbf{Etiket ve kaynak}
\end{minipage} \\
\midrule\noalign{}
\endhead
\bottomrule\noalign{}
\endlastfoot
Üretim değeri; katma değer & \$7.803 mn; \$1.245 mn (GSYH'nin \%0,09) & TÜİK YSHİ 2024 \etiketV{} \\
Çalışan; girişim (0 / 1--9 / 10--49 / 50--249 / 250+) & 26.004; 2.657 (172 / 2.271 / 383 / 75 / 8) & TÜİK YSHİ ve İş Kayıtları 2024 \etiketV{}; gizli büyüklük hücreleri \etiketT{} \\
250+ girişimlerin istihdam payı; AB'ye göre kişi başı KD & \%10; 0,60 & \etiketT{} (AB27 SBS 2023, nominal kur) \\
İlk 3 İBBS-2 bölgesi & Konya, Karaman \%10; Tekirdağ, Edirne, Kırklareli \%8; Balıkesir, Çanakkale \%7 & \etiketT{} iş kayıtları + SGK 2025, RAS \\
Yabancı sermaye; çapa yatırımcılar & Yüzlerce orta ölçekli değirmen (Gaziantep, Konya, Mersin, Tekirdağ); yabancı çapa yok \etiketH{} & \etiketH{} (firma bilgisi; pay verisi yok: VB-30) \\
\end{longtable}

\endgroup

\textbf{3. Dış ticaret}

\begingroup\gbtablo\small

\begin{longtable}[]{@{}
  >{\raggedright\arraybackslash}p{(\columnwidth - 2\tabcolsep) * \real{0.4766}}
  >{\raggedright\arraybackslash}p{(\columnwidth - 2\tabcolsep) * \real{0.5234}}@{}}
\toprule\noalign{}
\begin{minipage}[b]{\linewidth}\raggedright
\textbf{Gösterge}
\end{minipage} & \begin{minipage}[b]{\linewidth}\raggedright
\textbf{Değer (BACI 2024, \$; \etiketV{}, eşleme \etiketT{})}
\end{minipage} \\
\midrule\noalign{}
\endhead
\bottomrule\noalign{}
\endlastfoot
İhracat (2024); 2020--24 değişimi & \$1.870 mn; \%25 \\
AB payı; ilk 5 AB pazarı & \%7; DEU \%41; BEL \%16; AUT \%7; NLD \%7; FRA \%6 \\
İlk 5 AB dışı pazar & IRQ \%31; SYR \%11; SOM \%3; USA \%3; DJI \%3 \\
İthalat; AB ve Çin payı; ilk kaynaklar & \$739 mn; AB \%17, Çin \%11; RUS \%43; CHN \%11; IND \%7; ITA \%6; BRA \%5 \\
Dış ticaret dengesi & \$+1.131 mn \\
İthal girdi içeriği (I-O, toplam) / KD÷üretim & \%20 / \%16 \\
RCA dünya (2020 → 2024) / AB pazarında & 2,85 → 2,40 / 0,95 \\
Birim değer ÷ AB pazarı ortalaması & 1,11 \\
AB pazar payı (AB dışı tedarik içinde, 2020 → 2024); TR'nin sırası & \%3,0 → \%3,8; 9. \\
\end{longtable}

\endgroup

\textbf{4. AB tarafı maruziyet}

\begin{itemize}
\tightlist
\item
  \kod{τ\_\allowbreak{}AB} (S4 artışı, rota A, TR→AB 2025 ağırlıklı): \%4,9; 1/98 kaybıyla \%10,3; tam MFN ortalaması \%24,5. Rejim: GB \%0, İTÜ \%64, tarım \%36 \etiketV{}/\etiketT{}.
\item
  Tepe satırlar (S4 artışı ≥ \%10; en büyük 3): yok.
\item
  AB pazarında AB dışı rakipler (tercih durumu): GBR \%22; PAK \%13; IND \%9; MMR \%9; THA \%8. Tercihli rakip payı \%65.
\item
  Menşe (S2 / S3): risk O/O; \kod{c\_\allowbreak{}RoO} \%2,5 / \%3,5; \kod{p\_\allowbreak{}uyum} 0,85 / 0,75 \etiketH{}. Kural: Fasıl 11 ve 1904: tahıl tamamen elde edilmiş (makarnalık buğday istisnası 1904'te).
\item
  EUDR kapsamındaki AB ihracatı payı: \%0; SKDM: \%0.
\item
  AB'ye ihracat (\$139 mn) un, bulgur ve kahvaltılık tahıl. İTÜ payının (\%64) S4 sanayi bileşeni \%4,9; AB'ye ihracatta S4'te −\$24 mn.
\item
  Menşe (S2/S3): Fasıl 11'de tahılın tamamen elde edilmiş olması isteniyor. DİR'de Rus buğdayından üretilen un menşe kazanmıyor (O).
\end{itemize}

\textbf{5. Türkiye tarafı maruziyet}

\begin{itemize}
\tightlist
\item
  Uygulanan tarife (2026, ithalat ağırlıklı): MFN \%19,2, İGV dahil \%19,2 \etiketV{}. S4'te AB menşeline (c2): \%0,9.
\item
  Kritik girdiler (HS4 1001, 1003, 1005), Türkiye'nin toplam ithalatı 2024: \$2.438 mn (AB \%5; RUS \%60, UKR \%30, BRA \%2) \etiketV{}.
\item
  Buğday ithalatının \%43'ü Rusya'dan. Karadeniz koridoru ve Rusya'nın ihracat kotası en büyük risk (2022, 2023) \etiketV{}.
\end{itemize}

\textbf{6. Kabiliyet}

\begin{itemize}
\tightlist
\item
  KS 48,1 (sıra 135/230; ağırlık duyarlılığında P10--P90 63--180). Bileşenler: RCA 60, PCI 7 (PCI −1,18), yurt içi KD 49, Ar-Ge 25, tedarikçi 84, verimlilik 75, firma 77 \etiketT{}.
\item
  Ürün uzayında atlanabilir ürünler (sınıf içi; yoğunluk üst çeyrek, PCI \textgreater{} sınıf): yok \etiketT{}.
\item
  \textbf{Kırmızı çizgi yetkinlikleri:} Değirmencilik ve paketleme kapasitesi; DİR ağıyla kurulan Afrika--Orta Doğu dağıtım ilişkileri.
\item
  \textbf{Geliştirme hedefi:} T+5: Batı Afrika ve Körfez'de pazar payı; 2035: Rusya payı ≤ \%35.
\end{itemize}

\textbf{7. Stratejik önem.} SS 70,8 (en yüksek boyut: gıda güvenliği 95; ağırlıklı ortalama 14,2) \etiketH{}. Ekmek zincirinin başı: un ve buğday arzı gıda güvenliğinin çekirdeği.

\textbf{8. Makro duyarlılık}

\begingroup\gbtablo\small

\begin{longtable}[]{@{}lr@{}}
\toprule\noalign{}
\textbf{Gösterge} & \textbf{Değer (\etiketT{})} \\
\midrule\noalign{}
\endhead
\bottomrule\noalign{}
\endlastfoot
TÜFE ağırlığı (doğrudan, COICOP 2018 eşlemesi); tüketimde ithal pay & \%0,86; \%6 \\
ÜFE vekili (imalat üretimindeki pay) & \%1,17 \\
EDE: 1 puan tarife, tüm menşe / yalnız AB (bp; ρ = 1) & 0,107 / 0,018 \\
Geçişkenlik ρ (senaryolarda) & 0,6--1,0 (P50 0,8) \\
Piyasa yoğunlaşması (250+ istihdam payı) & \%10 \\
CKE: tarife (\$ mn / bp) ve mali (\$ mn / mlr TL); 0--100 puan & 1.403,9 ve 9,03; 37 \\
\end{longtable}

\endgroup

\textbf{9. Senaryo etkileri, T+5, P50 (P10 / P90)}

\begingroup\gbtablo\small

\begin{longtable}[]{@{}
  >{\raggedright\arraybackslash}p{(\columnwidth - 6\tabcolsep) * \real{0.3895}}
  >{\raggedleft\arraybackslash}p{(\columnwidth - 6\tabcolsep) * \real{0.1831}}
  >{\raggedleft\arraybackslash}p{(\columnwidth - 6\tabcolsep) * \real{0.2076}}
  >{\raggedleft\arraybackslash}p{(\columnwidth - 6\tabcolsep) * \real{0.2198}}@{}}
\toprule\noalign{}
\begin{minipage}[b]{\linewidth}\raggedright
\textbf{Değişken}
\end{minipage} & \begin{minipage}[b]{\linewidth}\raggedleft
\textbf{S2}
\end{minipage} & \begin{minipage}[b]{\linewidth}\raggedleft
\textbf{S3}
\end{minipage} & \begin{minipage}[b]{\linewidth}\raggedleft
\textbf{S4}
\end{minipage} \\
\midrule\noalign{}
\endhead
\bottomrule\noalign{}
\endlastfoot
AB'ye ihracat (\$ mn) & −12 (−29 / +3) & −16 (−30 / −3) & −24 (−43 / −10) \\
Toplam ihracat (\$ mn) & −9 (−20 / +2) & −11 (−21 / −2) & −19 (−34 / −8) \\
Nihai mal ithalatı (\$ mn) & −1 (−2 / +8) & −1 (−3 / +5) & −2 (−4 / −1) \\
Cari denge (\$ mn) & −6 (−15 / −4) & −8 (−15 / −5) & −13 (−24 / −6) \\
Üretim (\$ mn) & −8 (−18 / −5) & −10 (−19 / −6) & −17 (−30 / −8) \\
İstihdam (kişi) & −19 (−43 / −11) & −23 (−44 / −14) & −39 (−71 / −18) \\
TÜFE (bp, birinci tur) & 0,00 (0,00 / 0,00) & +0,01 (0,00 / +0,01) & +0,01 (+0,01 / +0,02) \\
Yatırım ve DYY (bütün senaryolar) & Etkisi yok & & \\
\end{longtable}

\endgroup

\emph{Kaynak:} \etiketT{} f3e: kısmi denge, toplamlar Faz 2B GE dağılımına kalibre; S4'te 1/98 kaybı olasılığı 0,25.

\textbf{10. Politika paketi}

\begingroup\gbtablo\scriptsize

\begin{longtable}[]{@{}
  >{\raggedright\arraybackslash}p{(\columnwidth - 16\tabcolsep) * \real{0.0983}}
  >{\raggedright\arraybackslash}p{(\columnwidth - 16\tabcolsep) * \real{0.1222}}
  >{\raggedright\arraybackslash}p{(\columnwidth - 16\tabcolsep) * \real{0.1498}}
  >{\raggedright\arraybackslash}p{(\columnwidth - 16\tabcolsep) * \real{0.0944}}
  >{\raggedright\arraybackslash}p{(\columnwidth - 16\tabcolsep) * \real{0.1323}}
  >{\raggedright\arraybackslash}p{(\columnwidth - 16\tabcolsep) * \real{0.0588}}
  >{\raggedright\arraybackslash}p{(\columnwidth - 16\tabcolsep) * \real{0.1488}}
  >{\raggedright\arraybackslash}p{(\columnwidth - 16\tabcolsep) * \real{0.1197}}
  >{\raggedright\arraybackslash}p{(\columnwidth - 16\tabcolsep) * \real{0.0758}}@{}}
\toprule\noalign{}
\begin{minipage}[b]{\linewidth}\raggedright
\textbf{Eylem ID}
\end{minipage} & \begin{minipage}[b]{\linewidth}\raggedright
\textbf{Araç}
\end{minipage} & \begin{minipage}[b]{\linewidth}\raggedright
\textbf{Parametre}
\end{minipage} & \begin{minipage}[b]{\linewidth}\raggedright
\textbf{Sorumlu}
\end{minipage} & \begin{minipage}[b]{\linewidth}\raggedright
\textbf{Hukuki araç}
\end{minipage} & \begin{minipage}[b]{\linewidth}\raggedright
\textbf{Bütçe}
\end{minipage} & \begin{minipage}[b]{\linewidth}\raggedright
\textbf{Başlangıç--bitiş; sunset}
\end{minipage} & \begin{minipage}[b]{\linewidth}\raggedright
\textbf{Performans koşulu}
\end{minipage} & \begin{minipage}[b]{\linewidth}\raggedright
\textbf{KPI}
\end{minipage} \\
\midrule\noalign{}
\endhead
\bottomrule\noalign{}
\endlastfoot
E-10.61-01 & Buğday kaynak çeşitlendirme & Kazakistan, Romanya, Bulgaristan ve Kanada'dan TMO alım çerçeveleri; Rusya payı ≤ \%35 & TMO, Ticaret Bakanlığı & TMO yönetim kurulu kararı; ikili mutabakatlar & İdari & T−18 → T+36; sunset: --- & Alım fiyatı ≤ Rus teklifi + \%3 & Tek kaynak payı ≤ \%35 \\
\end{longtable}

\endgroup

\textbf{12. Erken uyarı ve tetikler}

\begingroup\gbtablo\small

\begin{longtable}[]{@{}
  >{\raggedright\arraybackslash}p{(\columnwidth - 4\tabcolsep) * \real{0.4746}}
  >{\raggedright\arraybackslash}p{(\columnwidth - 4\tabcolsep) * \real{0.0720}}
  >{\raggedright\arraybackslash}p{(\columnwidth - 4\tabcolsep) * \real{0.4533}}@{}}
\toprule\noalign{}
\begin{minipage}[b]{\linewidth}\raggedright
\textbf{Gösterge}
\end{minipage} & \begin{minipage}[b]{\linewidth}\raggedright
\textbf{Eşik}
\end{minipage} & \begin{minipage}[b]{\linewidth}\raggedright
\textbf{Tetiklenecek önlem}
\end{minipage} \\
\midrule\noalign{}
\endhead
\bottomrule\noalign{}
\endlastfoot
Rusya'nın buğday ihracat vergisi/kotası ya da Karadeniz kesintisi & Duyuru & TMO stok satışı; DİR'de alternatif kaynak izni hızlandırılır \\
\end{longtable}

\endgroup

\textbf{13. Kazanan ve kaybedenler; bölgesel ve sosyal etki; telafi.} GB çıkışından küçük kayıp (S4 P50: −39 kişi); bölgeler Konya--Karaman, Trakya, Balıkesir.

\textbf{14. Riskler ve azaltım}

\begin{itemize}
\tightlist
\item
  Karadeniz arz şoku (GB'den bağımsız) → kaynak çeşitlendirme.
\end{itemize}

\textbf{15. Güven düzeyi; veri boşlukları; \etiketD{}.} Yüksek.

\hypertarget{nace_kartlar_10-12-313-1062-niux15fasta-ve-niux15fastalux131-uxfcruxfcnlerin-imalatux131--gb-statuxfcsuxfc-karma--birincilikincil-arketip-b--duruux15f-geliux15ftir}{%
\subsection{3.13 10.62 Nişasta ve nişastalı ürünlerin imalatı --- GB statüsü: KARMA \textbar{} Birincil/ikincil arketip: B \textbar{} Duruş: Geliştir}\label{nace_kartlar_10-12-313-1062-niux15fasta-ve-niux15fastalux131-uxfcruxfcnlerin-imalatux131--gb-statuxfcsuxfc-karma--birincilikincil-arketip-b--duruux15f-geliux15ftir}}

\emph{Fruktozda dünya birincisi (\%35). AB'ye \$96 mn ihracat; AB MFN'si yüksek (glikozda \%16), S4'te AB'ye ihracat −\$31 mn. S2'de menşe sorunu yok, çünkü ithal mısırdan glikoz başlık değişikliğiyle menşe kazanıyor.}

\textbf{1. Kimlik.} Rev. 2.1: 10.62 (n:1); CPA 10.62. Başlıca HS4 (TR ihracatı 2024): 1702 \%47, 1108 \%25, 1515 \%20, 3505 \%7. BEC payları (ticaret ağırlıklı): ara \%89 / sermaye \%0 / tüketim \%11.

\textbf{2. Ölçek}

\begingroup\gbtablo\small

\begin{longtable}[]{@{}
  >{\raggedright\arraybackslash}p{(\columnwidth - 4\tabcolsep) * \real{0.2964}}
  >{\raggedright\arraybackslash}p{(\columnwidth - 4\tabcolsep) * \real{0.3696}}
  >{\raggedright\arraybackslash}p{(\columnwidth - 4\tabcolsep) * \real{0.3340}}@{}}
\toprule\noalign{}
\begin{minipage}[b]{\linewidth}\raggedright
\textbf{Gösterge}
\end{minipage} & \begin{minipage}[b]{\linewidth}\raggedright
\textbf{Değer}
\end{minipage} & \begin{minipage}[b]{\linewidth}\raggedright
\textbf{Etiket ve kaynak}
\end{minipage} \\
\midrule\noalign{}
\endhead
\bottomrule\noalign{}
\endlastfoot
Üretim değeri; katma değer & \$1.638 mn; \$357 mn (GSYH'nin \%0,03) & TÜİK YSHİ 2024 \etiketV{} \\
Çalışan; girişim (0 / 1--9 / 10--49 / 50--249 / 250+) & 3.117; 33 (3 / 18 / 4 / 6 / 5) & TÜİK YSHİ ve İş Kayıtları 2024 \etiketV{}; gizli büyüklük hücreleri \etiketT{} \\
250+ girişimlerin istihdam payı; AB'ye göre kişi başı KD & \%75; 0,58 & \etiketT{} (AB27 SBS 2023, nominal kur) \\
İlk 3 İBBS-2 bölgesi & Adana, Mersin \%45; İstanbul \%35; Kocaeli, Sakarya, Düzce, Bolu, Yalova \%6 & \etiketT{} iş kayıtları + SGK 2025, RAS \\
Yabancı sermaye; çapa yatırımcılar & Cargill (Orhangazi), Tat Nişasta, Sunar Mısır, Pendik Nişasta \etiketH{} & \etiketH{} (firma bilgisi; pay verisi yok: VB-30) \\
\end{longtable}

\endgroup

\textbf{3. Dış ticaret}

\begingroup\gbtablo\small

\begin{longtable}[]{@{}
  >{\raggedright\arraybackslash}p{(\columnwidth - 2\tabcolsep) * \real{0.4802}}
  >{\raggedright\arraybackslash}p{(\columnwidth - 2\tabcolsep) * \real{0.5198}}@{}}
\toprule\noalign{}
\begin{minipage}[b]{\linewidth}\raggedright
\textbf{Gösterge}
\end{minipage} & \begin{minipage}[b]{\linewidth}\raggedright
\textbf{Değer (BACI 2024, \$; \etiketV{}, eşleme \etiketT{})}
\end{minipage} \\
\midrule\noalign{}
\endhead
\bottomrule\noalign{}
\endlastfoot
İhracat (2024); 2020--24 değişimi & \$528 mn; \%66 \\
AB payı; ilk 5 AB pazarı & \%18; NLD \%36; ITA \%15; DEU \%14; ESP \%9; FRA \%4 \\
İlk 5 AB dışı pazar & LBY \%12; NGA \%11; IRQ \%9; USA \%6; EGY \%4 \\
İthalat; AB ve Çin payı; ilk kaynaklar & \$206 mn; AB \%68, Çin \%7; NLD \%17; DEU \%15; CHN \%7; BEL \%6; FRA \%6 \\
Dış ticaret dengesi & \$+322 mn \\
İthal girdi içeriği (I-O, toplam) / KD÷üretim & \%20 / \%22 \\
RCA dünya (2020 → 2024) / AB pazarında & 1,70 → 1,84 / 3,34 \\
Birim değer ÷ AB pazarı ortalaması & 0,85 \\
AB pazar payı (AB dışı tedarik içinde, 2020 → 2024); TR'nin sırası & \%7,2 → \%13,3; 2. \\
\end{longtable}

\endgroup

\textbf{4. AB tarafı maruziyet}

\begin{itemize}
\tightlist
\item
  \kod{τ\_\allowbreak{}AB} (S4 artışı, rota A, TR→AB 2025 ağırlıklı): \%14,2; 1/98 kaybıyla \%15,4; tam MFN ortalaması \%54,3. Rejim: GB \%0, İTÜ \%95, tarım \%5 \etiketV{}/\etiketT{}.
\item
  Tepe satırlar (S4 artışı ≥ \%10; en büyük 3): 1 satır, AB ihracatının \%83'i: KN 1702 50 00 \%16,0 (€67,0 mn).
\item
  AB pazarında AB dışı rakipler (tercih durumu): USA \%20; CHN \%11; GBR \%11; THA \%8; MEX \%7. Tercihli rakip payı \%34.
\item
  Menşe (S2 / S3): risk D/O; \kod{c\_\allowbreak{}RoO} \%1,5 / \%2,5; \kod{p\_\allowbreak{}uyum} 0,90 / 0,85 \etiketH{}. Kural: 1108: tahıl tamamen elde edilmiş; 1702 glikoz: başlık değişikliği (ithal mısırdan nişasta → glikoz menşe kazanır).
\item
  EUDR kapsamındaki AB ihracatı payı: \%0; SKDM: \%0.
\item
  Değer zinciri: AB gıda ve içecek sanayisine ara girdi (glikoz ve fruktoz şurubu).
\item
  AB'ye ihracat kaybı S2'de −\$11 mn, S4'te −\$31 mn; çapa yatırımcının (Cargill) AB'ye arzı AB'deki tesislerinden karşılayabilir \etiketH{}.
\end{itemize}

\textbf{5. Türkiye tarafı maruziyet}

\begin{itemize}
\tightlist
\item
  Uygulanan tarife (2026, ithalat ağırlıklı): MFN \%26,4, İGV dahil \%26,4 \etiketV{}. S4'te AB menşeline (c2): \%5,3.
\item
  Kritik girdiler (HS4 1005), Türkiye'nin toplam ithalatı 2024: \$1.114 mn (AB \%10; UKR \%54, RUS \%28, BRA \%5) \etiketV{}.
\item
  Nişasta bazlı şekerin (NBŞ) yurt içi satışı Şeker Kanunu'ndaki kotayla sınırlı; ihracat kotasız \etiketV{}.
\end{itemize}

\textbf{6. Kabiliyet}

\begin{itemize}
\tightlist
\item
  KS 52,6 (sıra 83/230; ağırlık duyarlılığında P10--P90 47--144). Bileşenler: RCA 65, PCI 52 (PCI 0,05), yurt içi KD 55, Ar-Ge 25, tedarikçi 84, verimlilik 74, firma 13 \etiketT{}.
\item
  Ürün uzayında atlanabilir ürünler (sınıf içi; yoğunluk üst çeyrek, PCI \textgreater{} sınıf): yok \etiketT{}.
\item
  \textbf{Kırmızı çizgi yetkinlikleri:} Mısır ıslak öğütme kapasitesi; fruktozda ihracat sertifikasyonu.
\item
  \textbf{Geliştirme hedefi:} T+5: AB dışı (Kuzey Afrika, Körfez) payı +10 puan.
\end{itemize}

\textbf{7. Stratejik önem.} SS 30,2 (en yüksek boyut: gıda güvenliği 40; ağırlıklı ortalama 7,5) \etiketH{}. Nişasta ve glikoz: gıda ve ilaç sanayisine ara girdi

\textbf{8. Makro duyarlılık}

\begingroup\gbtablo\small

\begin{longtable}[]{@{}lr@{}}
\toprule\noalign{}
\textbf{Gösterge} & \textbf{Değer (\etiketT{})} \\
\midrule\noalign{}
\endhead
\bottomrule\noalign{}
\endlastfoot
TÜFE ağırlığı (doğrudan, COICOP 2018 eşlemesi); tüketimde ithal pay & \%0,00; \%6 \\
ÜFE vekili (imalat üretimindeki pay) & \%0,25 \\
EDE: 1 puan tarife, tüm menşe / yalnız AB (bp; ρ = 1) & 0,016 / 0,011 \\
Geçişkenlik ρ (senaryolarda) & 0,6--1,0 (P50 0,8) \\
Piyasa yoğunlaşması (250+ istihdam payı) & \%75 \\
CKE: tarife (\$ mn / bp) ve mali (\$ mn / mlr TL); 0--100 puan & 2.899,1 ve 13,27; 50 \\
\end{longtable}

\endgroup

\textbf{9. Senaryo etkileri, T+5, P50 (P10 / P90)}

\begingroup\gbtablo\small

\begin{longtable}[]{@{}
  >{\raggedright\arraybackslash}p{(\columnwidth - 6\tabcolsep) * \real{0.2499}}
  >{\raggedleft\arraybackslash}p{(\columnwidth - 6\tabcolsep) * \real{0.4170}}
  >{\raggedleft\arraybackslash}p{(\columnwidth - 6\tabcolsep) * \real{0.1665}}
  >{\raggedleft\arraybackslash}p{(\columnwidth - 6\tabcolsep) * \real{0.1665}}@{}}
\toprule\noalign{}
\begin{minipage}[b]{\linewidth}\raggedright
\textbf{Değişken}
\end{minipage} & \begin{minipage}[b]{\linewidth}\raggedleft
\textbf{S2}
\end{minipage} & \begin{minipage}[b]{\linewidth}\raggedleft
\textbf{S3}
\end{minipage} & \begin{minipage}[b]{\linewidth}\raggedleft
\textbf{S4}
\end{minipage} \\
\midrule\noalign{}
\endhead
\bottomrule\noalign{}
\endlastfoot
AB'ye ihracat (\$ mn) & −11 (−26 / +2) & −15 (−28 / −3) & −31 (−49 / −13) \\
Toplam ihracat (\$ mn) & −8 (−19 / +1) & −11 (−20 / −2) & −25 (−39 / −10) \\
Nihai mal ithalatı (\$ mn) & −4 (−13 / +10) & −6 (−15 / +6) & −14 (−25 / −3) \\
Cari denge (\$ mn) & −4 (−8 / −1) & −4 (−7 / −2) & −9 (−13 / −6) \\
Üretim (\$ mn) & −5 (−10 / −2) & −5 (−9 / −2) & −11 (−17 / −7) \\
İstihdam (kişi) & −7 (−13 / −2) & −7 (−12 / −3) & −14 (−23 / −9) \\
TÜFE (bp, birinci tur) & +0,01 (+0,01 / +0,01) & +0,02 (+0,01 / +0,02) & +0,05 (+0,04 / +0,06) \\
Yatırım ve DYY (bütün senaryolar) & S4'te çapa yatırımcının AB'ye yönelik kapasite kararı riski \etiketH{} & & \\
\end{longtable}

\endgroup

\emph{Kaynak:} \etiketT{} f3e: kısmi denge, toplamlar Faz 2B GE dağılımına kalibre; S4'te 1/98 kaybı olasılığı 0,25.

\textbf{11. Müzakere pozisyonu}

\begin{itemize}
\tightlist
\item
  AB ile savunma: S2'de İTÜ'de sanayi bileşeninin sıfırı (bugünkü durumun sürmesi); 1702 için başlık değişikliği kuralı.
\end{itemize}

\textbf{12. Erken uyarı ve tetikler}

\begingroup\gbtablo\small

\begin{longtable}[]{@{}lll@{}}
\toprule\noalign{}
\textbf{Gösterge} & \textbf{Eşik} & \textbf{Tetiklenecek önlem} \\
\midrule\noalign{}
\endhead
\bottomrule\noalign{}
\endlastfoot
AB'ye glikoz ve fruktoz ihracatı (aylık) & Yıllık −\%20 & E-10-05 tamponu \\
\end{longtable}

\endgroup

\textbf{13. Kazanan ve kaybedenler; bölgesel ve sosyal etki; telafi.} Kaybeden (S4): Bursa (Orhangazi), Adana; P50'de −14 kişi (sermaye yoğun).

\textbf{14. Riskler ve azaltım}

\begin{itemize}
\tightlist
\item
  Çapa yatırımcının S4'te AB arzını kaydırması → S2 önceliği.
\end{itemize}

\textbf{15. Güven düzeyi; veri boşlukları; \etiketD{}.} Orta (çapa yatırımcı davranışı \etiketH{}).

\hypertarget{nace_kartlar_10-12-314-1071-ekmek-taze-pastane-uxfcruxfcnleri-ve-taze-kek-imalatux131--gb-statuxfcsuxfc-iux307tuxfc--birincilikincil-arketip-d--duruux15f-koru}{%
\subsection{3.14 10.71 Ekmek, taze pastane ürünleri ve taze kek imalatı --- GB statüsü: İTÜ \textbar{} Birincil/ikincil arketip: ---/D \textbar{} Duruş: Koru}\label{nace_kartlar_10-12-314-1071-ekmek-taze-pastane-uxfcruxfcnleri-ve-taze-kek-imalatux131--gb-statuxfcsuxfc-iux307tuxfc--birincilikincil-arketip-d--duruux15f-koru}}

\emph{İç pazara dönük en büyük gıda işvereni (190 bin kişi, çoğu küçük fırın). Dondurulmuş hamur ürünü ihracatı AB'ye \$50 mn; S4'te \%9 sanayi bileşeni, S2'de menşe (un) sorunu.}

\textbf{1. Kimlik.} Rev. 2.1: 10.71 (n:1); CPA 10.71. Başlıca HS4 (TR ihracatı 2024): 1905 \%100. BEC payları (ticaret ağırlıklı): ara \%0 / sermaye \%0 / tüketim \%100.

\textbf{2. Ölçek}

\begingroup\gbtablo\small

\begin{longtable}[]{@{}
  >{\raggedright\arraybackslash}p{(\columnwidth - 4\tabcolsep) * \real{0.2945}}
  >{\raggedright\arraybackslash}p{(\columnwidth - 4\tabcolsep) * \real{0.3739}}
  >{\raggedright\arraybackslash}p{(\columnwidth - 4\tabcolsep) * \real{0.3316}}@{}}
\toprule\noalign{}
\begin{minipage}[b]{\linewidth}\raggedright
\textbf{Gösterge}
\end{minipage} & \begin{minipage}[b]{\linewidth}\raggedright
\textbf{Değer}
\end{minipage} & \begin{minipage}[b]{\linewidth}\raggedright
\textbf{Etiket ve kaynak}
\end{minipage} \\
\midrule\noalign{}
\endhead
\bottomrule\noalign{}
\endlastfoot
Üretim değeri; katma değer & \$6.566 mn; \$2.097 mn (GSYH'nin \%0,15) & TÜİK YSHİ 2024 \etiketV{} \\
Çalışan; girişim (0 / 1--9 / 10--49 / 50--249 / 250+) & 189.869; 33.744 (524 / 31.889 / 3.178 / 304 / 40) & TÜİK YSHİ ve İş Kayıtları 2024 \etiketV{}; gizli büyüklük hücreleri \etiketT{} \\
250+ girişimlerin istihdam payı; AB'ye göre kişi başı KD & \%5; --- & \etiketT{} (AB27 SBS 2023, nominal kur) \\
İlk 3 İBBS-2 bölgesi & İstanbul \%11; İzmir \%8; Kocaeli, Sakarya, Düzce, Bolu, Yalova \%7 & \etiketT{} iş kayıtları + SGK 2025, RAS \\
Yabancı sermaye; çapa yatırımcılar & Uno, Simit Sarayı, Ülker (Bizim), dondurulmuş hamurda Hacıbaba, Yuvam \etiketH{} & \etiketH{} (firma bilgisi; pay verisi yok: VB-30) \\
\end{longtable}

\endgroup

\textbf{3. Dış ticaret}

\begingroup\gbtablo\small

\begin{longtable}[]{@{}
  >{\raggedright\arraybackslash}p{(\columnwidth - 2\tabcolsep) * \real{0.4824}}
  >{\raggedright\arraybackslash}p{(\columnwidth - 2\tabcolsep) * \real{0.5176}}@{}}
\toprule\noalign{}
\begin{minipage}[b]{\linewidth}\raggedright
\textbf{Gösterge}
\end{minipage} & \begin{minipage}[b]{\linewidth}\raggedright
\textbf{Değer (BACI 2024, \$; \etiketV{}, eşleme \etiketT{})}
\end{minipage} \\
\midrule\noalign{}
\endhead
\bottomrule\noalign{}
\endlastfoot
İhracat (2024); 2020--24 değişimi & \$264 mn; \%62 \\
AB payı; ilk 5 AB pazarı & \%19; DEU \%40; NLD \%11; BEL \%10; ROU \%9; BGR \%8 \\
İlk 5 AB dışı pazar & IRQ \%32; USA \%8; AZE \%5; YEM \%5; LBY \%4 \\
İthalat; AB ve Çin payı; ilk kaynaklar & \$29 mn; AB \%61, Çin \%1; MYS \%30; POL \%21; CZE \%8; BEL \%7; DEU \%6 \\
Dış ticaret dengesi & \$+235 mn \\
İthal girdi içeriği (I-O, toplam) / KD÷üretim & \%20 / \%32 \\
RCA dünya (2020 → 2024) / AB pazarında & 1,54 → 1,45 / 2,06 \\
Birim değer ÷ AB pazarı ortalaması & 0,89 \\
AB pazar payı (AB dışı tedarik içinde, 2020 → 2024); TR'nin sırası & \%6,0 → \%8,2; 2. \\
\end{longtable}

\endgroup

\textbf{4. AB tarafı maruziyet}

\begin{itemize}
\tightlist
\item
  \kod{τ\_\allowbreak{}AB} (S4 artışı, rota A, TR→AB 2025 ağırlıklı): \%9,0; 1/98 kaybıyla \%9,0; tam MFN ortalaması \%9,0. Rejim: GB \%0, İTÜ \%100, tarım \%0 \etiketV{}/\etiketT{}.
\item
  Tepe satırlar (S4 artışı ≥ \%10; en büyük 3): yok.
\item
  AB pazarında AB dışı rakipler (tercih durumu): GBR \%43; CHN \%6; CHE \%5; SRB \%4; THA \%3. Tercihli rakip payı \%66.
\item
  Menşe (S2 / S3): risk O/O; \kod{c\_\allowbreak{}RoO} \%3,0 / \%3,5; \kod{p\_\allowbreak{}uyum} 0,80 / 0,75 \etiketH{}. Kural: 1905: Fasıl 11 malzemesi kullanılamaz (un menşeli olmalı; DİR'deki Rus buğdayı unu menşesiz).
\item
  EUDR kapsamındaki AB ihracatı payı: \%0; SKDM: \%0.
\item
  S2'de menşe: Fasıl 11 malzemesi kullanılamıyor; DİR'deki Rus buğdayı unu menşesiz. S2'nin kaybı (−\$11 mn AB ihracatı) S4'ünkinin (−\$16 mn) üçte ikisinden fazla.
\end{itemize}

\textbf{5. Türkiye tarafı maruziyet}

\begin{itemize}
\tightlist
\item
  Uygulanan tarife (2026, ithalat ağırlıklı): MFN \%7,1, İGV dahil \%7,1 \etiketV{}. S4'te AB menşeline (c2): \%7,1.
\item
  Kritik girdiler (HS4 1001, 1101), Türkiye'nin toplam ithalatı 2024: \$1.302 mn (AB \%1; RUS \%87, UKR \%10, SYR \%1) \etiketV{}.
\end{itemize}

\textbf{6. Kabiliyet}

\begin{itemize}
\tightlist
\item
  KS 55,5 (sıra 57/230; ağırlık duyarlılığında P10--P90 19--122). Bileşenler: RCA 47, PCI 38 (PCI −0,18), yurt içi KD 79, Ar-Ge 25, tedarikçi 84, verimlilik ---, firma 90 \etiketT{}.
\item
  Ürün uzayında atlanabilir ürünler (sınıf içi; yoğunluk üst çeyrek, PCI \textgreater{} sınıf): yok \etiketT{}.
\item
  \textbf{Kırmızı çizgi yetkinlikleri:} Ekmeklik buğday unu arzı ve fiyat istikrarı (Halk Ekmek ve belediye fırınları).
\end{itemize}

\textbf{7. Stratejik önem.} SS 67,0 (en yüksek boyut: gıda güvenliği 90; ağırlıklı ortalama 13,5) \etiketH{}. Ekmek en geniş tabanlı temel gıda; fiyatı belediye düzenlemesine tabi.

\textbf{8. Makro duyarlılık}

\begingroup\gbtablo\small

\begin{longtable}[]{@{}lr@{}}
\toprule\noalign{}
\textbf{Gösterge} & \textbf{Değer (\etiketT{})} \\
\midrule\noalign{}
\endhead
\bottomrule\noalign{}
\endlastfoot
TÜFE ağırlığı (doğrudan, COICOP 2018 eşlemesi); tüketimde ithal pay & \%2,37; \%2 \\
ÜFE vekili (imalat üretimindeki pay) & \%0,98 \\
EDE: 1 puan tarife, tüm menşe / yalnız AB (bp; ρ = 1) & 0,052 / 0,032 \\
Geçişkenlik ρ (senaryolarda) & 0,6--1,0 (P50 0,8) \\
Piyasa yoğunlaşması (250+ istihdam payı) & \%5 \\
CKE: tarife (\$ mn / bp) ve mali (\$ mn / mlr TL); 0--100 puan & 184,8 ve 0,56; 13 \\
\end{longtable}

\endgroup

\textbf{9. Senaryo etkileri, T+5, P50 (P10 / P90)}

\begingroup\gbtablo\small

\begin{longtable}[]{@{}
  >{\raggedright\arraybackslash}p{(\columnwidth - 6\tabcolsep) * \real{0.3667}}
  >{\raggedleft\arraybackslash}p{(\columnwidth - 6\tabcolsep) * \real{0.2111}}
  >{\raggedleft\arraybackslash}p{(\columnwidth - 6\tabcolsep) * \real{0.2111}}
  >{\raggedleft\arraybackslash}p{(\columnwidth - 6\tabcolsep) * \real{0.2111}}@{}}
\toprule\noalign{}
\begin{minipage}[b]{\linewidth}\raggedright
\textbf{Değişken}
\end{minipage} & \begin{minipage}[b]{\linewidth}\raggedleft
\textbf{S2}
\end{minipage} & \begin{minipage}[b]{\linewidth}\raggedleft
\textbf{S3}
\end{minipage} & \begin{minipage}[b]{\linewidth}\raggedleft
\textbf{S4}
\end{minipage} \\
\midrule\noalign{}
\endhead
\bottomrule\noalign{}
\endlastfoot
AB'ye ihracat (\$ mn) & −11 (−25 / +1) & −12 (−22 / −2) & −16 (−25 / −7) \\
Toplam ihracat (\$ mn) & −8 (−18 / +1) & −8 (−16 / −2) & −12 (−20 / −5) \\
Nihai mal ithalatı (\$ mn) & −1 (−4 / +1) & −2 (−4 / +1) & −3 (−5 / −1) \\
Cari denge (\$ mn) & −5 (−11 / 0) & −5 (−9 / −2) & −8 (−12 / −4) \\
Üretim (\$ mn) & −7 (−14 / −1) & −6 (−12 / −2) & −9 (−15 / −5) \\
İstihdam (kişi) & −131 (−294 / −12) & −130 (−237 / −47) & −189 (−310 / −92) \\
TÜFE (bp, birinci tur) & +0,04 (+0,03 / +0,05) & +0,08 (+0,06 / +0,09) & +0,21 (+0,18 / +0,24) \\
Yatırım ve DYY (bütün senaryolar) & Etkisi yok & & \\
\end{longtable}

\endgroup

\emph{Kaynak:} \etiketT{} f3e: kısmi denge, toplamlar Faz 2B GE dağılımına kalibre; S4'te 1/98 kaybı olasılığı 0,25.

\textbf{12. Erken uyarı ve tetikler}

\begingroup\gbtablo\small

\begin{longtable}[]{@{}
  >{\raggedright\arraybackslash}p{(\columnwidth - 4\tabcolsep) * \real{0.3054}}
  >{\raggedright\arraybackslash}p{(\columnwidth - 4\tabcolsep) * \real{0.3113}}
  >{\raggedright\arraybackslash}p{(\columnwidth - 4\tabcolsep) * \real{0.3833}}@{}}
\toprule\noalign{}
\begin{minipage}[b]{\linewidth}\raggedright
\textbf{Gösterge}
\end{minipage} & \begin{minipage}[b]{\linewidth}\raggedright
\textbf{Eşik}
\end{minipage} & \begin{minipage}[b]{\linewidth}\raggedright
\textbf{Tetiklenecek önlem}
\end{minipage} \\
\midrule\noalign{}
\endhead
\bottomrule\noalign{}
\endlastfoot
Ekmek ve tahıl TÜFE alt endeksi & 3 ay üst üste genel TÜFE + 5 puan & TMO'dan değirmenlere satış fiyatı ayarı \\
\end{longtable}

\endgroup

\textbf{13. Kazanan ve kaybedenler; bölgesel ve sosyal etki; telafi.} S4'te küçük kayıp (P50 −189 kişi; emek yoğun, kişi başı üretim düşük). Bölge: İstanbul (\%11), İzmir.

\textbf{14. Riskler ve azaltım}

\begin{itemize}
\tightlist
\item
  Menşe uyumlu DİR uygulanmazsa S2'de dondurulmuş ürün ihracatı kaybı → E-10-02.
\end{itemize}

\textbf{15. Güven düzeyi; veri boşlukları; \etiketD{}.} Orta (sınıf üretimi YSHİ'de geniş: küçük fırınlar).

\hypertarget{nace_kartlar_10-12-315-1072-peksimet-ve-biskuxfcvi-imalatux131-dayanux131klux131-pastane-uxfcruxfcnleri-ve-dayanux131klux131-kek-imalatux131--gb-statuxfcsuxfc-iux307tuxfc--birincilikincil-arketip-b--duruux15f-geliux15ftir}{%
\subsection{3.15 10.72 Peksimet ve bisküvi imalatı; dayanıklı pastane ürünleri ve dayanıklı kek imalatı --- GB statüsü: İTÜ \textbar{} Birincil/ikincil arketip: B \textbar{} Duruş: Geliştir}\label{nace_kartlar_10-12-315-1072-peksimet-ve-biskuxfcvi-imalatux131-dayanux131klux131-pastane-uxfcruxfcnleri-ve-dayanux131klux131-kek-imalatux131--gb-statuxfcsuxfc-iux307tuxfc--birincilikincil-arketip-b--duruux15f-geliux15ftir}}

\emph{Küresel rekabetçi (bisküvide dünya sekizincisi), AB payı \%14. S4'te \%9 sanayi bileşeniyle AB'ye ihracat −\$51 mn. S2'de menşe kuralı (yerli un) kaybın üçte ikisinden fazlasını bırakıyor; çözüm DİR'de eşdeğer eşya ve PSR.}

\textbf{1. Kimlik.} Rev. 2.1: 10.72 (n:1); CPA 10.72. Başlıca HS4 (TR ihracatı 2024): 1905 \%100. BEC payları (ticaret ağırlıklı): ara \%0 / sermaye \%0 / tüketim \%100.

\textbf{2. Ölçek}

\begingroup\gbtablo\small

\begin{longtable}[]{@{}
  >{\raggedright\arraybackslash}p{(\columnwidth - 4\tabcolsep) * \real{0.2963}}
  >{\raggedright\arraybackslash}p{(\columnwidth - 4\tabcolsep) * \real{0.3697}}
  >{\raggedright\arraybackslash}p{(\columnwidth - 4\tabcolsep) * \real{0.3340}}@{}}
\toprule\noalign{}
\begin{minipage}[b]{\linewidth}\raggedright
\textbf{Gösterge}
\end{minipage} & \begin{minipage}[b]{\linewidth}\raggedright
\textbf{Değer}
\end{minipage} & \begin{minipage}[b]{\linewidth}\raggedright
\textbf{Etiket ve kaynak}
\end{minipage} \\
\midrule\noalign{}
\endhead
\bottomrule\noalign{}
\endlastfoot
Üretim değeri; katma değer & \$5.738 mn; \$1.548 mn (GSYH'nin \%0,11) & TÜİK YSHİ 2024 \etiketV{} \\
Çalışan; girişim (0 / 1--9 / 10--49 / 50--249 / 250+) & 42.160; 5.211 (70 / 5.139 / 242 / 32 / 11) & TÜİK YSHİ ve İş Kayıtları 2024 \etiketV{}; gizli büyüklük hücreleri \etiketT{} \\
250+ girişimlerin istihdam payı; AB'ye göre kişi başı KD & \%43; 0,78 & \etiketT{} (AB27 SBS 2023, nominal kur) \\
İlk 3 İBBS-2 bölgesi & İstanbul \%16; Konya, Karaman \%15; Bursa, Eskişehir, Bilecik \%15 & \etiketT{} iş kayıtları + SGK 2025, RAS \\
Yabancı sermaye; çapa yatırımcılar & Ülker (Yıldız Holding/pladis), Eti, Şölen, Tayaş, Elvan, Bifa, Torku \etiketH{} & \etiketH{} (firma bilgisi; pay verisi yok: VB-30) \\
\end{longtable}

\endgroup

\textbf{3. Dış ticaret}

\begingroup\gbtablo\small

\begin{longtable}[]{@{}
  >{\raggedright\arraybackslash}p{(\columnwidth - 2\tabcolsep) * \real{0.4787}}
  >{\raggedright\arraybackslash}p{(\columnwidth - 2\tabcolsep) * \real{0.5213}}@{}}
\toprule\noalign{}
\begin{minipage}[b]{\linewidth}\raggedright
\textbf{Gösterge}
\end{minipage} & \begin{minipage}[b]{\linewidth}\raggedright
\textbf{Değer (BACI 2024, \$; \etiketV{}, eşleme \etiketT{})}
\end{minipage} \\
\midrule\noalign{}
\endhead
\bottomrule\noalign{}
\endlastfoot
İhracat (2024); 2020--24 değişimi & \$1.174 mn; \%45 \\
AB payı; ilk 5 AB pazarı & \%14; DEU \%32; BEL \%19; ROU \%11; NLD \%9; POL \%7 \\
İlk 5 AB dışı pazar & IRQ \%33; SAU \%6; USA \%5; LBY \%4; GBR \%4 \\
İthalat; AB ve Çin payı; ilk kaynaklar & \$89 mn; AB \%75, Çin \%0; ITA \%19; DEU \%15; POL \%13; UKR \%8; MYS \%7 \\
Dış ticaret dengesi & \$+1.085 mn \\
İthal girdi içeriği (I-O, toplam) / KD÷üretim & \%20 / \%27 \\
RCA dünya (2020 → 2024) / AB pazarında & 3,52 → 3,10 / 3,15 \\
Birim değer ÷ AB pazarı ortalaması & 0,95 \\
AB pazar payı (AB dışı tedarik içinde, 2020 → 2024); TR'nin sırası & \%9,4 → \%12,6; 2. \\
\end{longtable}

\endgroup

\textbf{4. AB tarafı maruziyet}

\begin{itemize}
\tightlist
\item
  \kod{τ\_\allowbreak{}AB} (S4 artışı, rota A, TR→AB 2025 ağırlıklı): \%9,0; 1/98 kaybıyla \%9,0; tam MFN ortalaması \%9,0. Rejim: GB \%0, İTÜ \%100, tarım \%0 \etiketV{}/\etiketT{}.
\item
  Tepe satırlar (S4 artışı ≥ \%10; en büyük 3): yok (AB ihracatının \%1'inden az).
\item
  AB pazarında AB dışı rakipler (tercih durumu): GBR \%40; UKR \%10; SRB \%5; CHE \%4; THA \%4. Tercihli rakip payı \%65.
\item
  Menşe (S2 / S3): risk O/O; \kod{c\_\allowbreak{}RoO} \%3,0 / \%3,5; \kod{p\_\allowbreak{}uyum} 0,80 / 0,75 \etiketH{}. Kural: 1905: Fasıl 11 malzemesi kullanılamaz (un menşeli olmalı); şeker ağırlık sınırı.
\item
  EUDR kapsamındaki AB ihracatı payı: \%0; SKDM: \%0.
\item
  Menşe kuralı: 1905'te Fasıl 11 malzemesi kullanılamaz ve şeker ağırlık sınırı var. İhracatçıların DİR'de ithal buğday ve şeker kullanması bugün sorun değil, çünkü GB'de serbest dolaşım için A.TR yeterli.
\item
  AB pazarındaki AB dışı rakiplerin hepsi tercihli: İngiltere (TCA), Ukrayna (DCFTA), Sırbistan (İOA). S4'te Türkiye bu grupta MFN ödeyen tek büyük tedarikçi olur.
\end{itemize}

\textbf{5. Türkiye tarafı maruziyet}

\begin{itemize}
\tightlist
\item
  Uygulanan tarife (2026, ithalat ağırlıklı): MFN \%8,5, İGV dahil \%9,7 \etiketV{}. S4'te AB menşeline (c2): \%9,8.
\item
  Kritik girdiler (HS4 1001, 1101, 1701, 1804), Türkiye'nin toplam ithalatı 2024: \$2.015 mn (AB \%15; RUS \%56, MAR \%11, BGR \%9) \etiketV{}.
\item
  Girdiler: buğday (DİR'de Rusya), şeker (DİR'de ham kamış şekeri), kakao yağı ve tozu (AB'den; S4'te \%7,7--9,6). S4'te kakao ara ürünlerinde MFN \%0 (E-10.82-01) bu sınıfa da yarıyor.
\end{itemize}

\textbf{6. Kabiliyet}

\begin{itemize}
\tightlist
\item
  KS 54,9 (sıra 64/230; ağırlık duyarlılığında P10--P90 23--121). Bileşenler: RCA 67, PCI 25 (PCI −0,49), yurt içi KD 66, Ar-Ge 25, tedarikçi 84, verimlilik 89, firma 55 \etiketT{}.
\item
  Ürün uzayında atlanabilir ürünler (sınıf içi; yoğunluk üst çeyrek, PCI \textgreater{} sınıf): yok \etiketT{}.
\item
  \textbf{Kırmızı çizgi yetkinlikleri:} Marka ve ambalaj tasarımı, raf ömrü teknolojisi, Irak ve Körfez'deki dağıtım ağları.
\item
  \textbf{Geliştirme hedefi:} T+5: AB ihracatı ≥ T−12 düzeyi (menşe çözümüyle); 2035: Kuzey Amerika'da marka payı.
\end{itemize}

\textbf{7. Stratejik önem.} SS 22,4 (en yüksek boyut: gıda güvenliği 30; ağırlıklı ortalama 4,5) \etiketH{}.

\textbf{8. Makro duyarlılık}

\begingroup\gbtablo\small

\begin{longtable}[]{@{}lr@{}}
\toprule\noalign{}
\textbf{Gösterge} & \textbf{Değer (\etiketT{})} \\
\midrule\noalign{}
\endhead
\bottomrule\noalign{}
\endlastfoot
TÜFE ağırlığı (doğrudan, COICOP 2018 eşlemesi); tüketimde ithal pay & \%1,02; \%2 \\
ÜFE vekili (imalat üretimindeki pay) & \%0,86 \\
EDE: 1 puan tarife, tüm menşe / yalnız AB (bp; ρ = 1) & 0,026 / 0,019 \\
Geçişkenlik ρ (senaryolarda) & 0,6--1,0 (P50 0,8) \\
Piyasa yoğunlaşması (250+ istihdam payı) & \%43 \\
CKE: tarife (\$ mn / bp) ve mali (\$ mn / mlr TL); 0--100 puan & 1.103,7 ve 2,27; 25 \\
\end{longtable}

\endgroup

\textbf{9. Senaryo etkileri, T+5, P50 (P10 / P90)}

\begingroup\gbtablo\small

\begin{longtable}[]{@{}
  >{\raggedright\arraybackslash}p{(\columnwidth - 6\tabcolsep) * \real{0.2112}}
  >{\raggedleft\arraybackslash}p{(\columnwidth - 6\tabcolsep) * \real{0.4853}}
  >{\raggedleft\arraybackslash}p{(\columnwidth - 6\tabcolsep) * \real{0.1518}}
  >{\raggedleft\arraybackslash}p{(\columnwidth - 6\tabcolsep) * \real{0.1518}}@{}}
\toprule\noalign{}
\begin{minipage}[b]{\linewidth}\raggedright
\textbf{Değişken}
\end{minipage} & \begin{minipage}[b]{\linewidth}\raggedleft
\textbf{S2}
\end{minipage} & \begin{minipage}[b]{\linewidth}\raggedleft
\textbf{S3}
\end{minipage} & \begin{minipage}[b]{\linewidth}\raggedleft
\textbf{S4}
\end{minipage} \\
\midrule\noalign{}
\endhead
\bottomrule\noalign{}
\endlastfoot
AB'ye ihracat (\$ mn) & −35 (−78 / +3) & −37 (−69 / −7) & −51 (−81 / −21) \\
Toplam ihracat (\$ mn) & −25 (−58 / +2) & −26 (−50 / −5) & −40 (−66 / −17) \\
Nihai mal ithalatı (\$ mn) & −5 (−15 / +5) & −7 (−17 / +3) & −13 (−24 / −3) \\
Cari denge (\$ mn) & −16 (−34 / −2) & −16 (−27 / −6) & −21 (−34 / −11) \\
Üretim (\$ mn) & −20 (−43 / −3) & −20 (−34 / −8) & −27 (−43 / −14) \\
İstihdam (kişi) & −104 (−218 / −13) & −99 (−174 / −40) & −135 (−221 / −70) \\
TÜFE (bp, birinci tur) & +0,03 (+0,02 / +0,03) & +0,05 (+0,05 / +0,06) & +0,17 (+0,15 / +0,19) \\
Yatırım ve DYY (bütün senaryolar) & Yıldız Holding'in yurt dışı (pladis) tesisleri AB arzını karşılayabilir: S4'te üretim kayması riski \etiketH{} & & \\
\end{longtable}

\endgroup

\emph{Kaynak:} \etiketT{} f3e: kısmi denge, toplamlar Faz 2B GE dağılımına kalibre; S4'te 1/98 kaybı olasılığı 0,25.

\textbf{10. Politika paketi}

\begingroup\gbtablo\footnotesize

\begin{longtable}[]{@{}
  >{\raggedright\arraybackslash}p{(\columnwidth - 16\tabcolsep) * \real{0.1117}}
  >{\raggedright\arraybackslash}p{(\columnwidth - 16\tabcolsep) * \real{0.0786}}
  >{\raggedright\arraybackslash}p{(\columnwidth - 16\tabcolsep) * \real{0.1383}}
  >{\raggedright\arraybackslash}p{(\columnwidth - 16\tabcolsep) * \real{0.1411}}
  >{\raggedright\arraybackslash}p{(\columnwidth - 16\tabcolsep) * \real{0.0949}}
  >{\raggedright\arraybackslash}p{(\columnwidth - 16\tabcolsep) * \real{0.0668}}
  >{\raggedright\arraybackslash}p{(\columnwidth - 16\tabcolsep) * \real{0.1654}}
  >{\raggedright\arraybackslash}p{(\columnwidth - 16\tabcolsep) * \real{0.1256}}
  >{\raggedright\arraybackslash}p{(\columnwidth - 16\tabcolsep) * \real{0.0778}}@{}}
\toprule\noalign{}
\begin{minipage}[b]{\linewidth}\raggedright
\textbf{Eylem ID}
\end{minipage} & \begin{minipage}[b]{\linewidth}\raggedright
\textbf{Araç}
\end{minipage} & \begin{minipage}[b]{\linewidth}\raggedright
\textbf{Parametre}
\end{minipage} & \begin{minipage}[b]{\linewidth}\raggedright
\textbf{Sorumlu}
\end{minipage} & \begin{minipage}[b]{\linewidth}\raggedright
\textbf{Hukuki araç}
\end{minipage} & \begin{minipage}[b]{\linewidth}\raggedright
\textbf{Bütçe}
\end{minipage} & \begin{minipage}[b]{\linewidth}\raggedright
\textbf{Başlangıç--bitiş; sunset}
\end{minipage} & \begin{minipage}[b]{\linewidth}\raggedright
\textbf{Performans koşulu}
\end{minipage} & \begin{minipage}[b]{\linewidth}\raggedright
\textbf{KPI}
\end{minipage} \\
\midrule\noalign{}
\endhead
\bottomrule\noalign{}
\endlastfoot
E-10.72-01 & Fasıl 19 PSR talebi & 1905: `başlık değişikliği; menşesiz tahıl ve un fabrika çıkış fiyatının \%15'ine kadar' (genel toleransla uyumlu) & Ticaret Bakanlığı (Uluslararası Anlaşmalar ve AB GM) & S2 menşe protokolü & İdari & T−18 → T; sunset: --- & --- & Kabul edilen PSR \\
\end{longtable}

\endgroup

\textbf{11. Müzakere pozisyonu}

\begin{itemize}
\tightlist
\item
  AB ile savunma: İTÜ'de sanayi bileşeninin sıfırı sürsün; Fasıl 19 PSR toleransı ya da eşdeğer eşya kabulü.
\item
  Üçüncü ülke: Körfez STA'sında bisküvi ve şekerlemede sıfır.
\end{itemize}

\textbf{12. Erken uyarı ve tetikler}

\begingroup\gbtablo\small

\begin{longtable}[]{@{}
  >{\raggedright\arraybackslash}p{(\columnwidth - 4\tabcolsep) * \real{0.4669}}
  >{\raggedright\arraybackslash}p{(\columnwidth - 4\tabcolsep) * \real{0.1162}}
  >{\raggedright\arraybackslash}p{(\columnwidth - 4\tabcolsep) * \real{0.4168}}@{}}
\toprule\noalign{}
\begin{minipage}[b]{\linewidth}\raggedright
\textbf{Gösterge}
\end{minipage} & \begin{minipage}[b]{\linewidth}\raggedright
\textbf{Eşik}
\end{minipage} & \begin{minipage}[b]{\linewidth}\raggedright
\textbf{Tetiklenecek önlem}
\end{minipage} \\
\midrule\noalign{}
\endhead
\bottomrule\noalign{}
\endlastfoot
AB'ye bisküvi ihracatında menşe ispatlı pay (S2'den sonra) & \textless{} \%80 & E-10-02 kapsamı genişletilir; TMO yerli buğday arzı \\
AB'ye bisküvi ihracatı (aylık; S4) & Yıllık −\%20 & E-10-05 \\
\end{longtable}

\endgroup

\textbf{13. Kazanan ve kaybedenler; bölgesel ve sosyal etki; telafi.} Kaybeden (S4 P50): −135 kişi; bölgeler İstanbul, Konya--Karaman (Torku, Eti Karaman \etiketH{}), Bursa--Eskişehir (Eti). Telafi: E-10-05.

\textbf{14. Riskler ve azaltım}

\begin{itemize}
\tightlist
\item
  Üretimin pladis'in AB içindeki tesislerine kayması → S2 önceliği ve E-10-02.
\end{itemize}

\textbf{15. Güven düzeyi; veri boşlukları; \etiketD{}.} Orta (c\_RoO ve yerli buğday fiyat farkı \etiketT{}; VB-33).

\hypertarget{nace_kartlar_10-12-316-1073-makarna-ux15fehriye-kuskus-ve-benzeri-unlu-mamullerin-imalatux131--gb-statuxfcsuxfc-iux307tuxfc--birincilikincil-arketip-bd--duruux15f-geliux15ftir}{%
\subsection{3.16 10.73 Makarna, şehriye, kuskus ve benzeri unlu mamullerin imalatı --- GB statüsü: İTÜ \textbar{} Birincil/ikincil arketip: B/D \textbar{} Duruş: Geliştir}\label{nace_kartlar_10-12-316-1073-makarna-ux15fehriye-kuskus-ve-benzeri-unlu-mamullerin-imalatux131--gb-statuxfcsuxfc-iux307tuxfc--birincilikincil-arketip-bd--duruux15f-geliux15ftir}}

\emph{Makarnada dünya ikincisi (\%12,3; tonajda dörtte bir). İhracatın \%98'i AB dışına (Afrika, Venezuela, Japonya). GB etkisi küçük; menşe kuralı ithal makarnalık buğdaya izin veriyor (D).}

\textbf{1. Kimlik.} Rev. 2.1: 10.73 (n:1); CPA 10.73. Başlıca HS4 (TR ihracatı 2024): 1902 \%100. BEC payları (ticaret ağırlıklı): ara \%0 / sermaye \%0 / tüketim \%100.

\textbf{2. Ölçek}

\begingroup\gbtablo\small

\begin{longtable}[]{@{}
  >{\raggedright\arraybackslash}p{(\columnwidth - 4\tabcolsep) * \real{0.2949}}
  >{\raggedright\arraybackslash}p{(\columnwidth - 4\tabcolsep) * \real{0.3729}}
  >{\raggedright\arraybackslash}p{(\columnwidth - 4\tabcolsep) * \real{0.3321}}@{}}
\toprule\noalign{}
\begin{minipage}[b]{\linewidth}\raggedright
\textbf{Gösterge}
\end{minipage} & \begin{minipage}[b]{\linewidth}\raggedright
\textbf{Değer}
\end{minipage} & \begin{minipage}[b]{\linewidth}\raggedright
\textbf{Etiket ve kaynak}
\end{minipage} \\
\midrule\noalign{}
\endhead
\bottomrule\noalign{}
\endlastfoot
Üretim değeri; katma değer & \$1.959 mn; \$409 mn (GSYH'nin \%0,03) & TÜİK YSHİ 2024 \etiketV{} \\
Çalışan; girişim (0 / 1--9 / 10--49 / 50--249 / 250+) & 7.040; 429 (40 / 391 / 28 / 11 / 7) & TÜİK YSHİ ve İş Kayıtları 2024 \etiketV{}; gizli büyüklük hücreleri \etiketT{} \\
250+ girişimlerin istihdam payı; AB'ye göre kişi başı KD & \%58; --- & \etiketT{} (AB27 SBS 2023, nominal kur) \\
İlk 3 İBBS-2 bölgesi & Konya, Karaman \%23; Adana, Mersin \%16; Gaziantep, Adıyaman, Kilis \%15 & \etiketT{} iş kayıtları + SGK 2025, RAS \\
Yabancı sermaye; çapa yatırımcılar & Filiz, Nuh'un Ankara, Oba, Selva, Barilla (Bursa), Pastavilla, Durum Gıda \etiketH{} & \etiketH{} (firma bilgisi; pay verisi yok: VB-30) \\
\end{longtable}

\endgroup

\textbf{3. Dış ticaret}

\begingroup\gbtablo\small

\begin{longtable}[]{@{}
  >{\raggedright\arraybackslash}p{(\columnwidth - 2\tabcolsep) * \real{0.4877}}
  >{\raggedright\arraybackslash}p{(\columnwidth - 2\tabcolsep) * \real{0.5123}}@{}}
\toprule\noalign{}
\begin{minipage}[b]{\linewidth}\raggedright
\textbf{Gösterge}
\end{minipage} & \begin{minipage}[b]{\linewidth}\raggedright
\textbf{Değer (BACI 2024, \$; \etiketV{}, eşleme \etiketT{})}
\end{minipage} \\
\midrule\noalign{}
\endhead
\bottomrule\noalign{}
\endlastfoot
İhracat (2024); 2020--24 değişimi & \$852 mn; \%20 \\
AB payı; ilk 5 AB pazarı & \%2; DEU \%35; NLD \%14; FRA \%9; AUT \%8; POL \%6 \\
İlk 5 AB dışı pazar & SOM \%15; GHA \%9; VEN \%9; JPN \%8; TGO \%6 \\
İthalat; AB ve Çin payı; ilk kaynaklar & \$17 mn; AB \%86, Çin \%7; ITA \%82; CHN \%7; ROU \%3; EGY \%2; LBN \%1 \\
Dış ticaret dengesi & \$+835 mn \\
İthal girdi içeriği (I-O, toplam) / KD÷üretim & \%20 / \%21 \\
RCA dünya (2020 → 2024) / AB pazarında & 10,68 → 9,09 / 2,63 \\
Birim değer ÷ AB pazarı ortalaması & 0,55 \\
AB pazar payı (AB dışı tedarik içinde, 2020 → 2024); TR'nin sırası & \%9,4 → \%10,5; 3. \\
\end{longtable}

\endgroup

\textbf{4. AB tarafı maruziyet}

\begin{itemize}
\tightlist
\item
  \kod{τ\_\allowbreak{}AB} (S4 artışı, rota A, TR→AB 2025 ağırlıklı): \%7,7; 1/98 kaybıyla \%7,7; tam MFN ortalaması \%34,2. Rejim: GB \%0, İTÜ \%100, tarım \%0 \etiketV{}/\etiketT{}.
\item
  Tepe satırlar (S4 artışı ≥ \%10; en büyük 3): yok.
\item
  AB pazarında AB dışı rakipler (tercih durumu): MAR \%19; VNM \%11; THA \%10; CHN \%10; GBR \%6. Tercihli rakip payı \%58.
\item
  Menşe (S2 / S3): risk D/D; \kod{c\_\allowbreak{}RoO} \%1,5 / \%2,0; \kod{p\_\allowbreak{}uyum} 0,90 / 0,85 \etiketH{}. Kural: 1902: tahıl tamamen elde edilmiş, makarnalık buğday ve türevleri hariç (ithal durum serbest).
\item
  EUDR kapsamındaki AB ihracatı payı: \%0; SKDM: \%0.
\end{itemize}

\textbf{5. Türkiye tarafı maruziyet}

\begin{itemize}
\tightlist
\item
  Uygulanan tarife (2026, ithalat ağırlıklı): MFN \%7,7, İGV dahil \%15,0 \etiketV{}. S4'te AB menşeline (c2): \%14,8.
\item
  Kritik girdiler (HS4 1001), Türkiye'nin toplam ithalatı 2024: \$1.293 mn (AB \%1; RUS \%87, UKR \%10, SYR \%1) \etiketV{}.
\item
  Makarnalık buğday ithalatı (Kanada, Rusya, Kazakistan) DİR'de. GB'den bağımsız.
\end{itemize}

\textbf{6. Kabiliyet}

\begin{itemize}
\tightlist
\item
  KS 42,5 (sıra 187/230; ağırlık duyarlılığında P10--P90 110--216). Bileşenler: RCA 73, PCI 3 (PCI −1,59), yurt içi KD 53, Ar-Ge 25, tedarikçi 84, verimlilik ---, firma 28 \etiketT{}.
\item
  Ürün uzayında atlanabilir ürünler (sınıf içi; yoğunluk üst çeyrek, PCI \textgreater{} sınıf): yok \etiketT{}.
\item
  \textbf{Kırmızı çizgi yetkinlikleri:} Durum işleme ve kurutma teknolojisi; Afrika dağıtım ağları.
\item
  \textbf{Geliştirme hedefi:} T+5: Afrika Kıtasal Serbest Ticaret Bölgesi üyeleriyle STA'lar; 2035: makarnada dünya birinciliği (tonaj).
\end{itemize}

\textbf{7. Stratejik önem.} SS 52,1 (en yüksek boyut: gıda güvenliği 70; ağırlıklı ortalama 10,5) \etiketH{}. Makarna: temel ve ucuz karbonhidrat

\textbf{8. Makro duyarlılık}

\begingroup\gbtablo\small

\begin{longtable}[]{@{}lr@{}}
\toprule\noalign{}
\textbf{Gösterge} & \textbf{Değer (\etiketT{})} \\
\midrule\noalign{}
\endhead
\bottomrule\noalign{}
\endlastfoot
TÜFE ağırlığı (doğrudan, COICOP 2018 eşlemesi); tüketimde ithal pay & \%0,21; \%2 \\
ÜFE vekili (imalat üretimindeki pay) & \%0,29 \\
EDE: 1 puan tarife, tüm menşe / yalnız AB (bp; ρ = 1) & 0,005 / 0,004 \\
Geçişkenlik ρ (senaryolarda) & 0,6--1,0 (P50 0,8) \\
Piyasa yoğunlaşması (250+ istihdam payı) & \%58 \\
CKE: tarife (\$ mn / bp) ve mali (\$ mn / mlr TL); 0--100 puan & 594,7 ve 0,99; 20 \\
\end{longtable}

\endgroup

\textbf{9. Senaryo etkileri, T+5, P50 (P10 / P90)}

\begingroup\gbtablo\small

\begin{longtable}[]{@{}
  >{\raggedright\arraybackslash}p{(\columnwidth - 6\tabcolsep) * \real{0.3667}}
  >{\raggedleft\arraybackslash}p{(\columnwidth - 6\tabcolsep) * \real{0.2111}}
  >{\raggedleft\arraybackslash}p{(\columnwidth - 6\tabcolsep) * \real{0.2111}}
  >{\raggedleft\arraybackslash}p{(\columnwidth - 6\tabcolsep) * \real{0.2111}}@{}}
\toprule\noalign{}
\begin{minipage}[b]{\linewidth}\raggedright
\textbf{Değişken}
\end{minipage} & \begin{minipage}[b]{\linewidth}\raggedleft
\textbf{S2}
\end{minipage} & \begin{minipage}[b]{\linewidth}\raggedleft
\textbf{S3}
\end{minipage} & \begin{minipage}[b]{\linewidth}\raggedleft
\textbf{S4}
\end{minipage} \\
\midrule\noalign{}
\endhead
\bottomrule\noalign{}
\endlastfoot
AB'ye ihracat (\$ mn) & −2 (−4 / 0) & −2 (−4 / 0) & −4 (−6 / −2) \\
Toplam ihracat (\$ mn) & −1 (−3 / 0) & −1 (−3 / 0) & −3 (−5 / −1) \\
Nihai mal ithalatı (\$ mn) & −1 (−2 / +1) & −1 (−3 / +1) & −2 (−4 / −1) \\
Cari denge (\$ mn) & 0 (−1 / 0) & 0 (−1 / 0) & 0 (−1 / 0) \\
Üretim (\$ mn) & −1 (−1 / 0) & 0 (−1 / 0) & −1 (−1 / 0) \\
İstihdam (kişi) & −1 (−2 / 0) & −1 (−2 / 0) & −1 (−2 / 0) \\
TÜFE (bp, birinci tur) & +0,01 (+0,01 / +0,01) & +0,01 (+0,01 / +0,02) & +0,06 (+0,05 / +0,06) \\
Yatırım ve DYY (bütün senaryolar) & Etkisi yok & & \\
\end{longtable}

\endgroup

\emph{Kaynak:} \etiketT{} f3e: kısmi denge, toplamlar Faz 2B GE dağılımına kalibre; S4'te 1/98 kaybı olasılığı 0,25.

\textbf{13. Kazanan ve kaybedenler; bölgesel ve sosyal etki; telafi.} GB'den etkilenen yok.

\textbf{15. Güven düzeyi; veri boşlukları; \etiketD{}.} Yüksek.

\hypertarget{nace_kartlar_10-12-317-1081-ux15feker-imalatux131--gb-statuxfcsuxfc-tarim-198--birincilikincil-arketip---duruux15f-koru}{%
\subsection{3.17 10.81 Şeker imalatı --- GB statüsü: TARIM-1/98 \textbar{} Birincil/ikincil arketip: --- \textbar{} Duruş: Koru}\label{nace_kartlar_10-12-317-1081-ux15feker-imalatux131--gb-statuxfcsuxfc-tarim-198--birincilikincil-arketip---duruux15f-koru}}

\emph{Kotalı iç pazar (Şeker Kanunu); ihracat küçük. GB etkisi yok; ham şeker ithalatı (Fas, Brezilya) DİR'de.}

\textbf{1. Kimlik.} Rev. 2.1: 10.81 (1:1); CPA 10.81. Başlıca HS4 (TR ihracatı 2024): 1701 \%86, 2303 \%13, 1703 \%0, 1702 \%0. BEC payları (ticaret ağırlıklı): ara \%15 / sermaye \%0 / tüketim \%85.

\textbf{2. Ölçek}

\begingroup\gbtablo\small

\begin{longtable}[]{@{}
  >{\raggedright\arraybackslash}p{(\columnwidth - 4\tabcolsep) * \real{0.2753}}
  >{\raggedright\arraybackslash}p{(\columnwidth - 4\tabcolsep) * \real{0.4181}}
  >{\raggedright\arraybackslash}p{(\columnwidth - 4\tabcolsep) * \real{0.3066}}@{}}
\toprule\noalign{}
\begin{minipage}[b]{\linewidth}\raggedright
\textbf{Gösterge}
\end{minipage} & \begin{minipage}[b]{\linewidth}\raggedright
\textbf{Değer}
\end{minipage} & \begin{minipage}[b]{\linewidth}\raggedright
\textbf{Etiket ve kaynak}
\end{minipage} \\
\midrule\noalign{}
\endhead
\bottomrule\noalign{}
\endlastfoot
Üretim değeri; katma değer & \$3.202 mn; \$803 mn (GSYH'nin \%0,06) & TÜİK YSHİ 2024 \etiketV{} \\
Çalışan; girişim (0 / 1--9 / 10--49 / 50--249 / 250+) & 15.900; 24 (0 / 8 / 0 / 3 / 30) & TÜİK YSHİ ve İş Kayıtları 2024 \etiketV{}; gizli büyüklük hücreleri \etiketT{} \\
250+ girişimlerin istihdam payı; AB'ye göre kişi başı KD & \%98; 0,26 & \etiketT{} (AB27 SBS 2023, nominal kur) \\
İlk 3 İBBS-2 bölgesi & Manisa, Afyon, Kütahya, Uşak \%21; Konya, Karaman \%18; Samsun, Tokat, Çorum, Amasya \%13 & \etiketT{} iş kayıtları + SGK 2025, RAS \\
Yabancı sermaye; çapa yatırımcılar & Türkiye Şeker Fabrikaları (özelleştirilmiş tesisler), Konya Şeker (Torku), Kayseri Şeker \etiketH{} & \etiketH{} (firma bilgisi; pay verisi yok: VB-30) \\
\end{longtable}

\endgroup

\textbf{3. Dış ticaret}

\begingroup\gbtablo\small

\begin{longtable}[]{@{}
  >{\raggedright\arraybackslash}p{(\columnwidth - 2\tabcolsep) * \real{0.4751}}
  >{\raggedright\arraybackslash}p{(\columnwidth - 2\tabcolsep) * \real{0.5249}}@{}}
\toprule\noalign{}
\begin{minipage}[b]{\linewidth}\raggedright
\textbf{Gösterge}
\end{minipage} & \begin{minipage}[b]{\linewidth}\raggedright
\textbf{Değer (BACI 2024, \$; \etiketV{}, eşleme \etiketT{})}
\end{minipage} \\
\midrule\noalign{}
\endhead
\bottomrule\noalign{}
\endlastfoot
İhracat (2024); 2020--24 değişimi & \$143 mn; \%426 \\
AB payı; ilk 5 AB pazarı & \%10; ESP \%44; GRC \%30; ITA \%15; IRL \%7; BGR \%2 \\
İlk 5 AB dışı pazar & SYR \%50; IRQ \%35; GEO \%5; MAR \%4; CYP \%2 \\
İthalat; AB ve Çin payı; ilk kaynaklar & \$458 mn; AB \%10, Çin \%0; MAR \%50; BRA \%15; RUS \%12; UKR \%8; POL \%3 \\
Dış ticaret dengesi & \$−315 mn \\
İthal girdi içeriği (I-O, toplam) / KD÷üretim & \%20 / \%25 \\
RCA dünya (2020 → 2024) / AB pazarında & 0,09 → 0,27 / 0,23 \\
Birim değer ÷ AB pazarı ortalaması & 1,00 \\
AB pazar payı (AB dışı tedarik içinde, 2020 → 2024); TR'nin sırası & \%0,0 → \%0,9; 21. \\
\end{longtable}

\endgroup

\textbf{4. AB tarafı maruziyet}

\begin{itemize}
\tightlist
\item
  \kod{τ\_\allowbreak{}AB} (S4 artışı, rota A, TR→AB 2025 ağırlıklı): \%0,0; 1/98 kaybıyla \%0,3; tam MFN ortalaması \%0,3. Rejim: GB \%0, İTÜ \%0, tarım \%100 \etiketV{}/\etiketT{}.
\item
  Tepe satırlar (S4 artışı ≥ \%10; en büyük 3): yok.
\item
  AB pazarında AB dışı rakipler (tercih durumu): BRA \%19; UKR \%14; GBR \%7; EGY \%7; CAN \%6. Tercihli rakip payı \%68.
\item
  Menşe (S2 / S3): risk D/D; \kod{c\_\allowbreak{}RoO} \%1,5 / \%2,0; \kod{p\_\allowbreak{}uyum} 0,90 / 0,85 \etiketH{}. Kural: 1701: Fasıl 17 malzemesi menşeli.
\item
  EUDR kapsamındaki AB ihracatı payı: \%0; SKDM: \%0.
\end{itemize}

\textbf{5. Türkiye tarafı maruziyet}

\begin{itemize}
\tightlist
\item
  Uygulanan tarife (2026, ithalat ağırlıklı): MFN \%116,8, İGV dahil \%116,8 \etiketV{}. S4'te AB menşeline (c2): \%0,0.
\item
  Kritik girdiler (HS4 1701), Türkiye'nin toplam ithalatı 2024: \$389 mn (AB \%12; MAR \%59, BRA \%18, UKR \%9) \etiketV{}.
\item
  Şeker fiyatı İTÜ ihracatçıları için menşe maliyeti: yerli pancar şekeri dünya fiyatının üzerinde. Bu yüzden 10.72 ve 10.82 DİR'de ham kamış şekeri kullanıyor.
\end{itemize}

\textbf{6. Kabiliyet}

\begin{itemize}
\tightlist
\item
  KS 33,5 (sıra 223/230; ağırlık duyarlılığında P10--P90 195--228). Bileşenler: RCA 39, PCI 11 (PCI −0,94), yurt içi KD 62, Ar-Ge 25, tedarikçi 84, verimlilik 8, firma 12 \etiketT{}.
\item
  Ürün uzayında atlanabilir ürünler (sınıf içi; yoğunluk üst çeyrek, PCI \textgreater{} sınıf): yok \etiketT{}.
\end{itemize}

\textbf{7. Stratejik önem.} SS 52,6 (en yüksek boyut: gıda güvenliği 70; ağırlıklı ortalama 12,0) \etiketH{}. Şeker; pancar üreticisi; melastan biyoetanol

\textbf{8. Makro duyarlılık}

\begingroup\gbtablo\small

\begin{longtable}[]{@{}lr@{}}
\toprule\noalign{}
\textbf{Gösterge} & \textbf{Değer (\etiketT{})} \\
\midrule\noalign{}
\endhead
\bottomrule\noalign{}
\endlastfoot
TÜFE ağırlığı (doğrudan, COICOP 2018 eşlemesi); tüketimde ithal pay & \%0,35; \%16 \\
ÜFE vekili (imalat üretimindeki pay) & \%0,48 \\
EDE: 1 puan tarife, tüm menşe / yalnız AB (bp; ρ = 1) & 0,092 / 0,009 \\
Geçişkenlik ρ (senaryolarda) & 0,6--1,0 (P50 0,8) \\
Piyasa yoğunlaşması (250+ istihdam payı) & \%98 \\
CKE: tarife (\$ mn / bp) ve mali (\$ mn / mlr TL); 0--100 puan & 306,3 ve 3,21; 20 \\
\end{longtable}

\endgroup

\textbf{9. Senaryo etkileri, T+5, P50 (P10 / P90)}

\begingroup\gbtablo\small

\begin{longtable}[]{@{}lrrr@{}}
\toprule\noalign{}
\textbf{Değişken} & \textbf{S2} & \textbf{S3} & \textbf{S4} \\
\midrule\noalign{}
\endhead
\bottomrule\noalign{}
\endlastfoot
AB'ye ihracat (\$ mn) & 0 (0 / 0) & 0 (0 / 0) & 0 (0 / 0) \\
Toplam ihracat (\$ mn) & 0 (0 / 0) & 0 (0 / 0) & 0 (0 / 0) \\
Nihai mal ithalatı (\$ mn) & 0 (0 / +3) & 0 (0 / +2) & 0 (0 / 0) \\
Cari denge (\$ mn) & 0 (−2 / 0) & 0 (−1 / 0) & 0 (0 / 0) \\
Üretim (\$ mn) & 0 (−3 / 0) & 0 (−2 / 0) & 0 (0 / 0) \\
İstihdam (kişi) & 0 (−9 / 0) & 0 (−6 / 0) & 0 (0 / 0) \\
TÜFE (bp, birinci tur) & 0,00 (0,00 / 0,00) & 0,00 (0,00 / 0,00) & 0,00 (0,00 / 0,00) \\
Yatırım ve DYY (bütün senaryolar) & --- & & \\
\end{longtable}

\endgroup

\emph{Kaynak:} \etiketT{} f3e: kısmi denge, toplamlar Faz 2B GE dağılımına kalibre; S4'te 1/98 kaybı olasılığı 0,25.

\textbf{13. Kazanan ve kaybedenler; bölgesel ve sosyal etki; telafi.} Etkilenen kesim yok.

\textbf{15. Güven düzeyi; veri boşlukları; \etiketD{}.} Yüksek.

\hypertarget{nace_kartlar_10-12-318-1082-kakao-uxe7ikolata-ve-ux15fekerleme-imalatux131--gb-statuxfcsuxfc-iux307tuxfc--birincilikincil-arketip-bg--duruux15f-geliux15ftir}{%
\subsection{3.18 10.82 Kakao, çikolata ve şekerleme imalatı --- GB statüsü: İTÜ \textbar{} Birincil/ikincil arketip: B/G \textbar{} Duruş: Geliştir}\label{nace_kartlar_10-12-318-1082-kakao-uxe7ikolata-ve-ux15fekerleme-imalatux131--gb-statuxfcsuxfc-iux307tuxfc--birincilikincil-arketip-bg--duruux15f-geliux15ftir}}

\emph{Küresel rekabetçi ihracatçı (\$2,0 mlr; AB payı \%14). S4'te AB'ye ihracat −\$84 mn, ama AB'den kakao ara ürünü ithalatı (\$474 mn) \%8 tarifeye girdiği için iç ikame artı yazıyor. Doğru politika girdide sıfır: kakao yağı, tozu ve hamurunda MFN \%0.}

\textbf{1. Kimlik.} Rev. 2.1: 10.82 (1:1); CPA 10.82. Başlıca HS4 (TR ihracatı 2024): 1806 \%46, 1704 \%45, 1805 \%5, 1804 \%3. BEC payları (ticaret ağırlıklı): ara \%33 / sermaye \%0 / tüketim \%67.

\textbf{2. Ölçek}

\begingroup\gbtablo\small

\begin{longtable}[]{@{}
  >{\raggedright\arraybackslash}p{(\columnwidth - 4\tabcolsep) * \real{0.2801}}
  >{\raggedright\arraybackslash}p{(\columnwidth - 4\tabcolsep) * \real{0.4071}}
  >{\raggedright\arraybackslash}p{(\columnwidth - 4\tabcolsep) * \real{0.3128}}@{}}
\toprule\noalign{}
\begin{minipage}[b]{\linewidth}\raggedright
\textbf{Gösterge}
\end{minipage} & \begin{minipage}[b]{\linewidth}\raggedright
\textbf{Değer}
\end{minipage} & \begin{minipage}[b]{\linewidth}\raggedright
\textbf{Etiket ve kaynak}
\end{minipage} \\
\midrule\noalign{}
\endhead
\bottomrule\noalign{}
\endlastfoot
Üretim değeri; katma değer & \$6.016 mn; \$1.645 mn (GSYH'nin \%0,12) & TÜİK YSHİ 2024 \etiketV{} \\
Çalışan; girişim (0 / 1--9 / 10--49 / 50--249 / 250+) & 36.835; 2.055 (111 / 1.700 / 283 / 64 / 32) & TÜİK YSHİ ve İş Kayıtları 2024 \etiketV{}; gizli büyüklük hücreleri \etiketT{} \\
250+ girişimlerin istihdam payı; AB'ye göre kişi başı KD & \%51; 0,69 & \etiketT{} (AB27 SBS 2023, nominal kur) \\
İlk 3 İBBS-2 bölgesi & İstanbul \%36; Konya, Karaman \%8; Manisa, Afyon, Kütahya, Uşak \%8 & \etiketT{} iş kayıtları + SGK 2025, RAS \\
Yabancı sermaye; çapa yatırımcılar & Ülker (Yıldız), Eti, Şölen, Kervan, Elvan, Tayaş, Nestlé (Bursa), Mondelez (Bursa \etiketD{}) \etiketH{} & \etiketH{} (firma bilgisi; pay verisi yok: VB-30) \\
\end{longtable}

\endgroup

\textbf{3. Dış ticaret}

\begingroup\gbtablo\small

\begin{longtable}[]{@{}
  >{\raggedright\arraybackslash}p{(\columnwidth - 2\tabcolsep) * \real{0.4766}}
  >{\raggedright\arraybackslash}p{(\columnwidth - 2\tabcolsep) * \real{0.5234}}@{}}
\toprule\noalign{}
\begin{minipage}[b]{\linewidth}\raggedright
\textbf{Gösterge}
\end{minipage} & \begin{minipage}[b]{\linewidth}\raggedright
\textbf{Değer (BACI 2024, \$; \etiketV{}, eşleme \etiketT{})}
\end{minipage} \\
\midrule\noalign{}
\endhead
\bottomrule\noalign{}
\endlastfoot
İhracat (2024); 2020--24 değişimi & \$2.032 mn; \%78 \\
AB payı; ilk 5 AB pazarı & \%14; DEU \%25; NLD \%11; ROU \%11; BGR \%8; GRC \%6 \\
İlk 5 AB dışı pazar & IRQ \%14; USA \%13; GBR \%7; SAU \%5; RUS \%5 \\
İthalat; AB ve Çin payı; ilk kaynaklar & \$963 mn; AB \%73, Çin \%1; BGR \%26; NLD \%16; DEU \%11; BEL \%9; GHA \%6 \\
Dış ticaret dengesi & \$+1.069 mn \\
İthal girdi içeriği (I-O, toplam) / KD÷üretim & \%20 / \%27 \\
RCA dünya (2020 → 2024) / AB pazarında & 2,03 → 1,97 / 0,79 \\
Birim değer ÷ AB pazarı ortalaması & 0,94 \\
AB pazar payı (AB dışı tedarik içinde, 2020 → 2024); TR'nin sırası & \%2,8 → \%3,1; 8. \\
\end{longtable}

\endgroup

\textbf{4. AB tarafı maruziyet}

\begin{itemize}
\tightlist
\item
  \kod{τ\_\allowbreak{}AB} (S4 artışı, rota A, TR→AB 2025 ağırlıklı): \%8,4; 1/98 kaybıyla \%8,4; tam MFN ortalaması \%8,8. Rejim: GB \%0, İTÜ \%100, tarım \%0 \etiketV{}/\etiketT{}.
\item
  Tepe satırlar (S4 artışı ≥ \%10; en büyük 3): yok (AB ihracatının \%1'inden az).
\item
  AB pazarında AB dışı rakipler (tercih durumu): CIV \%29; GBR \%14; CHE \%11; GHA \%9; CMR \%6. Tercihli rakip payı \%79.
\item
  Menşe (S2 / S3): risk O/O; \kod{c\_\allowbreak{}RoO} \%2,5 / \%3,5; \kod{p\_\allowbreak{}uyum} 0,85 / 0,75 \etiketH{}. Kural: 1704 ve 1806: başlık değişikliği; menşesiz şeker ve Fasıl 4 malzemesi ağırlıkça \%40 (birlikte \%60) sınırı.
\item
  EUDR kapsamındaki AB ihracatı payı: \%62; SKDM: \%0.
\item
  EUDR: AB ihracatının \%62'si EUDR kapsamında (1806, 1803--1805). 30.12.2026'dan itibaren orta ve büyük AB operatörleri, kakaonun yetiştirildiği parselin coğrafi konumunu istiyor \etiketV{} (Tüzük 2025/2650). GB çıkışından bağımsız, ama hasmane AB'de uygulama sıkılığı artabilir \etiketH{}.
\item
  AB pazarında rakipler Fildişi Sahili ve Gana (EPA; kakao ürünleri), İngiltere ve İsviçre (çikolata). Hepsi tercihli; tercihli rakip payı \%79.
\item
  Menşe: başlık değişikliği, menşesiz şeker ve Fasıl 4 ağırlık sınırı \%40. DİR'deki ham şekerle yapılan şekerlemelerde sınır aşılıyor (O).
\end{itemize}

\textbf{5. Türkiye tarafı maruziyet}

\begin{itemize}
\tightlist
\item
  Uygulanan tarife (2026, ithalat ağırlıklı): MFN \%8,5, İGV dahil \%8,6 \etiketV{}. S4'te AB menşeline (c2): \%8,5.
\item
  Kritik girdiler (HS4 1801, 1803, 1804, 1805, 1701), Türkiye'nin toplam ithalatı 2024: \$1.539 mn (AB \%34; CIV \%21, BGR \%16, MAR \%15) \etiketV{}.
\item
  S4'te AB'den kakao yağı (1804; \$254 mn), kakao tozu (1805; \$149 mn) ve kakao hamurunun (1803; \$70 mn) tarifesi \%7,7--9,6. Bunlar ihracatçının girdisi. Yerli öğütme kapasitesi bu talebin tamamını karşılayamaz \etiketD{}. Matristeki üretim artışı (S4 P50 +56 \$ mn) bu yüzden üst sınır.
\item
  Fildişi Sahili ve Gana'dan doğrudan kakao ürünü ithalatı bugün \%7,7--9,6 MFN ödüyor (AB'nin EPA'sı var, Türkiye'nin yok). Çıkış girdide sıfırı mümkün kılıyor: özerklik temettüsü.
\end{itemize}

\textbf{6. Kabiliyet}

\begin{itemize}
\tightlist
\item
  KS 55,5 (sıra 56/230; ağırlık duyarlılığında P10--P90 24--109). Bileşenler: RCA 62, PCI 34 (PCI −0,27), yurt içi KD 67, Ar-Ge 25, tedarikçi 84, verimlilik 86, firma 54 \etiketT{}.
\item
  Ürün uzayında atlanabilir ürünler (sınıf içi; yoğunluk üst çeyrek, PCI \textgreater{} sınıf): yok \etiketT{}.
\item
  \textbf{Kırmızı çizgi yetkinlikleri:} Fındık--kakao harman bilgisi ve marka; çikolata kalıplama ve kaplama hatları; helal ve Körfez dağıtımı.
\item
  \textbf{Geliştirme hedefi:} T+5: kakaoyu doğrudan kaynaktan (Batı Afrika) alarak maliyette −\%3; 2035: premium segmentte AB dışı (ABD, Körfez) marka payı.
\end{itemize}

\textbf{7. Stratejik önem.} SS 7,5 (en yüksek boyut: gıda güvenliği 10; ağırlıklı ortalama 1,5) \etiketH{}.

\textbf{8. Makro duyarlılık}

\begingroup\gbtablo\small

\begin{longtable}[]{@{}lr@{}}
\toprule\noalign{}
\textbf{Gösterge} & \textbf{Değer (\etiketT{})} \\
\midrule\noalign{}
\endhead
\bottomrule\noalign{}
\endlastfoot
TÜFE ağırlığı (doğrudan, COICOP 2018 eşlemesi); tüketimde ithal pay & \%0,61; \%7 \\
ÜFE vekili (imalat üretimindeki pay) & \%0,90 \\
EDE: 1 puan tarife, tüm menşe / yalnız AB (bp; ρ = 1) & 0,114 / 0,084 \\
Geçişkenlik ρ (senaryolarda) & 0,6--1,0 (P50 0,8) \\
Piyasa yoğunlaşması (250+ istihdam payı) & \%51 \\
CKE: tarife (\$ mn / bp) ve mali (\$ mn / mlr TL); 0--100 puan & 2.661,9 ve 24,12; 57 \\
\end{longtable}

\endgroup

\textbf{9. Senaryo etkileri, T+5, P50 (P10 / P90)}

\begingroup\gbtablo\small

\begin{longtable}[]{@{}
  >{\raggedright\arraybackslash}p{(\columnwidth - 6\tabcolsep) * \real{0.2452}}
  >{\raggedleft\arraybackslash}p{(\columnwidth - 6\tabcolsep) * \real{0.4253}}
  >{\raggedleft\arraybackslash}p{(\columnwidth - 6\tabcolsep) * \real{0.1647}}
  >{\raggedleft\arraybackslash}p{(\columnwidth - 6\tabcolsep) * \real{0.1647}}@{}}
\toprule\noalign{}
\begin{minipage}[b]{\linewidth}\raggedright
\textbf{Değişken}
\end{minipage} & \begin{minipage}[b]{\linewidth}\raggedleft
\textbf{S2}
\end{minipage} & \begin{minipage}[b]{\linewidth}\raggedleft
\textbf{S3}
\end{minipage} & \begin{minipage}[b]{\linewidth}\raggedleft
\textbf{S4}
\end{minipage} \\
\midrule\noalign{}
\endhead
\bottomrule\noalign{}
\endlastfoot
AB'ye ihracat (\$ mn) & −52 (−123 / +5) & −64 (−121 / −13) & −84 (−137 / −36) \\
Toplam ihracat (\$ mn) & −38 (−90 / +4) & −45 (−86 / −9) & −67 (−110 / −28) \\
Nihai mal ithalatı (\$ mn) & −44 (−147 / +48) & −63 (−164 / +28) & −123 (−225 / −26) \\
Cari denge (\$ mn) & +3 (−35 / +50) & +13 (−30 / +64) & +44 (−2 / +94) \\
Üretim (\$ mn) & +3 (−44 / +63) & +16 (−37 / +80) & +56 (−2 / +117) \\
İstihdam (kişi) & +15 (−185 / +266) & +71 (−159 / +341) & +233 (−11 / +494) \\
TÜFE (bp, birinci tur) & +0,11 (+0,10 / +0,14) & +0,22 (+0,18 / +0,26) & +0,64 (+0,55 / +0,73) \\
Yatırım ve DYY (bütün senaryolar) & Nestlé ve Mondelez'in AB'ye yönelik üretim kararı S2'ye duyarlı \etiketH{} & & \\
\end{longtable}

\endgroup

\emph{Kaynak:} \etiketT{} f3e: kısmi denge, toplamlar Faz 2B GE dağılımına kalibre; S4'te 1/98 kaybı olasılığı 0,25.

\textbf{10. Politika paketi}

\begingroup\gbtablo\scriptsize

\begin{longtable}[]{@{}
  >{\raggedright\arraybackslash}p{(\columnwidth - 16\tabcolsep) * \real{0.0983}}
  >{\raggedright\arraybackslash}p{(\columnwidth - 16\tabcolsep) * \real{0.0752}}
  >{\raggedright\arraybackslash}p{(\columnwidth - 16\tabcolsep) * \real{0.1360}}
  >{\raggedright\arraybackslash}p{(\columnwidth - 16\tabcolsep) * \real{0.0952}}
  >{\raggedright\arraybackslash}p{(\columnwidth - 16\tabcolsep) * \real{0.1427}}
  >{\raggedright\arraybackslash}p{(\columnwidth - 16\tabcolsep) * \real{0.0839}}
  >{\raggedright\arraybackslash}p{(\columnwidth - 16\tabcolsep) * \real{0.1461}}
  >{\raggedright\arraybackslash}p{(\columnwidth - 16\tabcolsep) * \real{0.1152}}
  >{\raggedright\arraybackslash}p{(\columnwidth - 16\tabcolsep) * \real{0.1073}}@{}}
\toprule\noalign{}
\begin{minipage}[b]{\linewidth}\raggedright
\textbf{Eylem ID}
\end{minipage} & \begin{minipage}[b]{\linewidth}\raggedright
\textbf{Araç}
\end{minipage} & \begin{minipage}[b]{\linewidth}\raggedright
\textbf{Parametre}
\end{minipage} & \begin{minipage}[b]{\linewidth}\raggedright
\textbf{Sorumlu}
\end{minipage} & \begin{minipage}[b]{\linewidth}\raggedright
\textbf{Hukuki araç}
\end{minipage} & \begin{minipage}[b]{\linewidth}\raggedright
\textbf{Bütçe}
\end{minipage} & \begin{minipage}[b]{\linewidth}\raggedright
\textbf{Başlangıç--bitiş; sunset}
\end{minipage} & \begin{minipage}[b]{\linewidth}\raggedright
\textbf{Performans koşulu}
\end{minipage} & \begin{minipage}[b]{\linewidth}\raggedright
\textbf{KPI}
\end{minipage} \\
\midrule\noalign{}
\endhead
\bottomrule\noalign{}
\endlastfoot
E-10.82-01 & Girdide sıfır: kakao ara ürünleri & KN 1803--1805'te MFN \%7,7--9,6 → \%0, tüm menşeler (T'den itibaren; S2'de de) & Ticaret Bakanlığı & İthalat Rejimi Kararı eki (Cumhurbaşkanı kararı) & Gelir kaybı ≈ 0,55 mlr TL/yıl (AB dışı kakao ürünü ithalatı × \%8) \etiketT{} & T → kalıcı; sunset: Yok & Yurt içi çikolata fiyatında geçiş (izleme) & TÜFE çikolata alt endeksi; S4'te −0,36 bp \\
E-10.82-02 & Menşe uyumlu şeker & AB'ye yönelik şekerlemede yerli pancar şekerinin DİR'de eşdeğer eşya olarak kullanımı (E-10-02) & Ticaret Bakanlığı, Şeker Kurumu & DİR Tebliği & ≈ 0 & T−12 → T+60; sunset: T+60 & --- & AB'ye şekerleme ihracatında menşe ispatlı pay ≥ \%85 \\
\end{longtable}

\endgroup

\textbf{11. Müzakere pozisyonu}

\begin{itemize}
\tightlist
\item
  AB ile savunma: İTÜ sanayi bileşeninin sıfırı; 1704 ve 1806'da şeker ağırlık sınırına \%15 tolerans.
\item
  Üçüncü ülke: Fildişi Sahili ve Gana ile kakao ürünlerinde sıfır tarifeli STA ya da tek taraflı sıfır; Körfez STA'sında çikolata.
\end{itemize}

\textbf{12. Erken uyarı ve tetikler}

\begingroup\gbtablo\small

\begin{longtable}[]{@{}lll@{}}
\toprule\noalign{}
\textbf{Gösterge} & \textbf{Eşik} & \textbf{Tetiklenecek önlem} \\
\midrule\noalign{}
\endhead
\bottomrule\noalign{}
\endlastfoot
AB'ye çikolata ve şekerleme ihracatı (aylık) & Yıllık −\%15 & E-10-05 \\
EUDR nedeniyle AB'de bekletilen sevkiyat & ≥ 3 vaka & E-10-03: tedarikçi veri havuzu \\
\end{longtable}

\endgroup

\textbf{13. Kazanan ve kaybedenler; bölgesel ve sosyal etki; telafi.} Kazanan (S4): kakao öğütücüleri ve iç pazara satan üreticiler (P50 +233 kişi). Kaybeden: AB'ye ihracatçılar. Bölgeler: İstanbul (\%36), Konya--Karaman, Manisa.

\textbf{14. Riskler ve azaltım}

\begin{itemize}
\tightlist
\item
  Girdide sıfırın yerli kakao öğütücülerine etkisi → kademeli geçiş (T+12'ye kadar \%4).
\item
  EUDR uyumsuzluğu → E-10-03.
\end{itemize}

\textbf{15. Güven düzeyi; veri boşlukları; \etiketD{}.} Orta (iç ikamenin kapasite sınırı \etiketD{}; EUDR uygulama sıkılığı \etiketH{}).

\hypertarget{nace_kartlar_10-12-319-1083-kahve-ve-uxe7ayux131n-iux15flenmesi--gb-statuxfcsuxfc-karma--birincilikincil-arketip---duruux15f-koru}{%
\subsection{3.19 10.83 Kahve ve çayın işlenmesi --- GB statüsü: KARMA \textbar{} Birincil/ikincil arketip: --- \textbar{} Duruş: Koru}\label{nace_kartlar_10-12-319-1083-kahve-ve-uxe7ayux131n-iux15flenmesi--gb-statuxfcsuxfc-karma--birincilikincil-arketip---duruux15f-koru}}

\emph{Çay (yerli; Çaykur) ve kahve (ithal). AB'den hazır kahve ithalatı (\$93 mn) S4'te \%9; TÜFE 0,4 bp. GB'nin ihracat etkisi küçük.}

\textbf{1. Kimlik.} Rev. 2.1: 10.83 (1:1); CPA 10.83. Başlıca HS4 (TR ihracatı 2024): 2101 \%51, 0901 \%31, 0902 \%18. BEC payları (ticaret ağırlıklı): ara \%0 / sermaye \%0 / tüketim \%100.

\textbf{2. Ölçek}

\begingroup\gbtablo\small

\begin{longtable}[]{@{}
  >{\raggedright\arraybackslash}p{(\columnwidth - 4\tabcolsep) * \real{0.2737}}
  >{\raggedright\arraybackslash}p{(\columnwidth - 4\tabcolsep) * \real{0.4217}}
  >{\raggedright\arraybackslash}p{(\columnwidth - 4\tabcolsep) * \real{0.3046}}@{}}
\toprule\noalign{}
\begin{minipage}[b]{\linewidth}\raggedright
\textbf{Gösterge}
\end{minipage} & \begin{minipage}[b]{\linewidth}\raggedright
\textbf{Değer}
\end{minipage} & \begin{minipage}[b]{\linewidth}\raggedright
\textbf{Etiket ve kaynak}
\end{minipage} \\
\midrule\noalign{}
\endhead
\bottomrule\noalign{}
\endlastfoot
Üretim değeri; katma değer & \$3.787 mn; \$1.042 mn (GSYH'nin \%0,08) & TÜİK YSHİ 2024 \etiketV{} \\
Çalışan; girişim (0 / 1--9 / 10--49 / 50--249 / 250+) & 19.999; 1.018 (99 / 850 / 129 / 18 / 9) & TÜİK YSHİ ve İş Kayıtları 2024 \etiketV{}; gizli büyüklük hücreleri \etiketT{} \\
250+ girişimlerin istihdam payı; AB'ye göre kişi başı KD & \%63; 0,57 & \etiketT{} (AB27 SBS 2023, nominal kur) \\
İlk 3 İBBS-2 bölgesi & İstanbul \%40; Trabzon, Ordu, Giresun, Rize, Artvin, Gümüşhane \%19; Adana, Mersin \%11 & \etiketT{} iş kayıtları + SGK 2025, RAS \\
Yabancı sermaye; çapa yatırımcılar & Çaykur (kamu), Doğuş Çay, Lipton Teas and Infusions, Kurukahveci Mehmet Efendi, Nestlé, JDE \etiketH{} & \etiketH{} (firma bilgisi; pay verisi yok: VB-30) \\
\end{longtable}

\endgroup

\textbf{3. Dış ticaret}

\begingroup\gbtablo\small

\begin{longtable}[]{@{}
  >{\raggedright\arraybackslash}p{(\columnwidth - 2\tabcolsep) * \real{0.4766}}
  >{\raggedright\arraybackslash}p{(\columnwidth - 2\tabcolsep) * \real{0.5234}}@{}}
\toprule\noalign{}
\begin{minipage}[b]{\linewidth}\raggedright
\textbf{Gösterge}
\end{minipage} & \begin{minipage}[b]{\linewidth}\raggedright
\textbf{Değer (BACI 2024, \$; \etiketV{}, eşleme \etiketT{})}
\end{minipage} \\
\midrule\noalign{}
\endhead
\bottomrule\noalign{}
\endlastfoot
İhracat (2024); 2020--24 değişimi & \$170 mn; \%109 \\
AB payı; ilk 5 AB pazarı & \%23; DEU \%50; BEL \%19; NLD \%13; FRA \%5; ROU \%3 \\
İlk 5 AB dışı pazar & IRQ \%17; BLR \%10; SYR \%8; IRN \%6; GBR \%6 \\
İthalat; AB ve Çin payı; ilk kaynaklar & \$364 mn; AB \%60, Çin \%0; ESP \%15; DEU \%12; CHE \%10; NLD \%9; IND \%8 \\
Dış ticaret dengesi & \$−194 mn \\
İthal girdi içeriği (I-O, toplam) / KD÷üretim & \%20 / \%28 \\
RCA dünya (2020 → 2024) / AB pazarında & 0,32 → 0,44 / 0,29 \\
Birim değer ÷ AB pazarı ortalaması & 0,74 \\
AB pazar payı (AB dışı tedarik içinde, 2020 → 2024); TR'nin sırası & \%0,8 → \%1,1; 11. \\
\end{longtable}

\endgroup

\textbf{4. AB tarafı maruziyet}

\begin{itemize}
\tightlist
\item
  \kod{τ\_\allowbreak{}AB} (S4 artışı, rota A, TR→AB 2025 ağırlıklı): \%3,2; 1/98 kaybıyla \%4,8; tam MFN ortalaması \%4,8. Rejim: GB \%0, İTÜ \%41, tarım \%59 \etiketV{}/\etiketT{}.
\item
  Tepe satırlar (S4 artışı ≥ \%10; en büyük 3): 4 satır, AB ihracatının \%13'i: KN 2101 12 92 \%11,5 (€4,2 mn); KN 2101 30 11 \%11,5 (€0,1 mn); KN 2101 30 91 \%14,1 (€0,1 mn).
\item
  AB pazarında AB dışı rakipler (tercih durumu): CHE \%47; GBR \%14; VNM \%9; BRA \%5; IND \%4. Tercihli rakip payı \%82.
\item
  Menşe (S2 / S3): risk D/D; \kod{c\_\allowbreak{}RoO} \%1,5 / \%2,0; \kod{p\_\allowbreak{}uyum} 0,90 / 0,85 \etiketH{}. Kural: 0901 ve 2101: başlık değişikliği; ithal kahve çekirdeğinden kavrulmuş/hazır kahve menşe kazanır.
\item
  EUDR kapsamındaki AB ihracatı payı: \%20; SKDM: \%0.
\item
  EUDR: AB ihracatının \%20'si (kavrulmuş kahve) kapsamda.
\end{itemize}

\textbf{5. Türkiye tarafı maruziyet}

\begin{itemize}
\tightlist
\item
  Uygulanan tarife (2026, ithalat ağırlıklı): MFN \%11,2, İGV dahil \%11,2 \etiketV{}. S4'te AB menşeline (c2): \%4,6.
\item
  Kritik girdiler (HS4 0901, 0902), Türkiye'nin toplam ithalatı 2024: \$578 mn (AB \%18; BRA \%58, LKA \%8, NLD \%5) \etiketV{}.
\end{itemize}

\textbf{6. Kabiliyet}

\begin{itemize}
\tightlist
\item
  KS 44,4 (sıra 172/230; ağırlık duyarlılığında P10--P90 115--197). Bileşenler: RCA 39, PCI 17 (PCI −0,76), yurt içi KD 68, Ar-Ge 25, tedarikçi 84, verimlilik 72, firma 37 \etiketT{}.
\item
  Ürün uzayında atlanabilir ürünler (sınıf içi; yoğunluk üst çeyrek, PCI \textgreater{} sınıf): yok \etiketT{}.
\end{itemize}

\textbf{7. Stratejik önem.} SS 22,4 (en yüksek boyut: gıda güvenliği 30; ağırlıklı ortalama 4,5) \etiketH{}. Çay (yerli) ve kahve (ithal)

\textbf{8. Makro duyarlılık}

\begingroup\gbtablo\small

\begin{longtable}[]{@{}lr@{}}
\toprule\noalign{}
\textbf{Gösterge} & \textbf{Değer (\etiketT{})} \\
\midrule\noalign{}
\endhead
\bottomrule\noalign{}
\endlastfoot
TÜFE ağırlığı (doğrudan, COICOP 2018 eşlemesi); tüketimde ithal pay & \%1,04; \%14 \\
ÜFE vekili (imalat üretimindeki pay) & \%0,57 \\
EDE: 1 puan tarife, tüm menşe / yalnız AB (bp; ρ = 1) & 0,178 / 0,106 \\
Geçişkenlik ρ (senaryolarda) & 0,6--1,0 (P50 0,8) \\
Piyasa yoğunlaşması (250+ istihdam payı) & \%63 \\
CKE: tarife (\$ mn / bp) ve mali (\$ mn / mlr TL); 0--100 puan & 388,1 ve 6,50; 25 \\
\end{longtable}

\endgroup

\textbf{9. Senaryo etkileri, T+5, P50 (P10 / P90)}

\begingroup\gbtablo\small

\begin{longtable}[]{@{}
  >{\raggedright\arraybackslash}p{(\columnwidth - 6\tabcolsep) * \real{0.3667}}
  >{\raggedleft\arraybackslash}p{(\columnwidth - 6\tabcolsep) * \real{0.2111}}
  >{\raggedleft\arraybackslash}p{(\columnwidth - 6\tabcolsep) * \real{0.2111}}
  >{\raggedleft\arraybackslash}p{(\columnwidth - 6\tabcolsep) * \real{0.2111}}@{}}
\toprule\noalign{}
\begin{minipage}[b]{\linewidth}\raggedright
\textbf{Değişken}
\end{minipage} & \begin{minipage}[b]{\linewidth}\raggedleft
\textbf{S2}
\end{minipage} & \begin{minipage}[b]{\linewidth}\raggedleft
\textbf{S3}
\end{minipage} & \begin{minipage}[b]{\linewidth}\raggedleft
\textbf{S4}
\end{minipage} \\
\midrule\noalign{}
\endhead
\bottomrule\noalign{}
\endlastfoot
AB'ye ihracat (\$ mn) & −1 (−3 / +1) & −2 (−4 / 0) & −4 (−7 / −2) \\
Toplam ihracat (\$ mn) & −1 (−3 / +1) & −1 (−3 / 0) & −3 (−6 / −1) \\
Nihai mal ithalatı (\$ mn) & −5 (−15 / +15) & −7 (−18 / +9) & −17 (−31 / −4) \\
Cari denge (\$ mn) & +3 (−12 / +10) & +5 (−7 / +12) & +11 (+2 / +20) \\
Üretim (\$ mn) & +4 (−15 / +13) & +6 (−9 / +16) & +13 (+2 / +25) \\
İstihdam (kişi) & +13 (−54 / +47) & +21 (−33 / +56) & +49 (+8 / +93) \\
TÜFE (bp, birinci tur) & +0,08 (+0,07 / +0,09) & +0,15 (+0,13 / +0,18) & +0,44 (+0,38 / +0,50) \\
Yatırım ve DYY (bütün senaryolar) & --- & & \\
\end{longtable}

\endgroup

\emph{Kaynak:} \etiketT{} f3e: kısmi denge, toplamlar Faz 2B GE dağılımına kalibre; S4'te 1/98 kaybı olasılığı 0,25.

\textbf{13. Kazanan ve kaybedenler; bölgesel ve sosyal etki; telafi.} S4'te iç ikame kazancı küçük (+49 kişi).

\textbf{15. Güven düzeyi; veri boşlukları; \etiketD{}.} Yüksek.

\hypertarget{nace_kartlar_10-12-320-1084-baharat-sos-sirke-ve-diux11fer-uxe7eux15fni-maddelerinin-imalatux131--gb-statuxfcsuxfc-karma--birincilikincil-arketip-a--duruux15f-geliux15ftir}{%
\subsection{3.20 10.84 Baharat, sos, sirke ve diğer çeşni maddelerinin imalatı --- GB statüsü: KARMA \textbar{} Birincil/ikincil arketip: A \textbar{} Duruş: Geliştir}\label{nace_kartlar_10-12-320-1084-baharat-sos-sirke-ve-diux11fer-uxe7eux15fni-maddelerinin-imalatux131--gb-statuxfcsuxfc-karma--birincilikincil-arketip-a--duruux15f-geliux15ftir}}

\emph{Baharat, sos ve sirke: AB payı \%34. AB'ye \$83 mn ihracat, S4'te \%4,9 (ketçapta \%10,2). SPS kontrolleri (kalıntı) GB'den bağımsız risk.}

\textbf{1. Kimlik.} Rev. 2.1: 10.84 (1:1); CPA 10.84. Başlıca HS4 (TR ihracatı 2024): 2103 \%54, 0910 \%20, 2501 \%9, 0904 \%9. BEC payları (ticaret ağırlıklı): ara \%7 / sermaye \%0 / tüketim \%93.

\textbf{2. Ölçek}

\begingroup\gbtablo\small

\begin{longtable}[]{@{}
  >{\raggedright\arraybackslash}p{(\columnwidth - 4\tabcolsep) * \real{0.2738}}
  >{\raggedright\arraybackslash}p{(\columnwidth - 4\tabcolsep) * \real{0.4215}}
  >{\raggedright\arraybackslash}p{(\columnwidth - 4\tabcolsep) * \real{0.3047}}@{}}
\toprule\noalign{}
\begin{minipage}[b]{\linewidth}\raggedright
\textbf{Gösterge}
\end{minipage} & \begin{minipage}[b]{\linewidth}\raggedright
\textbf{Değer}
\end{minipage} & \begin{minipage}[b]{\linewidth}\raggedright
\textbf{Etiket ve kaynak}
\end{minipage} \\
\midrule\noalign{}
\endhead
\bottomrule\noalign{}
\endlastfoot
Üretim değeri; katma değer & \$1.244 mn; \$296 mn (GSYH'nin \%0,02) & TÜİK YSHİ 2024 \etiketV{} \\
Çalışan; girişim (0 / 1--9 / 10--49 / 50--249 / 250+) & 7.151; 659 (29 / 558 / 96 / 24 / 2) & TÜİK YSHİ ve İş Kayıtları 2024 \etiketV{}; gizli büyüklük hücreleri \etiketT{} \\
250+ girişimlerin istihdam payı; AB'ye göre kişi başı KD & \%11; 0,55 & \etiketT{} (AB27 SBS 2023, nominal kur) \\
İlk 3 İBBS-2 bölgesi & İstanbul \%14; İzmir \%13; Gaziantep, Adıyaman, Kilis \%8 & \etiketT{} iş kayıtları + SGK 2025, RAS \\
Yabancı sermaye; çapa yatırımcılar & Tat (ketçap ve sos), Kühne (Almanya; sos), Bağdat Baharat; defne ve kekikte Ege ihracatçıları \etiketH{} & \etiketH{} (firma bilgisi; pay verisi yok: VB-30) \\
\end{longtable}

\endgroup

\textbf{3. Dış ticaret}

\begingroup\gbtablo\small

\begin{longtable}[]{@{}
  >{\raggedright\arraybackslash}p{(\columnwidth - 2\tabcolsep) * \real{0.4839}}
  >{\raggedright\arraybackslash}p{(\columnwidth - 2\tabcolsep) * \real{0.5161}}@{}}
\toprule\noalign{}
\begin{minipage}[b]{\linewidth}\raggedright
\textbf{Gösterge}
\end{minipage} & \begin{minipage}[b]{\linewidth}\raggedright
\textbf{Değer (BACI 2024, \$; \etiketV{}, eşleme \etiketT{})}
\end{minipage} \\
\midrule\noalign{}
\endhead
\bottomrule\noalign{}
\endlastfoot
İhracat (2024); 2020--24 değişimi & \$242 mn; \%38 \\
AB payı; ilk 5 AB pazarı & \%34; DEU \%34; BEL \%16; NLD \%9; FRA \%8; ROU \%5 \\
İlk 5 AB dışı pazar & USA \%17; GBR \%13; IRQ \%11; SYR \%3; CAN \%3 \\
İthalat; AB ve Çin payı; ilk kaynaklar & \$90 mn; AB \%40, Çin \%5; IDN \%16; GBR \%11; ITA \%8; NLD \%7; HUN \%7 \\
Dış ticaret dengesi & \$+153 mn \\
İthal girdi içeriği (I-O, toplam) / KD÷üretim & \%20 / \%24 \\
RCA dünya (2020 → 2024) / AB pazarında & 0,85 → 0,78 / 1,01 \\
Birim değer ÷ AB pazarı ortalaması & 0,98 \\
AB pazar payı (AB dışı tedarik içinde, 2020 → 2024); TR'nin sırası & \%3,0 → \%4,0; 7. \\
\end{longtable}

\endgroup

\textbf{4. AB tarafı maruziyet}

\begin{itemize}
\tightlist
\item
  \kod{τ\_\allowbreak{}AB} (S4 artışı, rota A, TR→AB 2025 ağırlıklı): \%4,9; 1/98 kaybıyla \%7,4; tam MFN ortalaması \%7,4. Rejim: GB \%11, İTÜ \%56, tarım \%33 \etiketV{}/\etiketT{}.
\item
  Tepe satırlar (S4 artışı ≥ \%10; en büyük 3): 1 satır, AB ihracatının \%23'i: KN 2103 20 00 \%10,2 (€13,6 mn).
\item
  AB pazarında AB dışı rakipler (tercih durumu): GBR \%23; CHN \%14; THA \%11; USA \%6; VNM \%5. Tercihli rakip payı \%49.
\item
  Menşe (S2 / S3): risk D/D; \kod{c\_\allowbreak{}RoO} \%1,5 / \%2,5; \kod{p\_\allowbreak{}uyum} 0,90 / 0,85 \etiketH{}. Kural: 2103 ve 2209: başlık değişikliği (hardal unu istisnası).
\item
  EUDR kapsamındaki AB ihracatı payı: \%0; SKDM: \%0.
\item
  Kekik, defne ve kimyonda kalıntı ve aflatoksin kaynaklı RASFF bildirimleri GB'den bağımsız bir engel \etiketD{}.
\end{itemize}

\textbf{5. Türkiye tarafı maruziyet}

\begin{itemize}
\tightlist
\item
  Uygulanan tarife (2026, ithalat ağırlıklı): MFN \%12,4, İGV dahil \%12,4 \etiketV{}. S4'te AB menşeline (c2): \%5,0.
\item
  Kritik girdiler (HS4 0904, 0909, 0910), Türkiye'nin toplam ithalatı 2024: \$108 mn (AB \%6; SYR \%31, VNM \%16, CHN \%14) \etiketV{}.
\end{itemize}

\textbf{6. Kabiliyet}

\begin{itemize}
\tightlist
\item
  KS 42,3 (sıra 192/230; ağırlık duyarlılığında P10--P90 133--208). Bileşenler: RCA 30, PCI 12 (PCI −0,92), yurt içi KD 59, Ar-Ge 25, tedarikçi 84, verimlilik 69, firma 60 \etiketT{}.
\item
  Ürün uzayında atlanabilir ürünler (sınıf içi; yoğunluk üst çeyrek, PCI \textgreater{} sınıf): yok \etiketT{}.
\end{itemize}

\textbf{7. Stratejik önem.} SS 22,4 (en yüksek boyut: gıda güvenliği 30; ağırlıklı ortalama 4,5) \etiketH{}.

\textbf{8. Makro duyarlılık}

\begingroup\gbtablo\small

\begin{longtable}[]{@{}lr@{}}
\toprule\noalign{}
\textbf{Gösterge} & \textbf{Değer (\etiketT{})} \\
\midrule\noalign{}
\endhead
\bottomrule\noalign{}
\endlastfoot
TÜFE ağırlığı (doğrudan, COICOP 2018 eşlemesi); tüketimde ithal pay & \%0,28; \%6 \\
ÜFE vekili (imalat üretimindeki pay) & \%0,19 \\
EDE: 1 puan tarife, tüm menşe / yalnız AB (bp; ρ = 1) & 0,023 / 0,009 \\
Geçişkenlik ρ (senaryolarda) & 0,6--1,0 (P50 0,8) \\
Piyasa yoğunlaşması (250+ istihdam payı) & \%11 \\
CKE: tarife (\$ mn / bp) ve mali (\$ mn / mlr TL); 0--100 puan & 861,7 ve 6,88; 30 \\
\end{longtable}

\endgroup

\textbf{9. Senaryo etkileri, T+5, P50 (P10 / P90)}

\begingroup\gbtablo\small

\begin{longtable}[]{@{}
  >{\raggedright\arraybackslash}p{(\columnwidth - 6\tabcolsep) * \real{0.3667}}
  >{\raggedleft\arraybackslash}p{(\columnwidth - 6\tabcolsep) * \real{0.2111}}
  >{\raggedleft\arraybackslash}p{(\columnwidth - 6\tabcolsep) * \real{0.2111}}
  >{\raggedleft\arraybackslash}p{(\columnwidth - 6\tabcolsep) * \real{0.2111}}@{}}
\toprule\noalign{}
\begin{minipage}[b]{\linewidth}\raggedright
\textbf{Değişken}
\end{minipage} & \begin{minipage}[b]{\linewidth}\raggedleft
\textbf{S2}
\end{minipage} & \begin{minipage}[b]{\linewidth}\raggedleft
\textbf{S3}
\end{minipage} & \begin{minipage}[b]{\linewidth}\raggedleft
\textbf{S4}
\end{minipage} \\
\midrule\noalign{}
\endhead
\bottomrule\noalign{}
\endlastfoot
AB'ye ihracat (\$ mn) & −6 (−14 / +2) & −8 (−16 / −2) & −14 (−24 / −6) \\
Toplam ihracat (\$ mn) & −4 (−10 / +1) & −6 (−11 / −1) & −11 (−19 / −5) \\
Nihai mal ithalatı (\$ mn) & −1 (−4 / +2) & −2 (−5 / +1) & −3 (−6 / −1) \\
Cari denge (\$ mn) & −2 (−5 / −1) & −3 (−5 / −2) & −6 (−10 / −3) \\
Üretim (\$ mn) & −3 (−6 / −1) & −4 (−7 / −2) & −8 (−13 / −4) \\
İstihdam (kişi) & −10 (−24 / −5) & −14 (−26 / −9) & −31 (−54 / −16) \\
TÜFE (bp, birinci tur) & +0,01 (+0,01 / +0,01) & +0,02 (+0,02 / +0,02) & +0,04 (+0,04 / +0,05) \\
Yatırım ve DYY (bütün senaryolar) & --- & & \\
\end{longtable}

\endgroup

\emph{Kaynak:} \etiketT{} f3e: kısmi denge, toplamlar Faz 2B GE dağılımına kalibre; S4'te 1/98 kaybı olasılığı 0,25.

\textbf{13. Kazanan ve kaybedenler; bölgesel ve sosyal etki; telafi.} Kaybeden (S4 P50): −31 kişi; İzmir ve İstanbul.

\textbf{15. Güven düzeyi; veri boşlukları; \etiketD{}.} Orta.

\hypertarget{nace_kartlar_10-12-321-1085-hazux131r-yemeklerin-imalatux131--gb-statuxfcsuxfc-iux307tuxfc--birincilikincil-arketip---duruux15f-koru}{%
\subsection{3.21 10.85 Hazır yemeklerin imalatı --- GB statüsü: İTÜ \textbar{} Birincil/ikincil arketip: --- \textbar{} Duruş: Koru}\label{nace_kartlar_10-12-321-1085-hazux131r-yemeklerin-imalatux131--gb-statuxfcsuxfc-iux307tuxfc--birincilikincil-arketip---duruux15f-koru}}

\emph{Küçük sınıf (hazır yemek): ihracat \$100 mn, AB payı \%10. Kurallar emek yoğun ve gerileyen (F) dedi; GB'nin etkisi küçük olduğu için duruş Koru.}

\textbf{1. Kimlik.} Rev. 2.1: 10.85 (1:1); CPA 10.85. Başlıca HS4 (TR ihracatı 2024): 1902 \%100. BEC payları (ticaret ağırlıklı): ara \%0 / sermaye \%0 / tüketim \%100.

\textbf{2. Ölçek}

\begingroup\gbtablo\small

\begin{longtable}[]{@{}
  >{\raggedright\arraybackslash}p{(\columnwidth - 4\tabcolsep) * \real{0.3166}}
  >{\raggedright\arraybackslash}p{(\columnwidth - 4\tabcolsep) * \real{0.3232}}
  >{\raggedright\arraybackslash}p{(\columnwidth - 4\tabcolsep) * \real{0.3603}}@{}}
\toprule\noalign{}
\begin{minipage}[b]{\linewidth}\raggedright
\textbf{Gösterge}
\end{minipage} & \begin{minipage}[b]{\linewidth}\raggedright
\textbf{Değer}
\end{minipage} & \begin{minipage}[b]{\linewidth}\raggedright
\textbf{Etiket ve kaynak}
\end{minipage} \\
\midrule\noalign{}
\endhead
\bottomrule\noalign{}
\endlastfoot
Üretim değeri; katma değer & \$815 mn; \$221 mn (GSYH'nin \%0,02) & TÜİK YSHİ 2024 \etiketV{} \\
Çalışan; girişim (0 / 1--9 / 10--49 / 50--249 / 250+) & 10.010; 932 (67 / 857 / 105 / 12 / 8) & TÜİK YSHİ ve İş Kayıtları 2024 \etiketV{}; gizli büyüklük hücreleri \etiketT{} \\
250+ girişimlerin istihdam payı; AB'ye göre kişi başı KD & \%38; 0,35 & \etiketT{} (AB27 SBS 2023, nominal kur) \\
İlk 3 İBBS-2 bölgesi & İstanbul \%15; İzmir \%15; Tekirdağ, Edirne, Kırklareli \%11 & \etiketT{} iş kayıtları + SGK 2025, RAS \\
Yabancı sermaye; çapa yatırımcılar & Yerli orta ölçekli üreticiler \etiketH{} & \etiketH{} (firma bilgisi; pay verisi yok: VB-30) \\
\end{longtable}

\endgroup

\textbf{3. Dış ticaret}

\begingroup\gbtablo\small

\begin{longtable}[]{@{}
  >{\raggedright\arraybackslash}p{(\columnwidth - 2\tabcolsep) * \real{0.4787}}
  >{\raggedright\arraybackslash}p{(\columnwidth - 2\tabcolsep) * \real{0.5213}}@{}}
\toprule\noalign{}
\begin{minipage}[b]{\linewidth}\raggedright
\textbf{Gösterge}
\end{minipage} & \begin{minipage}[b]{\linewidth}\raggedright
\textbf{Değer (BACI 2024, \$; \etiketV{}, eşleme \etiketT{})}
\end{minipage} \\
\midrule\noalign{}
\endhead
\bottomrule\noalign{}
\endlastfoot
İhracat (2024); 2020--24 değişimi & \$100 mn; \%26 \\
AB payı; ilk 5 AB pazarı & \%10; DEU \%23; DNK \%15; FRA \%11; NLD \%8; ROU \%7 \\
İlk 5 AB dışı pazar & IRQ \%21; SYR \%19; MYS \%14; NER \%11; PSE \%2 \\
İthalat; AB ve Çin payı; ilk kaynaklar & \$2 mn; AB \%65, Çin \%18; ITA \%46; CHN \%18; DEU \%10; HUN \%5; VNM \%4 \\
Dış ticaret dengesi & \$+98 mn \\
İthal girdi içeriği (I-O, toplam) / KD÷üretim & \%20 / \%27 \\
RCA dünya (2020 → 2024) / AB pazarında & 1,30 → 1,00 / 0,25 \\
Birim değer ÷ AB pazarı ortalaması & 0,62 \\
AB pazar payı (AB dışı tedarik içinde, 2020 → 2024); TR'nin sırası & \%0,4 → \%1,0; 12. \\
\end{longtable}

\endgroup

\textbf{4. AB tarafı maruziyet}

\begin{itemize}
\tightlist
\item
  \kod{τ\_\allowbreak{}AB} (S4 artışı, rota A, TR→AB 2025 ağırlıklı): \%7,0; 1/98 kaybıyla \%7,0; tam MFN ortalaması \%17,7. Rejim: GB \%0, İTÜ \%100, tarım \%0 \etiketV{}/\etiketT{}.
\item
  Tepe satırlar (S4 artışı ≥ \%10; en büyük 3): yok.
\item
  AB pazarında AB dışı rakipler (tercih durumu): CHN \%19; KOR \%17; THA \%15; CHE \%13; VNM \%13. Tercihli rakip payı \%57.
\item
  Menşe (S2 / S3): risk O/O; \kod{c\_\allowbreak{}RoO} \%3,0 / \%4,0; \kod{p\_\allowbreak{}uyum} 0,80 / 0,70 \etiketH{}. Kural: 1902 20 ve 2104: Fasıl 2, 3 ve 16 malzemeleri tamamen elde edilmiş; tahıl kuralı.
\item
  EUDR kapsamındaki AB ihracatı payı: \%0; SKDM: \%0.
\end{itemize}

\textbf{5. Türkiye tarafı maruziyet}

\begin{itemize}
\tightlist
\item
  Uygulanan tarife (2026, ithalat ağırlıklı): MFN \%15,9, İGV dahil \%20,9 \etiketV{}. S4'te AB menşeline (c2): \%23,3.
\end{itemize}

\textbf{6. Kabiliyet}

\begin{itemize}
\tightlist
\item
  KS 37,5 (sıra 216/230; ağırlık duyarlılığında P10--P90 170--222). Bileşenler: RCA 34, PCI 8 (PCI −1,08), yurt içi KD 66, Ar-Ge 25, tedarikçi 84, verimlilik 25, firma 46 \etiketT{}.
\item
  Ürün uzayında atlanabilir ürünler (sınıf içi; yoğunluk üst çeyrek, PCI \textgreater{} sınıf): yok \etiketT{}.
\end{itemize}

\textbf{7. Stratejik önem.} SS 23,6 (en yüksek boyut: gıda güvenliği 30; ağırlıklı ortalama 8,5) \etiketH{}. Hazır yemek; askerî ve afet tayınları

\textbf{8. Makro duyarlılık}

\begingroup\gbtablo\small

\begin{longtable}[]{@{}lr@{}}
\toprule\noalign{}
\textbf{Gösterge} & \textbf{Değer (\etiketT{})} \\
\midrule\noalign{}
\endhead
\bottomrule\noalign{}
\endlastfoot
TÜFE ağırlığı (doğrudan, COICOP 2018 eşlemesi); tüketimde ithal pay & \%0,26; \%2 \\
ÜFE vekili (imalat üretimindeki pay) & \%0,12 \\
EDE: 1 puan tarife, tüm menşe / yalnız AB (bp; ρ = 1) & 0,005 / 0,003 \\
Geçişkenlik ρ (senaryolarda) & 0,6--1,0 (P50 0,8) \\
Piyasa yoğunlaşması (250+ istihdam payı) & \%38 \\
CKE: tarife (\$ mn / bp) ve mali (\$ mn / mlr TL); 0--100 puan & 35,3 ve 0,09; 8 \\
\end{longtable}

\endgroup

\textbf{9. Senaryo etkileri, T+5, P50 (P10 / P90)}

\begingroup\gbtablo\small

\begin{longtable}[]{@{}
  >{\raggedright\arraybackslash}p{(\columnwidth - 6\tabcolsep) * \real{0.3667}}
  >{\raggedleft\arraybackslash}p{(\columnwidth - 6\tabcolsep) * \real{0.2111}}
  >{\raggedleft\arraybackslash}p{(\columnwidth - 6\tabcolsep) * \real{0.2111}}
  >{\raggedleft\arraybackslash}p{(\columnwidth - 6\tabcolsep) * \real{0.2111}}@{}}
\toprule\noalign{}
\begin{minipage}[b]{\linewidth}\raggedright
\textbf{Değişken}
\end{minipage} & \begin{minipage}[b]{\linewidth}\raggedleft
\textbf{S2}
\end{minipage} & \begin{minipage}[b]{\linewidth}\raggedleft
\textbf{S3}
\end{minipage} & \begin{minipage}[b]{\linewidth}\raggedleft
\textbf{S4}
\end{minipage} \\
\midrule\noalign{}
\endhead
\bottomrule\noalign{}
\endlastfoot
AB'ye ihracat (\$ mn) & −1 (−1 / 0) & −1 (−1 / 0) & −1 (−1 / 0) \\
Toplam ihracat (\$ mn) & 0 (−1 / 0) & −1 (−1 / 0) & −1 (−1 / 0) \\
Nihai mal ithalatı (\$ mn) & 0 (0 / 0) & 0 (0 / 0) & 0 (0 / 0) \\
Cari denge (\$ mn) & 0 (−1 / 0) & 0 (−1 / 0) & 0 (−1 / 0) \\
Üretim (\$ mn) & 0 (−1 / 0) & 0 (−1 / 0) & −1 (−1 / 0) \\
İstihdam (kişi) & −3 (−8 / 0) & −4 (−7 / −1) & −5 (−8 / −2) \\
TÜFE (bp, birinci tur) & +0,01 (+0,01 / +0,01) & +0,01 (+0,01 / +0,02) & +0,07 (+0,06 / +0,07) \\
Yatırım ve DYY (bütün senaryolar) & --- & & \\
\end{longtable}

\endgroup

\emph{Kaynak:} \etiketT{} f3e: kısmi denge, toplamlar Faz 2B GE dağılımına kalibre; S4'te 1/98 kaybı olasılığı 0,25.

\textbf{13. Kazanan ve kaybedenler; bölgesel ve sosyal etki; telafi.} Etkilenen kesim küçük (S4 P50 −5 kişi).

\textbf{15. Güven düzeyi; veri boşlukları; \etiketD{}.} Orta.

\hypertarget{nace_kartlar_10-12-322-1086-homojenize-gux131da-muxfcstahzarlarux131-ve-diyetetik-gux131dalarux131n-imalatux131--gb-statuxfcsuxfc-karma--birincilikincil-arketip-dc--duruux15f-geliux15ftir}{%
\subsection{3.22 10.86 Homojenize gıda müstahzarları ve diyetetik gıdaların imalatı --- GB statüsü: KARMA \textbar{} Birincil/ikincil arketip: D/C \textbar{} Duruş: Geliştir}\label{nace_kartlar_10-12-322-1086-homojenize-gux131da-muxfcstahzarlarux131-ve-diyetetik-gux131dalarux131n-imalatux131--gb-statuxfcsuxfc-karma--birincilikincil-arketip-dc--duruux15f-geliux15ftir}}

\emph{Bebek maması ve diyet gıdaları: ithalatın \%91'i AB'den (Polonya, Almanya, İrlanda). S4'te AB'ye c2 (\%7,6) fiyatı artırır. Sağlık açısından hassas; tarife yanlış araç, yerlileştirme için yatırım teşviki gerekiyor.}

\textbf{1. Kimlik.} Rev. 2.1: 10.61, 10.62, 10.71, 10.72, 10.73, 10.86 (1:6); CPA 10.86. Başlıca HS4 (TR ihracatı 2024): 1905 \%85, 1901 \%7, 2007 \%5, 2104 \%3. BEC payları (ticaret ağırlıklı): ara \%0 / sermaye \%0 / tüketim \%100.

\textbf{2. Ölçek}

\begingroup\gbtablo\small

\begin{longtable}[]{@{}
  >{\raggedright\arraybackslash}p{(\columnwidth - 4\tabcolsep) * \real{0.2717}}
  >{\raggedright\arraybackslash}p{(\columnwidth - 4\tabcolsep) * \real{0.4264}}
  >{\raggedright\arraybackslash}p{(\columnwidth - 4\tabcolsep) * \real{0.3019}}@{}}
\toprule\noalign{}
\begin{minipage}[b]{\linewidth}\raggedright
\textbf{Gösterge}
\end{minipage} & \begin{minipage}[b]{\linewidth}\raggedright
\textbf{Değer}
\end{minipage} & \begin{minipage}[b]{\linewidth}\raggedright
\textbf{Etiket ve kaynak}
\end{minipage} \\
\midrule\noalign{}
\endhead
\bottomrule\noalign{}
\endlastfoot
Üretim değeri; katma değer & \$172 mn; \$61 mn (GSYH'nin \%0,00) & TÜİK YSHİ 2024 \etiketV{} \\
Çalışan; girişim (0 / 1--9 / 10--49 / 50--249 / 250+) & 1.321; 124 (16 / 102 / 11 / 4 / 1) & TÜİK YSHİ ve İş Kayıtları 2024 \etiketV{}; gizli büyüklük hücreleri \etiketT{} \\
250+ girişimlerin istihdam payı; AB'ye göre kişi başı KD & \%26; 0,45 & \etiketT{} (AB27 SBS 2023, nominal kur) \\
İlk 3 İBBS-2 bölgesi & İstanbul \%41; Ankara \%15; Mardin, Batman, Şırnak, Siirt \%9 & \etiketT{} iş kayıtları + SGK 2025, RAS \\
Yabancı sermaye; çapa yatırımcılar & Nutricia (Danone), Nestlé, Abbott ve HiPP ithalatçıları; yerli üretimde Hero Baby (Ülker--Hero) \etiketH{} & \etiketH{} (firma bilgisi; pay verisi yok: VB-30) \\
\end{longtable}

\endgroup

\textbf{3. Dış ticaret}

\begingroup\gbtablo\small

\begin{longtable}[]{@{}
  >{\raggedright\arraybackslash}p{(\columnwidth - 2\tabcolsep) * \real{0.4730}}
  >{\raggedright\arraybackslash}p{(\columnwidth - 2\tabcolsep) * \real{0.5270}}@{}}
\toprule\noalign{}
\begin{minipage}[b]{\linewidth}\raggedright
\textbf{Gösterge}
\end{minipage} & \begin{minipage}[b]{\linewidth}\raggedright
\textbf{Değer (BACI 2024, \$; \etiketV{}, eşleme \etiketT{})}
\end{minipage} \\
\midrule\noalign{}
\endhead
\bottomrule\noalign{}
\endlastfoot
İhracat (2024); 2020--24 değişimi & \$214 mn; \%62 \\
AB payı; ilk 5 AB pazarı & \%20; DEU \%46; NLD \%10; BEL \%8; ROU \%7; BGR \%7 \\
İlk 5 AB dışı pazar & IRQ \%29; USA \%7; AZE \%5; YEM \%4; SYR \%4 \\
İthalat; AB ve Çin payı; ilk kaynaklar & \$173 mn; AB \%91, Çin \%0; POL \%32; DEU \%25; FRA \%12; IRL \%10; NLD \%5 \\
Dış ticaret dengesi & \$+41 mn \\
İthal girdi içeriği (I-O, toplam) / KD÷üretim & \%20 / \%35 \\
RCA dünya (2020 → 2024) / AB pazarında & 0,64 → 0,80 / 2,14 \\
Birim değer ÷ AB pazarı ortalaması & 1,26 \\
AB pazar payı (AB dışı tedarik içinde, 2020 → 2024); TR'nin sırası & \%4,5 → \%8,5; 3. \\
\end{longtable}

\endgroup

\textbf{4. AB tarafı maruziyet}

\begin{itemize}
\tightlist
\item
  \kod{τ\_\allowbreak{}AB} (S4 artışı, rota A, TR→AB 2025 ağırlıklı): \%7,7; 1/98 kaybıyla \%11,0; tam MFN ortalaması \%11,2. Rejim: GB \%0, İTÜ \%86, tarım \%14 \etiketV{}/\etiketT{}.
\item
  Tepe satırlar (S4 artışı ≥ \%10; en büyük 3): yok (AB ihracatının \%1'inden az).
\item
  AB pazarında AB dışı rakipler (tercih durumu): GBR \%39; CHE \%9; CHN \%6; SRB \%5; USA \%3. Tercihli rakip payı \%67.
\item
  Menşe (S2 / S3): risk O/O; \kod{c\_\allowbreak{}RoO} \%2,5 / \%3,5; \kod{p\_\allowbreak{}uyum} 0,85 / 0,75 \etiketH{}. Kural: 1901 10 ve 2104 20: Fasıl 4 ve tahıl malzemeleri ağırlık/menşe sınırları.
\item
  EUDR kapsamındaki AB ihracatı payı: \%0; SKDM: \%0.
\end{itemize}

\textbf{5. Türkiye tarafı maruziyet}

\begin{itemize}
\tightlist
\item
  Uygulanan tarife (2026, ithalat ağırlıklı): MFN \%8,4, İGV dahil \%8,4 \etiketV{}. S4'te AB menşeline (c2): \%7,5.
\item
  Kritik girdiler (HS4 0402, 0404), Türkiye'nin toplam ithalatı 2024: \$15 mn (AB \%12; BLR \%32, NZL \%17, USA \%13) \etiketV{}.
\item
  Tüketimde ithal pay matristeki I-O tahmininden (\%9) çok daha yüksek; bebek mamasında yerli üretim sınırlı \etiketD{}. YSHİ üretimi (\$172 mn) ihracatın (\$214 mn) altında: sınıflama farkı var (VB-35).
\end{itemize}

\textbf{6. Kabiliyet}

\begin{itemize}
\tightlist
\item
  KS 50,1 (sıra 117/230; ağırlık duyarlılığında P10--P90 59--160). Bileşenler: RCA 46, PCI 40 (PCI −0,17), yurt içi KD 85, Ar-Ge 25, tedarikçi 84, verimlilik 46, firma 34 \etiketT{}.
\item
  Ürün uzayında atlanabilir ürünler (sınıf içi; yoğunluk üst çeyrek, PCI \textgreater{} sınıf): yok \etiketT{}.
\item
  \textbf{Kırmızı çizgi yetkinlikleri:} Bebek maması formülasyonu ve gıda güvenliği laboratuvarları.
\item
  \textbf{Geliştirme hedefi:} T+5: yerli bebek maması kapasitesi iç talebin \%30'u; 2035: \%50.
\end{itemize}

\textbf{7. Stratejik önem.} SS 54,9 (en yüksek boyut: sağlık 70; ağırlıklı ortalama 19,5) \etiketH{}. Bebek maması ve tıbbi gıdalar sağlık için kritik; ithalat bağımlılığı yüksek.

\textbf{8. Makro duyarlılık}

\begingroup\gbtablo\small

\begin{longtable}[]{@{}lr@{}}
\toprule\noalign{}
\textbf{Gösterge} & \textbf{Değer (\etiketT{})} \\
\midrule\noalign{}
\endhead
\bottomrule\noalign{}
\endlastfoot
TÜFE ağırlığı (doğrudan, COICOP 2018 eşlemesi); tüketimde ithal pay & \%0,09; \%9 \\
ÜFE vekili (imalat üretimindeki pay) & \%0,03 \\
EDE: 1 puan tarife, tüm menşe / yalnız AB (bp; ρ = 1) & 0,021 / 0,019 \\
Geçişkenlik ρ (senaryolarda) & 0,6--1,0 (P50 0,8) \\
Piyasa yoğunlaşması (250+ istihdam payı) & \%26 \\
CKE: tarife (\$ mn / bp) ve mali (\$ mn / mlr TL); 0--100 puan & 1.850,8 ve 101,62; 72 \\
\end{longtable}

\endgroup

\textbf{9. Senaryo etkileri, T+5, P50 (P10 / P90)}

\begingroup\gbtablo\small

\begin{longtable}[]{@{}
  >{\raggedright\arraybackslash}p{(\columnwidth - 6\tabcolsep) * \real{0.2366}}
  >{\raggedleft\arraybackslash}p{(\columnwidth - 6\tabcolsep) * \real{0.4404}}
  >{\raggedleft\arraybackslash}p{(\columnwidth - 6\tabcolsep) * \real{0.1615}}
  >{\raggedleft\arraybackslash}p{(\columnwidth - 6\tabcolsep) * \real{0.1615}}@{}}
\toprule\noalign{}
\begin{minipage}[b]{\linewidth}\raggedright
\textbf{Değişken}
\end{minipage} & \begin{minipage}[b]{\linewidth}\raggedleft
\textbf{S2}
\end{minipage} & \begin{minipage}[b]{\linewidth}\raggedleft
\textbf{S3}
\end{minipage} & \begin{minipage}[b]{\linewidth}\raggedleft
\textbf{S4}
\end{minipage} \\
\midrule\noalign{}
\endhead
\bottomrule\noalign{}
\endlastfoot
AB'ye ihracat (\$ mn) & −7 (−17 / +1) & −9 (−17 / −2) & −13 (−21 / −5) \\
Toplam ihracat (\$ mn) & −5 (−12 / +1) & −6 (−12 / −1) & −10 (−17 / −4) \\
Nihai mal ithalatı (\$ mn) & −7 (−23 / +8) & −10 (−26 / +5) & −19 (−34 / −4) \\
Cari denge (\$ mn) & +1 (−6 / +9) & +3 (−5 / +11) & +7 (0 / +14) \\
Üretim (\$ mn) & +1 (−7 / +11) & +3 (−6 / +14) & +9 (0 / +18) \\
İstihdam (kişi) & +7 (−39 / +60) & +18 (−31 / +76) & +45 (−2 / +96) \\
TÜFE (bp, birinci tur) & +0,03 (+0,02 / +0,03) & +0,05 (+0,04 / +0,06) & +0,13 (+0,12 / +0,15) \\
Yatırım ve DYY (bütün senaryolar) & Yatırım teşviki (öncelikli yatırım) ile bir çok uluslu üretim tesisi hedefi \etiketH{} & & \\
\end{longtable}

\endgroup

\emph{Kaynak:} \etiketT{} f3e: kısmi denge, toplamlar Faz 2B GE dağılımına kalibre; S4'te 1/98 kaybı olasılığı 0,25.

\textbf{10. Politika paketi}

\begingroup\gbtablo\scriptsize

\begin{longtable}[]{@{}
  >{\raggedright\arraybackslash}p{(\columnwidth - 16\tabcolsep) * \real{0.0983}}
  >{\raggedright\arraybackslash}p{(\columnwidth - 16\tabcolsep) * \real{0.1037}}
  >{\raggedright\arraybackslash}p{(\columnwidth - 16\tabcolsep) * \real{0.1274}}
  >{\raggedright\arraybackslash}p{(\columnwidth - 16\tabcolsep) * \real{0.0933}}
  >{\raggedright\arraybackslash}p{(\columnwidth - 16\tabcolsep) * \real{0.1261}}
  >{\raggedright\arraybackslash}p{(\columnwidth - 16\tabcolsep) * \real{0.0945}}
  >{\raggedright\arraybackslash}p{(\columnwidth - 16\tabcolsep) * \real{0.1598}}
  >{\raggedright\arraybackslash}p{(\columnwidth - 16\tabcolsep) * \real{0.1135}}
  >{\raggedright\arraybackslash}p{(\columnwidth - 16\tabcolsep) * \real{0.0836}}@{}}
\toprule\noalign{}
\begin{minipage}[b]{\linewidth}\raggedright
\textbf{Eylem ID}
\end{minipage} & \begin{minipage}[b]{\linewidth}\raggedright
\textbf{Araç}
\end{minipage} & \begin{minipage}[b]{\linewidth}\raggedright
\textbf{Parametre}
\end{minipage} & \begin{minipage}[b]{\linewidth}\raggedright
\textbf{Sorumlu}
\end{minipage} & \begin{minipage}[b]{\linewidth}\raggedright
\textbf{Hukuki araç}
\end{minipage} & \begin{minipage}[b]{\linewidth}\raggedright
\textbf{Bütçe}
\end{minipage} & \begin{minipage}[b]{\linewidth}\raggedright
\textbf{Başlangıç--bitiş; sunset}
\end{minipage} & \begin{minipage}[b]{\linewidth}\raggedright
\textbf{Performans koşulu}
\end{minipage} & \begin{minipage}[b]{\linewidth}\raggedright
\textbf{KPI}
\end{minipage} \\
\midrule\noalign{}
\endhead
\bottomrule\noalign{}
\endlastfoot
E-10.86-01 & S4'te bebek mamasında MFN \%0 & KN 1901 10: \%7,6 → \%0, tüm menşeler (yalnız S4'te; S2'de AB'ye zaten \%0) & Ticaret Bakanlığı & İthalat Rejimi Kararı eki & Gelir kaybı ≈ 0,06 mlr TL/yıl (AB dışı ithalat) \etiketT{} & T → T+36; sunset: T+36 (yerli kapasite devreye girince) & --- & Bebek maması TÜFE alt kalemi; −0,12 bp \\
E-10.86-02 & Bebek maması yatırım teşviki & Öncelikli yatırım statüsü; en az 30 bin ton/yıl kapasite & Sanayi ve Teknoloji Bakanlığı & Yatırımlarda Devlet Yardımları Kararı & 0,2 mlr TL (bir kez, vergi indirimi) \etiketT{} & T−12 → T+36; sunset: T+60 & Üretim başlangıcı T+36 & Yerli pay ≥ \%30 \\
\end{longtable}

\endgroup

\textbf{12. Erken uyarı ve tetikler}

\begingroup\gbtablo\small

\begin{longtable}[]{@{}lll@{}}
\toprule\noalign{}
\textbf{Gösterge} & \textbf{Eşik} & \textbf{Tetiklenecek önlem} \\
\midrule\noalign{}
\endhead
\bottomrule\noalign{}
\endlastfoot
Bebek maması perakende fiyatı (aylık) & 3 ayda +\%10 & Stok ve ithalat izni hızlandırma \\
\end{longtable}

\endgroup

\textbf{13. Kazanan ve kaybedenler; bölgesel ve sosyal etki; telafi.} Kazanan (S4): yerli üretici; tüketici MFN \%0 ile korunur.

\textbf{14. Riskler ve azaltım}

\begin{itemize}
\tightlist
\item
  Tek tedarikçi kesintisi → stok izleme.
\end{itemize}

\textbf{15. Güven düzeyi; veri boşlukları; \etiketD{}.} Düşük (sınıf verisi tutarsız; VB-35).

\hypertarget{nace_kartlar_10-12-323-1089-baux15fka-yerde-sux131nux131flandux131rux131lmamux131ux15f-diux11fer-gux131da-maddelerinin-imalatux131--gb-statuxfcsuxfc-iux307tuxfc--birincilikincil-arketip-ac--duruux15f-geliux15ftir}{%
\subsection{3.23 10.89 Başka yerde sınıflandırılmamış diğer gıda maddelerinin imalatı --- GB statüsü: İTÜ \textbar{} Birincil/ikincil arketip: A/C \textbar{} Duruş: Geliştir}\label{nace_kartlar_10-12-323-1089-baux15fka-yerde-sux131nux131flandux131rux131lmamux131ux15f-diux11fer-gux131da-maddelerinin-imalatux131--gb-statuxfcsuxfc-iux307tuxfc--birincilikincil-arketip-ac--duruux15f-geliux15ftir}}

\emph{İki yönlü maruziyet. AB'ye ihracat \$267 mn, en büyük satır 2106 90 92 (\%12,8): S4'te −\$30 mn. AB'den ithalat \$396 mn, gıda takviyesi ve müstahzarı: S4'te c2 \%30, TÜFE 0,7 bp. Mayada dünya ikincisi (\%10,6).}

\textbf{1. Kimlik.} Rev. 2.1: 10.89, 11.07 (1:2); CPA 10.89. Başlıca HS4 (TR ihracatı 2024): 2106 \%70, 2102 \%23, 2104 \%3, 1302 \%2. BEC payları (ticaret ağırlıklı): ara \%25 / sermaye \%0 / tüketim \%75.

\textbf{2. Ölçek}

\begingroup\gbtablo\small

\begin{longtable}[]{@{}
  >{\raggedright\arraybackslash}p{(\columnwidth - 4\tabcolsep) * \real{0.3051}}
  >{\raggedright\arraybackslash}p{(\columnwidth - 4\tabcolsep) * \real{0.3495}}
  >{\raggedright\arraybackslash}p{(\columnwidth - 4\tabcolsep) * \real{0.3454}}@{}}
\toprule\noalign{}
\begin{minipage}[b]{\linewidth}\raggedright
\textbf{Gösterge}
\end{minipage} & \begin{minipage}[b]{\linewidth}\raggedright
\textbf{Değer}
\end{minipage} & \begin{minipage}[b]{\linewidth}\raggedright
\textbf{Etiket ve kaynak}
\end{minipage} \\
\midrule\noalign{}
\endhead
\bottomrule\noalign{}
\endlastfoot
Üretim değeri; katma değer & \$1.372 mn; \$467 mn (GSYH'nin \%0,03) & TÜİK YSHİ 2024 \etiketV{} \\
Çalışan; girişim (0 / 1--9 / 10--49 / 50--249 / 250+) & 10.858; 1.167 (174 / 988 / 93 / 31 / 8) & TÜİK YSHİ ve İş Kayıtları 2024 \etiketV{}; gizli büyüklük hücreleri \etiketT{} \\
250+ girişimlerin istihdam payı; AB'ye göre kişi başı KD & \%22; 0,46 & \etiketT{} (AB27 SBS 2023, nominal kur) \\
İlk 3 İBBS-2 bölgesi & İstanbul \%34; Kocaeli, Sakarya, Düzce, Bolu, Yalova \%10; İzmir \%10 & \etiketT{} iş kayıtları + SGK 2025, RAS \\
Yabancı sermaye; çapa yatırımcılar & Pakmaya (maya), Dr. Oetker, Unilever (Knorr), Nestlé \etiketH{} & \etiketH{} (firma bilgisi; pay verisi yok: VB-30) \\
\end{longtable}

\endgroup

\textbf{3. Dış ticaret}

\begingroup\gbtablo\small

\begin{longtable}[]{@{}
  >{\raggedright\arraybackslash}p{(\columnwidth - 2\tabcolsep) * \real{0.4802}}
  >{\raggedright\arraybackslash}p{(\columnwidth - 2\tabcolsep) * \real{0.5198}}@{}}
\toprule\noalign{}
\begin{minipage}[b]{\linewidth}\raggedright
\textbf{Gösterge}
\end{minipage} & \begin{minipage}[b]{\linewidth}\raggedright
\textbf{Değer (BACI 2024, \$; \etiketV{}, eşleme \etiketT{})}
\end{minipage} \\
\midrule\noalign{}
\endhead
\bottomrule\noalign{}
\endlastfoot
İhracat (2024); 2020--24 değişimi & \$1.114 mn; \%59 \\
AB payı; ilk 5 AB pazarı & \%24; CZE \%26; DEU \%19; NLD \%14; ITA \%9; ESP \%7 \\
İlk 5 AB dışı pazar & IRQ \%15; USA \%7; DZA \%4; SAU \%4; AZE \%3 \\
İthalat; AB ve Çin payı; ilk kaynaklar & \$639 mn; AB \%62, Çin \%6; DEU \%15; ITA \%12; USA \%8; NLD \%8; FRA \%6 \\
Dış ticaret dengesi & \$+475 mn \\
İthal girdi içeriği (I-O, toplam) / KD÷üretim & \%20 / \%34 \\
RCA dünya (2020 → 2024) / AB pazarında & 0,89 → 0,99 / 1,10 \\
Birim değer ÷ AB pazarı ortalaması & 1,14 \\
AB pazar payı (AB dışı tedarik içinde, 2020 → 2024); TR'nin sırası & \%3,3 → \%4,4; 5. \\
\end{longtable}

\endgroup

\textbf{4. AB tarafı maruziyet}

\begin{itemize}
\tightlist
\item
  \kod{τ\_\allowbreak{}AB} (S4 artışı, rota A, TR→AB 2025 ağırlıklı): \%11,8; 1/98 kaybıyla \%12,0; tam MFN ortalaması \%13,0. Rejim: GB \%1, İTÜ \%98, tarım \%1 \etiketV{}/\etiketT{}.
\item
  Tepe satırlar (S4 artışı ≥ \%10; en büyük 3): 10 satır, AB ihracatının \%83'i: KN 2106 90 92 \%12,8 (€174,7 mn); KN 2102 10 31 \%12,0 (€9,4 mn); KN 2104 10 00 \%11,5 (€5,6 mn).
\item
  AB pazarında AB dışı rakipler (tercih durumu): GBR \%22; CHN \%15; USA \%14; CHE \%10; IND \%4. Tercihli rakip payı \%50.
\item
  Menşe (S2 / S3): risk D/O; \kod{c\_\allowbreak{}RoO} \%2,0 / \%3,0; \kod{p\_\allowbreak{}uyum} 0,85 / 0,80 \etiketH{}. Kural: 2106 ve 2102: başlık değişikliği; şeker ve Fasıl 4 ağırlık sınırları.
\item
  EUDR kapsamındaki AB ihracatı payı: \%0; SKDM: \%0.
\item
  AB'ye ihracatta en büyük satır KN 2106 90 92 (€175 mn; süt yağı ve şeker içermeyen müstahzarlar). İçerik \etiketD{}: aroma ve tatlandırıcı karışımları olabilir.
\end{itemize}

\textbf{5. Türkiye tarafı maruziyet}

\begin{itemize}
\tightlist
\item
  Uygulanan tarife (2026, ithalat ağırlıklı): MFN \%25,7, İGV dahil \%25,9 \etiketV{}. S4'te AB menşeline (c2): \%25,4.
\item
  2106 90'da Türkiye'nin MFN'si yüksek (satırlarda \%17--58,5). S4'te AB'den gelen takviye ve müstahzarlara c2 olarak yansıyor. İGV kaldırma ölçütü (AB baskın, yerli zayıf) bazı satırlarda karşılanıyor.
\end{itemize}

\textbf{6. Kabiliyet}

\begin{itemize}
\tightlist
\item
  KS 53,1 (sıra 76/230; ağırlık duyarlılığında P10--P90 39--130). Bileşenler: RCA 48, PCI 43 (PCI −0,13), yurt içi KD 83, Ar-Ge 25, tedarikçi 84, verimlilik 49, firma 56 \etiketT{}.
\item
  Ürün uzayında atlanabilir ürünler (sınıf içi; yoğunluk üst çeyrek, PCI \textgreater{} sınıf): yok \etiketT{}.
\item
  \textbf{Geliştirme hedefi:} T+5: maya ve maya ekstraktında Afrika ve Asya payı.
\end{itemize}

\textbf{7. Stratejik önem.} SS 38,1 (en yüksek boyut: gıda güvenliği 50; ağırlıklı ortalama 10,5) \etiketH{}. Maya (ekmek zinciri) ve gıda müstahzarları

\textbf{8. Makro duyarlılık}

\begingroup\gbtablo\small

\begin{longtable}[]{@{}lr@{}}
\toprule\noalign{}
\textbf{Gösterge} & \textbf{Değer (\etiketT{})} \\
\midrule\noalign{}
\endhead
\bottomrule\noalign{}
\endlastfoot
TÜFE ağırlığı (doğrudan, COICOP 2018 eşlemesi); tüketimde ithal pay & \%0,07; \%8 \\
ÜFE vekili (imalat üretimindeki pay) & \%0,21 \\
EDE: 1 puan tarife, tüm menşe / yalnız AB (bp; ρ = 1) & 0,054 / 0,033 \\
Geçişkenlik ρ (senaryolarda) & 0,6--1,0 (P50 0,8) \\
Piyasa yoğunlaşması (250+ istihdam payı) & \%22 \\
CKE: tarife (\$ mn / bp) ve mali (\$ mn / mlr TL); 0--100 puan & 438,5 ve 17,26; 36 \\
\end{longtable}

\endgroup

\textbf{9. Senaryo etkileri, T+5, P50 (P10 / P90)}

\begingroup\gbtablo\small

\begin{longtable}[]{@{}
  >{\raggedright\arraybackslash}p{(\columnwidth - 6\tabcolsep) * \real{0.3667}}
  >{\raggedleft\arraybackslash}p{(\columnwidth - 6\tabcolsep) * \real{0.2111}}
  >{\raggedleft\arraybackslash}p{(\columnwidth - 6\tabcolsep) * \real{0.2111}}
  >{\raggedleft\arraybackslash}p{(\columnwidth - 6\tabcolsep) * \real{0.2111}}@{}}
\toprule\noalign{}
\begin{minipage}[b]{\linewidth}\raggedright
\textbf{Değişken}
\end{minipage} & \begin{minipage}[b]{\linewidth}\raggedleft
\textbf{S2}
\end{minipage} & \begin{minipage}[b]{\linewidth}\raggedleft
\textbf{S3}
\end{minipage} & \begin{minipage}[b]{\linewidth}\raggedleft
\textbf{S4}
\end{minipage} \\
\midrule\noalign{}
\endhead
\bottomrule\noalign{}
\endlastfoot
AB'ye ihracat (\$ mn) & −12 (−28 / +5) & −15 (−30 / −3) & −30 (−49 / −13) \\
Toplam ihracat (\$ mn) & −9 (−21 / +4) & −11 (−22 / −2) & −24 (−40 / −10) \\
Nihai mal ithalatı (\$ mn) & −5 (−17 / +15) & −8 (−20 / +9) & −25 (−45 / −5) \\
Cari denge (\$ mn) & −4 (−9 / 0) & −3 (−9 / +1) & 0 (−4 / +5) \\
Üretim (\$ mn) & −5 (−11 / 0) & −4 (−11 / +1) & 0 (−5 / +7) \\
İstihdam (kişi) & −25 (−63 / 0) & −21 (−63 / +6) & +2 (−28 / +37) \\
TÜFE (bp, birinci tur) & +0,08 (+0,07 / +0,09) & +0,15 (+0,13 / +0,18) & +0,71 (+0,61 / +0,80) \\
Yatırım ve DYY (bütün senaryolar) & Etkisi küçük & & \\
\end{longtable}

\endgroup

\emph{Kaynak:} \etiketT{} f3e: kısmi denge, toplamlar Faz 2B GE dağılımına kalibre; S4'te 1/98 kaybı olasılığı 0,25.

\textbf{12. Erken uyarı ve tetikler}

\begingroup\gbtablo\small

\begin{longtable}[]{@{}
  >{\raggedright\arraybackslash}p{(\columnwidth - 4\tabcolsep) * \real{0.2536}}
  >{\raggedright\arraybackslash}p{(\columnwidth - 4\tabcolsep) * \real{0.1140}}
  >{\raggedright\arraybackslash}p{(\columnwidth - 4\tabcolsep) * \real{0.6325}}@{}}
\toprule\noalign{}
\begin{minipage}[b]{\linewidth}\raggedright
\textbf{Gösterge}
\end{minipage} & \begin{minipage}[b]{\linewidth}\raggedright
\textbf{Eşik}
\end{minipage} & \begin{minipage}[b]{\linewidth}\raggedright
\textbf{Tetiklenecek önlem}
\end{minipage} \\
\midrule\noalign{}
\endhead
\bottomrule\noalign{}
\endlastfoot
2106 90 ithalat fiyatı (S4) & Yıllık +\%15 & İGV'nin bu satırlarda kaldırılması (c2 kuralının gözden geçirilmesi) \\
\end{longtable}

\endgroup

\textbf{13. Kazanan ve kaybedenler; bölgesel ve sosyal etki; telafi.} S4'te iç ikame AB ihracat kaybını dengeliyor (P50 +2 kişi).

\textbf{15. Güven düzeyi; veri boşlukları; \etiketD{}.} Orta (en büyük satırın içeriği \etiketD{}).

\hypertarget{nace_kartlar_10-12-324-1091-uxe7iftlik-hayvanlarux131-iuxe7in-hazux131r-yem-imalatux131--gb-statuxfcsuxfc-tarim-198--birincilikincil-arketip-db--duruux15f-geliux15ftir}{%
\subsection{3.24 10.91 Çiftlik hayvanları için hazır yem imalatı --- GB statüsü: TARIM-1/98 \textbar{} Birincil/ikincil arketip: D/B \textbar{} Duruş: Geliştir}\label{nace_kartlar_10-12-324-1091-uxe7iftlik-hayvanlarux131-iuxe7in-hazux131r-yem-imalatux131--gb-statuxfcsuxfc-tarim-198--birincilikincil-arketip-db--duruux15f-geliux15ftir}}

\emph{Karma yem: hayvancılığın ana girdisi. İhracat AB dışına (Irak, Suriye). GB etkisi yok; soya küspesi, mısır ve DDGS ithalatına bağımlı.}

\textbf{1. Kimlik.} Rev. 2.1: 10.91 (1:1); CPA 10.91. Başlıca HS4 (TR ihracatı 2024): 2309 \%100, 1214 \%0. BEC payları (ticaret ağırlıklı): ara \%100 / sermaye \%0 / tüketim \%0.

\textbf{2. Ölçek}

\begingroup\gbtablo\small

\begin{longtable}[]{@{}
  >{\raggedright\arraybackslash}p{(\columnwidth - 4\tabcolsep) * \real{0.3004}}
  >{\raggedright\arraybackslash}p{(\columnwidth - 4\tabcolsep) * \real{0.3603}}
  >{\raggedright\arraybackslash}p{(\columnwidth - 4\tabcolsep) * \real{0.3393}}@{}}
\toprule\noalign{}
\begin{minipage}[b]{\linewidth}\raggedright
\textbf{Gösterge}
\end{minipage} & \begin{minipage}[b]{\linewidth}\raggedright
\textbf{Değer}
\end{minipage} & \begin{minipage}[b]{\linewidth}\raggedright
\textbf{Etiket ve kaynak}
\end{minipage} \\
\midrule\noalign{}
\endhead
\bottomrule\noalign{}
\endlastfoot
Üretim değeri; katma değer & \$6.390 mn; \$846 mn (GSYH'nin \%0,06) & TÜİK YSHİ 2024 \etiketV{} \\
Çalışan; girişim (0 / 1--9 / 10--49 / 50--249 / 250+) & 14.598; 757 (82 / 512 / 162 / 49 / 9) & TÜİK YSHİ ve İş Kayıtları 2024 \etiketV{}; gizli büyüklük hücreleri \etiketT{} \\
250+ girişimlerin istihdam payı; AB'ye göre kişi başı KD & \%29; 0,60 & \etiketT{} (AB27 SBS 2023, nominal kur) \\
İlk 3 İBBS-2 bölgesi & Manisa, Afyon, Kütahya, Uşak \%13; Balıkesir, Çanakkale \%11; İzmir \%10 & \etiketT{} iş kayıtları + SGK 2025, RAS \\
Yabancı sermaye; çapa yatırımcılar & Bolca orta ölçekli üretici; Cargill, Kocatepe Yem, Abalıoğlu \etiketH{} & \etiketH{} (firma bilgisi; pay verisi yok: VB-30) \\
\end{longtable}

\endgroup

\textbf{3. Dış ticaret}

\begingroup\gbtablo\small

\begin{longtable}[]{@{}
  >{\raggedright\arraybackslash}p{(\columnwidth - 2\tabcolsep) * \real{0.4802}}
  >{\raggedright\arraybackslash}p{(\columnwidth - 2\tabcolsep) * \real{0.5198}}@{}}
\toprule\noalign{}
\begin{minipage}[b]{\linewidth}\raggedright
\textbf{Gösterge}
\end{minipage} & \begin{minipage}[b]{\linewidth}\raggedright
\textbf{Değer (BACI 2024, \$; \etiketV{}, eşleme \etiketT{})}
\end{minipage} \\
\midrule\noalign{}
\endhead
\bottomrule\noalign{}
\endlastfoot
İhracat (2024); 2020--24 değişimi & \$352 mn; \%128 \\
AB payı; ilk 5 AB pazarı & \%1; NLD \%27; ESP \%10; ROU \%9; GRC \%9; BEL \%8 \\
İlk 5 AB dışı pazar & IRQ \%36; SYR \%25; RUS \%6; TUN \%4; URY \%4 \\
İthalat; AB ve Çin payı; ilk kaynaklar & \$279 mn; AB \%50, Çin \%9; FRA \%13; RUS \%10; CHN \%9; ITA \%7; USA \%6 \\
Dış ticaret dengesi & \$+73 mn \\
İthal girdi içeriği (I-O, toplam) / KD÷üretim & \%20 / \%13 \\
RCA dünya (2020 → 2024) / AB pazarında & 0,82 → 1,40 / 0,10 \\
Birim değer ÷ AB pazarı ortalaması & 2,59 \\
AB pazar payı (AB dışı tedarik içinde, 2020 → 2024); TR'nin sırası & \%0,4 → \%0,4; 13. \\
\end{longtable}

\endgroup

\textbf{4. AB tarafı maruziyet}

\begin{itemize}
\tightlist
\item
  \kod{τ\_\allowbreak{}AB} (S4 artışı, rota A, TR→AB 2025 ağırlıklı): \%0,0; 1/98 kaybıyla \%9,3; tam MFN ortalaması \%9,3. Rejim: GB \%0, İTÜ \%0, tarım \%100 \etiketV{}/\etiketT{}.
\item
  Tepe satırlar (S4 artışı ≥ \%10; en büyük 3): yok.
\item
  AB pazarında AB dışı rakipler (tercih durumu): CHN \%31; GBR \%25; USA \%12; NOR \%9; RUS \%4. Tercihli rakip payı \%42.
\item
  Menşe (S2 / S3): risk D/D; \kod{c\_\allowbreak{}RoO} \%1,5 / \%2,0; \kod{p\_\allowbreak{}uyum} 0,90 / 0,85 \etiketH{}. Kural: 2309: tahıl, şeker pancarı ve Fasıl 2--4 malzemeleri tamamen elde edilmiş; ağırlık sınırı.
\item
  EUDR kapsamındaki AB ihracatı payı: \%0; SKDM: \%0.
\end{itemize}

\textbf{5. Türkiye tarafı maruziyet}

\begin{itemize}
\tightlist
\item
  Uygulanan tarife (2026, ithalat ağırlıklı): MFN \%6,8, İGV dahil \%6,8 \etiketV{}. S4'te AB menşeline (c2): \%0,0.
\item
  Kritik girdiler (HS4 2304, 1005, 1201, 2303), Türkiye'nin toplam ithalatı 2024: \$3.768 mn (AB \%4; BRA \%33, UKR \%23, USA \%13) \etiketV{}.
\end{itemize}

\textbf{6. Kabiliyet}

\begin{itemize}
\tightlist
\item
  KS 49,3 (sıra 125/230; ağırlık duyarlılığında P10--P90 65--165). Bileşenler: RCA 66, PCI 18 (PCI −0,76), yurt içi KD 48, Ar-Ge 25, tedarikçi 84, verimlilik 76, firma 58 \etiketT{}.
\item
  Ürün uzayında atlanabilir ürünler (sınıf içi; yoğunluk üst çeyrek, PCI \textgreater{} sınıf): yok \etiketT{}.
\end{itemize}

\textbf{7. Stratejik önem.} SS 63,3 (en yüksek boyut: gıda güvenliği 85; ağırlıklı ortalama 12,8) \etiketH{}. Yem, kanatlı ve süt zincirinin maliyet tabanı.

\textbf{8. Makro duyarlılık}

\begingroup\gbtablo\small

\begin{longtable}[]{@{}lr@{}}
\toprule\noalign{}
\textbf{Gösterge} & \textbf{Değer (\etiketT{})} \\
\midrule\noalign{}
\endhead
\bottomrule\noalign{}
\endlastfoot
TÜFE ağırlığı (doğrudan, COICOP 2018 eşlemesi); tüketimde ithal pay & \%0,00; \%9 \\
ÜFE vekili (imalat üretimindeki pay) & \%0,96 \\
EDE: 1 puan tarife, tüm menşe / yalnız AB (bp; ρ = 1) & 0,021 / 0,011 \\
Geçişkenlik ρ (senaryolarda) & 0,6--1,0 (P50 0,8) \\
Piyasa yoğunlaşması (250+ istihdam payı) & \%29 \\
CKE: tarife (\$ mn / bp) ve mali (\$ mn / mlr TL); 0--100 puan & 1.600,7 ve 1,96; 30 \\
\end{longtable}

\endgroup

\textbf{9. Senaryo etkileri, T+5, P50 (P10 / P90)}

\begingroup\gbtablo\small

\begin{longtable}[]{@{}lrrr@{}}
\toprule\noalign{}
\textbf{Değişken} & \textbf{S2} & \textbf{S3} & \textbf{S4} \\
\midrule\noalign{}
\endhead
\bottomrule\noalign{}
\endlastfoot
AB'ye ihracat (\$ mn) & 0 (0 / 0) & 0 (0 / 0) & 0 (−1 / 0) \\
Toplam ihracat (\$ mn) & 0 (0 / 0) & 0 (0 / 0) & 0 (−1 / 0) \\
Nihai mal ithalatı (\$ mn) & 0 (0 / +10) & 0 (0 / +6) & 0 (0 / 0) \\
Cari denge (\$ mn) & 0 (−8 / 0) & 0 (−5 / 0) & 0 (−1 / 0) \\
Üretim (\$ mn) & 0 (−10 / 0) & 0 (−6 / 0) & 0 (−1 / 0) \\
İstihdam (kişi) & 0 (−16 / 0) & 0 (−9 / 0) & 0 (−1 / 0) \\
TÜFE (bp, birinci tur) & 0,00 (0,00 / 0,00) & 0,00 (0,00 / 0,00) & 0,00 (0,00 / 0,00) \\
Yatırım ve DYY (bütün senaryolar) & --- & & \\
\end{longtable}

\endgroup

\emph{Kaynak:} \etiketT{} f3e: kısmi denge, toplamlar Faz 2B GE dağılımına kalibre; S4'te 1/98 kaybı olasılığı 0,25.

\textbf{13. Kazanan ve kaybedenler; bölgesel ve sosyal etki; telafi.} Etkilenen kesim yok.

\textbf{15. Güven düzeyi; veri boşlukları; \etiketD{}.} Yüksek.

\hypertarget{nace_kartlar_10-12-325-1092-ev-hayvanlarux131-iuxe7in-hazux131r-gux131da-imalatux131--gb-statuxfcsuxfc-tarim-198--birincilikincil-arketip-c--duruux15f-geliux15ftir}{%
\subsection{3.25 10.92 Ev hayvanları için hazır gıda imalatı --- GB statüsü: TARIM-1/98 \textbar{} Birincil/ikincil arketip: C \textbar{} Duruş: Geliştir}\label{nace_kartlar_10-12-325-1092-ev-hayvanlarux131-iuxe7in-hazux131r-gux131da-imalatux131--gb-statuxfcsuxfc-tarim-198--birincilikincil-arketip-c--duruux15f-geliux15ftir}}

\emph{Evcil hayvan maması: ithalatın \%87'si AB'den (Fransa, Polonya). Yerli üretim hızla büyüyor (RCA 0,17 → 0,49). Tarım rejiminde; GB çıkışı tarifeyi değiştirmiyor.}

\textbf{1. Kimlik.} Rev. 2.1: 10.92 (1:1); CPA 10.92. Başlıca HS4 (TR ihracatı 2024): 2309 \%100. BEC payları (ticaret ağırlıklı): ara \%0 / sermaye \%0 / tüketim \%100.

\textbf{2. Ölçek}

\begingroup\gbtablo\small

\begin{longtable}[]{@{}
  >{\raggedright\arraybackslash}p{(\columnwidth - 4\tabcolsep) * \real{0.3075}}
  >{\raggedright\arraybackslash}p{(\columnwidth - 4\tabcolsep) * \real{0.3439}}
  >{\raggedright\arraybackslash}p{(\columnwidth - 4\tabcolsep) * \real{0.3485}}@{}}
\toprule\noalign{}
\begin{minipage}[b]{\linewidth}\raggedright
\textbf{Gösterge}
\end{minipage} & \begin{minipage}[b]{\linewidth}\raggedright
\textbf{Değer}
\end{minipage} & \begin{minipage}[b]{\linewidth}\raggedright
\textbf{Etiket ve kaynak}
\end{minipage} \\
\midrule\noalign{}
\endhead
\bottomrule\noalign{}
\endlastfoot
Üretim değeri; katma değer & \$347 mn; \$81 mn (GSYH'nin \%0,01) & TÜİK YSHİ 2024 \etiketV{} \\
Çalışan; girişim (0 / 1--9 / 10--49 / 50--249 / 250+) & 1.768; 140 (21 / 117 / 18 / --- / ---) & TÜİK YSHİ ve İş Kayıtları 2024 \etiketV{}; gizli büyüklük hücreleri \etiketT{} \\
250+ girişimlerin istihdam payı; AB'ye göre kişi başı KD & \%21; 0,44 & \etiketT{} (AB27 SBS 2023, nominal kur) \\
İlk 3 İBBS-2 bölgesi & İzmir \%38; Manisa, Afyon, Kütahya, Uşak \%16; Adana, Mersin \%11 & \etiketT{} iş kayıtları + SGK 2025, RAS \\
Yabancı sermaye; çapa yatırımcılar & Mars (Royal Canin) ve Nestlé Purina ithalatı; yerli üreticiler \etiketH{} & \etiketH{} (firma bilgisi; pay verisi yok: VB-30) \\
\end{longtable}

\endgroup

\textbf{3. Dış ticaret}

\begingroup\gbtablo\small

\begin{longtable}[]{@{}
  >{\raggedright\arraybackslash}p{(\columnwidth - 2\tabcolsep) * \real{0.4766}}
  >{\raggedright\arraybackslash}p{(\columnwidth - 2\tabcolsep) * \real{0.5234}}@{}}
\toprule\noalign{}
\begin{minipage}[b]{\linewidth}\raggedright
\textbf{Gösterge}
\end{minipage} & \begin{minipage}[b]{\linewidth}\raggedright
\textbf{Değer (BACI 2024, \$; \etiketV{}, eşleme \etiketT{})}
\end{minipage} \\
\midrule\noalign{}
\endhead
\bottomrule\noalign{}
\endlastfoot
İhracat (2024); 2020--24 değişimi & \$147 mn; \%382 \\
AB payı; ilk 5 AB pazarı & \%14; GRC \%32; BGR \%20; DEU \%13; ROU \%12; ITA \%6 \\
İlk 5 AB dışı pazar & MYS \%15; IRQ \%12; SAU \%6; USA \%5; ISR \%5 \\
İthalat; AB ve Çin payı; ilk kaynaklar & \$214 mn; AB \%87, Çin \%4; FRA \%39; POL \%14; HUN \%10; ITA \%6; NLD \%5 \\
Dış ticaret dengesi & \$−67 mn \\
İthal girdi içeriği (I-O, toplam) / KD÷üretim & \%20 / \%23 \\
RCA dünya (2020 → 2024) / AB pazarında & 0,17 → 0,49 / 0,25 \\
Birim değer ÷ AB pazarı ortalaması & 0,50 \\
AB pazar payı (AB dışı tedarik içinde, 2020 → 2024); TR'nin sırası & \%0,4 → \%1,0; 10. \\
\end{longtable}

\endgroup

\textbf{4. AB tarafı maruziyet}

\begin{itemize}
\tightlist
\item
  \kod{τ\_\allowbreak{}AB} (S4 artışı, rota A, TR→AB 2025 ağırlıklı): \%0,0; 1/98 kaybıyla \%5,7; tam MFN ortalaması \%5,7. Rejim: GB \%0, İTÜ \%0, tarım \%100 \etiketV{}/\etiketT{}.
\item
  Tepe satırlar (S4 artışı ≥ \%10; en büyük 3): yok.
\item
  AB pazarında AB dışı rakipler (tercih durumu): CHN \%26; THA \%18; CHE \%18; GBR \%17; CAN \%5. Tercihli rakip payı \%44.
\item
  Menşe (S2 / S3): risk D/D; \kod{c\_\allowbreak{}RoO} \%1,5 / \%2,0; \kod{p\_\allowbreak{}uyum} 0,90 / 0,85 \etiketH{}. Kural: 2309 10: aynı.
\item
  EUDR kapsamındaki AB ihracatı payı: \%0; SKDM: \%0.
\end{itemize}

\textbf{5. Türkiye tarafı maruziyet}

\begin{itemize}
\tightlist
\item
  Uygulanan tarife (2026, ithalat ağırlıklı): MFN \%8,5, İGV dahil \%8,5 \etiketV{}. S4'te AB menşeline (c2): \%0,0.
\end{itemize}

\textbf{6. Kabiliyet}

\begin{itemize}
\tightlist
\item
  KS 48,3 (sıra 133/230; ağırlık duyarlılığında P10--P90 95--167). Bileşenler: RCA 46, PCI 51 (PCI 0,03), yurt içi KD 58, Ar-Ge 25, tedarikçi 84, verimlilik 43, firma 38 \etiketT{}.
\item
  Ürün uzayında atlanabilir ürünler (sınıf içi; yoğunluk üst çeyrek, PCI \textgreater{} sınıf): yok \etiketT{}.
\end{itemize}

\textbf{7. Stratejik önem.} SS 3,7 (en yüksek boyut: gıda güvenliği 5; ağırlıklı ortalama 0,8) \etiketH{}.

\textbf{8. Makro duyarlılık}

\begingroup\gbtablo\small

\begin{longtable}[]{@{}lr@{}}
\toprule\noalign{}
\textbf{Gösterge} & \textbf{Değer (\etiketT{})} \\
\midrule\noalign{}
\endhead
\bottomrule\noalign{}
\endlastfoot
TÜFE ağırlığı (doğrudan, COICOP 2018 eşlemesi); tüketimde ithal pay & \%0,19; \%12 \\
ÜFE vekili (imalat üretimindeki pay) & \%0,05 \\
EDE: 1 puan tarife, tüm menşe / yalnız AB (bp; ρ = 1) & 0,040 / 0,035 \\
Geçişkenlik ρ (senaryolarda) & 0,6--1,0 (P50 0,8) \\
Piyasa yoğunlaşması (250+ istihdam payı) & \%21 \\
CKE: tarife (\$ mn / bp) ve mali (\$ mn / mlr TL); 0--100 puan & 585,6 ve 22,92; 41 \\
\end{longtable}

\endgroup

\textbf{9. Senaryo etkileri, T+5, P50 (P10 / P90)}

\begingroup\gbtablo\small

\begin{longtable}[]{@{}lrrr@{}}
\toprule\noalign{}
\textbf{Değişken} & \textbf{S2} & \textbf{S3} & \textbf{S4} \\
\midrule\noalign{}
\endhead
\bottomrule\noalign{}
\endlastfoot
AB'ye ihracat (\$ mn) & 0 (0 / 0) & 0 (0 / 0) & 0 (−2 / 0) \\
Toplam ihracat (\$ mn) & 0 (0 / 0) & 0 (0 / 0) & 0 (−2 / 0) \\
Nihai mal ithalatı (\$ mn) & 0 (0 / +12) & 0 (0 / +7) & 0 (0 / 0) \\
Cari denge (\$ mn) & 0 (−9 / 0) & 0 (−6 / 0) & 0 (−2 / 0) \\
Üretim (\$ mn) & 0 (−11 / 0) & 0 (−7 / 0) & 0 (−2 / 0) \\
İstihdam (kişi) & 0 (−41 / 0) & 0 (−24 / 0) & 0 (−7 / 0) \\
TÜFE (bp, birinci tur) & 0,00 (0,00 / 0,00) & 0,00 (0,00 / 0,00) & 0,00 (0,00 / 0,00) \\
Yatırım ve DYY (bütün senaryolar) & --- & & \\
\end{longtable}

\endgroup

\emph{Kaynak:} \etiketT{} f3e: kısmi denge, toplamlar Faz 2B GE dağılımına kalibre; S4'te 1/98 kaybı olasılığı 0,25.

\textbf{13. Kazanan ve kaybedenler; bölgesel ve sosyal etki; telafi.} Etkilenen kesim yok.

\textbf{15. Güven düzeyi; veri boşlukları; \etiketD{}.} Orta.

\hypertarget{nace_kartlar_10-12-326-1101-alkolluxfc-iuxe7eceklerin-damux131tux131lmasux131-arux131tux131lmasux131-ve-harmanlanmasux131--gb-statuxfcsuxfc-iux307tuxfc--birincilikincil-arketip---duruux15f-koru}{%
\subsection{3.26 11.01 Alkollü içeceklerin damıtılması, arıtılması ve harmanlanması --- GB statüsü: İTÜ \textbar{} Birincil/ikincil arketip: --- \textbar{} Duruş: Koru}\label{nace_kartlar_10-12-326-1101-alkolluxfc-iuxe7eceklerin-damux131tux131lmasux131-arux131tux131lmasux131-ve-harmanlanmasux131--gb-statuxfcsuxfc-iux307tuxfc--birincilikincil-arketip---duruux15f-koru}}

\emph{Rakı (Mey\textbar Diageo) ve diğer alkollü içkiler. AB'nin MFN'si sıfıra yakın; GB etkisi yok. İthalatın \%63'ü İngiltere'den (viski; TR--BK STA'sı).}

\textbf{1. Kimlik.} Rev. 2.1: 11.01 (1:1); CPA 11.01. Başlıca HS4 (TR ihracatı 2024): 2208 \%100. BEC payları (ticaret ağırlıklı): ara \%0 / sermaye \%0 / tüketim \%100.

\textbf{2. Ölçek}

\begingroup\gbtablo\small

\begin{longtable}[]{@{}
  >{\raggedright\arraybackslash}p{(\columnwidth - 4\tabcolsep) * \real{0.2925}}
  >{\raggedright\arraybackslash}p{(\columnwidth - 4\tabcolsep) * \real{0.3784}}
  >{\raggedright\arraybackslash}p{(\columnwidth - 4\tabcolsep) * \real{0.3290}}@{}}
\toprule\noalign{}
\begin{minipage}[b]{\linewidth}\raggedright
\textbf{Gösterge}
\end{minipage} & \begin{minipage}[b]{\linewidth}\raggedright
\textbf{Değer}
\end{minipage} & \begin{minipage}[b]{\linewidth}\raggedright
\textbf{Etiket ve kaynak}
\end{minipage} \\
\midrule\noalign{}
\endhead
\bottomrule\noalign{}
\endlastfoot
Üretim değeri; katma değer & \$1.210 mn; \$393 mn (GSYH'nin \%0,03) & TÜİK YSHİ 2024 \etiketV{} \\
Çalışan; girişim (0 / 1--9 / 10--49 / 50--249 / 250+) & 2.398; 21 (1 / 7 / 9 / 1 / 4) & TÜİK YSHİ ve İş Kayıtları 2024 \etiketV{}; gizli büyüklük hücreleri \etiketT{} \\
250+ girişimlerin istihdam payı; AB'ye göre kişi başı KD & \%85; --- & \etiketT{} (AB27 SBS 2023, nominal kur) \\
İlk 3 İBBS-2 bölgesi & İstanbul \%50; Manisa, Afyon, Kütahya, Uşak \%26; Antalya, Isparta, Burdur \%11 & \etiketT{} iş kayıtları + SGK 2025, RAS \\
Yabancı sermaye; çapa yatırımcılar & Mey/Diageo (2011'den beri Diageo) \etiketV{}, Efe Rakı, Tekirdağ (Mey) \etiketH{} & \etiketH{} (firma bilgisi; pay verisi yok: VB-30) \\
\end{longtable}

\endgroup

\textbf{3. Dış ticaret}

\begingroup\gbtablo\small

\begin{longtable}[]{@{}
  >{\raggedright\arraybackslash}p{(\columnwidth - 2\tabcolsep) * \real{0.4824}}
  >{\raggedright\arraybackslash}p{(\columnwidth - 2\tabcolsep) * \real{0.5176}}@{}}
\toprule\noalign{}
\begin{minipage}[b]{\linewidth}\raggedright
\textbf{Gösterge}
\end{minipage} & \begin{minipage}[b]{\linewidth}\raggedright
\textbf{Değer (BACI 2024, \$; \etiketV{}, eşleme \etiketT{})}
\end{minipage} \\
\midrule\noalign{}
\endhead
\bottomrule\noalign{}
\endlastfoot
İhracat (2024); 2020--24 değişimi & \$68 mn; \%92 \\
AB payı; ilk 5 AB pazarı & \%48; DEU \%69; BGR \%7; NLD \%7; AUT \%5; FRA \%3 \\
İlk 5 AB dışı pazar & CYP \%62; BLR \%13; USA \%6; IRQ \%4; GBR \%4 \\
İthalat; AB ve Çin payı; ilk kaynaklar & \$425 mn; AB \%25, Çin \%0; GBR \%63; DEU \%7; MEX \%5; IRL \%4; USA \%4 \\
Dış ticaret dengesi & \$−358 mn \\
İthal girdi içeriği (I-O, toplam) / KD÷üretim & \%20 / \%33 \\
RCA dünya (2020 → 2024) / AB pazarında & 0,11 → 0,14 / 0,17 \\
Birim değer ÷ AB pazarı ortalaması & 1,20 \\
AB pazar payı (AB dışı tedarik içinde, 2020 → 2024); TR'nin sırası & \%0,7 → \%0,7; 11. \\
\end{longtable}

\endgroup

\textbf{4. AB tarafı maruziyet}

\begin{itemize}
\tightlist
\item
  \kod{τ\_\allowbreak{}AB} (S4 artışı, rota A, TR→AB 2025 ağırlıklı): \%0,0; 1/98 kaybıyla \%0,4; tam MFN ortalaması \%0,4. Rejim: GB \%0, İTÜ \%100, tarım \%0 \etiketV{}/\etiketT{}.
\item
  Tepe satırlar (S4 artışı ≥ \%10; en büyük 3): yok.
\item
  AB pazarında AB dışı rakipler (tercih durumu): GBR \%54; USA \%23; MEX \%4; JPN \%2; CHN \%2. Tercihli rakip payı \%70.
\item
  Menşe (S2 / S3): risk D/D; \kod{c\_\allowbreak{}RoO} \%1,5 / \%2,0; \kod{p\_\allowbreak{}uyum} 0,90 / 0,85 \etiketH{}. Kural: 2208: başlık değişikliği (2207--2208 hariç); üzüm türevleri tamamen elde edilmiş (rakı: yerli suma).
\item
  EUDR kapsamındaki AB ihracatı payı: \%0; SKDM: \%0.
\end{itemize}

\textbf{5. Türkiye tarafı maruziyet}

\begin{itemize}
\tightlist
\item
  Uygulanan tarife (2026, ithalat ağırlıklı): MFN \%0,6, İGV dahil \%0,6 \etiketV{}. S4'te AB menşeline (c2): \%0,7.
\item
  ÖTV ile gümrük etkileşimi: alkollü içkide ÖTV maktu (litre başına) ve fiyatın çoğunu oluşturuyor. Gümrük vergisinin perakende fiyata oransal etkisi düşük.
\end{itemize}

\textbf{6. Kabiliyet}

\begin{itemize}
\tightlist
\item
  KS 41,0 (sıra 201/230; ağırlık duyarlılığında P10--P90 131--218). Bileşenler: RCA 29, PCI 29 (PCI −0,36), yurt içi KD 81, Ar-Ge 25, tedarikçi 84, verimlilik ---, firma 11 \etiketT{}.
\item
  Ürün uzayında atlanabilir ürünler (sınıf içi; yoğunluk üst çeyrek, PCI \textgreater{} sınıf): yok \etiketT{}.
\end{itemize}

\textbf{7. Stratejik önem.} SS 0,0 (en yüksek boyut: --- 0; ağırlıklı ortalama 0,0) \etiketH{}.

\textbf{8. Makro duyarlılık}

\begingroup\gbtablo\small

\begin{longtable}[]{@{}lr@{}}
\toprule\noalign{}
\textbf{Gösterge} & \textbf{Değer (\etiketT{})} \\
\midrule\noalign{}
\endhead
\bottomrule\noalign{}
\endlastfoot
TÜFE ağırlığı (doğrudan, COICOP 2018 eşlemesi); tüketimde ithal pay & \%0,28; \%18 \\
ÜFE vekili (imalat üretimindeki pay) & \%0,18 \\
EDE: 1 puan tarife, tüm menşe / yalnız AB (bp; ρ = 1) & 0,083 / 0,021 \\
Geçişkenlik ρ (senaryolarda) & 0,6--1,0 (P50 0,8) \\
Piyasa yoğunlaşması (250+ istihdam payı) & \%85 \\
CKE: tarife (\$ mn / bp) ve mali (\$ mn / mlr TL); 0--100 puan & 279,0 ve 6,69; 23 \\
\end{longtable}

\endgroup

\textbf{9. Senaryo etkileri, T+5, P50 (P10 / P90)}

\begingroup\gbtablo\small

\begin{longtable}[]{@{}
  >{\raggedright\arraybackslash}p{(\columnwidth - 6\tabcolsep) * \real{0.3667}}
  >{\raggedleft\arraybackslash}p{(\columnwidth - 6\tabcolsep) * \real{0.2111}}
  >{\raggedleft\arraybackslash}p{(\columnwidth - 6\tabcolsep) * \real{0.2111}}
  >{\raggedleft\arraybackslash}p{(\columnwidth - 6\tabcolsep) * \real{0.2111}}@{}}
\toprule\noalign{}
\begin{minipage}[b]{\linewidth}\raggedright
\textbf{Değişken}
\end{minipage} & \begin{minipage}[b]{\linewidth}\raggedleft
\textbf{S2}
\end{minipage} & \begin{minipage}[b]{\linewidth}\raggedleft
\textbf{S3}
\end{minipage} & \begin{minipage}[b]{\linewidth}\raggedleft
\textbf{S4}
\end{minipage} \\
\midrule\noalign{}
\endhead
\bottomrule\noalign{}
\endlastfoot
AB'ye ihracat (\$ mn) & −1 (−2 / +1) & −1 (−2 / 0) & 0 (−1 / 0) \\
Toplam ihracat (\$ mn) & −1 (−1 / 0) & −1 (−1 / 0) & 0 (−1 / 0) \\
Nihai mal ithalatı (\$ mn) & −1 (−3 / +6) & −1 (−3 / +4) & −1 (−2 / 0) \\
Cari denge (\$ mn) & 0 (−4 / +1) & 0 (−3 / +1) & 0 (0 / +1) \\
Üretim (\$ mn) & 0 (−5 / +1) & 0 (−4 / +2) & +1 (0 / +1) \\
İstihdam (kişi) & 0 (−7 / +2) & +1 (−5 / +2) & +1 (0 / +2) \\
TÜFE (bp, birinci tur) & +0,02 (+0,01 / +0,02) & +0,03 (+0,03 / +0,04) & +0,03 (+0,02 / +0,04) \\
Yatırım ve DYY (bütün senaryolar) & --- & & \\
\end{longtable}

\endgroup

\emph{Kaynak:} \etiketT{} f3e: kısmi denge, toplamlar Faz 2B GE dağılımına kalibre; S4'te 1/98 kaybı olasılığı 0,25.

\textbf{11. Müzakere pozisyonu}

\begin{itemize}
\tightlist
\item
  AB ile hücum: S2'de `Türk Rakısı' coğrafi işaretinin karşılıklı korunması \etiketD{} (AB'deki tescil durumu).
\end{itemize}

\textbf{13. Kazanan ve kaybedenler; bölgesel ve sosyal etki; telafi.} Etkilenen kesim yok.

\textbf{15. Güven düzeyi; veri boşlukları; \etiketD{}.} Orta (YSHİ gizli; bölüm kalanından tahmin).

\hypertarget{nace_kartlar_10-12-327-1102-uxfczuxfcmden-ux15farap-imalatux131--gb-statuxfcsuxfc-tarim-198--birincilikincil-arketip---duruux15f-koru}{%
\subsection{3.27 11.02 Üzümden şarap imalatı --- GB statüsü: TARIM-1/98 \textbar{} Birincil/ikincil arketip: --- \textbar{} Duruş: Koru}\label{nace_kartlar_10-12-327-1102-uxfczuxfcmden-ux15farap-imalatux131--gb-statuxfcsuxfc-tarim-198--birincilikincil-arketip---duruux15f-koru}}

\emph{Küçük şarap sınıfı (ihracat \$24 mn). Tarım rejiminde; GB etkisi yok.}

\textbf{1. Kimlik.} Rev. 2.1: 11.02 (1:1); CPA 11.02. Başlıca HS4 (TR ihracatı 2024): 2204 \%100. BEC payları (ticaret ağırlıklı): ara \%0 / sermaye \%0 / tüketim \%100.

\textbf{2. Ölçek}

\begingroup\gbtablo\small

\begin{longtable}[]{@{}
  >{\raggedright\arraybackslash}p{(\columnwidth - 4\tabcolsep) * \real{0.2995}}
  >{\raggedright\arraybackslash}p{(\columnwidth - 4\tabcolsep) * \real{0.3624}}
  >{\raggedright\arraybackslash}p{(\columnwidth - 4\tabcolsep) * \real{0.3381}}@{}}
\toprule\noalign{}
\begin{minipage}[b]{\linewidth}\raggedright
\textbf{Gösterge}
\end{minipage} & \begin{minipage}[b]{\linewidth}\raggedright
\textbf{Değer}
\end{minipage} & \begin{minipage}[b]{\linewidth}\raggedright
\textbf{Etiket ve kaynak}
\end{minipage} \\
\midrule\noalign{}
\endhead
\bottomrule\noalign{}
\endlastfoot
Üretim değeri; katma değer & \$278 mn; \$109 mn (GSYH'nin \%0,01) & TÜİK YSHİ 2024 \etiketV{} \\
Çalışan; girişim (0 / 1--9 / 10--49 / 50--249 / 250+) & 1.968; 196 (21 / 145 / 35 / 8 / 0) & TÜİK YSHİ ve İş Kayıtları 2024 \etiketV{}; gizli büyüklük hücreleri \etiketT{} \\
250+ girişimlerin istihdam payı; AB'ye göre kişi başı KD & \%0; 0,59 & \etiketT{} (AB27 SBS 2023, nominal kur) \\
İlk 3 İBBS-2 bölgesi & Aydın, Denizli, Muğla \%16; Ankara \%14; Tekirdağ, Edirne, Kırklareli \%14 & \etiketT{} iş kayıtları + SGK 2025, RAS \\
Yabancı sermaye; çapa yatırımcılar & Kavaklıdere, Doluca, Kayra \etiketH{} & \etiketH{} (firma bilgisi; pay verisi yok: VB-30) \\
\end{longtable}

\endgroup

\textbf{3. Dış ticaret}

\begingroup\gbtablo\small

\begin{longtable}[]{@{}
  >{\raggedright\arraybackslash}p{(\columnwidth - 2\tabcolsep) * \real{0.4824}}
  >{\raggedright\arraybackslash}p{(\columnwidth - 2\tabcolsep) * \real{0.5176}}@{}}
\toprule\noalign{}
\begin{minipage}[b]{\linewidth}\raggedright
\textbf{Gösterge}
\end{minipage} & \begin{minipage}[b]{\linewidth}\raggedright
\textbf{Değer (BACI 2024, \$; \etiketV{}, eşleme \etiketT{})}
\end{minipage} \\
\midrule\noalign{}
\endhead
\bottomrule\noalign{}
\endlastfoot
İhracat (2024); 2020--24 değişimi & \$24 mn; \%213 \\
AB payı; ilk 5 AB pazarı & \%61; GRC \%23; DEU \%21; NLD \%15; BEL \%12; FRA \%12 \\
İlk 5 AB dışı pazar & BLR \%34; GBR \%21; CYP \%14; USA \%11; RUS \%9 \\
İthalat; AB ve Çin payı; ilk kaynaklar & \$53 mn; AB \%86, Çin \%0; FRA \%43; ITA \%24; GRC \%7; ESP \%5; MDA \%4 \\
Dış ticaret dengesi & \$−29 mn \\
İthal girdi içeriği (I-O, toplam) / KD÷üretim & \%20 / \%39 \\
RCA dünya (2020 → 2024) / AB pazarında & 0,02 → 0,05 / 0,23 \\
Birim değer ÷ AB pazarı ortalaması & 1,72 \\
AB pazar payı (AB dışı tedarik içinde, 2020 → 2024); TR'nin sırası & \%0,2 → \%0,9; 12. \\
\end{longtable}

\endgroup

\textbf{4. AB tarafı maruziyet}

\begin{itemize}
\tightlist
\item
  \kod{τ\_\allowbreak{}AB} (S4 artışı, rota A, TR→AB 2025 ağırlıklı): \%0,0; 1/98 kaybıyla \%0,3; tam MFN ortalaması \%0,3. Rejim: GB \%0, İTÜ \%0, tarım \%100 \etiketV{}/\etiketT{}.
\item
  Tepe satırlar (S4 artışı ≥ \%10; en büyük 3): yok.
\item
  AB pazarında AB dışı rakipler (tercih durumu): CHL \%21; USA \%19; ZAF \%14; AUS \%12; NZL \%8. Tercihli rakip payı \%60.
\item
  Menşe (S2 / S3): risk D/D; \kod{c\_\allowbreak{}RoO} \%1,5 / \%2,0; \kod{p\_\allowbreak{}uyum} 0,90 / 0,85 \etiketH{}. Kural: 2204: üzüm ve türevleri tamamen elde edilmiş.
\item
  EUDR kapsamındaki AB ihracatı payı: \%0; SKDM: \%0.
\end{itemize}

\textbf{5. Türkiye tarafı maruziyet}

\begin{itemize}
\tightlist
\item
  Uygulanan tarife (2026, ithalat ağırlıklı): MFN \%69,9, İGV dahil \%69,9 \etiketV{}. S4'te AB menşeline (c2): \%0,0.
\end{itemize}

\textbf{6. Kabiliyet}

\begin{itemize}
\tightlist
\item
  KS 50,3 (sıra 110/230; ağırlık duyarlılığında P10--P90 45--168). Bileşenler: RCA 32, PCI 23 (PCI −0,56), yurt içi KD 90, Ar-Ge 25, tedarikçi 84, verimlilik 74, firma 62 \etiketT{}.
\item
  Ürün uzayında atlanabilir ürünler (sınıf içi; yoğunluk üst çeyrek, PCI \textgreater{} sınıf): 220422 (PCI -0.15, RCA 0.00) \etiketT{}.
\end{itemize}

\textbf{7. Stratejik önem.} SS 0,0 (en yüksek boyut: --- 0; ağırlıklı ortalama 0,0) \etiketH{}.

\textbf{8. Makro duyarlılık}

\begingroup\gbtablo\small

\begin{longtable}[]{@{}lr@{}}
\toprule\noalign{}
\textbf{Gösterge} & \textbf{Değer (\etiketT{})} \\
\midrule\noalign{}
\endhead
\bottomrule\noalign{}
\endlastfoot
TÜFE ağırlığı (doğrudan, COICOP 2018 eşlemesi); tüketimde ithal pay & \%0,05; \%15 \\
ÜFE vekili (imalat üretimindeki pay) & \%0,04 \\
EDE: 1 puan tarife, tüm menşe / yalnız AB (bp; ρ = 1) & 0,011 / 0,009 \\
Geçişkenlik ρ (senaryolarda) & 0,6--1,0 (P50 0,8) \\
Piyasa yoğunlaşması (250+ istihdam payı) & \%0 \\
CKE: tarife (\$ mn / bp) ve mali (\$ mn / mlr TL); 0--100 puan & 773,8 ve 10,18; 33 \\
\end{longtable}

\endgroup

\textbf{9. Senaryo etkileri, T+5, P50 (P10 / P90)}

\begingroup\gbtablo\small

\begin{longtable}[]{@{}lrrr@{}}
\toprule\noalign{}
\textbf{Değişken} & \textbf{S2} & \textbf{S3} & \textbf{S4} \\
\midrule\noalign{}
\endhead
\bottomrule\noalign{}
\endlastfoot
AB'ye ihracat (\$ mn) & 0 (0 / 0) & 0 (0 / 0) & 0 (0 / 0) \\
Toplam ihracat (\$ mn) & 0 (0 / 0) & 0 (0 / 0) & 0 (0 / 0) \\
Nihai mal ithalatı (\$ mn) & 0 (0 / +3) & 0 (0 / +2) & 0 (0 / 0) \\
Cari denge (\$ mn) & 0 (−2 / 0) & 0 (−1 / 0) & 0 (0 / 0) \\
Üretim (\$ mn) & 0 (−3 / 0) & 0 (−2 / 0) & 0 (0 / 0) \\
İstihdam (kişi) & 0 (−15 / 0) & 0 (−9 / 0) & 0 (−1 / 0) \\
TÜFE (bp, birinci tur) & 0,00 (0,00 / 0,00) & 0,00 (0,00 / 0,00) & 0,00 (0,00 / 0,00) \\
Yatırım ve DYY (bütün senaryolar) & --- & & \\
\end{longtable}

\endgroup

\emph{Kaynak:} \etiketT{} f3e: kısmi denge, toplamlar Faz 2B GE dağılımına kalibre; S4'te 1/98 kaybı olasılığı 0,25.

\textbf{13. Kazanan ve kaybedenler; bölgesel ve sosyal etki; telafi.} Etkilenen kesim yok.

\textbf{15. Güven düzeyi; veri boşlukları; \etiketD{}.} Yüksek.

\hypertarget{nace_kartlar_10-12-328-1103-elma-ux15farabux131-ve-diux11fer-meyve-ux15faraplarux131nux131n-imalatux131--gb-statuxfcsuxfc-tarim-198--birincilikincil-arketip---duruux15f-koru}{%
\subsection{3.28 11.03 Elma şarabı ve diğer meyve şaraplarının imalatı --- GB statüsü: TARIM-1/98 \textbar{} Birincil/ikincil arketip: --- \textbar{} Duruş: Koru}\label{nace_kartlar_10-12-328-1103-elma-ux15farabux131-ve-diux11fer-meyve-ux15faraplarux131nux131n-imalatux131--gb-statuxfcsuxfc-tarim-198--birincilikincil-arketip---duruux15f-koru}}

\emph{İhmal edilebilir ticaret (elma şarabı); YSHİ verisi yok.}

\textbf{1. Kimlik.} Rev. 2.1: 11.03 (1:1); CPA 11.03. Başlıca HS4 (TR ihracatı 2024): 2206 \%100. BEC payları (ticaret ağırlıklı): ara \%0 / sermaye \%0 / tüketim \%100.

\textbf{2. Ölçek}

\begingroup\gbtablo\small

\begin{longtable}[]{@{}
  >{\raggedright\arraybackslash}p{(\columnwidth - 4\tabcolsep) * \real{0.3634}}
  >{\raggedright\arraybackslash}p{(\columnwidth - 4\tabcolsep) * \real{0.2154}}
  >{\raggedright\arraybackslash}p{(\columnwidth - 4\tabcolsep) * \real{0.4211}}@{}}
\toprule\noalign{}
\begin{minipage}[b]{\linewidth}\raggedright
\textbf{Gösterge}
\end{minipage} & \begin{minipage}[b]{\linewidth}\raggedright
\textbf{Değer}
\end{minipage} & \begin{minipage}[b]{\linewidth}\raggedright
\textbf{Etiket ve kaynak}
\end{minipage} \\
\midrule\noalign{}
\endhead
\bottomrule\noalign{}
\endlastfoot
Üretim değeri; katma değer & \$--- mn; \$--- mn (GSYH'nin ---) & TÜİK YSHİ 2024 \etiketV{} \\
Çalışan; girişim (0 / 1--9 / 10--49 / 50--249 / 250+) & ---; --- (--- / --- / --- / --- / ---) & TÜİK YSHİ ve İş Kayıtları 2024 \etiketV{}; gizli büyüklük hücreleri \etiketT{} \\
250+ girişimlerin istihdam payı; AB'ye göre kişi başı KD & ---; --- & \etiketT{} (AB27 SBS 2023, nominal kur) \\
İlk 3 İBBS-2 bölgesi & --- & \etiketT{} iş kayıtları + SGK 2025, RAS \\
Yabancı sermaye; çapa yatırımcılar & --- & \etiketH{} (firma bilgisi; pay verisi yok: VB-30) \\
\end{longtable}

\endgroup

\textbf{3. Dış ticaret}

\begingroup\gbtablo\small

\begin{longtable}[]{@{}
  >{\raggedright\arraybackslash}p{(\columnwidth - 2\tabcolsep) * \real{0.4821}}
  >{\raggedright\arraybackslash}p{(\columnwidth - 2\tabcolsep) * \real{0.5179}}@{}}
\toprule\noalign{}
\begin{minipage}[b]{\linewidth}\raggedright
\textbf{Gösterge}
\end{minipage} & \begin{minipage}[b]{\linewidth}\raggedright
\textbf{Değer (BACI 2024, \$; \etiketV{}, eşleme \etiketT{})}
\end{minipage} \\
\midrule\noalign{}
\endhead
\bottomrule\noalign{}
\endlastfoot
İhracat (2024); 2020--24 değişimi & \$0 mn; \%−61 \\
AB payı; ilk 5 AB pazarı & \%10; DEU \%66; BEL \%25; CZE \%10; IRL \%0 \\
İlk 5 AB dışı pazar & CYP \%60; AUS \%34; CHE \%3; GBR \%1; ZAF \%1 \\
İthalat; AB ve Çin payı; ilk kaynaklar & \$1 mn; AB \%47, Çin \%0; JPN \%30; SWE \%23; BEL \%11; CAN \%6; USA \%6 \\
Dış ticaret dengesi & \$−1 mn \\
İthal girdi içeriği (I-O, toplam) / KD÷üretim & \%20 / --- \\
RCA dünya (2020 → 2024) / AB pazarında & 0,02 → 0,01 / 0,01 \\
Birim değer ÷ AB pazarı ortalaması & 1,07 \\
AB pazar payı (AB dışı tedarik içinde, 2020 → 2024); TR'nin sırası & \%0,0 → \%0,0; 33. \\
\end{longtable}

\endgroup

\textbf{4. AB tarafı maruziyet}

\begin{itemize}
\tightlist
\item
  \kod{τ\_\allowbreak{}AB} (S4 artışı, rota A, TR→AB 2025 ağırlıklı): \%0,0; 1/98 kaybıyla \%1,1; tam MFN ortalaması \%1,1. Rejim: GB \%0, İTÜ \%0, tarım \%100 \etiketV{}/\etiketT{}.
\item
  Tepe satırlar (S4 artışı ≥ \%10; en büyük 3): yok.
\item
  AB pazarında AB dışı rakipler (tercih durumu): GBR \%38; JPN \%23; USA \%11; CHN \%10; ARM \%3. Tercihli rakip payı \%73.
\item
  Menşe (S2 / S3): risk D/D; \kod{c\_\allowbreak{}RoO} \%1,5 / \%2,0; \kod{p\_\allowbreak{}uyum} 0,90 / 0,85 \etiketH{}. Kural: 2206: başlık değişikliği; meyve tamamen elde edilmiş.
\item
  EUDR kapsamındaki AB ihracatı payı: \%0; SKDM: \%0.
\end{itemize}

\textbf{5. Türkiye tarafı maruziyet}

\begin{itemize}
\tightlist
\item
  Uygulanan tarife (2026, ithalat ağırlıklı): MFN \%70,0, İGV dahil \%70,0 \etiketV{}. S4'te AB menşeline (c2): \%0,0.
\end{itemize}

\textbf{6. Kabiliyet}

\begin{itemize}
\tightlist
\item
  KS 42,5 (sıra 186/230; ağırlık duyarlılığında P10--P90 104--221). Bileşenler: RCA 2, PCI 54 (PCI 0,07), yurt içi KD 72, Ar-Ge 25, tedarikçi 84, verimlilik ---, firma --- \etiketT{}.
\item
  Ürün uzayında atlanabilir ürünler (sınıf içi; yoğunluk üst çeyrek, PCI \textgreater{} sınıf): yok \etiketT{}.
\end{itemize}

\textbf{7. Stratejik önem.} SS 0,0 (en yüksek boyut: --- 0; ağırlıklı ortalama 0,0) \etiketH{}.

\textbf{8. Makro duyarlılık}

\begingroup\gbtablo\small

\begin{longtable}[]{@{}ll@{}}
\toprule\noalign{}
\textbf{Gösterge} & \textbf{Değer (\etiketT{})} \\
\midrule\noalign{}
\endhead
\bottomrule\noalign{}
\endlastfoot
TÜFE ağırlığı (doğrudan, COICOP 2018 eşlemesi); tüketimde ithal pay & \%0,00; \%19 \\
ÜFE vekili (imalat üretimindeki pay) & --- \\
EDE: 1 puan tarife, tüm menşe / yalnız AB (bp; ρ = 1) & 0,000 / 0,000 \\
Geçişkenlik ρ (senaryolarda) & 0,6--1,0 (P50 0,8) \\
Piyasa yoğunlaşması (250+ istihdam payı) & --- \\
CKE: tarife (\$ mn / bp) ve mali (\$ mn / mlr TL); 0--100 puan & 0,0 ve ---; 3 \\
\end{longtable}

\endgroup

\textbf{9. Senaryo etkileri, T+5, P50 (P10 / P90)}

\begingroup\gbtablo\small

\begin{longtable}[]{@{}lrrr@{}}
\toprule\noalign{}
\textbf{Değişken} & \textbf{S2} & \textbf{S3} & \textbf{S4} \\
\midrule\noalign{}
\endhead
\bottomrule\noalign{}
\endlastfoot
AB'ye ihracat (\$ mn) & 0 (0 / 0) & 0 (0 / 0) & 0 (0 / 0) \\
Toplam ihracat (\$ mn) & 0 (0 / 0) & 0 (0 / 0) & 0 (0 / 0) \\
Nihai mal ithalatı (\$ mn) & 0 (0 / 0) & 0 (0 / 0) & 0 (0 / 0) \\
Cari denge (\$ mn) & 0 (0 / 0) & 0 (0 / 0) & 0 (0 / 0) \\
Üretim (\$ mn) & 0 (0 / 0) & 0 (0 / 0) & 0 (0 / 0) \\
İstihdam (kişi) & 0 (0 / 0) & 0 (0 / 0) & 0 (0 / 0) \\
TÜFE (bp, birinci tur) & 0,00 (0,00 / 0,00) & 0,00 (0,00 / 0,00) & 0,00 (0,00 / 0,00) \\
Yatırım ve DYY (bütün senaryolar) & --- & & \\
\end{longtable}

\endgroup

\emph{Kaynak:} \etiketT{} f3e: kısmi denge, toplamlar Faz 2B GE dağılımına kalibre; S4'te 1/98 kaybı olasılığı 0,25.

\textbf{15. Güven düzeyi; veri boşlukları; \etiketD{}.} Düşük (YSHİ verisi yok).

\hypertarget{nace_kartlar_10-12-329-1104-diux11fer-damux131tux131lmamux131ux15f-fermente-iuxe7eceklerin-imalatux131--gb-statuxfcsuxfc-iux307tuxfc--birincilikincil-arketip---duruux15f-koru}{%
\subsection{3.29 11.04 Diğer damıtılmamış fermente içeceklerin imalatı --- GB statüsü: İTÜ \textbar{} Birincil/ikincil arketip: --- \textbar{} Duruş: Koru}\label{nace_kartlar_10-12-329-1104-diux11fer-damux131tux131lmamux131ux15f-fermente-iuxe7eceklerin-imalatux131--gb-statuxfcsuxfc-iux307tuxfc--birincilikincil-arketip---duruux15f-koru}}

\emph{İhmal edilebilir ticaret (vermut ve benzeri); YSHİ verisi yok.}

\textbf{1. Kimlik.} Rev. 2.1: 11.04 (1:1); CPA 11.04. Başlıca HS4 (TR ihracatı 2024): 2205 \%100. BEC payları (ticaret ağırlıklı): ara \%0 / sermaye \%0 / tüketim \%100.

\textbf{2. Ölçek}

\begingroup\gbtablo\small

\begin{longtable}[]{@{}
  >{\raggedright\arraybackslash}p{(\columnwidth - 4\tabcolsep) * \real{0.3634}}
  >{\raggedright\arraybackslash}p{(\columnwidth - 4\tabcolsep) * \real{0.2154}}
  >{\raggedright\arraybackslash}p{(\columnwidth - 4\tabcolsep) * \real{0.4211}}@{}}
\toprule\noalign{}
\begin{minipage}[b]{\linewidth}\raggedright
\textbf{Gösterge}
\end{minipage} & \begin{minipage}[b]{\linewidth}\raggedright
\textbf{Değer}
\end{minipage} & \begin{minipage}[b]{\linewidth}\raggedright
\textbf{Etiket ve kaynak}
\end{minipage} \\
\midrule\noalign{}
\endhead
\bottomrule\noalign{}
\endlastfoot
Üretim değeri; katma değer & \$--- mn; \$--- mn (GSYH'nin ---) & TÜİK YSHİ 2024 \etiketV{} \\
Çalışan; girişim (0 / 1--9 / 10--49 / 50--249 / 250+) & ---; --- (--- / --- / --- / --- / ---) & TÜİK YSHİ ve İş Kayıtları 2024 \etiketV{}; gizli büyüklük hücreleri \etiketT{} \\
250+ girişimlerin istihdam payı; AB'ye göre kişi başı KD & ---; --- & \etiketT{} (AB27 SBS 2023, nominal kur) \\
İlk 3 İBBS-2 bölgesi & --- & \etiketT{} iş kayıtları + SGK 2025, RAS \\
Yabancı sermaye; çapa yatırımcılar & --- & \etiketH{} (firma bilgisi; pay verisi yok: VB-30) \\
\end{longtable}

\endgroup

\textbf{3. Dış ticaret}

\begingroup\gbtablo\small

\begin{longtable}[]{@{}
  >{\raggedright\arraybackslash}p{(\columnwidth - 2\tabcolsep) * \real{0.4839}}
  >{\raggedright\arraybackslash}p{(\columnwidth - 2\tabcolsep) * \real{0.5161}}@{}}
\toprule\noalign{}
\begin{minipage}[b]{\linewidth}\raggedright
\textbf{Gösterge}
\end{minipage} & \begin{minipage}[b]{\linewidth}\raggedright
\textbf{Değer (BACI 2024, \$; \etiketV{}, eşleme \etiketT{})}
\end{minipage} \\
\midrule\noalign{}
\endhead
\bottomrule\noalign{}
\endlastfoot
İhracat (2024); 2020--24 değişimi & \$0 mn; \%696 \\
AB payı; ilk 5 AB pazarı & \%15; DEU \%100 \\
İlk 5 AB dışı pazar & NGA \%93; GBR \%5; OMN \%1; GEO \%1 \\
İthalat; AB ve Çin payı; ilk kaynaklar & \$2 mn; AB \%100, Çin \%0; ITA \%71; DEU \%12; ESP \%6; ROU \%5; NLD \%2 \\
Dış ticaret dengesi & \$−2 mn \\
İthal girdi içeriği (I-O, toplam) / KD÷üretim & \%20 / --- \\
RCA dünya (2020 → 2024) / AB pazarında & 0,00 → 0,01 / 0,02 \\
Birim değer ÷ AB pazarı ortalaması & 1,74 \\
AB pazar payı (AB dışı tedarik içinde, 2020 → 2024); TR'nin sırası & \%0,0 → \%0,1; 28. \\
\end{longtable}

\endgroup

\textbf{4. AB tarafı maruziyet}

\begin{itemize}
\tightlist
\item
  \kod{τ\_\allowbreak{}AB} (S4 artışı, rota A, TR→AB 2025 ağırlıklı): \%0,0; 1/98 kaybıyla \%0,0; tam MFN ortalaması \%0,7. Rejim: GB \%0, İTÜ \%100, tarım \%0 \etiketV{}/\etiketT{}.
\item
  Tepe satırlar (S4 artışı ≥ \%10; en büyük 3): yok.
\item
  AB pazarında AB dışı rakipler (tercih durumu): ARG \%48; USA \%20; CHL \%10; GEO \%6; GBR \%5. Tercihli rakip payı \%30.
\item
  Menşe (S2 / S3): risk D/D; \kod{c\_\allowbreak{}RoO} \%1,5 / \%2,0; \kod{p\_\allowbreak{}uyum} 0,90 / 0,85 \etiketH{}. Kural: 2205: başlık değişikliği; üzüm türevleri tamamen elde edilmiş.
\item
  EUDR kapsamındaki AB ihracatı payı: \%0; SKDM: \%0.
\end{itemize}

\textbf{5. Türkiye tarafı maruziyet}

\begin{itemize}
\tightlist
\item
  Uygulanan tarife (2026, ithalat ağırlıklı): MFN \%4,2, İGV dahil \%4,2 \etiketV{}. S4'te AB menşeline (c2): \%4,2.
\end{itemize}

\textbf{6. Kabiliyet}

\begin{itemize}
\tightlist
\item
  KS 43,1 (sıra 180/230; ağırlık duyarlılığında P10--P90 114--207). Bileşenler: RCA 31, PCI 27 (PCI −0,42), yurt içi KD 72, Ar-Ge 25, tedarikçi 84, verimlilik ---, firma --- \etiketT{}.
\item
  Ürün uzayında atlanabilir ürünler (sınıf içi; yoğunluk üst çeyrek, PCI \textgreater{} sınıf): yok \etiketT{}.
\end{itemize}

\textbf{7. Stratejik önem.} SS 0,0 (en yüksek boyut: --- 0; ağırlıklı ortalama 0,0) \etiketH{}.

\textbf{8. Makro duyarlılık}

\begingroup\gbtablo\small

\begin{longtable}[]{@{}ll@{}}
\toprule\noalign{}
\textbf{Gösterge} & \textbf{Değer (\etiketT{})} \\
\midrule\noalign{}
\endhead
\bottomrule\noalign{}
\endlastfoot
TÜFE ağırlığı (doğrudan, COICOP 2018 eşlemesi); tüketimde ithal pay & \%0,00; \%20 \\
ÜFE vekili (imalat üretimindeki pay) & --- \\
EDE: 1 puan tarife, tüm menşe / yalnız AB (bp; ρ = 1) & 0,000 / 0,000 \\
Geçişkenlik ρ (senaryolarda) & 0,6--1,0 (P50 0,8) \\
Piyasa yoğunlaşması (250+ istihdam payı) & --- \\
CKE: tarife (\$ mn / bp) ve mali (\$ mn / mlr TL); 0--100 puan & --- ve ---; --- \\
\end{longtable}

\endgroup

\textbf{9. Senaryo etkileri, T+5, P50 (P10 / P90)}

\begingroup\gbtablo\small

\begin{longtable}[]{@{}lrrr@{}}
\toprule\noalign{}
\textbf{Değişken} & \textbf{S2} & \textbf{S3} & \textbf{S4} \\
\midrule\noalign{}
\endhead
\bottomrule\noalign{}
\endlastfoot
AB'ye ihracat (\$ mn) & 0 (0 / 0) & 0 (0 / 0) & 0 (0 / 0) \\
Toplam ihracat (\$ mn) & 0 (0 / 0) & 0 (0 / 0) & 0 (0 / 0) \\
Nihai mal ithalatı (\$ mn) & 0 (0 / 0) & 0 (0 / 0) & 0 (0 / 0) \\
Cari denge (\$ mn) & 0 (0 / 0) & 0 (0 / 0) & 0 (0 / 0) \\
Üretim (\$ mn) & 0 (0 / 0) & 0 (0 / 0) & 0 (0 / 0) \\
İstihdam (kişi) & 0 (0 / 0) & 0 (0 / 0) & 0 (0 / 0) \\
TÜFE (bp, birinci tur) & 0,00 (0,00 / 0,00) & 0,00 (0,00 / 0,00) & 0,00 (0,00 / 0,00) \\
Yatırım ve DYY (bütün senaryolar) & --- & & \\
\end{longtable}

\endgroup

\emph{Kaynak:} \etiketT{} f3e: kısmi denge, toplamlar Faz 2B GE dağılımına kalibre; S4'te 1/98 kaybı olasılığı 0,25.

\textbf{15. Güven düzeyi; veri boşlukları; \etiketD{}.} Düşük (YSHİ verisi yok).

\hypertarget{nace_kartlar_10-12-330-1105-bira-imalatux131--gb-statuxfcsuxfc-iux307tuxfc--birincilikincil-arketip---duruux15f-koru}{%
\subsection{3.30 11.05 Bira imalatı --- GB statüsü: İTÜ \textbar{} Birincil/ikincil arketip: --- \textbar{} Duruş: Koru}\label{nace_kartlar_10-12-330-1105-bira-imalatux131--gb-statuxfcsuxfc-iux307tuxfc--birincilikincil-arketip---duruux15f-koru}}

\emph{Bira (Anadolu Efes, Türk Tuborg). AB'nin bira MFN'si sıfır; GB etkisi yok. İthalatın (\$348 mn) \%92'si bira değil, ABD'den damıtma artığı yem (KN 2303 30, DDGS).}

\textbf{1. Kimlik.} Rev. 2.1: 11.05 (1:1); CPA 11.05. Başlıca HS4 (TR ihracatı 2024): 2203 \%90, 2202 \%7, 2303 \%3. BEC payları (ticaret ağırlıklı): ara \%78 / sermaye \%0 / tüketim \%22. CPA 11.05.20 (bira ve damıtma artığı) DDGS'yi bu sınıfa koyuyor; ekonomik olarak bir yem girdisi (10.91).

\textbf{2. Ölçek}

\begingroup\gbtablo\small

\begin{longtable}[]{@{}
  >{\raggedright\arraybackslash}p{(\columnwidth - 4\tabcolsep) * \real{0.3162}}
  >{\raggedright\arraybackslash}p{(\columnwidth - 4\tabcolsep) * \real{0.3241}}
  >{\raggedright\arraybackslash}p{(\columnwidth - 4\tabcolsep) * \real{0.3597}}@{}}
\toprule\noalign{}
\begin{minipage}[b]{\linewidth}\raggedright
\textbf{Gösterge}
\end{minipage} & \begin{minipage}[b]{\linewidth}\raggedright
\textbf{Değer}
\end{minipage} & \begin{minipage}[b]{\linewidth}\raggedright
\textbf{Etiket ve kaynak}
\end{minipage} \\
\midrule\noalign{}
\endhead
\bottomrule\noalign{}
\endlastfoot
Üretim değeri; katma değer & \$199 mn; \$65 mn (GSYH'nin \%0,00) & TÜİK YSHİ 2024 \etiketV{} \\
Çalışan; girişim (0 / 1--9 / 10--49 / 50--249 / 250+) & 394; 15 (3 / 12 / 0 / 0 / 1) & TÜİK YSHİ ve İş Kayıtları 2024 \etiketV{}; gizli büyüklük hücreleri \etiketT{} \\
250+ girişimlerin istihdam payı; AB'ye göre kişi başı KD & \%90; 1,50 & \etiketT{} (AB27 SBS 2023, nominal kur) \\
İlk 3 İBBS-2 bölgesi & İzmir \%59; İstanbul \%15; Adana, Mersin \%9 & \etiketT{} iş kayıtları + SGK 2025, RAS \\
Yabancı sermaye; çapa yatırımcılar & Anadolu Efes (AB InBev ortaklığı \etiketD{}), Türk Tuborg \etiketH{} & \etiketH{} (firma bilgisi; pay verisi yok: VB-30) \\
\end{longtable}

\endgroup

\textbf{3. Dış ticaret}

\begingroup\gbtablo\small

\begin{longtable}[]{@{}
  >{\raggedright\arraybackslash}p{(\columnwidth - 2\tabcolsep) * \real{0.4787}}
  >{\raggedright\arraybackslash}p{(\columnwidth - 2\tabcolsep) * \real{0.5213}}@{}}
\toprule\noalign{}
\begin{minipage}[b]{\linewidth}\raggedright
\textbf{Gösterge}
\end{minipage} & \begin{minipage}[b]{\linewidth}\raggedright
\textbf{Değer (BACI 2024, \$; \etiketV{}, eşleme \etiketT{})}
\end{minipage} \\
\midrule\noalign{}
\endhead
\bottomrule\noalign{}
\endlastfoot
İhracat (2024); 2020--24 değişimi & \$61 mn; \%29 \\
AB payı; ilk 5 AB pazarı & \%8; BGR \%32; ROU \%22; FRA \%11; SVK \%8; NLD \%6 \\
İlk 5 AB dışı pazar & IRQ \%39; CYP \%29; GBR \%8; CHN \%7; AZE \%3 \\
İthalat; AB ve Çin payı; ilk kaynaklar & \$348 mn; AB \%10, Çin \%0; USA \%64; BRA \%12; RUS \%8; BGR \%3; MEX \%3 \\
Dış ticaret dengesi & \$−287 mn \\
İthal girdi içeriği (I-O, toplam) / KD÷üretim & \%20 / \%33 \\
RCA dünya (2020 → 2024) / AB pazarında & 0,23 → 0,23 / 0,15 \\
Birim değer ÷ AB pazarı ortalaması & 0,79 \\
AB pazar payı (AB dışı tedarik içinde, 2020 → 2024); TR'nin sırası & \%0,7 → \%0,6; 11. \\
\end{longtable}

\endgroup

\textbf{4. AB tarafı maruziyet}

\begin{itemize}
\tightlist
\item
  \kod{τ\_\allowbreak{}AB} (S4 artışı, rota A, TR→AB 2025 ağırlıklı): \%0,0; 1/98 kaybıyla \%0,0; tam MFN ortalaması \%0,0. Rejim: GB \%0, İTÜ \%100, tarım \%0 \etiketV{}/\etiketT{}.
\item
  Tepe satırlar (S4 artışı ≥ \%10; en büyük 3): yok.
\item
  AB pazarında AB dışı rakipler (tercih durumu): GBR \%43; USA \%23; MEX \%16; BRA \%3; CAN \%3. Tercihli rakip payı \%67.
\item
  Menşe (S2 / S3): risk D/D; \kod{c\_\allowbreak{}RoO} \%1,5 / \%2,0; \kod{p\_\allowbreak{}uyum} 0,90 / 0,85 \etiketH{}. Kural: 2203: başlık değişikliği (ithal malttan bira menşe kazanır).
\item
  EUDR kapsamındaki AB ihracatı payı: \%0; SKDM: \%0.
\end{itemize}

\textbf{5. Türkiye tarafı maruziyet}

\begin{itemize}
\tightlist
\item
  Uygulanan tarife (2026, ithalat ağırlıklı): MFN \%0,1, İGV dahil \%1,6 \etiketV{}. S4'te AB menşeline (c2): \%8,5.
\item
  İthalatın asıl kalemi DDGS (\$319 mn, ABD). Tarım rejiminde (Fasıl 23) ve GB'den bağımsız; yem maliyeti için önemli.
\end{itemize}

\textbf{6. Kabiliyet}

\begin{itemize}
\tightlist
\item
  KS 42,4 (sıra 189/230; ağırlık duyarlılığında P10--P90 105--220). Bileşenler: RCA 19, PCI 20 (PCI −0,64), yurt içi KD 80, Ar-Ge 25, tedarikçi 84, verimlilik 98, firma 4 \etiketT{}.
\item
  Ürün uzayında atlanabilir ürünler (sınıf içi; yoğunluk üst çeyrek, PCI \textgreater{} sınıf): yok \etiketT{}.
\end{itemize}

\textbf{7. Stratejik önem.} SS 0,0 (en yüksek boyut: --- 0; ağırlıklı ortalama 0,0) \etiketH{}.

\textbf{8. Makro duyarlılık}

\begingroup\gbtablo\small

\begin{longtable}[]{@{}lr@{}}
\toprule\noalign{}
\textbf{Gösterge} & \textbf{Değer (\etiketT{})} \\
\midrule\noalign{}
\endhead
\bottomrule\noalign{}
\endlastfoot
TÜFE ağırlığı (doğrudan, COICOP 2018 eşlemesi); tüketimde ithal pay & \%0,27; \%18 \\
ÜFE vekili (imalat üretimindeki pay) & \%0,03 \\
EDE: 1 puan tarife, tüm menşe / yalnız AB (bp; ρ = 1) & 0,075 / 0,008 \\
Geçişkenlik ρ (senaryolarda) & 0,6--1,0 (P50 0,8) \\
Piyasa yoğunlaşması (250+ istihdam payı) & \%90 \\
CKE: tarife (\$ mn / bp) ve mali (\$ mn / mlr TL); 0--100 puan & 121,4 ve 21,81; 33 \\
\end{longtable}

\endgroup

\textbf{9. Senaryo etkileri, T+5, P50 (P10 / P90)}

\begingroup\gbtablo\small

\begin{longtable}[]{@{}
  >{\raggedright\arraybackslash}p{(\columnwidth - 6\tabcolsep) * \real{0.3667}}
  >{\raggedleft\arraybackslash}p{(\columnwidth - 6\tabcolsep) * \real{0.2111}}
  >{\raggedleft\arraybackslash}p{(\columnwidth - 6\tabcolsep) * \real{0.2111}}
  >{\raggedleft\arraybackslash}p{(\columnwidth - 6\tabcolsep) * \real{0.2111}}@{}}
\toprule\noalign{}
\begin{minipage}[b]{\linewidth}\raggedright
\textbf{Değişken}
\end{minipage} & \begin{minipage}[b]{\linewidth}\raggedleft
\textbf{S2}
\end{minipage} & \begin{minipage}[b]{\linewidth}\raggedleft
\textbf{S3}
\end{minipage} & \begin{minipage}[b]{\linewidth}\raggedleft
\textbf{S4}
\end{minipage} \\
\midrule\noalign{}
\endhead
\bottomrule\noalign{}
\endlastfoot
AB'ye ihracat (\$ mn) & 0 (0 / 0) & 0 (0 / 0) & 0 (0 / 0) \\
Toplam ihracat (\$ mn) & 0 (0 / 0) & 0 (0 / 0) & 0 (0 / 0) \\
Nihai mal ithalatı (\$ mn) & 0 (−1 / +1) & 0 (−1 / 0) & −1 (−1 / 0) \\
Cari denge (\$ mn) & 0 (−1 / 0) & 0 (0 / 0) & 0 (0 / +1) \\
Üretim (\$ mn) & 0 (−1 / 0) & 0 (0 / 0) & +1 (0 / +1) \\
İstihdam (kişi) & 0 (−1 / 0) & 0 (−1 / 0) & +1 (0 / +2) \\
TÜFE (bp, birinci tur) & +0,01 (+0,01 / +0,01) & +0,02 (+0,01 / +0,02) & +0,06 (+0,05 / +0,06) \\
Yatırım ve DYY (bütün senaryolar) & --- & & \\
\end{longtable}

\endgroup

\emph{Kaynak:} \etiketT{} f3e: kısmi denge, toplamlar Faz 2B GE dağılımına kalibre; S4'te 1/98 kaybı olasılığı 0,25.

\textbf{13. Kazanan ve kaybedenler; bölgesel ve sosyal etki; telafi.} Etkilenen kesim yok.

\textbf{15. Güven düzeyi; veri boşlukları; \etiketD{}.} Orta (YSHİ gizli; sınıflama notu).

\hypertarget{nace_kartlar_10-12-331-1106-malt-imalatux131--gb-statuxfcsuxfc-tarim-198--birincilikincil-arketip---duruux15f-koru}{%
\subsection{3.31 11.06 Malt imalatı --- GB statüsü: TARIM-1/98 \textbar{} Birincil/ikincil arketip: --- \textbar{} Duruş: Koru}\label{nace_kartlar_10-12-331-1106-malt-imalatux131--gb-statuxfcsuxfc-tarim-198--birincilikincil-arketip---duruux15f-koru}}

\emph{Malt: ithalat \$72 mn (Almanya, Danimarka, Fransa). YSHİ verisi yok. Tarım rejiminde; GB etkisi yok.}

\textbf{1. Kimlik.} Rev. 2.1: 11.06 (1:1); CPA 11.06. Başlıca HS4 (TR ihracatı 2024): 1107 \%100. BEC payları (ticaret ağırlıklı): ara \%100 / sermaye \%0 / tüketim \%0.

\textbf{2. Ölçek}

\begingroup\gbtablo\small

\begin{longtable}[]{@{}
  >{\raggedright\arraybackslash}p{(\columnwidth - 4\tabcolsep) * \real{0.3629}}
  >{\raggedright\arraybackslash}p{(\columnwidth - 4\tabcolsep) * \real{0.2167}}
  >{\raggedright\arraybackslash}p{(\columnwidth - 4\tabcolsep) * \real{0.4204}}@{}}
\toprule\noalign{}
\begin{minipage}[b]{\linewidth}\raggedright
\textbf{Gösterge}
\end{minipage} & \begin{minipage}[b]{\linewidth}\raggedright
\textbf{Değer}
\end{minipage} & \begin{minipage}[b]{\linewidth}\raggedright
\textbf{Etiket ve kaynak}
\end{minipage} \\
\midrule\noalign{}
\endhead
\bottomrule\noalign{}
\endlastfoot
Üretim değeri; katma değer & \$--- mn; \$--- mn (GSYH'nin ---) & TÜİK YSHİ 2024 \etiketV{} \\
Çalışan; girişim (0 / 1--9 / 10--49 / 50--249 / 250+) & ---; --- (--- / --- / --- / --- / ---) & TÜİK YSHİ ve İş Kayıtları 2024 \etiketV{}; gizli büyüklük hücreleri \etiketT{} \\
250+ girişimlerin istihdam payı; AB'ye göre kişi başı KD & ---; --- & \etiketT{} (AB27 SBS 2023, nominal kur) \\
İlk 3 İBBS-2 bölgesi & --- & \etiketT{} iş kayıtları + SGK 2025, RAS \\
Yabancı sermaye; çapa yatırımcılar & Efes Malt (Afyon, Konya) \etiketH{} & \etiketH{} (firma bilgisi; pay verisi yok: VB-30) \\
\end{longtable}

\endgroup

\textbf{3. Dış ticaret}

\begingroup\gbtablo\small

\begin{longtable}[]{@{}
  >{\raggedright\arraybackslash}p{(\columnwidth - 2\tabcolsep) * \real{0.4766}}
  >{\raggedright\arraybackslash}p{(\columnwidth - 2\tabcolsep) * \real{0.5234}}@{}}
\toprule\noalign{}
\begin{minipage}[b]{\linewidth}\raggedright
\textbf{Gösterge}
\end{minipage} & \begin{minipage}[b]{\linewidth}\raggedright
\textbf{Değer (BACI 2024, \$; \etiketV{}, eşleme \etiketT{})}
\end{minipage} \\
\midrule\noalign{}
\endhead
\bottomrule\noalign{}
\endlastfoot
İhracat (2024); 2020--24 değişimi & \$0 mn; \%−66 \\
AB payı; ilk 5 AB pazarı & \%2; NLD \%90; DEU \%6; BEL \%4; ITA \%0; POL \%0 \\
İlk 5 AB dışı pazar & CAF \%86; CYP \%5; UKR \%3; IRQ \%2; CAN \%1 \\
İthalat; AB ve Çin payı; ilk kaynaklar & \$83 mn; AB \%86, Çin \%0; DEU \%29; DNK \%23; FRA \%23; RUS \%10; ROU \%7 \\
Dış ticaret dengesi & \$−83 mn \\
İthal girdi içeriği (I-O, toplam) / KD÷üretim & \%20 / --- \\
RCA dünya (2020 → 2024) / AB pazarında & 0,03 → 0,01 / 0,01 \\
Birim değer ÷ AB pazarı ortalaması & 2,98 \\
AB pazar payı (AB dışı tedarik içinde, 2020 → 2024); TR'nin sırası & \%0,0 → \%0,0; 20. \\
\end{longtable}

\endgroup

\textbf{4. AB tarafı maruziyet}

\begin{itemize}
\tightlist
\item
  \kod{τ\_\allowbreak{}AB} (S4 artışı, rota A, TR→AB 2025 ağırlıklı): \%0,0; 1/98 kaybıyla \%8,2; tam MFN ortalaması \%8,2. Rejim: GB \%0, İTÜ \%0, tarım \%100 \etiketV{}/\etiketT{}.
\item
  Tepe satırlar (S4 artışı ≥ \%10; en büyük 3): yok.
\item
  AB pazarında AB dışı rakipler (tercih durumu): GBR \%71; UKR \%23; COD \%2; TGO \%1; SRB \%1. Tercihli rakip payı \%95.
\item
  Menşe (S2 / S3): risk D/D; \kod{c\_\allowbreak{}RoO} \%1,5 / \%2,0; \kod{p\_\allowbreak{}uyum} 0,90 / 0,85 \etiketH{}. Kural: 1107: tahıl tamamen elde edilmiş.
\item
  EUDR kapsamındaki AB ihracatı payı: \%0; SKDM: \%0.
\end{itemize}

\textbf{5. Türkiye tarafı maruziyet}

\begin{itemize}
\tightlist
\item
  Uygulanan tarife (2026, ithalat ağırlıklı): MFN \%18,7, İGV dahil \%18,7 \etiketV{}. S4'te AB menşeline (c2): \%0,0.
\end{itemize}

\textbf{6. Kabiliyet}

\begin{itemize}
\tightlist
\item
  KS 43,7 (sıra 176/230; ağırlık duyarlılığında P10--P90 92--218). Bileşenler: RCA 1, PCI 59 (PCI 0,16), yurt içi KD 72, Ar-Ge 25, tedarikçi 84, verimlilik ---, firma --- \etiketT{}.
\item
  Ürün uzayında atlanabilir ürünler (sınıf içi; yoğunluk üst çeyrek, PCI \textgreater{} sınıf): yok \etiketT{}.
\end{itemize}

\textbf{7. Stratejik önem.} SS 14,9 (en yüksek boyut: gıda güvenliği 20; ağırlıklı ortalama 3,0) \etiketH{}. Malt; bira ve gıda girdisi

\textbf{8. Makro duyarlılık}

\begingroup\gbtablo\small

\begin{longtable}[]{@{}ll@{}}
\toprule\noalign{}
\textbf{Gösterge} & \textbf{Değer (\etiketT{})} \\
\midrule\noalign{}
\endhead
\bottomrule\noalign{}
\endlastfoot
TÜFE ağırlığı (doğrudan, COICOP 2018 eşlemesi); tüketimde ithal pay & \%0,00; \%21 \\
ÜFE vekili (imalat üretimindeki pay) & --- \\
EDE: 1 puan tarife, tüm menşe / yalnız AB (bp; ρ = 1) & 0,006 / 0,005 \\
Geçişkenlik ρ (senaryolarda) & 0,6--1,0 (P50 0,8) \\
Piyasa yoğunlaşması (250+ istihdam payı) & --- \\
CKE: tarife (\$ mn / bp) ve mali (\$ mn / mlr TL); 0--100 puan & 0,0 ve ---; 3 \\
\end{longtable}

\endgroup

\textbf{9. Senaryo etkileri, T+5, P50 (P10 / P90)}

\begingroup\gbtablo\small

\begin{longtable}[]{@{}lrrr@{}}
\toprule\noalign{}
\textbf{Değişken} & \textbf{S2} & \textbf{S3} & \textbf{S4} \\
\midrule\noalign{}
\endhead
\bottomrule\noalign{}
\endlastfoot
AB'ye ihracat (\$ mn) & 0 (0 / 0) & 0 (0 / 0) & 0 (0 / 0) \\
Toplam ihracat (\$ mn) & 0 (0 / 0) & 0 (0 / 0) & 0 (0 / 0) \\
Nihai mal ithalatı (\$ mn) & 0 (0 / +5) & 0 (0 / +3) & 0 (0 / 0) \\
Cari denge (\$ mn) & 0 (−4 / 0) & 0 (−2 / 0) & 0 (0 / 0) \\
Üretim (\$ mn) & 0 (−5 / 0) & 0 (−3 / 0) & 0 (0 / 0) \\
İstihdam (kişi) & 0 (−10 / 0) & 0 (−6 / 0) & 0 (0 / 0) \\
TÜFE (bp, birinci tur) & 0,00 (0,00 / 0,00) & 0,00 (0,00 / 0,00) & 0,00 (0,00 / 0,00) \\
Yatırım ve DYY (bütün senaryolar) & --- & & \\
\end{longtable}

\endgroup

\emph{Kaynak:} \etiketT{} f3e: kısmi denge, toplamlar Faz 2B GE dağılımına kalibre; S4'te 1/98 kaybı olasılığı 0,25.

\textbf{15. Güven düzeyi; veri boşlukları; \etiketD{}.} Düşük (YSHİ verisi yok).

\hypertarget{nace_kartlar_10-12-332-1107-alkolsuxfcz-iuxe7eceklerin-imalatux131-maden-sularux131nux131n-ve-diux11fer-ux15fiux15felenmiux15f-sularux131n-uxfcretimi--gb-statuxfcsuxfc-karma--birincilikincil-arketip---duruux15f-koru}{%
\subsection{3.32 11.07 Alkolsüz içeceklerin imalatı; maden sularının ve diğer şişelenmiş suların üretimi --- GB statüsü: KARMA \textbar{} Birincil/ikincil arketip: --- \textbar{} Duruş: Koru}\label{nace_kartlar_10-12-332-1107-alkolsuxfcz-iuxe7eceklerin-imalatux131-maden-sularux131nux131n-ve-diux11fer-ux15fiux15felenmiux15f-sularux131n-uxfcretimi--gb-statuxfcsuxfc-karma--birincilikincil-arketip---duruux15f-koru}}

\emph{Alkolsüz içecek ve su. AB'ye ihracat \$111 mn, S4'te \%6,6. AB'den ithalat \$258 mn, enerji içecekleri (Avusturya \%42): S4'te c2 \%8,6, TÜFE 0,8 bp. Etki küçük.}

\textbf{1. Kimlik.} Rev. 2.1: 11.07 (n:1); CPA 11.07. Başlıca HS4 (TR ihracatı 2024): 2202 \%79, 2201 \%21. BEC payları (ticaret ağırlıklı): ara \%0 / sermaye \%0 / tüketim \%100.

\textbf{2. Ölçek}

\begingroup\gbtablo\small

\begin{longtable}[]{@{}
  >{\raggedright\arraybackslash}p{(\columnwidth - 4\tabcolsep) * \real{0.2812}}
  >{\raggedright\arraybackslash}p{(\columnwidth - 4\tabcolsep) * \real{0.4044}}
  >{\raggedright\arraybackslash}p{(\columnwidth - 4\tabcolsep) * \real{0.3144}}@{}}
\toprule\noalign{}
\begin{minipage}[b]{\linewidth}\raggedright
\textbf{Gösterge}
\end{minipage} & \begin{minipage}[b]{\linewidth}\raggedright
\textbf{Değer}
\end{minipage} & \begin{minipage}[b]{\linewidth}\raggedright
\textbf{Etiket ve kaynak}
\end{minipage} \\
\midrule\noalign{}
\endhead
\bottomrule\noalign{}
\endlastfoot
Üretim değeri; katma değer & \$3.641 mn; \$764 mn (GSYH'nin \%0,06) & TÜİK YSHİ 2024 \etiketV{} \\
Çalışan; girişim (0 / 1--9 / 10--49 / 50--249 / 250+) & 14.765; 537 (72 / 343 / 105 / 44 / 13) & TÜİK YSHİ ve İş Kayıtları 2024 \etiketV{}; gizli büyüklük hücreleri \etiketT{} \\
250+ girişimlerin istihdam payı; AB'ye göre kişi başı KD & \%45; --- & \etiketT{} (AB27 SBS 2023, nominal kur) \\
İlk 3 İBBS-2 bölgesi & Bursa, Eskişehir, Bilecik \%17; Kocaeli, Sakarya, Düzce, Bolu, Yalova \%14; Adana, Mersin \%12 & \etiketT{} iş kayıtları + SGK 2025, RAS \\
Yabancı sermaye; çapa yatırımcılar & Coca-Cola İçecek (Anadolu Grubu), PepsiCo şişeleyicisi, Uludağ, Erikli ve Hayat Su \etiketH{} & \etiketH{} (firma bilgisi; pay verisi yok: VB-30) \\
\end{longtable}

\endgroup

\textbf{3. Dış ticaret}

\begingroup\gbtablo\small

\begin{longtable}[]{@{}
  >{\raggedright\arraybackslash}p{(\columnwidth - 2\tabcolsep) * \real{0.4766}}
  >{\raggedright\arraybackslash}p{(\columnwidth - 2\tabcolsep) * \real{0.5234}}@{}}
\toprule\noalign{}
\begin{minipage}[b]{\linewidth}\raggedright
\textbf{Gösterge}
\end{minipage} & \begin{minipage}[b]{\linewidth}\raggedright
\textbf{Değer (BACI 2024, \$; \etiketV{}, eşleme \etiketT{})}
\end{minipage} \\
\midrule\noalign{}
\endhead
\bottomrule\noalign{}
\endlastfoot
İhracat (2024); 2020--24 değişimi & \$385 mn; \%61 \\
AB payı; ilk 5 AB pazarı & \%29; DEU \%42; NLD \%14; AUT \%7; FRA \%7; BEL \%7 \\
İlk 5 AB dışı pazar & GBR \%17; SYR \%13; LBY \%8; ARE \%7; CYP \%7 \\
İthalat; AB ve Çin payı; ilk kaynaklar & \$330 mn; AB \%78, Çin \%0; AUT \%42; NLD \%19; CHE \%16; DEU \%4; ITA \%4 \\
Dış ticaret dengesi & \$+55 mn \\
İthal girdi içeriği (I-O, toplam) / KD÷üretim & \%20 / \%21 \\
RCA dünya (2020 → 2024) / AB pazarında & 0,89 → 0,93 / 1,41 \\
Birim değer ÷ AB pazarı ortalaması & 1,29 \\
AB pazar payı (AB dışı tedarik içinde, 2020 → 2024); TR'nin sırası & \%5,7 → \%5,6; 4. \\
\end{longtable}

\endgroup

\textbf{4. AB tarafı maruziyet}

\begin{itemize}
\tightlist
\item
  \kod{τ\_\allowbreak{}AB} (S4 artışı, rota A, TR→AB 2025 ağırlıklı): \%6,6; 1/98 kaybıyla \%6,6; tam MFN ortalaması \%7,0. Rejim: GB \%28, İTÜ \%72, tarım \%0 \etiketV{}/\etiketT{}.
\item
  Tepe satırlar (S4 artışı ≥ \%10; en büyük 3): yok.
\item
  AB pazarında AB dışı rakipler (tercih durumu): CHE \%33; GBR \%23; SRB \%7; USA \%3; KOR \%3. Tercihli rakip payı \%73.
\item
  Menşe (S2 / S3): risk D/O; \kod{c\_\allowbreak{}RoO} \%1,5 / \%2,5; \kod{p\_\allowbreak{}uyum} 0,90 / 0,85 \etiketH{}. Kural: 2201--2202: başlık değişikliği; şeker \%40 ağırlık sınırı; meyve suyu menşeli.
\item
  EUDR kapsamındaki AB ihracatı payı: \%0; SKDM: \%0.
\end{itemize}

\textbf{5. Türkiye tarafı maruziyet}

\begin{itemize}
\tightlist
\item
  Uygulanan tarife (2026, ithalat ağırlıklı): MFN \%8,6, İGV dahil \%8,6 \etiketV{}. S4'te AB menşeline (c2): \%8,6.
\item
  ÖTV ile gümrük etkileşimi: kolalı içeceklerde ÖTV (III) oransal. S4'te AB'den gelen enerji içeceğindeki gümrük vergisi ÖTV matrahına giriyor, ama etki küçük (0,8 bp).
\end{itemize}

\textbf{6. Kabiliyet}

\begin{itemize}
\tightlist
\item
  KS 41,3 (sıra 200/230; ağırlık duyarlılığında P10--P90 150--207). Bileşenler: RCA 42, PCI 20 (PCI −0,64), yurt içi KD 54, Ar-Ge 25, tedarikçi 84, verimlilik ---, firma 45 \etiketT{}.
\item
  Ürün uzayında atlanabilir ürünler (sınıf içi; yoğunluk üst çeyrek, PCI \textgreater{} sınıf): yok \etiketT{}.
\end{itemize}

\textbf{7. Stratejik önem.} SS 30,2 (en yüksek boyut: gıda güvenliği 40; ağırlıklı ortalama 7,5) \etiketH{}. İçme suyu arzı (afet ve kriz)

\textbf{8. Makro duyarlılık}

\begingroup\gbtablo\small

\begin{longtable}[]{@{}lr@{}}
\toprule\noalign{}
\textbf{Gösterge} & \textbf{Değer (\etiketT{})} \\
\midrule\noalign{}
\endhead
\bottomrule\noalign{}
\endlastfoot
TÜFE ağırlığı (doğrudan, COICOP 2018 eşlemesi); tüketimde ithal pay & \%1,11; \%10 \\
ÜFE vekili (imalat üretimindeki pay) & \%0,55 \\
EDE: 1 puan tarife, tüm menşe / yalnız AB (bp; ρ = 1) & 0,134 / 0,104 \\
Geçişkenlik ρ (senaryolarda) & 0,6--1,0 (P50 0,8) \\
Piyasa yoğunlaşması (250+ istihdam payı) & \%45 \\
CKE: tarife (\$ mn / bp) ve mali (\$ mn / mlr TL); 0--100 puan & 567,9 ve 7,85; 29 \\
\end{longtable}

\endgroup

\textbf{9. Senaryo etkileri, T+5, P50 (P10 / P90)}

\begingroup\gbtablo\small

\begin{longtable}[]{@{}
  >{\raggedright\arraybackslash}p{(\columnwidth - 6\tabcolsep) * \real{0.3667}}
  >{\raggedleft\arraybackslash}p{(\columnwidth - 6\tabcolsep) * \real{0.2111}}
  >{\raggedleft\arraybackslash}p{(\columnwidth - 6\tabcolsep) * \real{0.2111}}
  >{\raggedleft\arraybackslash}p{(\columnwidth - 6\tabcolsep) * \real{0.2111}}@{}}
\toprule\noalign{}
\begin{minipage}[b]{\linewidth}\raggedright
\textbf{Değişken}
\end{minipage} & \begin{minipage}[b]{\linewidth}\raggedleft
\textbf{S2}
\end{minipage} & \begin{minipage}[b]{\linewidth}\raggedleft
\textbf{S3}
\end{minipage} & \begin{minipage}[b]{\linewidth}\raggedleft
\textbf{S4}
\end{minipage} \\
\midrule\noalign{}
\endhead
\bottomrule\noalign{}
\endlastfoot
AB'ye ihracat (\$ mn) & −10 (−25 / +2) & −15 (−28 / −3) & −23 (−38 / −10) \\
Toplam ihracat (\$ mn) & −7 (−18 / +2) & −10 (−20 / −2) & −18 (−30 / −8) \\
Nihai mal ithalatı (\$ mn) & −12 (−39 / +18) & −18 (−44 / +11) & −37 (−67 / −8) \\
Cari denge (\$ mn) & +3 (−13 / +17) & +5 (−10 / +20) & +15 (0 / +30) \\
Üretim (\$ mn) & +4 (−17 / +22) & +7 (−13 / +26) & +18 (0 / +38) \\
İstihdam (kişi) & +10 (−46 / +62) & +19 (−37 / +73) & +52 (0 / +106) \\
TÜFE (bp, birinci tur) & +0,14 (+0,12 / +0,17) & +0,27 (+0,23 / +0,33) & +0,81 (+0,70 / +0,93) \\
Yatırım ve DYY (bütün senaryolar) & --- & & \\
\end{longtable}

\endgroup

\emph{Kaynak:} \etiketT{} f3e: kısmi denge, toplamlar Faz 2B GE dağılımına kalibre; S4'te 1/98 kaybı olasılığı 0,25.

\textbf{13. Kazanan ve kaybedenler; bölgesel ve sosyal etki; telafi.} S4'te iç ikame kazancı (+52 kişi); AB ihracatçısı küçük kayıp.

\textbf{15. Güven düzeyi; veri boşlukları; \etiketD{}.} Düşük (AB verimlilik verisi yok; YSHİ gizli).

\hypertarget{nace_kartlar_10-12-333-1200-tuxfctuxfcn-uxfcruxfcnleri-imalatux131--gb-statuxfcsuxfc-karma--birincilikincil-arketip-b--duruux15f-koru}{%
\subsection{3.33 12.00 Tütün ürünleri imalatı --- GB statüsü: KARMA \textbar{} Birincil/ikincil arketip: B \textbar{} Duruş: Koru}\label{nace_kartlar_10-12-333-1200-tuxfctuxfcn-uxfcruxfcnleri-imalatux131--gb-statuxfcsuxfc-karma--birincilikincil-arketip-b--duruux15f-koru}}

\emph{Çok uluslu üretim üssü. AB'ye \$116 mn ihracat; sigarada AB MFN'si \%57,6, S4'te AB pazarı kapanır. S2'de de \%70 menşeli yaprak kuralı Amerikan harmanını dışarıda bırakıyor. İthalatın \%48'i AB'den, ısıtılmış tütün ve sigara: S4 TÜFE etkisi 2,7 bp.}

\textbf{1. Kimlik.} Rev. 2.1: 12.00 (1:1); CPA 12.00. Başlıca HS4 (TR ihracatı 2024): 2402 \%62, 2403 \%27, 2401 \%11. BEC payları (ticaret ağırlıklı): ara \%45 / sermaye \%0 / tüketim \%55. NACE 12.00 yaprak tütün işlemeyi de (sapından ayırma, kurutma) içeriyor; AB pazarındaki rakip listesi (Brezilya, Hindistan, Malavi) bu yüzden yaprak tütün ihracatçıları.

\textbf{2. Ölçek}

\begingroup\gbtablo\small

\begin{longtable}[]{@{}
  >{\raggedright\arraybackslash}p{(\columnwidth - 4\tabcolsep) * \real{0.2617}}
  >{\raggedright\arraybackslash}p{(\columnwidth - 4\tabcolsep) * \real{0.4493}}
  >{\raggedright\arraybackslash}p{(\columnwidth - 4\tabcolsep) * \real{0.2890}}@{}}
\toprule\noalign{}
\begin{minipage}[b]{\linewidth}\raggedright
\textbf{Gösterge}
\end{minipage} & \begin{minipage}[b]{\linewidth}\raggedright
\textbf{Değer}
\end{minipage} & \begin{minipage}[b]{\linewidth}\raggedright
\textbf{Etiket ve kaynak}
\end{minipage} \\
\midrule\noalign{}
\endhead
\bottomrule\noalign{}
\endlastfoot
Üretim değeri; katma değer & \$2.816 mn; \$656 mn (GSYH'nin \%0,05) & TÜİK YSHİ 2024 \etiketV{} \\
Çalışan; girişim (0 / 1--9 / 10--49 / 50--249 / 250+) & 5.769; 80 (8 / 43 / 21 / 8 / 5) & TÜİK YSHİ ve İş Kayıtları 2024 \etiketV{}; gizli büyüklük hücreleri \etiketT{} \\
250+ girişimlerin istihdam payı; AB'ye göre kişi başı KD & \%75; 0,48 & \etiketT{} (AB27 SBS 2023, nominal kur) \\
İlk 3 İBBS-2 bölgesi & İzmir \%59; Samsun, Tokat, Çorum, Amasya \%11; Adana, Mersin \%6 & \etiketT{} iş kayıtları + SGK 2025, RAS \\
Yabancı sermaye; çapa yatırımcılar & Philip Morris (Torbalı; 2024'te üretiminin \%11'inden fazlasını ihraç etti) \etiketV{}, JTI (Torbalı), BAT \etiketD{}, KT\&G \etiketH{} & \etiketH{} (firma bilgisi; pay verisi yok: VB-30) \\
\end{longtable}

\endgroup

\textbf{3. Dış ticaret}

\begingroup\gbtablo\small

\begin{longtable}[]{@{}
  >{\raggedright\arraybackslash}p{(\columnwidth - 2\tabcolsep) * \real{0.4802}}
  >{\raggedright\arraybackslash}p{(\columnwidth - 2\tabcolsep) * \real{0.5198}}@{}}
\toprule\noalign{}
\begin{minipage}[b]{\linewidth}\raggedright
\textbf{Gösterge}
\end{minipage} & \begin{minipage}[b]{\linewidth}\raggedright
\textbf{Değer (BACI 2024, \$; \etiketV{}, eşleme \etiketT{})}
\end{minipage} \\
\midrule\noalign{}
\endhead
\bottomrule\noalign{}
\endlastfoot
İhracat (2024); 2020--24 değişimi & \$876 mn; \%19 \\
AB payı; ilk 5 AB pazarı & \%13; ROU \%36; NLD \%15; DEU \%12; BEL \%10; MLT \%7 \\
İlk 5 AB dışı pazar & IRQ \%20; IRN \%11; AZE \%7; KWT \%7; USA \%7 \\
İthalat; AB ve Çin payı; ilk kaynaklar & \$958 mn; AB \%48, Çin \%1; DEU \%14; NLD \%11; BEL \%9; IND \%6; BRA \%6 \\
Dış ticaret dengesi & \$−81 mn \\
İthal girdi içeriği (I-O, toplam) / KD÷üretim & \%20 / \%23 \\
RCA dünya (2020 → 2024) / AB pazarında & 1,64 → 1,52 / 0,71 \\
Birim değer ÷ AB pazarı ortalaması & 0,88 \\
AB pazar payı (AB dışı tedarik içinde, 2020 → 2024); TR'nin sırası & \%3,9 → \%2,8; 12. \\
\end{longtable}

\endgroup

\textbf{4. AB tarafı maruziyet}

\begin{itemize}
\tightlist
\item
  \kod{τ\_\allowbreak{}AB} (S4 artışı, rota A, TR→AB 2025 ağırlıklı): \%51,3; 1/98 kaybıyla \%51,8; tam MFN ortalaması \%51,8. Rejim: GB \%0, İTÜ \%95, tarım \%5 \etiketV{}/\etiketT{}.
\item
  Tepe satırlar (S4 artışı ≥ \%10; en büyük 3): 8 satır, AB ihracatının \%95'i: KN 2402 20 90 \%57,6 (€98,8 mn); KN 2403 99 90 \%16,6 (€15,6 mn); KN 2403 11 00 \%74,9 (€9,0 mn).
\item
  AB pazarında AB dışı rakipler (tercih durumu): BRA \%21; IND \%12; USA \%7; MWI \%6; CHN \%5. Tercihli rakip payı \%29.
\item
  Menşe (S2 / S3): risk Y/Y; \kod{c\_\allowbreak{}RoO} \%5,0 / \%6,0; \kod{p\_\allowbreak{}uyum} 0,40 / 0,30 \etiketH{}. Kural: 2402: başlık 2401 tütününün ağırlıkça en az \%70'i menşeli (Amerikan harmanı ithal Virginia/Burley içerir).
\item
  EUDR kapsamındaki AB ihracatı payı: \%0; SKDM: \%0.
\item
  Menşe: 2402 için `başlık 2401 tütününün ağırlıkça en az \%70'i menşeli'. Türk sigaralarında Virginia ve Burley ithal; şark tütünü yerli. Uyum olasılığı S2'de 0,40 \etiketH{}.
\item
  Çok uluslu şirket AB pazarını AB fabrikalarından besleyebilir. PMI'ın Yunanistan'daki Papastratos fabrikası 2018'de ısıtılmış tütüne döndü \etiketV{}.
\end{itemize}

\textbf{5. Türkiye tarafı maruziyet}

\begin{itemize}
\tightlist
\item
  Uygulanan tarife (2026, ithalat ağırlıklı): MFN \%29,1, İGV dahil \%29,1 \etiketV{}. S4'te AB menşeline (c2): \%22,3.
\item
  Kritik girdiler (HS4 2401), Türkiye'nin toplam ithalatı 2024: \$519 mn (AB \%29; BEL \%13, IND \%11, DEU \%11) \etiketV{}.
\item
  ÖTV ile gümrük etkileşimi: sigarada ÖTV oransal + asgari maktu. S4'te AB'den gelen ürüne MFN (sigarada \%41,7, ısıtılmış tütün ve diğer mamul tütünde \%29,1) perakende fiyatın vergi öncesi kısmını artırıyor. Isıtılmış tütünde (KN 2404 11; HS2017'de 2403 99) MFN \%0, S4 etkisinin \%39'unu kaldırır.
\item
  Isıtılmış tütünün Türkiye'de üretilip üretilmediği \etiketD{} (PMI Torbalı'da ısıtılmış tütün hattı doğrulanmadı).
\end{itemize}

\textbf{6. Kabiliyet}

\begin{itemize}
\tightlist
\item
  KS 39,8 (sıra 206/230; ağırlık duyarlılığında P10--P90 157--219). Bileşenler: RCA 49, PCI 8 (PCI −1,08), yurt içi KD 58, Ar-Ge 25, tedarikçi 84, verimlilik 57, firma 17 \etiketT{}.
\item
  Ürün uzayında atlanabilir ürünler (sınıf içi; yoğunluk üst çeyrek, PCI \textgreater{} sınıf): yok \etiketT{}.
\item
  \textbf{Kırmızı çizgi yetkinlikleri:} Şark tütünü harman bilgisi ve yaprak işleme (Ege, Karadeniz); fabrikaların bölgesel ihracat merkezi rolü.
\item
  \textbf{Geliştirme hedefi:} T+5: AB'ye kaybın AB dışı (Irak, Körfez, Kafkasya) ihracatla telafisi.
\end{itemize}

\textbf{7. Stratejik önem.} SS 0,0 (en yüksek boyut: --- 0; ağırlıklı ortalama 0,0) \etiketH{}. Stratejik değil; kamu sağlığı ve vergi geliri açısından izlenir.

\textbf{8. Makro duyarlılık}

\begingroup\gbtablo\small

\begin{longtable}[]{@{}lr@{}}
\toprule\noalign{}
\textbf{Gösterge} & \textbf{Değer (\etiketT{})} \\
\midrule\noalign{}
\endhead
\bottomrule\noalign{}
\endlastfoot
TÜFE ağırlığı (doğrudan, COICOP 2018 eşlemesi); tüketimde ithal pay & \%2,16; \%11 \\
ÜFE vekili (imalat üretimindeki pay) & \%0,42 \\
EDE: 1 puan tarife, tüm menşe / yalnız AB (bp; ρ = 1) & 0,310 / 0,148 \\
Geçişkenlik ρ (senaryolarda) & 0,6--1,0 (P50 0,8) \\
Piyasa yoğunlaşması (250+ istihdam payı) & \%75 \\
CKE: tarife (\$ mn / bp) ve mali (\$ mn / mlr TL); 0--100 puan & 215,2 ve 10,30; 26 \\
\end{longtable}

\endgroup

\textbf{9. Senaryo etkileri, T+5, P50 (P10 / P90)}

\begingroup\gbtablo\small

\begin{longtable}[]{@{}
  >{\raggedright\arraybackslash}p{(\columnwidth - 6\tabcolsep) * \real{0.2477}}
  >{\raggedleft\arraybackslash}p{(\columnwidth - 6\tabcolsep) * \real{0.4209}}
  >{\raggedleft\arraybackslash}p{(\columnwidth - 6\tabcolsep) * \real{0.1657}}
  >{\raggedleft\arraybackslash}p{(\columnwidth - 6\tabcolsep) * \real{0.1657}}@{}}
\toprule\noalign{}
\begin{minipage}[b]{\linewidth}\raggedright
\textbf{Değişken}
\end{minipage} & \begin{minipage}[b]{\linewidth}\raggedleft
\textbf{S2}
\end{minipage} & \begin{minipage}[b]{\linewidth}\raggedleft
\textbf{S3}
\end{minipage} & \begin{minipage}[b]{\linewidth}\raggedleft
\textbf{S4}
\end{minipage} \\
\midrule\noalign{}
\endhead
\bottomrule\noalign{}
\endlastfoot
AB'ye ihracat (\$ mn) & −35 (−78 / +2) & −34 (−63 / −7) & −39 (−62 / −16) \\
Toplam ihracat (\$ mn) & −25 (−57 / +2) & −24 (−46 / −5) & −31 (−50 / −13) \\
Nihai mal ithalatı (\$ mn) & −9 (−30 / +26) & −14 (−36 / +15) & −46 (−84 / −10) \\
Cari denge (\$ mn) & −15 (−26 / −10) & −9 (−17 / −5) & +11 (−3 / +28) \\
Üretim (\$ mn) & −19 (−33 / −13) & −11 (−21 / −6) & +14 (−3 / +35) \\
İstihdam (kişi) & −27 (−48 / −17) & −16 (−30 / −9) & +20 (−5 / +49) \\
TÜFE (bp, birinci tur) & +0,29 (+0,25 / +0,33) & +0,55 (+0,47 / +0,63) & +2,73 (+2,34 / +3,09) \\
Yatırım ve DYY (bütün senaryolar) & AB'ye yönelik üretim hattının AB'ye kayması olası (S4 ve S2) \etiketH{} & & \\
\end{longtable}

\endgroup

\emph{Kaynak:} \etiketT{} f3e: kısmi denge, toplamlar Faz 2B GE dağılımına kalibre; S4'te 1/98 kaybı olasılığı 0,25.

\textbf{10. Politika paketi}

\begingroup\gbtablo\footnotesize

\begin{longtable}[]{@{}
  >{\raggedright\arraybackslash}p{(\columnwidth - 16\tabcolsep) * \real{0.1117}}
  >{\raggedright\arraybackslash}p{(\columnwidth - 16\tabcolsep) * \real{0.1010}}
  >{\raggedright\arraybackslash}p{(\columnwidth - 16\tabcolsep) * \real{0.1212}}
  >{\raggedright\arraybackslash}p{(\columnwidth - 16\tabcolsep) * \real{0.1067}}
  >{\raggedright\arraybackslash}p{(\columnwidth - 16\tabcolsep) * \real{0.0944}}
  >{\raggedright\arraybackslash}p{(\columnwidth - 16\tabcolsep) * \real{0.0845}}
  >{\raggedright\arraybackslash}p{(\columnwidth - 16\tabcolsep) * \real{0.1644}}
  >{\raggedright\arraybackslash}p{(\columnwidth - 16\tabcolsep) * \real{0.1246}}
  >{\raggedright\arraybackslash}p{(\columnwidth - 16\tabcolsep) * \real{0.0916}}@{}}
\toprule\noalign{}
\begin{minipage}[b]{\linewidth}\raggedright
\textbf{Eylem ID}
\end{minipage} & \begin{minipage}[b]{\linewidth}\raggedright
\textbf{Araç}
\end{minipage} & \begin{minipage}[b]{\linewidth}\raggedright
\textbf{Parametre}
\end{minipage} & \begin{minipage}[b]{\linewidth}\raggedright
\textbf{Sorumlu}
\end{minipage} & \begin{minipage}[b]{\linewidth}\raggedright
\textbf{Hukuki araç}
\end{minipage} & \begin{minipage}[b]{\linewidth}\raggedright
\textbf{Bütçe}
\end{minipage} & \begin{minipage}[b]{\linewidth}\raggedright
\textbf{Başlangıç--bitiş; sunset}
\end{minipage} & \begin{minipage}[b]{\linewidth}\raggedright
\textbf{Performans koşulu}
\end{minipage} & \begin{minipage}[b]{\linewidth}\raggedright
\textbf{KPI}
\end{minipage} \\
\midrule\noalign{}
\endhead
\bottomrule\noalign{}
\endlastfoot
E-12.00-01 & S2'de PSR talebi & 2402: `başlık değişikliği; menşesiz 2401 tütünü ağırlıkça ≤ \%50' & Ticaret Bakanlığı & S2 menşe protokolü & İdari & T−18 → T; sunset: --- & --- & Kabul (olasılık düşük \etiketH{}) \\
E-12.00-02 & S4'te ısıtılmış tütünde MFN \%0 & KN 2404 11 ve 2403 99: \%29,1 → \%0, tüm menşeler; ÖTV yapısı değişmez & Ticaret Bakanlığı, Hazine ve Maliye Bakanlığı & İthalat Rejimi Kararı eki & Gelir kaybı ≈ 0,5 mlr TL/yıl (AB dışı ithalat) \etiketT{} & T → T+36; sunset: T+36 & --- & TÜFE tütün alt endeksi; −1,07 bp \\
\end{longtable}

\endgroup

\textbf{11. Müzakere pozisyonu}

\begin{itemize}
\tightlist
\item
  AB ile savunma: PSR talebi; kabul görmezse AB'ye ihracat kaybını kabul et, karşılığında başka bir başlıkta taviz iste.
\end{itemize}

\textbf{12. Erken uyarı ve tetikler}

\begingroup\gbtablo\small

\begin{longtable}[]{@{}
  >{\raggedright\arraybackslash}p{(\columnwidth - 4\tabcolsep) * \real{0.2574}}
  >{\raggedright\arraybackslash}p{(\columnwidth - 4\tabcolsep) * \real{0.1165}}
  >{\raggedright\arraybackslash}p{(\columnwidth - 4\tabcolsep) * \real{0.6261}}@{}}
\toprule\noalign{}
\begin{minipage}[b]{\linewidth}\raggedright
\textbf{Gösterge}
\end{minipage} & \begin{minipage}[b]{\linewidth}\raggedright
\textbf{Eşik}
\end{minipage} & \begin{minipage}[b]{\linewidth}\raggedright
\textbf{Tetiklenecek önlem}
\end{minipage} \\
\midrule\noalign{}
\endhead
\bottomrule\noalign{}
\endlastfoot
AB'ye sigara ihracatı (aylık) & Yıllık −\%30 & Çok uluslu şirketlerle üretim hattı kararı görüşmesi; AB dışı pazar desteği \\
Kayıt dışı sigara payı (S4) & +2 puan & Gümrük ve Muhafaza denetimi \\
\end{longtable}

\endgroup

\textbf{13. Kazanan ve kaybedenler; bölgesel ve sosyal etki; telafi.} Kaybeden: AB'ye ihracat yapan hat (P50'de istihdam etkisi küçük; iç ikame dengeliyor). Bölge: İzmir (\%59).

\textbf{14. Riskler ve azaltım}

\begin{itemize}
\tightlist
\item
  AB'ye yönelik üretimin AB'ye kayması → kabullenilen kayıp; AB dışı pazar.
\item
  Kaçakçılık (fiyat artışı) → MFN \%0 (ısıtılmış tütün) ve denetim.
\end{itemize}

\textbf{15. Güven düzeyi; veri boşlukları; \etiketD{}.} Orta. \etiketD{}: ısıtılmış tütün üretiminin yeri; BAT'ın Türkiye'deki üretim durumu.

\hypertarget{nace_kartlar_10-12-4-ortak-eylemler-1012}{%
\section{4. Ortak eylemler (10--12)}\label{nace_kartlar_10-12-4-ortak-eylemler-1012}}

\begin{landscape}

\textbf{Tablo 4.1 --- Bölüm düzeyi eylemler}

\nopagebreak[4]

\begingroup\gbtablo\footnotesize

\begin{longtable}[]{@{}
  >{\raggedright\arraybackslash}p{(\columnwidth - 18\tabcolsep) * \real{0.0564}}
  >{\raggedright\arraybackslash}p{(\columnwidth - 18\tabcolsep) * \real{0.0760}}
  >{\raggedright\arraybackslash}p{(\columnwidth - 18\tabcolsep) * \real{0.1911}}
  >{\raggedright\arraybackslash}p{(\columnwidth - 18\tabcolsep) * \real{0.1218}}
  >{\raggedright\arraybackslash}p{(\columnwidth - 18\tabcolsep) * \real{0.0855}}
  >{\raggedright\arraybackslash}p{(\columnwidth - 18\tabcolsep) * \real{0.0815}}
  >{\raggedright\arraybackslash}p{(\columnwidth - 18\tabcolsep) * \real{0.1027}}
  >{\raggedright\arraybackslash}p{(\columnwidth - 18\tabcolsep) * \real{0.0937}}
  >{\raggedright\arraybackslash}p{(\columnwidth - 18\tabcolsep) * \real{0.0855}}
  >{\raggedright\arraybackslash}p{(\columnwidth - 18\tabcolsep) * \real{0.1058}}@{}}
\toprule\noalign{}
\begin{minipage}[b]{\linewidth}\raggedright
\textbf{Eylem ID}
\end{minipage} & \begin{minipage}[b]{\linewidth}\raggedright
\textbf{Araç}
\end{minipage} & \begin{minipage}[b]{\linewidth}\raggedright
\textbf{Parametre}
\end{minipage} & \begin{minipage}[b]{\linewidth}\raggedright
\textbf{Sorumlu}
\end{minipage} & \begin{minipage}[b]{\linewidth}\raggedright
\textbf{Hukuki araç}
\end{minipage} & \begin{minipage}[b]{\linewidth}\raggedright
\textbf{Bütçe}
\end{minipage} & \begin{minipage}[b]{\linewidth}\raggedright
\textbf{Başlangıç--bitiş}
\end{minipage} & \begin{minipage}[b]{\linewidth}\raggedright
\textbf{Performans koşulu}
\end{minipage} & \begin{minipage}[b]{\linewidth}\raggedright
\textbf{Sunset}
\end{minipage} & \begin{minipage}[b]{\linewidth}\raggedright
\textbf{KPI}
\end{minipage} \\
\midrule\noalign{}
\endhead
\bottomrule\noalign{}
\endlastfoot
E-10-01 & 1/98 sürekliliği ve S2'ye aktarımı & (i) 1/98'in Ankara Anlaşması'na dayanan bağımsız hukuki dosyası; (ii) S2 metninde 1/98 tavizlerinin ve İTÜ rejiminin aktarımı; (iii) S4'te 1/98'i askıya alma girişimine karşı DTÖ ve Ortaklık Konseyi yolu & Ticaret Bakanlığı (Uluslararası Anlaşmalar ve AB GM), Dışişleri (AB Başkanlığı), Tarım ve Orman Bakanlığı (TOB) & Ortaklık Konseyi kararı; S2 anlaşma metni (Ek: tarım tavizleri) & İdari & T−18 → T & S2 metninde 1/98 tavizleri en az bugünkü kapsamda & Yok (anlaşmayla kalıcı) & 1/98 kapsamındaki AB ihracatı (€4,0 mlr civarı) T+1'de ≥ \%95 \\
E-10-02 & Menşe uyumlu DİR: yerli eşdeğer eşya & AB'ye yönelik İTÜ ihracatında (Fasıl 17--21) DİR izin belgesiyle yerli (menşeli) buğday ve şeker kullanımı; ithal eşya serbest dolaşıma vergisiz girer (eşdeğer eşya). S2'de PSR talebi: Fasıl 19'da menşesiz tahıla fabrika çıkış fiyatının \%15'ine kadar tolerans & Ticaret Bakanlığı (İhracat GM), TMO, Gümrükler GM & DİR Kararı ve Tebliği değişikliği; S2 menşe protokolü & Gelir etkisi ≈ 0 (vergi zaten alınmıyor) \etiketT{} & T−12 → T+60 & Belge sahibinin AB ihracatı ≥ yerli eşdeğer miktar & T+60 (PSR yürürlüğe girince gözden geçirilir) & Bisküvi/ekmek AB ihracatında menşe ispatlı pay ≥ \%85 \\
E-10-03 & EUDR uyum masası & Kakao (10.82), kahve (10.83), soya/palm (10.41, 10.91) ve sığır (10.11, 15.11) için coğrafi konum verisi ve özen beyanı desteği; ithalatçı birlikleriyle ortak veri havuzu & TOB, Ticaret Bakanlığı, TİM (ilgili ihracatçı birlikleri) & Bakanlık genelgesi; TİM yönetim kurulu kararı & ≈ 0,05 mlr TL/yıl \etiketT{} & T−18 → T+24 & Kullanan firma sayısı & T+24 & EUDR nedeniyle AB'de reddedilen sevkiyat = 0 \\
E-10-04 & Ticaret savunması erken uyarı & Su ürünleri (10.20) ve işlenmiş meyve-sebzede (10.39) AB anti-damping ve sübvansiyon soruşturmalarına hazırlık: destek programlarının ASCM uyum dosyası, maliyet verisi & Ticaret Bakanlığı (İthalat GM ve İhracat GM), TOB, ilgili ihracatçı birlikleri & İç talimat & İdari & T−12 → T+36 & --- & T+36 & Soruşturma açılırsa 30 gün içinde tam yanıt \\
E-10-05 & S4 geçiş tamponu (yalnız S4) & AB'ye İTÜ ihracatçılarına (10.62, 10.72, 10.82, 10.89, 11.07) ihracat kredi sigortası ve pazar çeşitlendirme desteği; yıllık tavan & Türk Eximbank, Ticaret Bakanlığı & Eximbank programı; İhracat Destekleri Kararı & 0,3 mlr TL/yıl × 3 yıl \etiketT{} & T → T+36 & AB dışı ihracat artışı & T+36 & Tampon kapsamındaki firmaların toplam ihracatı ≥ T−12 düzeyi \\
E-10-06 & Tarımda hücum gündemi (S2) & Meyve-sebzede giriş fiyatı sistemi istisnası, domates ve işlenmiş ürün kotalarının genişletilmesi, su ürünlerinde kota yerine sıfır tarife, kanatlı ve süt ürünlerinde SPS listelenme & Ticaret Bakanlığı, TOB & S2 anlaşma metni (tarım ve SPS bölümleri) & İdari & T−18 → T & AB'nin tarım hassas listesine karşılıklı taviz & --- & Tavizli satırların AB ihracatı T+5'te +\%10 \\
\end{longtable}

\endgroup

\end{landscape}

\emph{Kaynak:} Bütçeler \etiketT{}, kurumsal sorumluluk \etiketV{} (mevcut görev dağılımı), tutarlar 2026 fiyatlarıyla; mali yükler defterde toplanıyor.

\hypertarget{nace_kartlar_10-12-5-kuxfcmuxfclatif-defter}{%
\section{5. Kümülatif defter}\label{nace_kartlar_10-12-5-kuxfcmuxfclatif-defter}}

\Needspace{7\baselineskip}

\textbf{Tablo 5.1 --- Kümülatif defter (Faz 3, parça 1 sonrası)}

\nopagebreak[4]

\begingroup\gbtablo\small

\begin{longtable}[]{@{}
  >{\raggedright\arraybackslash}p{(\columnwidth - 2\tabcolsep) * \real{0.1515}}
  >{\raggedright\arraybackslash}p{(\columnwidth - 2\tabcolsep) * \real{0.8485}}@{}}
\toprule\noalign{}
\begin{minipage}[b]{\linewidth}\raggedright
\textbf{Kalem}
\end{minipage} & \begin{minipage}[b]{\linewidth}\raggedright
\textbf{Değer}
\end{minipage} \\
\midrule\noalign{}
\endhead
\bottomrule\noalign{}
\endlastfoot
Tamamlanan / kalan sınıf & 33 / 197 C sınıfı; ek kapsam (B, D, E, hizmet) 0 / 29 \\
Kullanılan / kalan enflasyon bütçesi (bp; birinci tur, P50) & S2 yolu: şok 0,9 + önlem −0,11 = 0,8 bp; S4 yolu: şok 6,2 + önlem −1,55 = 4,7 bp. Kümülatif 50 bp bütçeden kalan: S2 49,2, S4 45,3 bp \\
Kümülatif cari denge katkısı (\$ mlr/yıl, T+5) & Senaryo P50: S2 −0,044, S4 +0,030 (S4 P10 −1,06); önlemler: S2 +0,071, S4 +0,010 \etiketT{}/\etiketH{} \\
Kümülatif mali maliyet (mlr TL/yıl; GSYH \%) & S2 yolu 0,83 (\%0,001); S4 yolu 1,69 (\%0,002) \\
Riskteki istihdam (kişi, T+5) & S2 P50 −284; S4 P50 −15, P10 −4.698. S4 P10'da kaybı \%5'i aşan sınıf: yok \\
Açık kalan veri boşlukları & VB-30: sınıf düzeyinde yabancı sermaye payı (TÜİK yabancı kontrollü girişim istatistikleri yalnız bölüm düzeyinde); VB-31: Türkiye'ye özgü sınıf düzeyi Ar-Ge yoğunluğu (OECD sınıflaması vekil); VB-32: DTÖ bağlı oranları, satır düzeyi (WTO CTS); 10--12'de tarife artışı önerilmediği için bağlayıcı değil, sonraki parçalarda gerekli; VB-33: TMO yerli buğday fiyatı ile ithal paritesi farkı (menşe maliyeti c\_RoO'nun kalibrasyonu); VB-34: EUDR kapsamındaki ihracatçıların tedarik zinciri verisi; VB-35: YSHİ ile BACI sınıf eşlemesindeki tutarsızlık (10.86; 11.03, 11.04, 11.06'da YSHİ verisi yok); VB-06 (sürüyor): 1/98 tercih oranları satır düzeyinde; 1/98 kaybı varyantı bu yüzden üst sınır \\
Varsayım kaydına eklenen yeni maddeler & V-30: Faz 3 senaryoları --- toplam GE'den, dağılım kısmi dengeden (ortak z, N = 2.000); V-31: S4'te 1/98'in askıya alınma olasılığı 0,25 (hasmane AB) \etiketH{}; V-32: ticaret esnekliği σ --- ProTEE HS6, sınıf içi AB ihracatı ağırlıklı medyan, {[}2, 15{]}; V-33: sınır sürtünmesi \%0,5--2 (yalnız GB ve İTÜ), S2'de × 0,5; AB ihracatçısının menşe maliyeti S2 \%0,5, S3 \%1; V-34: arketip eşikleri (KS, AME, ithalat penetrasyonu, EDE ≥ 0,3 bp/puan, SS ≥ 60) \etiketH{}; V-35: SS boyut puanları ve ağırlıkları \etiketH{} (veri/f3\_ss\_boyutlar.csv; 10--12 gözden geçirildi, diğerleri ilk tahmin); V-36: CKE tanımı --- 1 puan tarife / \%1 üretim desteği, 5 yıl, kapasite yatırımının ithalatı (sermaye/hasıla 1, ithal payı 0,6) düşülü; V-37: KN8 → NACE eşlemesinde CPA 2.2 → 2.1 kimlik bağlantısı önceliği (D-1) \\
\end{longtable}

\endgroup

\emph{Kaynak:} Bu tablo \kod{kod/\allowbreak{}f3g\_\allowbreak{}kartlar.\allowbreak{}py} ile matristen ve kartlardaki önlem kayıtlarından üretildi. Senaryo şokları politika seçimi değil, çıkışın kendisi; önlemler kartların eylem tablolarından. Makro mutabakat Faz 6'da.

\textbf{Kontrol listesi (parça 1).}

\begin{enumerate}
\def\labelenumi{\arabic{enumi}.}
\tightlist
\item
  Etiketler: her sayının etiketi, birimi ve dönemi var. Sınıf değerleri 2024, AB tarifeleri 2026, senaryolar T+5.
\item
  Toplam ve çift sayım: bölüm toplamları matristen (230 sınıf); EDE toplamı I-O ile tutarlı (22,0 / 23,5 bp). KKTC ve KN8 düzeltmeleri yapıldı (D-1, D-2).
\item
  Kısıtlar:

  \begin{itemize}
  \tightlist
  \item
    S2 yolunda 10--12'nin birinci tur TÜFE etkisi 0,9 bp (önlem değil, sürtünme şoku).
  \item
    S4 yolunda 6,2 bp; önlemler (E-10.82-01, E-10.86-01, E-12.00-02) bunu 1,55 bp düşürüyor (Tablo 5.1).
  \item
    Mali kısıt: önlemler S2 yolunda 0,8, S4 yolunda 1,7 mlr TL/yıl (GSYH'nin \%0,002'si; Tablo 5.1).
  \item
    İstihdam: S4'ün P10 değerinde 10.39 −\%4,9 ile \%5 eşiğine yakın; tetik tanımlandı.
  \end{itemize}
\item
  Tarife önerileri: 10--12'de bağlı oranı aşan öneri yok. Öneriler ya MFN'yi koruyor ya indiriyor; ITA bu bölümleri kapsamıyor; STA kontrolü kartlarda.
\item
  DTÖ: E-10-02, DİR'in mevcut eşdeğer eşya rejimi içinde. Tarım ürünlerinde ihracat sübvansiyonu yasağına (Nairobi 2015) takılmamak için TMO satışı ihracat koşullu indirimle yapılmıyor.
\item
  ``Pazartesi sabahı'': eylemlerde kurum, hukuki araç, bütçe, tarih ve KPI var. Bütçe ve tarih kalibrasyonu Faz 5--8'de.
\item
  AB tepkisi: 1/98'in askıya alınması (0,25) ve ticaret savunması (10.20, 10.39) fiyatlandı.
\item
  Zaman yolu: eylemler T−18 → T+60.
\item
  Kaçınılacak hatalar: tarım satırlarına GB tarifesi eklenmedi (hata 15); spesifik vergiler AVE'ye çevrildi.
\item
  Kaçamak ifade bırakılmadı. Belirsiz olgular \etiketD{} olarak işaretli.
\item
  Olumsuz bulgular girişte (sigara, fındık kaldıracının sınırı).
\end{enumerate}

\begin{gbilerleme}

İLERLEME: Faz 3 parça 1 / 16 tamamlandı (33 / 230 C sınıfı kartı) --- SONRAKİ ADIM: Faz 3 parça 2 (13--15: tekstil, giyim, deri); DEVAM bekleniyor.

\end{gbilerleme}
